\documentclass[a4paper,UTF8,zihao=-4,openany]{ctexbook}
\usepackage{amsmath,amssymb,amsthm,mathtools,bm}
\usepackage[a4paper,top=2.5cm,bottom=2.5cm,left=2.65cm,right=2.65cm,headheight=15pt]{geometry}
\usepackage{booktabs,array,graphicx,xcolor,enumitem,hyperref}
\usepackage{caption}
\usepackage{tikz}
\usetikzlibrary{arrows.meta,positioning,fit,calc,shapes.geometric,decorations.pathreplacing}
\tikzset{
	idi node/.style={draw, rounded corners=2pt, align=center, minimum height=0.72cm, inner sep=5pt, font=\small},
	idi input/.style={idi node, fill=blue!8, draw=blue!55!black},
	idi process/.style={idi node, fill=orange!12, draw=orange!65!black},
	idi state/.style={idi node, fill=green!10, draw=green!55!black},
	idi output/.style={idi node, fill=red!8, draw=red!55!black},
	idi arrow/.style={-{Latex[length=2mm]}, thick},
	idi feedback/.style={-{Latex[length=2mm]}, dashed, thick, gray!75!black}
}
\usepackage{algorithm}
\usepackage{algpseudocode}
\usepackage{microtype}
\usepackage{listings}
\lstset{basicstyle=\ttfamily\small,breaklines=true,columns=fullflexible,frame=single}
\captionsetup{font=small,labelfont=bf}

\hypersetup{colorlinks=true,linkcolor=blue,citecolor=teal,urlcolor=blue}
\setcounter{secnumdepth}{3}
\setcounter{tocdepth}{1}
\newcommand{\R}{\mathbb{R}}
\newcommand{\E}{\mathbb{E}}
\newcommand{\Var}{\operatorname{Var}}
\newcommand{\KL}{\operatorname{KL}}
\newcommand{\argmin}{\mathop{\mathrm{arg\,min}}}
\newcommand{\argmax}{\mathop{\mathrm{arg\,max}}}
\newcommand{\softmax}{\operatorname{softmax}}
\newcommand{\diag}{\operatorname{diag}}
\newtheorem{definition}{定义}[chapter]
\newtheorem{proposition}{命题}[chapter]
\newtheorem{example}{例}[chapter]
\newtheorem{remark}{说明}[chapter]
\newtheorem{theorem}{定理}[chapter]
\newtheorem{lemma}{引理}[chapter]

\title{工业数据智能入门指南 \\ Introduction to Industrial Data Intelligence}
\author{张恺}
\date{\today}

\begin{document}
\maketitle
\frontmatter
\tableofcontents
\chapter*{写在前面}
本指南面向刚开始接触数据智能、机器学习或深度学习的读者。工业数据智能包括多类任务：预测设备故障或风功率，根据有限观测重建风场，通过缺失值填充恢复不完整数据，定位图像中的缺陷，以及根据手册和工单回答问题。本书以设备维护作为贯穿案例之一，同时结合这些任务介绍数据、模型、优化、推理和不确定性。公式用来说明计算怎样成立，例子用来说明结果会怎样影响现场判断。

\section*{先从一个不需要算法名称的问题开始}
假设每台设备每小时记录一次温度和振动。现在是周一上午八点，值班人员想知道：从此刻起的二十四小时内，这台设备是否可能发生故障，以便决定是否安排复检。这里已经有三个不同的对象：温度和振动是\textbf{输入}，未来二十四小时内是否故障是\textbf{预测目标}，是否复检是人员或系统需要作出的\textbf{决策}。目标与决策有关，却不是同一件事：同样的故障概率，在复检便宜和复检昂贵时，可能对应不同的行动。

过去的记录可以帮助建立输入与目标之间的关系。例如，把每个历史预测时刻之前已经收到的温度和振动作为输入，等该时刻之后的二十四小时结束、故障记录确认后，再补上实际结果。这样一对输入和实际结果称为一个\textbf{带标签样本}。利用许多历史样本调整计算规则的过程叫\textbf{训练}；把调整好的规则用于当前输入、估计尚未知晓的目标，叫\textbf{预测}。实际运行模型得到输出，通常叫\textbf{推理}；不过在概率章节中，“推断”还可能指根据已知观测计算未知量的条件分布，需要按上下文区分。模型算出的是估计，不是提前观察到了未来事实。训练时可以使用后来确认的结果来学习，但当前预测不能偷看未来。

本书中的\textbf{模型}就是这样一条可以计算的规则。例如，“温度超过某个值就告警”是一条规则；把温度和振动分别乘以一个系数、相加后再输出风险，也是一条规则。规则中需要从数据确定的系数叫\textbf{参数}。机器学习研究怎样根据数据选择或调整这些规则，并检验它们对没见过的数据是否仍然有效。深度学习是机器学习的一类方法，它用多层可调整的计算来形成表示，并不意味着可以跳过数据检查和结果验证。

当模型不仅给出一个数，还描述不同结果出现的可能性时，需要用到\textbf{概率分布}。例如，预测故障概率为$0.2$，并不是说这台设备会发生“五分之一次故障”，也不是保证五台这样的设备恰好坏一台；它是一种需要用后续数据检验的风险表达。如何计算、解释和检验这样的概率，是本书第二部分的主题。

\section*{公式怎样读，读到什么程度再往下走}
不要求读者在开篇就认识全部数学符号。第一章会逐一建立它们的含义；以后遇到公式，可以按“对象是什么、输入是什么、要计算什么、结果怎样使用”的顺序阅读。例如，$\hat y=f_\theta(x)$只是在压缩一句话：把输入$x$交给参数为$\theta$的规则$f$，得到预测$\hat y$。字母上方的小帽子用来提醒我们，这是预测值，不是后来观测到的真实结果$y$。参数下标$\theta$说明规则可以随训练而改变，不表示再乘一次$\theta$。

除特别说明外，向量默认为列向量，即把几个数竖着排成一列；$\|\cdot\|_2$表示欧氏范数，可先理解为由勾股定理推广得到的长度；$\mathbb{E}$表示期望，即按照各结果的概率加权平均。它们的计算方法、适用条件和例子都在数学预备知识章中展开。第一次读到较长推导时，可以先保留“这段在解决什么问题”的主线，再回头逐步计算，不必把暂时不会推导误认为完全不懂这个任务。

设备维护案例联系了其中几类学习任务：根据运行记录预测未来二十四小时内的故障，利用巡检图像定位缺陷，再结合有效手册和历史工单组织复检依据。这些任务共享设备、时间和证据来源，但并不共享同一种标签。预测未来二十四小时内是否故障，也不等于保证提前二十四小时告警。各章中的数值设置是教学设置，不代表已经完成的工业实验。阅读时应追问：观测在何时可用，目标如何定义，模型保留了什么结构，损失优化了什么，错误如何影响决策，以及系统需要采取什么动作。

因此，本书的章节不是按模型名称组成的目录。前两部分先说明学习什么、可以用哪些数据、错误有什么代价，以及结果有多不确定；第三部分分别处理按时间排列的读数、图像、文本、不同类型观测之间的联系，以及新数据的生成，再以风场重建说明物理方程如何参与学习；第四部分讨论旧设备上的经验怎样用于新设备、不同机构怎样共同训练，以及新任务只有少量样本时怎样适应。每次引入新模型时，都应把它放回明确的案例设定中，并与简单方法比较。“基线”就是这种用于比较的参考方法，例如始终预测历史平均值；它并不是故意设置的弱对手，而是检查复杂方法是否值得使用的起点。

贯穿全书的核心观点是：先把现场能看到的东西和要做的决定说清楚，再谈模型。传感器晚到五分钟，不能当成实时输入；拆机后才拍到的裂纹照片，不能拿来证明拆机前预警有效；一次误报会占用检修人员的时间，一次漏报可能损坏设备。数据、假设、损失、优化和验证，最终都要对这些具体后果负责。它们不是彼此独立的技术清单。例如，为什么常把“预测值减去真实值”之后的差再平方？这既与如何衡量偏差有关，也与我们怎样描述观测误差有关。又如，为什么训练时有时要惩罚过大的参数？这既可以限制规则对训练数据的过度迎合，也可以表达训练之前对参数的偏好。后文会把这些做法分别命名为平方损失、正则化和先验，并解释它们在什么条件下能够联系起来。这里先记住问题，不必提前记住术语。

阅读每章时，可以先沿着开篇的几个问题建立方向：上一章已经解决了什么，仍留下什么困难，本章通过增加哪一种假设或结构处理这个困难。进入推导后，应先确认随机变量与具体观测、待学习参数与已知常量的区别，再检查矩阵维度和概率分布的条件。等式说明在给定假设下严格成立的关系，近似式通常意味着舍弃了某些项、引入了采样或限制了函数族；二者不能混为一谈。遇到较长推导，不必先记结论，可以用一个标量、两个样本或一维高斯分布重算一遍，再回到一般形式。

本书中的数学展开承担三项任务：从问题假设推出目标，从目标推出求解步骤，再从求解结果返回实际含义。因此，得到梯度或闭式解不是讨论的终点，还要追问解是否存在、是否唯一、算法是否能到达它，以及所需假设在工业数据中是否成立。章末的总结段落把这些结论重新放回相应的工业任务，并指出下一章需要改变的条件。读者可以据此区分已经得到证明的结论、依赖模型的解释和仍须通过实验检验的判断。

\section*{贯穿全书的参考书}
除了各章末的辅助学习材料，下面三本书可以作为阅读整本指南时的常备参考。遇到推导看不懂、概念分不清的地方，可按问题查阅对应章节，不必先从头读完这些书再开始本指南。
\begin{itemize}
 \item Marc Peter Deisenroth、A. Aldo Faisal、Cheng Soon Ong：\textit{Mathematics for Machine Learning}。这本书适合补齐线性代数、解析几何、矩阵分解、向量微积分、概率与连续优化基础，并通过线性回归、主成分分析等方法说明这些数学工具如何使用。阅读本指南的数学基础和模型推导时，可用它核对矩阵运算、求导步骤及概率计算。\href{https://mml-book.com}{作者维护的书籍网站}提供公开电子版。
 \item Sergios Theodoridis：\textit{Machine Learning: A Bayesian and Optimization Perspective}。这本书从贝叶斯推断和优化两个角度组织机器学习方法，适合进一步理解概率模型、参数估计、正则化、核方法与学习算法之间的联系。阅读本指南中的统计学习、贝叶斯建模和优化推导时，可对照查阅，重点看模型假设如何导出目标函数，以及目标函数如何决定求解方法。
 \item Ian Goodfellow、Yoshua Bengio、Aaron Courville：\textit{Deep Learning}，即常说的“花书”。这本书系统介绍深度前馈网络、正则化、优化、卷积网络、序列模型以及概率与生成建模。阅读本指南的深度学习部分时，可用它补充基础原理和推导。\href{https://www.deeplearningbook.org/}{官方书籍网站}提供免费在线阅读。它适合打基础，但不涵盖后来的 Transformer、Mamba 和当前大语言模型训练方法；这些内容需结合相应章节的补充材料学习。
\end{itemize}
三本书的用途各有侧重：数学运算不熟时查第一本，想理清概率假设与优化目标的关系时查第二本，想系统学习神经网络原理时查第三本。本指南中的工业例子用于连接这些原理与具体任务，不替代对数据条件和模型假设的检查。

\mainmatter
\part{机器学习的基本语言与完整流程}
本部分先补齐后文要用的数学语言，再依次回答八个问题：公式里的向量和概率怎样读；机器学习究竟学什么；简单模型怎样计算；一次学习实验怎样组织；参数怎样一步步更新；为什么训练成绩不能代表新数据上的成绩；有限数据能支持多强的结论；原始记录怎样变成预测时刻确实可用的输入。这里的“泛化”指模型在未见数据上的表现，“优化”指按目标调整参数，“部署”指把训练好的系统放到实际流程中使用。读者应看到，机器学习不是一条训练命令，而是一个完整的流程：先明确问题，再选择模型，最后验证结果。
\chapter{数学与概率预备知识}

\begin{quote}\small
\textit{本章摘要。}本章建立全书共同使用的数学语言。先从向量、矩阵和微积分说明模型如何表示对象、变化与优化，再引入概率、分布和期望，最后解释熵、距离、降维和数值稳定性。阅读顺序对应后文的依赖关系：先能读懂确定性模型，再能处理随机性，最后才能理解贝叶斯推断、深度学习和流形表示中的公式。
\end{quote}

本书从数学开始，并不是要求读者先成为数学家，而是要建立一种可核查的阅读习惯：每个运算必须作用在合适的对象上，每个结论必须有成立条件。设备的一组传感器读数先成为向量，多个样本组成矩阵；预测误差成为标量函数，其变化由梯度描述；同一工况下仍会出现不同结果，则需要条件概率。以下工具沿着这条路径出现，而不是彼此孤立的公式表。

可以始终带着同一组设备读数来理解这些工具。把温度和振动画成一个点，是几何表示；调整权重让预测更接近测量值，需要导数；重复运行后读数仍然有波动，就要用概率和期望。先看清每一步输入和输出的维度，再跟着推导计算。后面遇到具体模型时，可以回到这里核对它用了哪些条件。

\section{为什么需要数学语言}
机器学习不是把公式堆在代码旁边，而是用一套精确语言说明四件事：对象是什么，假设是什么，优化什么，以及结果在哪些条件下成立。本章只介绍后文反复使用的工具。读者不必一次记住所有细节，但应能在看到一个公式时识别其输入、输出、随机性与求和范围。初读时先掌握向量乘法、单变量求导和离散概率加权平均；高斯最大熵、PCA优化证明与期望求导可在对应模型用到时回读，不必先完成这些证明才能进入下一章。

\paragraph{先约定怎样读符号。}
标量是一个数，向量是一列数，矩阵是按行列排列的一张数表。$\R$表示实数集合，$\in$读作“属于”，$\R^d$表示含$d$个实数的向量空间；上标$\mathsf T$把列变成行，不是乘方。下标$i$常编号样本，下标$j$常编号特征；$\sum_{j=1}^d$表示把第$1$项到第$d$项相加。$\hat y$是预测值，未加帽的$y$是观测标签；$\argmin$返回让目标最小的参数，$\min$返回最小目标值。例如$(w-2)^2$的最小值是$0$，取得它的参数是$w=2$。

\section{向量、矩阵与几何}
先把两个传感器的读数按固定次序排成列，再用加权和生成输出。若$\bm x=(2,3)^{\mathsf T}$，$A=\begin{pmatrix}1&2\\0&1\end{pmatrix}$，则$A\bm x=(8,3)^{\mathsf T}$：第一项是$1\times2+2\times3$。列空间是矩阵所有列的线性组合能到达的集合；零空间是被矩阵乘成零的输入集合；秩是独立输出方向的个数。线性组合就是给各向量乘系数再相加。

向量$\bm{x}\in\mathbb{R}^d$可以表示一个样本，也可以表示参数、梯度或隐状态。内积$\bm{x}^{\mathsf T}\bm{y}$衡量方向相似性，范数$\|\bm{x}\|_2=(\sum_jx_j^2)^{1/2}$衡量长度。矩阵$A\in\mathbb{R}^{m\times d}$是从$d$维空间到$m$维空间的线性变换。其列空间、零空间和秩分别描述可表达方向、被压缩方向与有效维数。

对称半正定矩阵是满足$A=A^{\mathsf T}$且对所有$\bm{v}\in\R^d$都有$\bm{v}^{\mathsf T}A\bm{v}\geq0$的方阵；这里的“半”允许某些非零方向的二次型等于零。协方差矩阵必须具有这一性质，因为方差不可能为负。二次型$\bm{v}^{\mathsf T}A\bm{v}$还可以看作沿方向$\bm{v}$观察曲率，这使 Hessian 矩阵成为优化和不确定性分析的核心对象。

\paragraph{维度、投影与半正定性的由来。}
若$A\in\R^{m\times d}$、$\bm x\in\R^d$，则$(A\bm x)_i=\sum_{j=1}^d A_{ij}x_j$，输出有$m$个分量。内积不是不受尺度影响的相似度：将一个向量放大两倍会使内积也放大两倍；只有先归一化，内积才等于夹角余弦。对单位向量$\bm u$，把$\bm x$投影到其张成的直线上，需要最小化
\begin{align}
 \|\bm x-a\bm u\|_2^2
 &=\|\bm x\|_2^2-2a\bm u^{\mathsf T}\bm x+a^2,\\
 \frac{d}{da}\|\bm x-a\bm u\|_2^2&=-2\bm u^{\mathsf T}\bm x+2a=0.\label{eq:math-projection-stationarity}
\end{align}
由式\eqref{eq:math-projection-stationarity}得$a=\bm u^{\mathsf T}\bm x$，投影为$\bm u\bm u^{\mathsf T}\bm x$，残差与$\bm u$正交。若$\bm u$不是单位向量且$\bm u\ne\bm0$，系数必须再除以$\bm u^{\mathsf T}\bm u$；零向量不张成一条有方向的投影直线，不能用于这个公式。

设随机向量$\bm X$有有限二阶矩，$\bm\mu=\E\bm X$，定义$\Sigma=\E[(\bm X-\bm\mu)(\bm X-\bm\mu)^{\mathsf T}]\in\R^{d\times d}$。对任意$\bm v\in\R^d$，
\begin{equation}
 \bm v^{\mathsf T}\Sigma\bm v
 =\E\bigl[(\bm v^{\mathsf T}(\bm X-\bm\mu))^2\bigr]
 =\Var(\bm v^{\mathsf T}\bm X)\geq0.\label{eq:math-covariance-psd}
\end{equation}
式\eqref{eq:math-covariance-psd}同时证明了协方差半正定，并说明二次型就是投影后的方差。半正定并不保证可逆：若某个传感器恰好是另一个的两倍，就存在非零方向上的方差为零。

\section{微积分与链式法则}
导数描述输入作微小改变时输出怎样变化；偏导数只改变一个参数，其他参数暂时固定。梯度把所有偏导数排成列。Jacobian（雅可比矩阵）把多个输出对多个输入的偏导数排成表；Hessian（海森矩阵）则收集标量函数的二阶偏导数，用来描述斜率怎样变化。
例如$x=2$、预测$\hat y=wx$、标签$y=3$，在$w=1$时误差是$-1$，损失$J=(\hat y-y)^2/2=1/2$。先求$\partial J/\partial\hat y=-1$，再求$\partial\hat y/\partial w=2$，相乘得$dJ/dw=-2$。用步长$0.1$更新得到$w=1.2$，新损失为$(2.4-3)^2/2=0.18$。

若标量函数$J(\bm{\theta})$对参数可微，梯度为
\begin{equation}
 \nabla_{\bm{\theta}}J=\left[\frac{\partial J}{\partial\theta_1},\ldots,\frac{\partial J}{\partial\theta_d}\right]^{\mathsf T}.
\end{equation}
梯度非零时，它指向欧氏单位方向中函数增长最快的方向，因此$-\nabla J$是局部下降方向；梯度为零时，仅凭一阶导数不能判断是否达到极小值。对复合函数$J(f(\bm{\theta}))$，链式法则给出
\begin{equation}
 \frac{\partial J}{\partial\theta_j}=\sum_i\frac{\partial J}{\partial f_i}\frac{\partial f_i}{\partial\theta_j}.
\end{equation}
反向传播正是以高效的图结构方式反复应用链式法则，而不是一种神秘的独立学习规则。

\paragraph{为什么负梯度是局部下降方向。}
令$\bm g=\nabla J(\bm\theta)$。可微性给出$J(\bm\theta+\epsilon\bm v)=J(\bm\theta)+\epsilon\bm g^{\mathsf T}\bm v+o(\epsilon)$。在$\|\bm v\|_2=1$的方向中，Cauchy--Schwarz不等式给出$\bm g^{\mathsf T}\bm v\geq-\|\bm g\|_2$，当$\bm g\ne0$时由$\bm v=-\bm g/\|\bm g\|_2$取到等号。这只说明足够小的步长能下降，不保证任意步长都有效，也不保证梯度为零时已经达到最小值。

以$\bm r=A\bm\theta-\bm y\in\R^m$和$J=\tfrac12\bm r^{\mathsf T}\bm r$为例，微分逐步给出
\begin{equation}
 dJ=\bm r^{\mathsf T}d\bm r
 =\bm r^{\mathsf T}A\,d\bm\theta
 =(A^{\mathsf T}\bm r)^{\mathsf T}d\bm\theta,
 \qquad \nabla J=A^{\mathsf T}(A\bm\theta-\bm y)\in\R^d.
\end{equation}
一般地，若$\bm f:\R^d\to\R^m$的Jacobian为$D\bm f\in\R^{m\times d}$，则$\nabla_{\bm\theta}J=(D\bm f)^{\mathsf T}\nabla_{\bm f}J$。转置把输出空间的误差信号拉回参数空间，维度检查可以帮助发现漏转置和广播错误。

\section{概率、条件概率与期望}
大写$X$表示还未确定的随机量，小写$x$表示一次具体取值。离散概率给每个可能值分配权重，所有权重之和为一；连续密度则要在区间上积分才得到概率，密度值本身可以大于一。条件概率是先限定已知条件，再在剩余情形中计算比例。假设同一工况下温度$Y$以$1/4$概率为$2$，以$3/4$概率为$4$，则均值为$2/4+3\times4/4=3.5$，方差为$(2-3.5)^2/4+3(4-3.5)^2/4=0.75$。均值不一定是能实际观测到的一个值。

随机变量$X$的分布描述不确定性。联合分布满足乘法法则
\begin{equation}
 p(x,y)=p(y\mid x)p(x),
 \qquad p(y\mid x)=\frac{p(x,y)}{p(x)}.\label{eq:math-conditional-probability}
\end{equation}
边缘化把未知变量积分掉：$p(x)=\int p(x,y)\,dy$。期望是随机量的加权平均，方差是
\begin{equation}
 \Var(X)=\mathbb{E}[(X-\mathbb{E}[X])^2].
\end{equation}
条件期望$\mathbb{E}[Y\mid X]$是给定输入后对输出的最佳平方损失预测，这解释了为什么均方误差回归在理想条件下学习条件均值。

\paragraph{条件均值为何最优。}
对离散变量使用概率质量，对连续变量使用密度；式\eqref{eq:math-conditional-probability}的条件概率比值只在分母为正处使用，连续情形不是把单点概率相除。假设$\E[Y^2]<\infty$，固定输入$x$，记$\mu(x)=\E[Y\mid X=x]$。对任意实数预测$a$，把误差拆为$(Y-\mu)+(\mu-a)$，得到
\begin{align}
 \E[(Y-a)^2\mid X=x]
 &=\E[(Y-\mu)^2\mid X=x]
   +2(\mu-a)\E[Y-\mu\mid X=x]+(\mu-a)^2\\
 &=\Var(Y\mid X=x)+(\mu(x)-a)^2.\label{eq:math-conditional-mse}
\end{align}
式\eqref{eq:math-conditional-mse}的第一项与$a$无关，交叉项因条件残差均值为零而消失，所以最优$a$为$\mu(x)$。这不是说有限样本拟合一定得到条件均值，而是说它是无限制平方损失预测的理想目标。

将$Y-\E Y=(Y-\E[Y\mid X])+(\E[Y\mid X]-\E Y)$平方后取期望。交叉项先对$X$条件化，其中$\E[Y-\E[Y\mid X]\mid X]=0$，故该项消失；两个平方项分别给出全方差公式
\begin{equation}
 \Var(Y)=\E[\Var(Y\mid X)]+\Var(\E[Y\mid X]).\label{eq:math-total-variance}
\end{equation}
式\eqref{eq:math-total-variance}将总波动分为给定输入后仍存在的波动，以及不同输入导致的均值变化。缺少关键传感器时，前一项可能很大；仅靠更精细地优化参数无法消除这种信息缺失。

\section{高斯分布与最大熵}
高斯分布把靠近中心的读数看得更常见，把远离中心的读数看得更少见。先看一维：均值$\mu$给中心，标准差$\sigma>0$给横向尺度，$p(x)=\exp[-(x-\mu)^2/(2\sigma^2)]/(\sqrt{2\pi}\sigma)$。多维时用向量$\bm\mu\in\R^d$替代中心，用协方差矩阵$\Sigma\in\R^{d\times d}$替代方差；$|\Sigma|$是行列式，描述体积缩放，不是逐元素绝对值。正定指任意非零$\bm v$都有$\bm v^{\mathsf T}\Sigma\bm v>0$，比半正定更强。
多元高斯分布写作
\begin{equation}
 p(\bm{x})=\frac{1}{(2\pi)^{d/2}|\Sigma|^{1/2}}
 \exp\left[-\frac12(\bm{x}-\bm{\mu})^{\mathsf T}\Sigma^{-1}(\bm{x}-\bm{\mu})\right].\label{eq:math-gaussian-density}
\end{equation}
均值决定中心，协方差决定椭球的方向和尺度。给定均值和协方差，高斯分布具有最大的微分熵，因此它常作为“只知道一二阶统计量时不过度添加假设”的基线。

\paragraph{密度公式与最大熵结论的条件。}
式\eqref{eq:math-gaussian-density}的普通密度要求$\Sigma$正定。写$\Sigma=U\Lambda U^{\mathsf T}$，令$\bm z=\Lambda^{-1/2}U^{\mathsf T}(\bm x-\bm\mu)$，则指数中的二次型为$\|\bm z\|_2^2$，变量替换的体积因子为$|\Sigma|^{1/2}$。因此它实际上是独立标准正态密度经旋转、伸缩和平移得到的分布。若$\Sigma$奇异，概率集中在较低维子空间上，不能直接套用含$\Sigma^{-1}$的$d$维密度。

设$q$为具有给定均值和正定协方差的高斯密度，$p$为具有相同矩且相关积分有限的另一密度。因为$\log q$只是常数加二次型，相同的一二阶矩意味着$\E_p\log q=\E_q\log q$。利用下一节证明的Kullback--Leibler divergence (KL)非负性，
\begin{equation}
 0\leq\KL(p\|q)=\int p\log p-\int p\log q
 =-h(p)+h(q),\qquad
 h(q)=\frac12\log\bigl((2\pi e)^d|\Sigma|\bigr).
\end{equation}
故$h(p)\leq h(q)$。最大熵说的是固定矩约束下不再添加额外分布结构，并不证明现实数据必为高斯。微分熵还会随坐标单位改变，不应直接当成离散事件数目。

\section{熵、交叉熵与 KL 散度}
一个必然发生的结果没有不确定性，两个等可能结果则较难事先猜中。熵把这种不确定性量化；交叉熵还计入使用了不准确概率的额外代价；KL散度正是这部分差额。本书无特别说明时$\log$是自然对数，信息单位为奈特；改用以$2$为底的对数才以比特计。对$p=(1/2,1/2)$、$q=(3/4,1/4)$，$H(p)=\log2\approx0.693$，$H(p,q)=-\tfrac12\log(3/4)-\tfrac12\log(1/4)\approx0.837$，差额约$0.144$。

熵的信息度量来自 Shannon 的信息论框架，相对熵的统计区分意义可追溯到 Kullback 与 Leibler 的工作\cite{shannon1948information,kullback1951information}。以下公式区分随机性、编码代价与分布差异，三者不能仅因为都使用对数就被视为同一个量。
离散分布$p$的熵为$H(p)=-\sum_xp(x)\log p(x)$。用$q$编码来自$p$的数据时，平均代价是交叉熵$H(p,q)=-\sum_xp(x)\log q(x)$。两者关系为
\begin{equation}
 H(p,q)=H(p)+\KL(p\|q),
 \qquad \KL(p\|q)=\sum_xp(x)\log\frac{p(x)}{q(x)}\geq0.\label{eq:math-cross-entropy-kl}
\end{equation}
由式\eqref{eq:math-cross-entropy-kl}可知，最小化交叉熵等价于让预测分布接近真实分布。KL 散度不对称，选择$\KL(p\|q)$还是$\KL(q\|p)$会带来不同的近似行为，这一点在变分推断中非常重要。

\begin{example}
 若故障类别的真实分布为$(0.9,0.1)$，模型给出$(0.6,0.4)$，交叉熵会惩罚模型对高概率类别的低估。若模型输出概率而不是只输出类别，系统还可以根据置信度设置不同告警阈值。
\end{example}

\paragraph{KL非负性与交叉熵分解。}
约定$0\log0=0$，若$p(x)>0$但$q(x)=0$，则KL为无穷大。其余情形利用$-\log u\geq1-u$，逐项得到
\begin{equation}
 \KL(p\|q)=\sum_{p(x)>0}p(x)\left[-\log\frac{q(x)}{p(x)}\right]
 \geq\sum_{p(x)>0}(p(x)-q(x))\geq0.
\end{equation}
等号要求两分布相同。再把$\log(p/q)=\log p-\log q$代回求和，立即得到$\KL=H(p,q)-H(p)$。训练时真实分布$p$固定，因而最小化交叉熵与最小化这个方向的KL有相同最优解；用有限样本平均替代期望后，还存在采样误差和模型表达能力限制。

\section{从几何距离到概率距离}
机器学习中常见的“距离”并不只有欧氏距离。欧氏距离比较两个向量在坐标空间中的位置，余弦距离比较方向，Mahalanobis 距离则利用协方差修正不同方向的尺度：
\begin{equation}
 d_M(x,y)=\sqrt{(x-y)^{\mathsf T}\Sigma^{-1}(x-y)}.
\end{equation}
若$\Sigma$刻画的是噪声协方差，这种缩放可降低高噪声方向的影响；若使用全部观测的协方差，大方差还可能来自有用信号，不能直接等同于噪声。下面先比较向量之间的距离，再转向分布之间的差异：

\begin{align}
 d_1(x,y)&=\|x-y\|_1=\sum_{j=1}^{d}|x_j-y_j|,\label{eq:math-manhattan-distance}\\
 d_2(x,y)&=\|x-y\|_2=\sqrt{\sum_{j=1}^{d}(x_j-y_j)^2},\label{eq:math-euclidean-distance}\\
 d_p(x,y)&=\|x-y\|_p=\left(\sum_{j=1}^{d}|x_j-y_j|^p\right)^{1/p},\quad p\geq1,\label{eq:math-minkowski-distance}\\
 d_\infty(x,y)&=\|x-y\|_\infty=\max_{1\leq j\leq d}|x_j-y_j|.\label{eq:math-chebyshev-distance}
\end{align}

式\eqref{eq:math-manhattan-distance}把各坐标的绝对差直接相加，适合把每个传感器的偏差看成可以分别累积的误差。式\eqref{eq:math-euclidean-distance}对较大的单项差异惩罚更重，常用于连续特征经过同尺度处理后的比较。式\eqref{eq:math-minkowski-distance}把前两者统一起来：$p=1$得到曼哈顿距离，$p=2$得到欧氏距离；当$p$增大时，最大坐标差异的影响逐渐增强。式\eqref{eq:math-chebyshev-distance}只关注最不一致的那个坐标，例如设备的多个安全指标中只要有一项超出容许范围，就把该项作为整体差异的代表。

如果只关心方向而不关心向量长度，可以使用余弦相似度及其对应的余弦距离：
\begin{equation}
 s_{\cos}(x,y)=\frac{x^{\mathsf T}y}{\|x\|_2\|y\|_2},\qquad
 d_{\cos}(x,y)=1-s_{\cos}(x,y),
 \quad x\ne0,\ y\ne0.\label{eq:math-cosine-distance}
\end{equation}
例如，两台设备的振动频谱一个整体强、一个整体弱，但频率峰的相对分布相似时，余弦距离可能很小，而欧氏距离仍然可能很大。相反，若要比较类别、开关状态或字符串中有多少位置不同，可以使用汉明距离：
\begin{equation}
 d_{\mathrm H}(x,y)=\sum_{j=1}^{d}\mathbb{1}\{x_j\ne y_j\},\label{eq:math-hamming-distance}
\end{equation}
它只统计不相同的位置，不表示数值差异的大小，因此不适合直接代替连续传感器变量的欧氏距离。

概率模型中的 KL 散度比较分布而非单个点：
\begin{equation}
 \KL(p\|q)=\sum_xp(x)\log\frac{p(x)}{q(x)},\label{eq:math-kl-distance}
\end{equation}
它会把概率质量放错位置的代价纳入比较，因此适合描述预测分布与真实分布之间的差异。不过，KL 散度一般不对称，也可能为无穷大。若希望比较两个分布搬运概率质量所需的“距离”，离散情形还可以使用 Wasserstein 距离：
\begin{equation}
 W_1(p,q)=\min_{\gamma\in\Gamma(p,q)}
 \sum_{a,b}\gamma(a,b)c(a,b),\label{eq:math-wasserstein-distance}
\end{equation}
其中$\Gamma(p,q)$是边缘分布分别为$p$和$q$的联合分布集合，$c(a,b)$是把质量从$a$搬到$b$的代价。举例来说，把预测的风速分布从“每秒 8 米”移到“每秒 9 米”，Wasserstein 距离会利用这两个数值相近这一事实；KL 散度则只看概率质量在各个取值上的比例。实际使用哪一种距离，取决于特征的尺度、是否关心方向、是否关心最坏坐标，以及分布取值之间是否有自然的距离。

\paragraph{手算距离并检查“距离”是否严格。}
取$x=(0,0)$、$y=(3,4)$，曼哈顿、欧氏和最大坐标距离分别是$7,5,4$。若$\Sigma=\diag(9,16)$，Mahalanobis距离为$\sqrt{3^2/9+4^2/16}=\sqrt2$。严格的度量要求非负、相同点且仅相同点距离为零、对称以及三角不等式；余弦距离一般不是这种度量，两个不同长度的同向向量也可距离为零。Wasserstein公式中的$c$若是底层空间的度量，才具有相应距离解释；任意成本表不自动保证这些性质。

\section{特征值、主方向与降维}
特征向量是乘矩阵后仍落在原向量所张成直线上的非零向量，特征值是相应倍数，即$A\bm u=\lambda\bm u$；负特征值会反向，零特征值会把该向量映到零。协方差的特征值非负。PCA（主成分分析）保留变化较大的少数坐标，以较少数值近似原始样本。若协方差是$\diag(4,1)$，保留第一坐标就保留$4/(4+1)=80\%$的总方差；丢弃第二坐标的平均平方重构误差为$1$。这里是方差比例，不是分类准确率。

对称矩阵可以正交对角化：$A=U\Lambda U^{\mathsf T}$，其中$U$的列是两两正交、长度为一的特征向量，$\Lambda$是特征值组成的对角矩阵。若$A$是协方差矩阵，按特征值从大到小排列后，第一个方向使投影方差最大，后续方向在与已选方向正交的约束下依次最大；并不是每个特征向量都取得同一个最大方差。Principal Component Analysis (PCA) 选择前$r$个方向的投影
\begin{equation}
 \bm z^{\mathrm{PCA}}=U_r^{\mathsf T}(\bm x-\bar{\bm x}),
 \qquad \hat{\bm x}=U_r\bm z^{\mathrm{PCA}}+\bar{\bm x}.
\end{equation}
这里的$\bm z^{\mathrm{PCA}}\in\R^r$是单个样本的降维坐标，不要与后文证明中表示中心化样本的$\bm z_i\in\R^d$混淆。训练后须保存均值$\bar{\bm x}$与投影基$U_r$；新样本沿用二者，不能用新样本自身替换均值，否则单样本中心化会把输入变成零。
它同时可以理解为最大化投影方差，或最小化线性子空间的重构误差。这个双重解释为后文的表示学习和流形学习提供了桥梁。

\paragraph{从白化到PCA的优化证明。}
当$\Sigma$正定，令$W=\Lambda^{-1/2}U^{\mathsf T}$，则$d_M(x,y)=\|W(x-y)\|_2$。变换后协方差为$W\Sigma W^{\mathsf T}=I$，这种使协方差成为单位阵的操作称为白化，不同于只缩放各原始特征的标准化。这是先把各主方向除以其标准差再测欧氏距离；但协方差同时含有信号变化和噪声变化，不能无条件把大方差方向都解释为低质量传感器。样本协方差奇异时可考虑$\Sigma+\epsilon I$，代价是改变了原距离。

对中心化样本$\bm z_i=\bm x_i-\bar{\bm x}\in\R^d$，定义$S=n^{-1}\sum_i\bm z_i\bm z_i^{\mathsf T}$。单个单位方向$\bm u$的平均投影平方为$\bm u^{\mathsf T}S\bm u$。拉格朗日函数$\bm u^{\mathsf T}S\bm u-\lambda(\bm u^{\mathsf T}\bm u-1)$的驻点满足$S\bm u=\lambda\bm u$；按特征向量展开$\bm u$可知该二次型是特征值的加权平均，最大值因此为最大特征值。

对$U_r\in\R^{d\times r}$且$U_r^{\mathsf T}U_r=I_r$，投影与残差正交，故
\begin{align}
 \frac1n\sum_i\|\bm z_i-U_rU_r^{\mathsf T}\bm z_i\|_2^2
 &=\frac1n\sum_i\bigl(\|\bm z_i\|_2^2-\|U_r^{\mathsf T}\bm z_i\|_2^2\bigr)\\
 &=\operatorname{tr}(S)-\operatorname{tr}(U_r^{\mathsf T}SU_r).\label{eq:math-pca-reconstruction}
\end{align}
式\eqref{eq:math-pca-reconstruction}中，$\operatorname{tr}$表示对角元素之和。在$S$的特征基中，后一迹是特征值按各方向保留程度的加权和，权重在$[0,1]$内且总和为$r$，所以保留最大的$r$个特征值达到最大值，最小重构误差为$\sum_{j>r}\lambda_j$。重根时最优基不唯一；大方差也不必与故障标签有关，PCA最优的是重构，不是分类。

\section{微分、积分与期望的对应}
连续概率密度的积分必须等于一；离散概率则通过求和归一化。对参数化密度$p_\theta(x)$，当分布本身也依赖参数时，采样分布的变化也必须被考虑：
\begin{equation}
 \nabla_\theta\mathbb{E}_{x\sim p_\theta}[f_\theta(x)]
 \neq \mathbb{E}_{x\sim p_\theta}[\nabla_\theta f_\theta(x)].\label{eq:math-expectation-gradient-warning}
\end{equation}
这里有两种变化要一起计算：参数改变后，同一个样本上的函数值会变，抽到各个样本的概率也可能变。普通监督学习通常把训练样本固定，只需处理前一种变化；变分推断中的采样分布也要学习，因此常用重参数化技巧或得分函数估计来处理后一种变化。

\paragraph{把缺少的梯度项补出来。}
式\eqref{eq:math-expectation-gradient-warning}的$\neq$应理解为一般不能直接交换，并非所有函数都严格不等。若积分支撑集不随$\theta$变化，密度与函数可微，且存在允许微分与积分交换的可积控制函数，则乘积法则给出
\begin{align}
 \nabla_\theta\int f_\theta(x)p_\theta(x)\,dx
 &=\int\bigl[p_\theta(x)\nabla_\theta f_\theta(x)
       +f_\theta(x)\nabla_\theta p_\theta(x)\bigr]dx\\
 &=\E_{p_\theta}[\nabla_\theta f_\theta(X)]
   +\E_{p_\theta}[f_\theta(X)\nabla_\theta\log p_\theta(X)].\label{eq:math-score-gradient}
\end{align}
式\eqref{eq:math-score-gradient}的第二项就是采样分布变化的贡献。若可写$X=\mu+\sigma\varepsilon$、$\varepsilon\sim\mathcal N(0,1)$且$\sigma>0$，也可对固定噪声分布下的$f_\theta(\mu+\sigma\varepsilon)$使用链式法则。若支撑集依赖参数，例如$X\sim\operatorname{Uniform}(0,\theta)$，则还要处理积分边界，不能机械使用式\eqref{eq:math-score-gradient}。

\paragraph{一个不能漏掉分布变化的例子。}
令$X\in\{0,1\}$、$P_\theta(X=1)=\theta$，$0<\theta<1$，且$f(X)=X$不显含参数。由于$\E_\theta X=\theta$，其导数为$1$；但固定$X$求导得$\partial f/\partial\theta=0$。缺少的$1$完全来自概率权重的变化。离散情形把积分换成求和即可看清这一点。

\section{数值稳定性}
计算机只能保存有限范围和精度的数。softmax把一组任意实数得分变成总和为一的概率：第$i$项为$\exp(a_i)/\sum_j\exp(a_j)$。所有得分减去同一个常数不会改变概率；例如得分$(1000,1000)$应给出$(1/2,1/2)$，减去$1000$后计算$(1,1)/2$，就不必先算可能溢出的$e^{1000}$。NaN表示运算没有得到可表示的有效数，而不是一种很大的正常损失。

概率模型经常使用对数概率。直接计算$\exp(a)$可能溢出，通常使用
\begin{equation}
 \log\sum_i\exp(a_i)=m+\log\sum_i\exp(a_i-m),\qquad m=\max_i a_i.
\end{equation}
这就是 log-sum-exp 技巧。softmax 的实现、交叉熵损失和扩散模型中的噪声权重都依赖类似的稳定计算。数值稳定性不是实现细节之外的问题：溢出的 loss 会让优化器得到 NaN，随后整个模型参数都可能失效。

\section{自测与参考答案}
\begin{enumerate}
\item 把$(3,4)^{\mathsf T}$投影到单位方向$(1,0)^{\mathsf T}$，投影和残差是什么？\textbf{答：}系数为$3$，投影为$(3,0)^{\mathsf T}$，残差为$(0,4)^{\mathsf T}$，内积为零。
\item 对本章取值$2,4$的随机温度，预测$3$的平均平方误差是多少？\textbf{答：}$1/4+3/4=1$；也可用$0.75+(3.5-3)^2=1$核对。最优预测$3.5$仍留下$0.75$误差。
\item 协方差特征值为$5,2,1$，PCA保留两个方向的平均平方重构误差和方差保留率是什么？\textbf{答：}丢弃值之和为$1$，保留率为$7/8$。不需要把误差再除以维数，除非另行定义按坐标平均的指标。
\end{enumerate}

\section{本章小结：看到公式时问什么}
面对一个新公式，先问变量是确定量还是随机量，再问每一项的维度和求和范围；然后检查它是在拟合条件均值、概率分布、几何结构还是约束；最后确认该量在训练和推理阶段是否都可计算。这个检查顺序把符号阅读转化为模型理解。

本章的几条推导可以连起来阅读：投影把逼近误差分解为可表达部分和残差；条件均值把预测误差分解为可约误差和条件波动；KL把交叉熵分解为真实分布自身的不确定性和模型失配。它们共同提醒读者，优化只能改变目标中的某些项，不能凭空补充缺失的信息，也不能消除模型假设的边界。

下一章把这些运算用到具体任务中：一段振动记录算一个样本，还是一次完整运行算一个样本？要预测平均温度、温度的某个分位数，还是故障概率？这些选择会决定标签和损失，也决定本章的公式该用在哪里。


\chapter{机器学习究竟在学习什么}

\begin{quote}\small
\textit{本章摘要。}本章定义样本、任务、模型、损失和风险，说明不同损失对应不同预测对象，并讨论正则化、可识别性及预测到决策的边界。
\end{quote}

上一章给出了向量、导数和概率的语言，但能计算一个梯度并不意味着已经定义了值得学习的问题。现在需要把设备观测、未来事件和维护动作分开：观测决定模型能看见什么，标签决定它试图回答什么，损失决定怎样的错误更值得避免。三者不一致时，即使训练完全收敛，也可能得到一个回答错误问题的模型。

\paragraph{本章导语。}
同一段振动记录，可以用来判断轴承是否故障，也可以用来预测下一时刻的振幅。这两个问题需要不同的标签，算错后的后果也不同。本章先把这些差别写清楚，再说明模型中的参数如何通过损失函数学出来。随后讨论正则化、可识别性和部署反馈：为什么训练数据上看起来很好，换一台设备却可能失效？后文介绍各种模型时，都要回到输入、目标和错误代价这几件事。

\section{学习对象：数据、任务与模型}

% auxiliary-resource-guide
本章末的“辅助学习材料”列出与核心方法对应的讲解和实践入口，可在阅读相关推导后，按所列小节对照学习。

\paragraph{从数据到任务。}
数据（data）是对对象、过程或环境的观测记录。一个样本（sample）是一次相对完整的观测，通常记为$\bm{x}\in\R^d$；监督学习中还伴随标签$y$，于是样本写为$(\bm{x},y)$。这里的“样本”不是数据表的一行这么简单：一次设备运行周期、一张图像、一个句子或一段传感器窗口都可以成为样本，关键在于它是否对应一个明确的可学习单位。

“任务”由输入空间$\mathcal{X}$、输出空间$\mathcal{Y}$以及评价准则共同决定。输入空间是模型允许接收的对象集合（例如时序数据、图像、文本等），输出空间是允许回答的结果集合，这两个集合先于可学习参数的确定。分类的输出是离散类别，回归的输出是连续数值，预测性维护可能输出剩余寿命，异常检测则可能没有逐样本标签。相同的数据可以支持不同任务。例如，振动信号既可以用于故障分类，也可以用于生成未来波形或估计设备状态。

\paragraph{把一个任务写成可核查的样本。}
假设每小时读一次温度和振动，把这两项读数排成输入向量：
\[
 \bm x=(\text{温度},\text{振动})^{\mathsf T}\in\R^2.
\]
标签为下一小时温度$y\in\R$时，这就是回归任务；改成“未来一天是否故障”的$y\in\{0,1\}$时，就成为二分类。

监督学习用已知输入和标签学习对应关系，无监督学习则在没有这些目标标签时寻找输入结构，即发现样本之间的相似性、变量之间的关联等规律。例如，不提供故障标签，仅根据温度和振动把相似的设备记录分组，但这些分组不自动代表正常或故障。

同一数据记录能面向不同任务，不同任务同样需要不同的标签或指标。训练是利用已有样本调整参数；推理是在参数训练固定后处理一个新输入，得到对新样本标签或指标的预测。


\paragraph{模型、参数与预测。}
模型根据对数据和标签的观测方式不同，可分为两类：\textbf{确定性模型}和\textbf{概率模型}。确定性模型（deterministic model）是带有可调参数的函数族
\begin{equation}
 f_{\bm{\theta}}:\mathcal{X}\rightarrow\mathcal{Y},
 \qquad \bm{\theta}\in\Theta.
\end{equation}
给定输入$\bm{x}$，确定性模型直接输出$\hat{y}=f_{\bm{\theta}}(\bm{x})$。参数$\bm{\theta}$不是人为凭空指定的数值，而是由训练数据和训练目标共同决定的可学习量。模型结构规定了“允许怎样形式的映射”（线性形式 or 非线性形式），参数则规定了“采用哪一个特定的映射”。通常，这类模型也被称为“\textbf{确定性判别模型}”，因为模型函数拟合的是观测数据到标签的映射。

至于函数族的形式，应是由数据观测$x$与标签$y$（任务）的形式决定的。例如，“图像+分类任务”适用的函数族可以是ResNet、ViT或MLP；“时序+回归任务”适用的函数族可以是RNN、LSTM或TCN。函数族的选择通常由数据类型、任务类型和计算资源共同决定。

对应地，有确定性判别模型，就有“\textbf{概率判别模型}”。概率判别模型输出条件分布$p_{\bm{\theta}}(y\mid\bm{x})$，而不是直接输出$\mathcal Y$中的一个标签。例如，$K$类分类任务的概率输出是$K$个非负且总和为一的数，回归任务的概率输出可以是预测均值和噪声大小（在对标签的噪声建模为高斯的情形下）。确定性输出可以看作概率模型的特殊情形，例如$p(y\mid\bm{x})$集中在单点$f_{\bm{\theta}}(\bm{x})$上。因此，确定性模型与概率模型不是互斥的两个世界，而是表达不确定性程度不同的两种建模方式。

需要注意的是，这里的例子只把随机性放在模型的最终输出端，并不意味着概率模型中的概率性仅存在于输出层。模型内部的隐藏状态可以是随机变量，例如用潜变量$z$表示未直接观测的设备状态，并通过$p_{\bm{\theta}}(y\mid\bm{x},z)$和$p_{\bm{\theta}}(z\mid\bm{x})$共同描述预测过程；模型参数本身也可以被视为随机变量，例如贝叶斯模型用参数分布$p(\bm{\theta}\mid\mathcal D)$表达训练数据给参数带来的不确定性。因而，概率建模可以作用于输出、隐藏变量或参数等不同层次，前面的输出分布只是其中最简单的两种情形。

同时，有一类区别于判别模型的模型，它们叫做“\textbf{生成模型}”。与建模$\bm{x}$到$y$的映射的判别模型不同，生成模型更关注观测$\bm{x}$的生成方式，即建模联合分布$p_{\bm{\theta}}(\bm{x},y)$。生成模型可以用来生成新的样本，也可以用来计算条件分布$p_{\bm{\theta}}(y\mid\bm{x})$，但它们的训练目标和假设通常与判别模型不同。

\section{让模型有点追求}

\paragraph{给定一个模型的训练目标——损失函数。}
训练的最终目的，是在模型允许的参数、隐藏状态或其他可训练变量中，找出使训练目标最优的取值或状态。“最优”并不是脱离任务而自然存在的性质，而是由目标函数规定的；损失函数正是把模型产生的输出与真实任务目标连接起来的媒介。一个好的损失函数应当能够充分反映任务真正关心的目标，并且对任务中的关键差异敏感。例如，分类任务关心的是类别概率及其相对关系，交叉熵通常比直接把“分对或分错”作为训练信号更适合优化。

损失函数的具体形式还取决于标签情况和模型的建模方式。对确定性判别模型，通常先用函数衡量真实标签与模型输出之间的差异，例如回归模型$f_{\bm{\theta}}(\bm{x})$可以采用
\begin{equation}
 \ell_{\mathrm{sq}}\bigl(f_{\bm{\theta}}(\bm{x}),y\bigr)=\frac12\bigl(f_{\bm{\theta}}(\bm{x})-y\bigr)^2,
\end{equation}
再对样本损失求和或取平均。对概率判别模型，若把参数视为随机变量并给定先验$p(\bm{\theta})$，则可以最大化后验
\begin{equation}
 \hat{\bm{\theta}}_{\mathrm{MAP}}
 =\argmax_{\bm{\theta}}p(\bm{\theta}\mid\mathcal D)
 =\argmax_{\bm{\theta}}\bigl[\log p(\mathcal D\mid\bm{\theta})+\log p(\bm{\theta})\bigr].
\end{equation}
等价地，也可以最小化负对数后验；不过在更一般的语境下，这类量通常称为目标函数（objective），不一定称为损失函数。对概率生成模型，则常通过最大化观测数据的对数似然
\begin{equation}
 \hat{\bm{\theta}}_{\mathrm{ML}}
 =\argmax_{\bm{\theta}}\log p_{\bm{\theta}}(\mathcal D)
 =\argmax_{\bm{\theta}}\sum_{i=1}^{n}\log p_{\bm{\theta}}(\bm{x}_i,y_i)
\end{equation}
来确定参数，或等价地最小化负对数似然。因而，损失或目标的构建必须同时考虑任务类型、标签的可观测性以及模型的概率或确定性建模方式。

\paragraph{经验风险与总体风险.}  
有限样本上的拟合与总体上的预测能力需要分开讨论，经验分布到真实分布的一致收敛是统计学习理论的基础问题之一\cite{vapnik1971uniform}。
设$(X,Y)\sim P$描述部署环境中的输入与标签，$X\in\R^d$，$Y$可为实数或类别。总体风险定义为
\begin{equation}
 R(f)=\E_{(X,Y)\sim P}[\ell(f(X),Y)].
\end{equation}
这里的期望包含尚未观察到的未来样本，因此通常不能直接计算。若训练样本独立同分布于$P$（independent and identically distributed，彼此独立且共享同一个分布），对一个事先固定且损失可积的$f$，有$\E[\hat R(f)]=R(f)$。训练前允许选择的模型参数属于参数空间$\Theta$；“事先固定”强调不能先看测试结果再挑选$f$而仍使用这条无偏性解释。若损失方差有限，独立性进一步给出
\begin{equation}
 \Var(\hat R(f))=\frac{1}{n^2}\sum_{i=1}^n\Var(\ell(f(X_i),Y_i))
 =\frac{\Var(\ell(f(X),Y))}{n}.
\end{equation}
这解释了平均损失为什么可用于估计风险，也暴露了限制：相邻设备窗口相关时有额外协方差项；$f$由同一训练集选择后，上述固定模型的无偏论证不能直接用于其训练误差。验证与测试的必要性由此产生，而不是来自惯例。

训练时通常最小化经验风险，真正希望最小化的却是部署分布上的总体风险。“\textbf{过拟合}”的核心表现，是模型在训练集上的经验风险很小，但在未参与训练的新样本上的总体风险仍然较大，因而出现较大的泛化差距。模型容量过大时，它不仅可以学习输入与标签之间可迁移的规律，也可能记住训练样本中的噪声、偶然相关性甚至标签错误；但模型容量过大并不是过拟合的必要条件，训练样本过少、标签噪声较强或反复利用验证集也可能造成过拟合。当训练集有限、包含噪声，或训练集与部署分布存在差异时，这些被模型记住的细节可能无法在新样本中重复出现。典型地，训练早期经验风险和总体风险可能同时下降；训练后期经验风险继续下降，而总体风险可能转而上升，因为模型开始更多地拟合训练数据中的偶然性。


\section{归纳偏置}

前文提到了过拟合，通俗来讲指的是训练误差很低而新样本误差较高。与之对应的一个概念叫做“\textbf{泛化性}”，指模型在“未参与训练和选择的数据”上的有效程度。使用损失函数或目标函数进行训练，直接下降的是经验风险；但我们真正关心的是部署分布上的总体风险。经验风险最小化并不自动意味着总体风险最小化，因此需要在模型构造和训练目标中加入能够约束学习过程的偏好。下面从模型角度和训练目标角度分别说明这种约束如何发挥作用。

正如前文所说，模型本质上是一个函数族；无论这个函数族是确定性的还是含有随机变量，它都不可能平等地偏好所有可能的输入输出关系。模型结构、参数化方式以及训练算法对哪些解更容易被找到、哪些解更容易被保留所作的假设，统称为“\textbf{归纳偏置}”（inductive bias）。归纳偏置的作用，是在训练样本不足以唯一决定规律时，帮助模型从许多能够解释训练数据的候选规律中选择更可能适用于新样本的规律。

例如，对同一图像分类任务，卷积神经网络把局部连接和参数共享写进了模型结构，因此偏好局部模式以及对平移具有一定稳定性的函数；全连接网络则没有这种结构性限制。若图像任务确实具有这些性质，卷积网络就能把有限样本用于学习更有用的变化，而不必分别记住每个像素位置的完全独立关系，因此通常更容易获得较好的泛化性能。这里的“更适合”是统计意义上的，不是保证在任何数据集上都更好。

归纳偏置与可解释性有关，但二者不是同一个概念。可解释性讨论人能否理解模型的预测依据；归纳偏置讨论模型在学习时预先接受了哪些结构假设。一个模型可以具有明确的归纳偏置，却仍然难以解释；也可以相对容易解释，但其归纳偏置并不适合当前任务。判断归纳偏置是否良好，关键在于它是否与数据生成过程和任务目标相匹配，并能否在有限样本下缩小训练风险与总体风险之间的差距。

正则化（regularization）提供了加入这种偏好的另一条途径：不一定更换模型结构，而是在训练目标中明确惩罚某些不希望出现的解。其基本形式是
\begin{equation}
 \min_{\bm{\theta}}\;\hat{R}(\bm{\theta})+\lambda\Omega(\bm{\theta}),
 \qquad \Omega(\bm{\theta})=\frac12\|\bm{\theta}\|_2^2.\label{eq:foundations-regularized-risk}
\end{equation}
这里第一项$\hat R(\bm{\theta})$要求模型解释训练数据，第二项$\Omega(\bm{\theta})$表达我们对模型复杂度的偏好。选择$\Omega(\bm{\theta})=\frac12\|\bm{\theta}\|_2^2$，表示在其他条件相近时偏好范数较小的参数；它并不声称真实参数一定接近零，而是把“大权重可能对应不稳定拟合”作为一种可调的建模偏好。系数$\lambda\geq0$决定这种偏好有多强：$\lambda=0$时只拟合训练数据；$\lambda$增大时，模型愿意牺牲一部分训练集拟合精度，以换取更小的参数和可能更好的泛化。正则化的目的不是让训练误差尽可能低，而是在经验风险和对复杂度的约束之间寻找适合部署的平衡。

\begin{example}
 对于温度传感器预测下一时刻温度，模型并不需要学习所有历史序列到未来数值的任意映射。若已知相邻时刻通常连续变化，可以采用平滑性偏置；若设备存在周期，则可以把周期特征加入输入。模型的“聪明”来自数据，也来自这些合理的限制。
\end{example}

\paragraph{惩罚形式与约束形式如何联系。}
为什么可以把正则化理解成“控制复杂度”？一种直接的写法是先规定允许的复杂度上限：
\begin{equation}
 \min_{\theta}\hat R(\theta)\quad\text{subject to}\quad \Omega(\theta)\leq c.
\end{equation}
这句话的意思是：在所有不超过复杂度预算$c$的模型中，选择训练风险最小的模型。它明确表达了两件事：训练数据需要拟合，但模型不能为了拟合少量样本而无限增加复杂度。

约束问题也可以通过拉格朗日乘子写成
\begin{equation}
 L(\theta,\lambda)=\hat R(\theta)+\lambda\bigl(\Omega(\theta)-c\bigr),\qquad \lambda\geq0.
\end{equation}
固定$\lambda$后，$-\lambda c$与$\theta$无关，不会改变参数最优解，于是得到惩罚形式$\hat R(\theta)+\lambda\Omega(\theta)$。因此，$c$和$\lambda$分别是两种不同的控制方式：前者规定复杂度上限，后者规定超出较简单解的代价。在线性凸问题并满足适当严格可行性条件时，二者可以通过拉格朗日最优性条件联系起来；对非凸网络，不能据此声称任意约束问题都与某个惩罚问题全局等价，也不能认为$\lambda$与$c$必然一一对应。

在约束问题的最优点$\hat\theta$处，若目标与约束满足相应的凸性条件，可以写出
\begin{equation}
 0\in\partial\hat R(\hat\theta)+\lambda\partial\Omega(\hat\theta),\qquad
 \Omega(\hat\theta)\leq c,\qquad
 \lambda(\Omega(\hat\theta)-c)=0.
\end{equation}
$\partial$是允许不可微凸函数的次梯度集合。最后一个等式是互补条件。若$\Omega(\hat\theta)<c$，说明复杂度预算没有用完，限制没有起作用，此时可以有$\lambda=0$；若$\Omega(\hat\theta)=c$，说明模型已经碰到复杂度边界，进一步降低训练风险通常需要增加复杂度，约束因此会通过正的$\lambda$影响最优解。这里的“补偿”不是说两个数必须相等，而是说训练风险的改善必须值得为增加复杂度所付出的代价。

最后用一维例子看清“收缩”究竟发生了什么。设无正则化时的训练风险为
\begin{equation}
 \hat R(w)=\frac12(w-a)^2,
\end{equation}
其中$a$是只根据训练数据得到的最优估计。它的最小值在$w=a$处，因为此时$w$与$a$完全一致。现在加入平方惩罚，训练目标变为
\begin{equation}
 J(w)=\frac12(w-a)^2+\frac{\lambda}{2}w^2.
\end{equation}
第一项希望$w$接近数据给出的估计$a$，第二项希望$w$接近零。对$w$求导得到
\begin{equation}
 J'(w)=(w-a)+\lambda w.
\end{equation}
令导数为零，得到$(1+\lambda)w=a$，因此
\begin{equation}
 \hat w=\frac{a}{1+\lambda}.
\end{equation}
由于$\lambda\geq0$，有$|\hat w|\leq|a|$；当$\lambda=0$时不收缩，随着$\lambda$增大，$\hat w$逐渐靠近零。这就是“收缩”的含义：正则化不会自动判断某个变量是噪声，也不会直接把所有不重要的系数删掉，而是把参数估计向预先选定的中心（这里是零）拉近。

收缩为什么可能改善总体风险？训练数据中的偶然波动会使无正则化估计$a$在不同训练集之间变化很大，这种变化称为估计的不稳定性或方差。平方惩罚牺牲了一部分对当前训练集的贴合，使估计更稳定；但如果真实系数确实远离零，收缩会带来偏差。因此，正则化是在偏差与方差之间作权衡：只有方差的降低足以抵消偏差增加，使用验证数据选择的$\lambda$才可能降低总体预测风险。它不是“惩罚越强越好”，也不是“参数越小越接近真实机制”。

模型选择至少同时受四类约束：数据规模与标注质量、任务中的依赖结构、训练与推理资源、错误的业务代价。复杂模型并不自动优于简单模型；它只有在表示了简单模型无法表示的结构时才有必要。后文讨论的序列、视觉和语言模型，实际上是在不同结构约束下对这四个问题的回答。


\section{任务的可学习性}
在选择模型之前，还应先问一个更根本的问题：给定的输入中，是否含有足以预测目标的稳定信息？例如，设输入是风电机组叶片的图像，标签是同一时刻的发电功率。凭借领域知识可以预期，功率还强烈依赖风速、风向、空气密度、机组运行状态和控制策略；一张只反映叶片外观的图像，通常不足以充分决定功率。图像中也许包含结冰、污损或损伤等与功率有关的间接线索，但这种联系必须由数据和机理共同证实，不能因为两个变量出现在同一张数据表中，就假定二者必然存在可用于预测的关系。

训练集上的拟合能力不能单独回答“任务是否可以被学习”这个问题。只要模型足够灵活，它就可以把输入空间划分为许多很小的区域，并为每个训练样本建立对应映射。以决策树回归为例，当训练输入彼此可区分，且树的深度、叶节点样本数等不受限制时，树可以把训练样本逐个分到不同叶节点，从而达到零训练误差；分类树也可以达到百分之百训练准确率。这说明模型有能力记住训练样本，却不说明图像中存在对未来功率稳定有效的预测线索。若新图像落入训练时未被充分覆盖的区域，模型只能依赖未经验证的划分和映射，训练集上的完美表现便可能无法迁移到新样本。这正是过拟合的一种典型表现。

其实，这个问题很简单，只需要判断一件事：“相似的”输入是否具有“相似的”标签。从统计角度看，一个监督任务在给定输入$X$时是否可学习，取决于真实条件分布$P(Y\mid X)$是否具有可利用的结构，以及这一结构能否在训练环境和部署环境之间保持足够稳定。这里的$P(Y\mid X)$描述“在给定输入时，标签可能如何变化”。若能够在输入空间的各个相关区域持续获得带标签样本，并且条件分布在邻近输入之间不会发生任意剧烈的变化，那么观察到的局部模式才有可能外推到附近的未见样本。连续性、平滑性或低维结构都可以成为这种外推的依据。那么，在足够的样本覆盖、合适的模型族的条件下，任务的训练才具备泛化性。

现实中的困难在于，输入空间往往高维，且样本真实分布不均匀；有限样本更不可能把整个空间都覆盖。因此，数据中是否存在规律通常不能仅靠观察几个样本来断言。一个有用的直觉是：若相同或非常相近的输入在输入空间的一个小邻域内反复对应差异很大的标签，那么只使用这些输入时就存在不可消除的不确定性；即使知道真实条件分布，也不可能准确预测每一个样本。相反，若$X$能够稳定地缩小$Y$的不确定性，使合适的预测函数在独立新样本上显著优于“不看输入、只预测平均值或多数类”的基线，则该任务具有可学习的信号。这里的“规律”不要求是线性的、显式的，或容易写成公式；它可以是高度非线性的统计关系，但必须能够在新的数据中重复出现。

有限数据无法从逻辑上证明这种规律一定存在，只能不断积累证据。实践中应结合四类检查：首先用领域机理审查输入到标签的路径，确认不是把结果发生之后的信息误当作输入；其次检查标签的测量质量、时间对齐和样本覆盖，避免真实信号被噪声或错误配对淹没；再次在严格隔离的验证集或测试集上，与不使用输入的简单基线比较；最后检验结果是否能跨时间、机组或运行条件复现。若模型只在训练集上优于基线，或一旦按时间、设备划分测试集就失效，更合理的结论通常是现有输入不足、分布发生了变化，或模型只记住了偶然性，而不是急于增加模型复杂度。

\paragraph{从训练到部署及反馈。}
沿图\ref{fig:foundations-learning-loop}从采集到上线的路径，可以依次看四个阶段。先把设备记录整理成输入和标签，再选择能处理这些输入的模型，通过训练求出参数，最后用这些参数预测新记录。上线后的错误又会提示我们补数据、改标签或重新训练。评价时除了训练损失，还要检查新工况下的表现、概率是否校准、计算是否及时，以及误报和漏报的后果。

\begin{figure}[htbp]
\centering
\begin{tikzpicture}[>=Latex, node distance=0.85cm, every node/.style={font=\small}]
 \node[draw, rounded corners, fill=blue!10, minimum width=2.15cm, minimum height=0.9cm] (data) {\shortstack{数据\\$\mathcal D$}};
 \node[draw, rounded corners, fill=orange!12, right=of data, minimum width=2.35cm, minimum height=0.9cm] (task) {\shortstack{任务定义\\$X,Y,\ell$}};
 \node[draw, rounded corners, fill=green!12, right=of task, minimum width=2.35cm, minimum height=0.9cm] (model) {\shortstack{模型\\$f_\theta$}};
 \node[draw, rounded corners, fill=yellow!20, right=of model, minimum width=2.35cm, minimum height=0.9cm] (train) {\shortstack{训练\\$\argmin_\theta$}};
 \node[draw, rounded corners, fill=red!10, right=of train, minimum width=2.15cm, minimum height=0.9cm] (deploy) {\shortstack{推理\\与部署}};
 \draw[->, thick] (data)--(task); \draw[->, thick] (task)--(model); \draw[->, thick] (model)--(train); \draw[->, thick] (train)--(deploy);
 \draw[->, thick, dashed] (deploy.north) .. controls +(0,0.72) and +(0,0.72) .. node[above,font=\scriptsize]{反馈、漂移、错误分析} (data.north);
\end{tikzpicture}
\caption{工业数据智能的训练、部署与反馈：部署反馈会重新影响数据和任务定义。}
\label{fig:foundations-learning-loop}
\end{figure}

\section{尚未解决的问题——优化}

本章讲解了数据、任务、目标函数、泛化性等概念，告诉了大家什么是一个好的模型，应该用怎样的目标函数训练一个模型，应该怎样评价一个模型。但是还缺一点，有了目标函数，模型中的可学习参数（或者是先验分布），如何根据目标函数一步一步被训练到“最优”状态。关于这个问题，我们后面会用一个单独的章节去介绍。
\section*{辅助学习材料}
\begingroup
\emergencystretch=3em
\begin{itemize}
\item \textbf{Linear Models}；作者/平台：scikit-learn Developers；英文；官方文档。重点阅读 Ordinary Least Squares、Ridge、Lasso 与 Logistic Regression 小节，配合本章的线性预测、正则化和分类损失理解参数约束与风险最小化。\href{https://scikit-learn.org/stable/modules/linear\_model.html}{在线阅读}。
\item \textbf{Cross-validation: evaluating estimator performance}；作者/平台：scikit-learn Developers；英文；官方文档。重点阅读训练/验证/测试划分、交叉验证和数据泄漏说明，帮助把本章的经验风险、部署风险与可复现实验边界落实为流程。\href{https://scikit-learn.org/stable/modules/cross\_validation.html}{在线阅读}。
\item \textbf{损失函数（Loss Function）}；知乎；中文解说。比较平方损失与交叉熵，核对本章预测对象；标题与主题据必应索引，原页受限。\href{https://zhuanlan.zhihu.com/p/261059231}{在线阅读}。
\item \textbf{损失函数（lossfunction）的全面介绍（简单易懂版）}；CSDN；中文博客。阅读 L1、MSE 与交叉熵说明，比较不同误差的代价；已核对公开页面标题与索引摘要。\href{https://blog.csdn.net/weixin\_57643648/article/details/122704657}{在线阅读}。
\item \textbf{2021 - 类神经网络训练不起来怎么办(四) 损失函数 (Loss) 也可能有影响}；李宏毅；bilibili；中文课程。选看第22集，对照本章损失与学习目标。2026年10月4日官方公开元数据显示整套课程累计播放约299.5万，并非本分集播放量；已核对目录与计数，未逐帧观看。\href{https://www.bilibili.com/video/BV1Wv411h7kN/?p=22}{观看分集}。
\end{itemize}
\endgroup


\chapter{从一个简单模型开始：线性回归}

\begin{quote}\small
\textit{本章摘要。}本章用线性回归，把上一章的数据、任务、模型、目标函数与正则化落实为可以计算的例子。我们从三个角度看同一条线性关系：直接预测数值；在固定参数下预测标签的条件分布；用参数分布表达有限数据带来的不确定性。平方损失、高斯似然、岭惩罚与高斯先验由此联系起来，同时区分最大后验点估计与完整的贝叶斯预测。
\end{quote}

上一章介绍了许多概念，但只知道名称还不足以建立一个模型。例如，输入和标签究竟怎样排成数据？参数怎样从目标函数中求出来？输出一个概率，是否就意味着参数也是随机的？本章沿着“问题形式化、参数估计、贝叶斯线性回归、投影解释”的线索组织推导\,\cite{deisenroth2020mml}，但用本书的工业预测任务重新表述这些关系。

我们始终围绕同一个计算形式：给输入的每个分量乘一个权重，再加上截距。先把它当作直接预测的函数，再把它当作条件分布中的均值，最后不再只保存一组权重，而是保存权重的后验分布。变化的是建模假设和答案的表达方式，不是突然换成了另一类复杂模型。

\section{先固定数据和任务}
一个带标签样本由输入$\bm x_i$和对应的标签$y_i$组成，其中$i$是样本编号。输入是预测时可用的信息，用$d$个数值特征组成的列向量表示；标签是希望根据这些信息预测的目标量。本章讨论实数值的回归任务，因此
\begin{equation}
 \bm x_i=(x_{i1},\ldots,x_{id})^{\mathsf T}\in\R^d,
 \qquad y_i\in\R.
\end{equation}
这里$x_{ij}$表示第$i$个样本的第$j$个特征。标签$y_i$是训练时已经观测到的目标值，模型给出的预测值则记为$\hat y_i$，两者需要区分。标签也可能含有观测噪声，并不一定是完全精确的真值。

将$n$个输入与各自的标签配对，得到训练集
\begin{equation}
 \mathcal D=\{(\bm x_i,y_i)\}_{i=1}^n.
\end{equation}
训练的目的是根据这些已知配对学习从输入到目标的关系，再对新输入预测尚未知的标签。这属于监督学习。对未来量的预测，输入只能包含预测时刻已经可用的信息，未来的目标观测只能作为历史样本的标签，不能提前放入输入。

\paragraph{例：根据当前读数预测下一小时温度。}
在每个历史预测时刻，记录设备当前的温度和振动，并将下一小时实际测得的温度作为标签：
\begin{equation}
 \bm x_i=(\text{当前温度}_i,\text{当前振动}_i)^{\mathsf T}\in\R^2,
 \qquad y_i=\text{下一小时温度}_i.
\end{equation}
这个例子中$d=2$，每一对$(\bm x_i,y_i)$对应一次历史预测及其后续观测。训练完成后，给出新的当前温度和振动读数，模型便输出对下一小时温度的预测。

\paragraph{模型结构与参数各是什么。}
回归模型先写成
\begin{equation}
 f_{\bm w,b}(\bm x)=\bm w^{\mathsf T}\bm x+b.
\end{equation}
“加权求和再加截距”是模型结构，$\bm w$和$b$是要从数据中学习的参数。输入$\bm x$不是参数，历史标签$y$也不是参数。例如，两个权重为$0.8,0.2$、截距为$1$时，预测就是$0.8x_1+0.2x_2+1$；训练要做的是决定这些数，而不是改变已经观测到的读数。

\paragraph{基函数：先表示输入，再在线性空间中学习。}
更一般地，可先将原始输入映射为$M$个预先规定的特征（basis functions，基函数）
\begin{equation}
 \bm\phi(\bm x)=(\phi_1(\bm x),\ldots,\phi_M(\bm x))^{\mathsf T},
 \qquad f_{\bm w}(\bm x)=\bm w^{\mathsf T}\bm\phi(\bm x).
 \label{eq:linear-basis-function-model}
\end{equation}
这时“线性”始终是指$f_{\bm w}$对参数$\bm w$线性，并不要求$\phi_j$对原始输入线性。对一维温度读数$x$，可取$\bm\phi(x)=(1,x,x^2)^{\mathsf T}$，得到二次曲线；对振动窗口，可将均方根、峭度、频带能量或某个频率分量作为不同的$\phi_j$。基函数负责规定模型能够看见和组合什么模式，权重负责根据标签决定这些模式各占多大作用。若基函数由训练数据本身估计，例如标准化均值、PCA 投影或特征选择规则，则它们也必须只在训练部分拟合，再应用于验证和测试数据。

给定$n$个输入后，式\eqref{eq:linear-basis-function-model}形成的设计矩阵为$\Phi\in\R^{n\times M}$，其中$\Phi_{ij}=\phi_j(\bm x_i)$。为避免符号重复，后续仍用$X$表示这个已经形成的设计矩阵。因此，改变原始特征或加入非线性基函数，改变的是$X$的列；最小二乘、岭回归和贝叶斯线性回归的参数推导保持不变。这正是线性模型可作为许多复杂表示的基础层的原因。

以下推导都可以自然包含截距。令增广特征为$\tilde{\bm x}_i=(\bm x_i^{\mathsf T},1)^{\mathsf T}$、增广参数为$\tilde{\bm w}=(\bm w^{\mathsf T},b)^{\mathsf T}$，则$\tilde{\bm w}^{\mathsf T}\tilde{\bm x}_i=\bm w^{\mathsf T}\bm x_i+b$。也就是说，只需在设计矩阵$X$中增加一列全为$1$的固定特征，后续的最小二乘、似然和后验推导无需改变。为使符号简洁，后文仍用$X,\bm w$表示已经增广后的矩阵和参数向量；若加入岭惩罚或高斯先验，通常将截距对应坐标排除在惩罚或收缩之外。本章的手算例子为突出主线，取截距固定为零。

\section{第一种视角：确定性判别模型}
\subsection{线性回归：直接输出一个数}
给定一组参数，模型直接输出$\hat y_i=\bm w^{\mathsf T}\bm x_i$。所谓“确定性”，是指输入和参数固定后，输出就是确定的；并不是说真实观测没有噪声，也不是说训练数据改变后参数仍然不变。

模型已经规定了允许怎样预测，还需要目标函数规定怎样比较预测。选择平方损失
\begin{equation}
 \ell_i=\frac12(\bm w^{\mathsf T}\bm x_i-y_i)^2,
 \qquad \hat R(\bm w)=\frac1n\sum_{i=1}^n\ell_i.
\end{equation}
单个样本的损失衡量一次预测偏差，经验风险是训练样本损失的平均值。平方让正负偏差不会抵消，也使较大的偏差受到更强惩罚。这是我们对错误代价的选择，不是线性模型结构自动规定的唯一损失。

\paragraph{用两个样本求一个权重。}
为便于手算，暂时只保留一个输入，取$(x_1,y_1)=(1,2)$、$(x_2,y_2)=(2,3)$，模型为$\hat y=wx$。本章用损失之和写训练目标：
\begin{equation}
 J(w)=\frac12[(w-2)^2+(2w-3)^2].
\end{equation}
求导得到$J'(w)=(w-2)+2(2w-3)=5w-8$。令导数为零，得到$\hat w=8/5$，两个预测分别为$1.6$和$3.2$。它没有逐点穿过观测，却使总平方误差最小。这就是从数据、模型和损失走到参数估计的完整过程。

\paragraph{从手算到矩阵。}
把每个$\bm x_i^{\mathsf T}$排成设计矩阵$X\in\R^{n\times d}$的一行，将标签排成$\bm y\in\R^n$，则
\begin{equation}
 J(\bm w)=\frac12\|X\bm w-\bm y\|_2^2.\label{eq:linear-least-squares}
\end{equation}
梯度为$X^{\mathsf T}(X\bm w-\bm y)$，令它为零得到\textbf{正规方程}(normal equations)
\begin{equation}
 X^{\mathsf T}X\hat{\bm w}=X^{\mathsf T}\bm y.\label{eq:linear-normal-equation}
\end{equation}
当$X$的列线性无关时，$X^{\mathsf T}X$可逆，解唯一。若两个特征列完全相同，数据只能确定它们系数的和，不能区分各自的系数，这就是上一章所说的可识别性问题。公式可以写成逆矩阵形式，但实际计算通常用矩阵分解求解，而不显式求逆。

这里的条件也说明了样本数与特征数的关系。对$X\in\R^{n\times d}$，$X^{\mathsf T}X$可逆当且仅当$X$有满列秩，即$\operatorname{rank}(X)=d$。因此$n\geq d$只是可逆的必要条件，并非充分条件：即使样本数不少于特征数，重复或线性相关的特征列仍会使矩阵不可逆。反过来，若$n<d$，则
\begin{equation}
 \operatorname{rank}(X)\leq n<d.
\end{equation}
于是$X$不可能满列秩，必存在非零向量$\bm v$满足$X\bm v=\bm0$，从而$X^{\mathsf T}X\bm v=\bm0$；所以$X^{\mathsf T}X$必定不可逆。直观地说，参数个数多于训练样本时，仅凭这些样本无法唯一确定每个权重。

\paragraph{不可逆与病态：解是否唯一，与解是否稳定。}
正规方程的矩阵不可逆，意味着参数不能由数据唯一确定，并不意味着最小二乘问题无解。若特征列完全共线，或样本数不足以约束所有参数，$X$就不满列秩。此时若$\hat{\bm w}$是一个最小二乘解，则对任意满足$X\bm v=\bm0$的向量$\bm v$，$\hat{\bm w}+\bm v$也给出相同的训练集预测。问题在于数据不能唯一识别参数，而非最小二乘目标没有最优解。可以删除冗余特征、重新参数化，或用 Moore--Penrose 伪逆选取最小范数解$X^+\bm y$.

与此同时，更容易被忽略的是：$X$满列秩，解在数学上唯一，但特征近似共线，使某些参数方向只受到数据很弱的约束。令$A=X^{\mathsf T}X$，其特征分解为
\begin{equation}
 A=Q\operatorname{diag}(\mu_1,\ldots,\mu_d)Q^{\mathsf T},
 \qquad \mu_j>0,\qquad Q^{\mathsf T}Q=I.
\end{equation}
记$Q$的第$j$列为$\bm q_j$，则
\begin{equation}
 A^{-1}=\sum_{j=1}^d\frac1{\mu_j}\bm q_j\bm q_j^{\mathsf T}.
\end{equation}
固定$X$，若正规方程的右端$\bm b=X^{\mathsf T}\bm y$受到扰动$\delta\bm b$，参数变化恰好为
\begin{equation}
 \delta\bm w=A^{-1}\delta\bm b
 =\sum_{j=1}^d\frac{\bm q_j^{\mathsf T}\delta\bm b}{\mu_j}\bm q_j.
 \label{eq:linear-ill-conditioned-perturbation}
\end{equation}
因此，小特征值方向上的右端扰动会被放大$1/\mu_j$倍。标签若发生小扰动$\delta\bm y$，右端就相应变化为$\delta\bm b=X^{\mathsf T}\delta\bm y$。当矩阵病态时，标签中的小扰动也可能使求出的权重发生很大变化。实际标签通常含有测量误差、记录误差或其他观测噪声，不能把它们当作完全精确的数值。因此，\emph{矩阵可逆只保证解唯一，不保证解可靠；病态矩阵会使参数估计对标签噪声敏感。}

\paragraph{为什么正规方程也是正交投影条件。}
所有可由该模型给出的训练集预测向量构成$X$的列空间
\begin{equation}
 \mathcal C(X)=\{X\bm w:\bm w\in\R^d\}\subseteq\R^n.
\end{equation}
最小二乘不是在$\R^n$中任意寻找一个接近$\bm y$的向量，而是只允许从$\mathcal C(X)$中选择预测向量。令残差为$\bm r=\bm y-X\hat{\bm w}$。从预测向量$X\hat{\bm w}$沿列空间中的任意方向$X\bm v$作一个微小移动，平方残差不能在一阶继续下降，因此
\begin{equation}
 0=\left.\frac{\mathrm d}{\mathrm d\epsilon}
 \frac12\|\bm r-\epsilon X\bm v\|_2^2\right|_{\epsilon=0}
 =-\bm v^{\mathsf T}X^{\mathsf T}\bm r
 \quad\text{对任意 }\bm v.
\end{equation}
故$X^{\mathsf T}\bm r=0$，这正是式\eqref{eq:linear-normal-equation}。换言之，拟合向量$X\hat{\bm w}$是$\bm y$到列空间的正交投影，残差垂直于每一列特征。若$X$秩不足，不同权重可生成同一个投影向量；因此训练集上的最佳预测唯一，但参数表示未必唯一。选择最小范数解$X^+\bm y$，是在这些同样拟合训练数据的参数中选取一个代表，不改变最优的训练拟合。到这里，我们已经区分了两个问题：数据是否能唯一确定参数，以及求出的参数是否对扰动稳定。接下来讨论怎样在拟合要求之外加入对参数的偏好，以缓解病态情况下参数对标签噪声敏感的问题。

\subsection{正则化：在拟合之外加入偏好}
只追求训练误差小，可能得到对数据扰动敏感的权重。岭回归在平方损失之外加入二次惩罚\cite{hoerl1970ridge}：
\begin{equation}
 J_\lambda(\bm w)=\frac12\|X\bm w-\bm y\|_2^2
 +\frac\lambda2\|\bm w\|_2^2,\qquad\lambda\geq0.\label{eq:linear-ridge-objective}
\end{equation}
第一项要求解释数据，第二项在拟合效果相近时偏好较小的权重。与前面用伪逆从原有最小二乘解中选取最小范数代表不同，有限强度的岭惩罚改变了优化目标，通常也改变最优的训练拟合。因此，它是为了控制参数而引入的建模选择，并不是矩阵不可逆时必须采取的步骤。对参数求导，得到
\begin{equation}
 (X^{\mathsf T}X+\lambda I)\hat{\bm w}_\lambda=X^{\mathsf T}\bm y,
\end{equation}
其中$I$是$d\times d$单位矩阵。$\lambda>0$时，无截距模型的这个解唯一，即使特征列相关也可以求出。

岭回归的作用不仅是让矩阵可逆，更是在数据约束薄弱的方向上抑制过大的权重。若所有坐标均受惩罚，$X^{\mathsf T}X+\lambda I$的特征值为$\mu_j+\lambda$，求解时相应的放大因子从$1/\mu_j$变为$1/(\mu_j+\lambda)$，从而减弱小特征值方向上的噪声放大。但它通过改变估计目标换取稳定性，会引入收缩偏差。若截距不受惩罚，上述统一平移所有特征值的结论不直接适用，可先中心化数据，再对特征权重讨论这一结论。

处理这类问题时，应区分三个目的：完全共线时，删除冗余列、重新参数化或选取伪逆解，可以处理参数表示不唯一；计算普通最小二乘时，QR 或 SVD 分解可以避免显式求逆及构造正规方程造成的额外数值误差，但无法消除数据本身的敏感性；若希望降低噪声引起的参数波动，则可依据验证结果选择岭惩罚、截断小奇异值，或补充能区分相关特征的数据。后两类收缩处理改变了估计规则，需要评估稳定性与偏差的取舍。故既不能从“不可逆”直接推出“必须使用岭回归”，也不能从“已经求出唯一解”直接推出“估计结果可靠”。

在刚才的例子中，$X^{\mathsf T}X=5$、$X^{\mathsf T}\bm y=8$，因此$\hat w_\lambda=8/(5+\lambda)$。取$\lambda=3$，权重从$1.6$缩为$1$，预测变成$1,2$。训练误差变大了，但权重也更小了；是否因此更适合新数据，要用未参与拟合的数据检验，不能只凭惩罚项判断。

$\lambda$是控制学习偏好的超参数，而$\bm w$是给定这个偏好后学出的参数。选择$\lambda$通常需要验证数据。惩罚还受特征单位影响，因此一般先按训练数据的统计量调整特征尺度。若改用平均损失，目标中的相同数值$\lambda$不再代表同样的惩罚强度：$\hat R+\lambda\|\bm w\|^2/2$乘以$n$后，对应的求解矩阵是$X^{\mathsf T}X+n\lambda I$。以下推导都沿用损失求和的约定。

\section{第二种视角：输出有分布，参数取一个确定值}
\subsection{概率回归：把线性输出作为均值}
现在不只说“预测为$\bm w^{\mathsf T}\bm x$”，还假设给定输入后，观测在这个均值附近波动：
\begin{equation}
 Y=\bm w^{\mathsf T}\bm x+\varepsilon,\qquad
 \varepsilon\sim\mathcal N(0,\sigma^2),\qquad
 Y\mid\bm x,\bm w\sim\mathcal N(\bm w^{\mathsf T}\bm x,\sigma^2).
\end{equation}
这里$\mathcal N$表示高斯分布，$\sigma^2>0$是噪声方差。本节先把它作为给定常数。随机性放在标签$Y$上；参数$\bm w$仍是待估计的一个向量，训练后保存的是确定的$\hat{\bm w}$。

假设给定输入和参数后，各样本的观测噪声独立。为了选出最能解释已知标签的参数，可以最大化条件似然
\begin{equation}
 L(\bm w)=\prod_{i=1}^n p(y_i\mid\bm x_i,\bm w).
\end{equation}
似然是把已经观测到的标签固定后，比较不同参数对这些数据的解释程度；它不是参数的概率分布。取负对数，将乘积变成求和，有
\begin{equation}
 -\log L(\bm w)=\frac n2\log(2\pi\sigma^2)
 +\frac1{2\sigma^2}\|X\bm w-\bm y\|_2^2.
\end{equation}
第一项与$\bm w$无关，第二项只是平方损失乘一个正常数。因此，在这些高斯噪声假设下，最大似然估计与最小二乘得到相同权重。平方损失可以直接作为错误度量，也可以从概率假设推出来；两种解释相通，但概率解释增加了具体的观测假设。

训练完成后，模型输出$\mathcal N(\hat{\bm w}^{\mathsf T}\bm x_*,\sigma^2)$。星号表示新输入。需要一个数值时可取其均值；需要描述观测波动时则保留分布。这里尚未将有限数据导致的参数不确定性纳入预测。

\section{第三种视角：参数也用分布表示}
\subsection{先验与后验：不再只选一组权重}
前面两种方法训练后都保存一组参数。现在问：只有少量数据时，是否还有其他权重也能合理解释观测？贝叶斯回归用参数分布保留这种不确定性\cite{bishop2006pattern,murphy2022pml}。这里的随机参数不意味着设备的真实系数每次预测都在改变，而是用分布表达我们对未知系数的认识。

仍采用高斯观测模型，并为权重指定零均值高斯先验：
\begin{equation}
 \bm w\sim\mathcal N(\bm0,\tau^2I),\qquad
 \bm y\mid X,\bm w\sim\mathcal N(X\bm w,\sigma^2I_n),
\end{equation}
其中$\tau^2>0$描述先验允许权重偏离零的程度，$I_n$是$n\times n$单位矩阵。本章把$\tau^2$和$\sigma^2$都作为给定的超参数，先看清参数分布怎样更新。训练输入$X$作为已知条件，不对输入的产生过程建模，因此这里仍是条件判别建模，而不是因为使用了贝叶斯公式就成了生成模型。

若权重维数为$d$，上述两个高斯分布写成密度分别为
\begin{align}
 p(\bm y\mid X,\bm w)
 &=\frac{1}{(2\pi\sigma^2)^{n/2}}
 \exp\left[-\frac{1}{2\sigma^2}
 \|\bm y-X\bm w\|_2^2\right],\label{eq:linear-gaussian-likelihood}\\
 p(\bm w)
 &=\frac{1}{(2\pi\tau^2)^{d/2}}
 \exp\left[-\frac{1}{2\tau^2}\|\bm w\|_2^2\right].
 \label{eq:linear-gaussian-prior-density}
\end{align}
似然的指数项衡量参数给出的均值与观测标签相差多远；先验的指数项衡量权重离零多远。两个前面的系数保证各自对相应变量积分为一。它们在比较不同$\bm w$时看似只是常数，但在写出完整后验或计算模型证据时不能随意丢弃。

将$X$视为已知条件，贝叶斯公式写成
\begin{equation}
 p(\bm w\mid\mathcal D)=
 \frac{p(\bm y\mid X,\bm w)p(\bm w)}
 {\int p(\bm y\mid X,\bm u)p(\bm u)\,\mathrm d\bm u}.
\end{equation}
分子把数据证据与原有偏好结合，分母保证后验积分为一。先验是看见本次训练标签之前的参数分布，后验是结合这些标签之后的参数分布。后验不是一个更复杂的权重向量，而是给许多可能的权重分配不同的概率密度。

\subsection{最大后验估计：从后验中选一个点}
若仍希望最终只保存一组权重，可以选择后验密度最大的点，称为最大后验估计（Maximum A Posteriori, MAP）。取负对数并舍去与权重无关的常数，得到
\begin{equation}
 \hat{\bm w}_{\mathrm{MAP}}=\argmin_{\bm w}
 \left\{\frac1{2\sigma^2}\|X\bm w-\bm y\|_2^2
 +\frac1{2\tau^2}\|\bm w\|_2^2\right\}.
\end{equation}
目标乘以$\sigma^2$不改变最优点。因此，在本章损失求和的约定下，它就是岭回归，且
\begin{equation}
 \lambda=\frac{\sigma^2}{\tau^2}.
\end{equation}
观测噪声越大，数据的相对作用越弱；先验越集中于零，收缩越强。这是正则化的概率解释：二次惩罚来自零均值高斯先验的负对数，并不是看见某个惩罚项就能不加条件地称它为贝叶斯模型。

\textbf{MAP 使用了先验和后验，但最终仍是点估计；它不等于保留了完整的参数分布。}若直接把 MAP 权重代入概率模型，仍会省略参数不确定性对预测的影响。

\subsection{完整后验：均值与方差都保留}
在线性高斯模型中，式\eqref{eq:linear-gaussian-likelihood}与式\eqref{eq:linear-gaussian-prior-density}的乘积仍可写成高斯形式。为看清这个结论，先把后验中依赖$\bm w$的指数项展开。忽略暂时与$\bm w$无关的因子，有
\begin{equation}
 p(\bm w\mid\mathcal D)\propto
 \exp\left\{-\frac12\left[
 \bm w^{\mathsf T}(\sigma^{-2}X^{\mathsf T}X+\tau^{-2}I)\bm w
 -2\bm w^{\mathsf T}\sigma^{-2}X^{\mathsf T}\bm y
 \right]\right\}.
\end{equation}
定义后验精度矩阵和信息向量
\begin{align}
 \Lambda&=\sigma^{-2}X^{\mathsf T}X+\tau^{-2}I,\qquad
 \bm h=\sigma^{-2}X^{\mathsf T}\bm y,\\
 S&=\Lambda^{-1},\qquad \bm m=S\bm h.
\end{align}
由于$\Lambda\bm m=\bm h$，二次项可完成平方：
\begin{equation}
 \bm w^{\mathsf T}\Lambda\bm w-2\bm w^{\mathsf T}\bm h
 =(\bm w-\bm m)^{\mathsf T}\Lambda(\bm w-\bm m)
 -\bm m^{\mathsf T}\Lambda\bm m.
 \label{eq:linear-posterior-completing-square}
\end{equation}
最后一项不依赖$\bm w$，会被贝叶斯公式分母中的证据吸收。于是归一化后的完整后验密度为
\begin{equation}
 p(\bm w\mid\mathcal D)=
 \frac{1}{(2\pi)^{d/2}|S|^{1/2}}
 \exp\left[-\frac12(\bm w-\bm m)^{\mathsf T}S^{-1}(\bm w-\bm m)\right],
 \label{eq:linear-gaussian-posterior-density}
\end{equation}
即$\bm w\mid\mathcal D\sim\mathcal N(\bm m,S)$。其中$\bm m$是后验均值，$S$是后验协方差，描述权重的不确定性及其相互关联；$|S|$是协方差矩阵的行列式，归一化系数保证式\eqref{eq:linear-gaussian-posterior-density}对所有$\bm w$的积分为一。高斯分布的均值恰好也是密度最高点，因此这里$\bm m=\hat{\bm w}_{\mathrm{MAP}}$；这不表示任意后验的均值与 MAP 都相等。

\paragraph{新数据怎样更新已有后验。}
后验也可作为下一批独立数据到来之前的先验。设已有后验为$\mathcal N(\bm m_0,S_0)$，新到一批设计矩阵为$X_1$、标签为$\bm y_1$，且噪声仍为$\mathcal N(\bm0,\sigma^2I)$，则更新后
\begin{align}
 S_1^{-1}&=S_0^{-1}+\sigma^{-2}X_1^{\mathsf T}X_1,\label{eq:linear-sequential-posterior-precision}\\
 \bm m_1&=S_1\left(S_0^{-1}\bm m_0+\sigma^{-2}X_1^{\mathsf T}\bm y_1\right).
 \label{eq:linear-sequential-posterior-mean}
\end{align}
式\eqref{eq:linear-sequential-posterior-precision}说明独立观测增加的是精度（协方差的逆），而不是简单把方差相加。新数据在某个权重方向上的信息由$X_1^{\mathsf T}X_1$决定：若新记录在该方向没有变化，就无法显著降低该方向的不确定性。将全部训练数据一次代入，或将它们分批按式\eqref{eq:linear-sequential-posterior-precision}--\eqref{eq:linear-sequential-posterior-mean}更新，在模型假设相同且不重复计数的条件下得到同一个后验。这个顺序解释也说明，重复复制同一条记录并不能制造独立证据。

仍用两个样本$(1,2),(2,3)$，取$\sigma^2=1$、$\tau^2=1/3$。此时$\lambda=3$，并且
\begin{equation}
 S=(5+3)^{-1}=\frac18,\qquad m=\frac18\times8=1,
 \qquad w\mid\mathcal D\sim\mathcal N(1,1/8).
\end{equation}
岭回归与 MAP 只留下$w=1$；完整后验还留下方差$1/8$。三者在这个例子中可以给出相同的中心预测，但保存的信息并不相同。

\subsection{后验预测：把参数不确定性带到新样本}
对新输入$\bm x_*$，贝叶斯预测不是只代入一组权重，而是对后验中的权重积分：
\begin{equation}
 p(y_*\mid\bm x_*,\mathcal D)=
 \int p(y_*\mid\bm x_*,\bm w)p(\bm w\mid\mathcal D)\,\mathrm d\bm w.
\end{equation}
在线性高斯模型中，这个结果可以直接算出：
\begin{equation}
 Y_*\mid\bm x_*,\mathcal D\sim
 \mathcal N(\bm x_*^{\mathsf T}\bm m,
 \underbrace{\sigma^2}_{\text{观测噪声}}
 +\underbrace{\bm x_*^{\mathsf T}S\bm x_*}_{\text{参数不确定性}}).
\end{equation}
第一项是在给定权重后仍存在的观测波动，第二项是不同合理权重在新输入处产生的预测差异。也可以先看没有观测噪声的线性函数值：它的后验方差只有第二项。二者分别回答“下一次观测有多不确定”和“我们对均值函数有多不确定”，不能混用。

在上述一维例子中，取$x_*=2$，后验预测的均值为$2$，方差为$1+2^2/8=3/2$。若只代入 MAP 权重，则得到$\mathcal N(2,1)$，省略了额外的$1/2$。这正是“输出是概率分布”和“参数也是分布”在预测结果上的具体差别。

本章的线性高斯假设使后验和预测分布都能写成闭式形式，便于看清参数不确定性怎样传递到预测中。更一般的模型可能需要近似推断，后面的贝叶斯章节再讨论这些方法。

\section{把上一章的概念逐项放回例子}
回顾这条线性模型主线，首先要区分几组容易混淆的对象。
\begin{itemize}
 \item \textbf{数据与任务：}输入$\bm x_i$是预测时可用的特征，标签$y_i$是对应的实数目标。在设备温度例子中，当前温度和振动是输入，下一小时温度是标签，任务是回归。
 \item \textbf{结构与参数：}加权求和是结构，权重和截距是参数。线性结构本身就是一种归纳偏置：不允许模型任意改变输入与输出的关系。
 \item \textbf{损失与目标：}平方损失衡量单样本的预测误差；求和后形成训练目标，加入惩罚后则同时考虑拟合和参数偏好。
 \item \textbf{训练与预测：}训练用历史标签求权重或后验；预测将结果用于新输入。点估计下代入固定参数，完整贝叶斯预测则对参数后验积分。
 \item \textbf{正则化与先验：}岭惩罚可以在给定高斯似然和高斯先验时解释为 MAP。惩罚强度、先验方差和噪声方差之间的对应依赖目标的尺度约定。
 \item \textbf{经验风险与泛化：}这些计算决定怎样拟合已有样本，不保证新设备、新时段上的表现。验证仍然需要独立于参数拟合的数据。
\end{itemize}

三种视角不是三项彼此排斥的任务。同一个回归任务，既能用确定性函数给出一个数，也能用固定参数的高斯模型给出条件分布，还能用参数后验给出更完整的预测分布。应该选择哪一种，取决于需要回答的问题，而不是“概率模型”这个名称听起来是否更高级。

本章只把上一章的判别模型分支具体算了一遍，并没有试图用线性回归代替所有模型类型。生成建模、隐藏变量和无监督学习改变了要描述的对象，需要另作假设。即使使用参数后验，模型结构和噪声假设也仍可能不适合数据；后验不确定性并不能自动覆盖所有建模错误。

\section*{辅助学习材料}
\begingroup
\emergencystretch=3em
\begin{itemize}
 \item \textbf{Linear Models}；scikit-learn Developers；英文官方文档。阅读 Ordinary Least Squares、Ridge 与 Bayesian Regression 小节，对照本章的点预测、惩罚目标与参数不确定性。软件中损失的归一化和超参数记法可能不同，需要先核对定义。\href{https://scikit-learn.org/stable/modules/linear\_model.html}{在线阅读}。
 \item \textbf{Mathematics for Machine Learning}；Deisenroth、Faisal、Ong；英文教材。阅读线性回归相关章节，对照设计矩阵、最大似然估计与贝叶斯回归。\href{https://mml-book.com}{书籍网站}。
\end{itemize}
\endgroup

\chapter{从数据到推理的完整流程}
\label{chap:pipeline}
\begin{quote}\small
\textit{本章摘要。}本章把前面定义的学习问题与经典模型落实为可执行流程：先定义事件和样本，再划分数据、训练模型、进行推理并评价部署风险。数据治理和时间边界贯穿全章，因为任何后续算法都无法修复样本构造或信息泄漏造成的错误。
\end{quote}
前一章说明了模型如何依据目标求解，却有意把设计矩阵当作已经给定。工业任务最容易出错的地方，恰恰发生在矩阵形成之前：一条记录何时进入系统，一个窗口对应哪个未来事件，哪些设备可以出现在训练和测试两侧。这些选择不是附属工程细节，而是在决定风险公式中的样本究竟来自什么分布。

本章会反复检查一个朴素的问题：作出预测的那一刻，这条信息真的已经到手了吗？标准化的均值从哪里算，参数用哪些记录训练，阈值在哪一批数据上选择，都要留下清楚的记录。评价时还要把相邻窗口合并回故障事件：同一次故障连续触发许多条告警，并不等于发现了许多次故障。否则离线分数可能很好看，现场却仍然漏报或频繁打扰检修人员。

\section{数据划分与数据治理}

% auxiliary-resource-guide
本章末的“辅助学习材料”列出与核心方法对应的讲解和实践入口，可在阅读相关推导后，按所列小节对照学习。

数据泄漏不仅发生在训练测试样本重叠时，也可能通过预处理、特征选择与标签构造进入整个流程\cite{kaufman2012leakage}。因此以下划分规则约束的是信息使用过程，而不只是一次随机切分。
训练集用于拟合参数，验证集用于选择超参数和比较方案，测试集只在最终报告时使用。若时间顺序具有因果意义，应采用按时间切分，而不是随机打乱。工业数据还需处理缺失、漂移、传感器失真、重复记录与标签延迟。数据清洗不是模型之外的琐事，因为它决定了模型究竟看到了什么分布。

\paragraph{先把三份数据的工作分开。}
例如用一月至六月的数据学习权重，用七月比较惩罚强度和告警阈值，全部选择完成后只用八月报告一次最终结果。测试结果不好当然可以启动下一轮研究，但此后八月不再是未被选择过程使用的测试集。交叉验证是把可用于开发的数据分成若干折，轮流留一折评价、其余折训练，用多次留出结果减少对某一次划分的依赖；它不额外创造数据，也不取代最终测试。时间任务必须改用只拿过去预测未来的滚动形式。数据泄漏指训练或选择过程使用了在目标评估规则下不该提前知道的信息，并不限于直接复制测试标签。

\paragraph{预处理也是需要拟合的模型。}
标准化先减去训练均值，再除以训练标准差，目的是避免不同单位造成过大的尺度差异。若训练值为$2,4$，均值为$3$，本节分母约定下方差为$1$；当尺度下限$0<\epsilon\leq1$时，验证值$6$应变成$(6-3)/1=3$，不是重新用验证值计算均值把它变成零。这里为变换数据而使用除以样本数的方差，不是在估计总体方差时作无偏校正。
设训练索引为$\mathcal T$，第$j$个特征的均值和尺度应只在训练集计算：
\begin{equation}
 \hat\mu_j=\frac1{|\mathcal T|}\sum_{i\in\mathcal T}x_{ij},\qquad
 \hat s_j^2=\frac1{|\mathcal T|}\sum_{i\in\mathcal T}(x_{ij}-\hat\mu_j)^2,
 \qquad z_{ij}=\frac{x_{ij}-\hat\mu_j}{\max(\hat s_j,\epsilon)}.
\end{equation}
这里$\hat s_j$是$\hat s_j^2$的非负平方根，$\epsilon>0$是与该特征单位一致的尺度下限，避免常数列除零；也可明确删除常数列。验证、测试与上线输入都使用同一组训练统计量。其原因可写为函数依赖：最终预测是$f_{\hat\theta}(T_{\hat\phi}(x))$，其中$\hat\phi=\operatorname{Fit}(\mathcal D_{\mathcal T})$表示只用训练数据估计变换参数，随后$T_{\hat\phi}$才把这些参数应用于训练、验证或测试输入；有些软件把这两步合称为FitTransform。若用全数据估计$\hat\phi$，测试集就参与了模型构造，即使没有使用测试标签，也改变了原定的归纳评估问题。

交叉验证时每折必须重新拟合$\hat\phi$与$\hat\theta$，而不是先全量标准化再切折。缺失值填补、PCA、异常值阈值和监督特征筛选也属于这条规则；若部署允许使用一批无标签目标数据适配，应将其明确写成另一个评估规则。

\section{训练、评价、部署与反馈}
可以把流程理解为三种不同的操作：训练改变参数，评价只记录固定规则下的表现，部署则把预测接入实际服务。延迟是单次请求花多久，吞吐是单位时间处理多少请求；二者不能互相替代。所谓冻结参数，是推理时不再因当前输入执行训练更新，但仍需加载训练时确定的预处理统计量。

\begin{figure}[htbp]
\centering
\begin{tikzpicture}[node distance=0.7cm and 0.65cm]
 \node[idi input, minimum width=2.0cm] (raw) {原始记录\\时间戳、设备};
 \node[idi process, right=of raw, minimum width=2.1cm] (window) {窗口与标签\\只用可用信息};
 \node[idi state, right=of window, minimum width=2.0cm] (split) {按时间/设备\\划分数据};
 \node[idi process, right=of split, minimum width=2.0cm] (fit) {训练与验证\\模型版本};
 \node[idi output, right=of fit, minimum width=2.0cm] (serve) {上线推理\\告警与回退};
 \draw[idi arrow] (raw)--(window); \draw[idi arrow] (window)--(split);
 \draw[idi arrow] (split)--(fit); \draw[idi arrow] (fit)--(serve);
 \draw[idi feedback] (serve.north) .. controls +(0,0.68) and +(0,0.68) .. node[above,font=\scriptsize]{延迟标签与审计} (window.north);
 \draw[idi feedback] (split.south) .. controls +(0,-0.42) and +(0,-0.42) .. node[below,font=\scriptsize]{测试集只用于最终报告} (fit.south);
\end{tikzpicture}
\caption{工业学习管道中的信息边界。切分、特征生成和标签构造共同决定是否发生泄漏。}
\label{fig:pipeline-information-boundary}
\end{figure}
图\ref{fig:pipeline-information-boundary}把窗口构造、时间与设备划分和上线反馈连接起来；其中延迟标签返回样本构造环节，而测试集边界限制了模型选择可使用的信息。
\paragraph{贯穿案例的固定设置。}
先以“预测未来二十四小时内的设备故障”建立维护流程的预测任务。设备$i$在预测时刻$t$的可用观测记为
\begin{equation}
 \mathcal{I}_{i,t}=\{x_{i,s}:s\leq t\}\cup\{m_{i,s}:s\leq t\},
\end{equation}
其中$s$表示记录进入系统并可查询的时刻，原始采集时刻另行保存；$x_{i,s}$表示传感器和控制量，$m_{i,s}$表示维护信息。二十四小时故障标签定义为
\begin{equation}
 y_{i,t}=\mathbb{1}\{\text{设备 }i\text{ 在 }(t,t+24\mathrm{h}]\text{ 内发生经确认的故障事件}\}.
\end{equation}
这里的预测时刻、事件发生时刻和记录可用时刻必须分开保存。计划保养并不自动构成故障标签，维修工单只是确认事件的证据来源。若未来二十四小时尚未完整观测，且还没有确认故障，该样本不能直接标为正常。二十四小时是预测时间窗，实际提前量则是故障发生时刻与有效告警时刻之差，需要单独评价。后文的剩余寿命回归、图像定位和维修问答是关联任务，各有标签与指标；共享的是证据可用性原则，而非强行统一输出。

后文比较模型时，都沿用这项任务定义，并逐项记录输入、结构、损失和测试划分。还要写清分数超过多少才告警，传感器缺失时是沿用备用规则还是请人工检查。例如，模型能算出故障概率，并不意味着它已经规定了谁来接收告警、何时安排复检。把这些安排与模型结果分别记录，才能复查成绩差异来自算法还是处理流程。

\paragraph{可用性、标签成熟与时间隔离的形式化。}
对记录$r$分别保存采集时间$e(r)$和可用时间$a(r)$。预测时刻$t$、回看长度$w>0$对应的输入集合应为
\begin{equation}
 \mathcal R_t=\{r:t-w\leq e(r)\leq t,\ a(r)\leq t\},\qquad
 x_t=\operatorname{Aggregate}(\mathcal R_t).
\end{equation}
因此仅比较采集时间不够；迟到数据不能反向修正历史在线输入而仍声称是实时评估。可将这一原则理解为$x_t$必须由$t$时刻已经掌握的信息计算出来。

令$h=24\mathrm h$，训练快照截止于$T$。若标签确认最多延迟$\delta$，且事件追踪完整，则要求$t+h+\delta\leq T$可保守保证一个二元标签成熟。已确认的早期正例虽可提前得知，但不能把尚未成熟的其余样本直接当负例；只保留这些早期正例又会改变样本分布。若延迟没有确定上界，应该依据逐样本的标签成熟状态筛选，或使用明确处理删失的模型。

若训练预测时刻为$t_a$，验证最早预测时刻为$t_b$，为使训练标签不跨入验证预测期，可要求$t_a+h<t_b$。若还要求原始窗口与标签支撑区间完全不重叠，则由$[t_a-w,t_a+h]$与$[t_b-w,t_b+h]$分离得到更保守条件$t_a+h<t_b-w$。后一条件不是所有滚动预测都必须采用：部署时合理地重用过去观测并不自动泄漏，但会造成统计依赖；应根据要模拟的部署方式选择规则并说明残留相关性。

\paragraph{沿时间轴手查一个样本。}
在周二中午预测未来$24$小时故障，回看长度为$6$小时。周二上午十点采集、下午一点才入库的记录不能用作输入；周二晚间发生、周三确认的故障可以事后成为正标签，但确认工单不能回填成周二中午的特征。如果数据只保存到周三上午十点且没有故障记录，这个样本仍未完整观察未来$24$小时，不能标为负例。未完整观察的状态称为删失，处理它需要额外规则而不是猜一个标签。

\paragraph{优化：如何找到参数。}
以可微目标$J(\bm{\theta})$为例，梯度下降更新为
\begin{equation}
 \bm{\theta}_{t+1}=\bm{\theta}_t-\eta_t\nabla_{\bm{\theta}}J(\bm{\theta}_t),
\end{equation}
其中$\eta_t$是学习率。随机梯度下降用第$t$个小批量集合$B_t$（其大小记为$|B_t|$）近似全量梯度；这里的下标$t$是优化步号，不是前文的时刻。
\begin{equation}
 \bm{\theta}_{t+1}=\bm{\theta}_t-\eta_t\frac{1}{|B_t|}\sum_{i\in B_t}\nabla_{\bm{\theta}}\ell_i(\bm{\theta}_t).
\end{equation}
动量、Adaptive Moment Estimation (Adam)和学习率调度都改变了优化轨迹，但不改变“依据目标函数更新参数”的基本逻辑。优化收敛并不自动意味着泛化良好，因为训练目标只是有限数据上的代理目标。

\paragraph{推理与部署。}
训练完成后冻结参数，在新样本上计算$f_{\hat{\bm{\theta}}}(\bm{x})$或$p_{\hat{\bm{\theta}}}(y\mid\bm{x})$。部署阶段还要考虑延迟、吞吐、内存、量化误差、异常输入和模型更新策略。一个在离线测试集上准确的模型，如果无法在传感器数据到达后的时间约束内完成推理，仍然不是有效的工业系统。

\paragraph{评价指标。}
分类可报告准确率、精确率、召回率、$F_1$和Area Under the Receiver Operating Characteristic Curve (ROC-AUC)；回归常用Mean Absolute Error (MAE)与Root Mean Squared Error (RMSE)。概率预测还应检查负对数似然和校准误差。指标必须对应决策目标。例如，安全监测通常优先控制漏报，故不能只使用准确率。报告结果时应同时给出数据划分、随机种子、置信区间或多次运行波动，以区分真实改进和偶然差异。

\paragraph{先用十个样本算指标。}
“正”表示本任务关心的故障，“真/假”表示预测是否正确。假设十个样本中有四个故障，模型报了三个告警，其中两个正确，则$TP=2,FP=1,FN=2,TN=5$。准确率为$(2+5)/10=0.7$，精确率为$2/3$（报出的有多少是真的），召回率为$2/4$（真实故障找到了多少），$F_1=4/7$。对连续预测，若两个误差为$1,-3$，MAE为$(1+3)/2=2$，RMSE为$\sqrt{(1+9)/2}=\sqrt5$；后者更重视大误差。校准指给出约$0.2$概率的一组可比预测中，长期事件比例是否约为$20\%$，不等于单个样本预测正确。

\paragraph{从混淆矩阵推导指标及其基率依赖。}
固定阈值后，真正例、假正例、假负例与真负例数量记为$TP,FP,FN,TN$。于是
\begin{equation}
 \operatorname{Precision}=\frac{TP}{TP+FP},\quad
 \operatorname{Recall}=\frac{TP}{TP+FN},\quad
 F_1=\frac{2TP}{2TP+FP+FN}.\label{eq:pipeline-classification-metrics}
\end{equation}
当精确率$P$与召回率$R$均为正时，调和平均为$2/(1/P+1/R)=2PR/(P+R)$，代入计数即得式\eqref{eq:pipeline-classification-metrics}中的$F_1$；分母为零时必须预先约定报告方式，不能把无正例设备的指标静默设成满分。

令部署故障比例为$\pi=P(Y=1)$，真正例率为$u=P(\hat Y=1\mid Y=1)$，假正例率为$v=P(\hat Y=1\mid Y=0)$。由条件概率和全概率公式，
\begin{equation}
 P(Y=1\mid\hat Y=1)=\frac{\pi u}{\pi u+(1-\pi)v}.
\end{equation}
例如$\pi=0.01,u=0.9,v=0.05$时，精确率约为$0.154$：高召回仍可能伴随大量无效告警。因此在故障比例不同的设备群间，不能只依据同一受试者工作特征（Receiver Operating Characteristic，ROC）工作点推断告警质量。ROC曲线让阈值从高到低变化，横轴为假正例率$v$，纵轴为真正例率$u$；曲线下面积就是ROC-AUC。它也可解释为随机取一正一负样本，正样本得分更高的概率，加上同分概率的一半。因此ROC-AUC考察排序，部署动作还需要固定阈值、基率与成本。只有一种类别的测试子集不能据此计算有意义的ROC-AUC。

\paragraph{窗口误差如何转换为事件指标。}
设第$e$次故障时间为$\tau_e$，要求最小提前量为$L_{\min}$，其中$0<L_{\min}\leq h$。只有位于$[\tau_e-h,\tau_e-L_{\min}]$的有效告警才算及时命中。定义$a_e$为该区间内按既定合并规则选取的首个有效告警，事件召回率为
\begin{equation}
 \frac{\sum_e\mathbb1\{a_e\text{存在}\}}{\text{可评价故障事件数}},\qquad
 L_e=\tau_e-a_e.
\end{equation}
未命中事件不能从召回分母移除；仅在命中事件上报告平均提前量也应明确这一条件。必须事先规定同一告警能否匹配多个事件、连续超阈值如何合并、告警冷却时间如何设置。无效告警频率宜用“未匹配告警数/有效监测设备天数”计算，而不是把重复窗口当成独立动作。

\paragraph{事件合并的手算例子。}
有两次可评价故障，提前量要求为$2$小时；第一次故障前$5,4,3$小时各有一个超过阈值的窗口，按“连续告警合并为一次”规则记一次命中，提前量为$5$小时。第二次故障仅在前$1$小时告警，不满足及时命中条件。因此事件召回率为$1/2$，不能把三个窗口算成三次命中。是否把过晚告警另列为延迟告警，必须在报告中预先定义。

\paragraph{误差条为何依赖抽样单位。}
标准误描述如果重新抽取测试样本，所报告的平均误差会波动多大；它不是单个预测误差。置信区间是由抽样规则构造的区间，其覆盖率含义针对重复抽样，不能仅凭一条曲线给出。
对固定模型的$n$个独立测试错误指示$E_i\in\{0,1\}$，若错误率为$q$，则$\hat q=n^{-1}\sum_iE_i$且$\Var(\hat q)=q(1-q)/n$，标准误可用$\sqrt{\hat q(1-\hat q)/n}$估计；小样本或极端比例应使用更合适的区间方法。若同一设备有$m$个窗口，设备间独立、设备内错误指示方差为$\sigma^2$且两两相关系数为$\rho$，共$G$台设备，则
\begin{equation}
 \Var(\bar E)=\frac{Gm\sigma^2+Gm(m-1)\rho\sigma^2}{(Gm)^2}
 =\frac{\sigma^2}{Gm}\{1+(m-1)\rho\}.\label{eq:pipeline-cluster-variance}
\end{equation}
式\eqref{eq:pipeline-cluster-variance}说明窗口数巨大不代表独立证据同样巨大。可按设备进行成组重采样，时间相关则使用适当的块重采样；重采样方案仍需匹配目标设备总体及依赖结构。多随机种子波动主要衡量训练随机性，不能替代测试抽样的不确定性。

\begin{remark}
 一条实用的诊断顺序是：先检查标签和切分，再检查训练与验证曲线，随后检查错误样本，最后才调整模型结构。否则，复杂模型很容易被用来掩盖数据问题。
\end{remark}

\paragraph{从原始记录到学习样本。}
工业数据通常不是天然的$(x,y)$对，而是多个系统中异步产生的事件。构造样本需要确定实体键、观测时间、预测时间、标签时间和允许使用的信息集合。实现前文的$\operatorname{Aggregate}(\mathcal R_t)$时，按预定规则计算均值、最大值或其他窗口特征，并同时保留采集时间$e(r)$与可用时间$a(r)$的筛选条件。标签则由部署时真正关心的未来区间生成。若把$t$之后才产生的维修结果加入$x_t$，离线指标会被人为抬高。

落实到日志时，输入必须同时满足“已经采集”和“已经可用”。维修工单可能在故障发生数小时后才录入，可以事后用于确认标签，却不能被反向放入预测时刻的输入特征。标签审计还要区分故障维修和计划保养，避免模型实际学习维护排班而非故障风险。

该问题至少需要三组切分：按时间向后滚动的验证集、完全留出的后期测试集，以及按设备留出的迁移测试集。训练阶段可以使用窗口统计、频谱特征和最近一次维护信息，但每个特征都要通过时间可用性检查。常见泄漏包括用全周期均值标准化、用故障后的完整轨迹计算斜率、把同一设备的相邻窗口随机分到不同集合，以及使用在全数据上拟合的异常阈值。每项检查都应有一个可失败的单元测试或样本审计报告。

上线验收不应只看 ROC-AUC。系统还要规定每台设备每天最多允许的告警数、最小提前量、连续超阈值窗口数和人工确认时限。传感器断流时，服务应返回“数据不足”状态并触发数据质量工单；模型超时则回退到上一版经过验证的模型或规则基线。这样，模型风险、数据风险和服务风险才能在运行记录中被区分，而不是都表现为一个模糊的“预测错误”。

\paragraph{训练、验证和测试的职责。}
训练集回答“参数如何拟合”；验证集回答“在候选方案中如何选择”；测试集回答“在选择完成后能否泛化”。当研究者根据测试集反复改模型时，测试集已经承担验证集职责，最终数字会有乐观偏差。时间序列还应使用滚动验证，设备任务则应按设备分组，避免同一实体的相邻窗口跨越划分边界。

\paragraph{数据管道的可追溯性。}
每个特征都应能追溯到原始字段、单位、时间窗口、清洗规则和版本。保存最终矩阵还不够，因为矩阵无法说明它是否使用了未来统计量。建议同时保存样本 Identifier (ID)、设备 ID、预测时间、标签版本、特征生成版本和划分哈希。这样错误样本才可以从模型输出反查到数据源。

\paragraph{评价不等于排序。}
准确率和 RMSE 适合比较某些预测任务，却不直接等于业务价值。一个故障检测器可能提高宏平均$F_1$，却增加大量无效告警；一个回归器可能降低平均误差，却在关键峰值时系统性低估。指标应由决策流程倒推，并报告阈值、事件合并规则、提前量、资源预算和错误成本。

\paragraph{上线前后的差异。}
离线评估使用固定数据和批量计算，线上系统还会遇到数据延迟、网络断开、字段缺失、模型版本不一致和标签延迟。部署验收应至少包括正常输入、缺失输入、极值输入、重复输入、乱序输入和服务超时。模型服务失败时必须有明确回退策略，而不是返回一个未校验的默认预测。

本章关于时间边界、泄漏和评价的规则，可用下面的官方文档作为实施检查表。
\section*{辅助学习材料}
\begingroup
\emergencystretch=3em
\begin{itemize}
\item \textbf{Cross-validation: evaluating estimator performance}；作者/平台：scikit-learn Developers；英文；官方文档。阅读 hold-out、validation、cross-validation 及 test-set overfitting 小节，帮助检查本章的数据切分和泄漏风险。\href{https://scikit-learn.org/stable/modules/cross\_validation.html}{在线阅读}。
\item \textbf{Pipelines and composite estimators}；作者/平台：scikit-learn Developers；英文；官方文档。重点看 Pipeline、ColumnTransformer 与交叉验证中的预处理，帮助落实“只在训练侧拟合变换”的流程约束。\href{https://scikit-learn.org/stable/modules/compose.html}{在线阅读}。
\item \textbf{【机器学习】Cross-Validation（交叉验证）详解}；知乎；中文解说。比较留出法与交叉验证，检查训练和测试的信息边界；标题与主题据必应索引，原页受限。\href{https://zhuanlan.zhihu.com/p/24825503}{在线阅读}。
\item \textbf{交叉验证（Cross-Validation）}；CSDN；中文博客。对照验证流程；时序数据仍需按时间或设备划分，不能直接随机切分；标题与主题据必应索引，原页受限。\href{https://blog.csdn.net/weixin\_42691585/article/details/113971857}{在线阅读}。
\item \textbf{第二节 2021 - 机器学习任务攻略}；李宏毅；bilibili；中文课程。选看第18集，按任务、训练与验证梳理流程。2026年10月4日官方公开元数据显示整套课程累计播放约299.5万，并非本分集播放量；已核对目录与计数，未逐帧观看。\href{https://www.bilibili.com/video/BV1Wv411h7kN/?p=18}{观看分集}。
\end{itemize}
\endgroup
\section{自测与参考答案}
\begin{enumerate}
\item 训练值$10,14$，测试值$16$，按本章标准化规则算测试值。\textbf{答：}训练均值$12$，方差$[(10-12)^2+(14-12)^2]/2=4$，标准差$2$，测试变换为$4/\max(2,\epsilon)$；常用的$0<\epsilon\leq2$时结果为$2$。测试值不参与前两步。
\item $h=24$小时、确认延迟上界$\delta=6$小时，截止时刻为周三中午，最晚哪个预测时刻的标签可按充分条件成熟？\textbf{答：}往前减$30$小时，得到周二上午六点；仍需满足事件追踪完整的前提。
\item 记录模型权重是否足以重现上线结果？\textbf{答：}不足，还需特征与标签版本、训练统计量、阈值和事件合并规则、输入可用时间、模型运行模式等。否则同一权重也会接收不同输入或产生不同动作。
\end{enumerate}

\paragraph{本章小结。}
完整流程不是“下载数据、训练模型、报告分数”，而是把现实过程转换为可复现的学习问题，再把模型输出转换为可审计的决策。数据定义、划分规则、损失函数、优化过程、评价指标和部署约束必须互相一致；其中任何一环含糊，后续的精度数字都难以解释。

本章最重要的边界有两条：一是信息边界，任何输入只能使用预测当时已经可获得的证据；二是评价边界，训练、选择和最终检验不能共用同一份证据而不承认选择偏差。标准化统计量、标签成熟规则、事件合并与置信区间都应服从这些边界，不能仅把模型权重当作完整实验。

确定样本、标签和数据用途之后，下一章开始逐步更新参数：先沿计算图求梯度，再检查步长、曲率、初始化、归一化与数值精度。优化器可以帮助降低已经定义好的损失，却不会发现工单录入晚了，也不会替我们纠正训练集和测试集的混用。

\chapter{优化算法与神经网络训练}

\begin{quote}\small
\textit{本章摘要。}本章解释神经网络如何把损失函数转化为参数更新。叙事依次经过计算图、梯度优化、曲率与尺度、初始化和正则化，最后回到训练故障诊断。这样读者能够把“训练不稳定”定位到目标函数、梯度传播、数值尺度或数据管线，而不是把问题简单归咎于模型容量。
\end{quote}

前一章已经固定样本、信息边界和评价规则，本章在这些条件不变的前提下研究求解。我们暂时不问模型能否应对新设备，而先问：损失是否被正确求导，参数是否沿着有效方向移动，信号是否在多层网络中保持可计算的尺度。把这些问题分开，才能避免用增加模型容量掩盖错误梯度，或用延长训练掩盖数值发散。

选好优化器以后，训练仍可能出问题。例如，同一批误差取总和还是取平均，会改变梯度的大小；初始权重过大，信号经过多层后可能迅速放大。下面先推导一步更新什么时候能降低损失，再分析随机梯度、曲率、初始化和归一化怎样影响训练。章末把这些推导用于排查故障，帮助读者决定该查数据、梯度还是学习率。

\section{目标、梯度与参数更新}

% auxiliary-resource-guide
本章末的“辅助学习材料”列出与核心方法对应的讲解和实践入口，可在阅读相关推导后，按所列小节对照学习。

随机近似奠定了含噪更新的基础，Adam 则将梯度一阶与二阶矩估计用于自适应更新\cite{robbins1951stochastic,kingma2015adam}。算法的更新规则、初始化和偏差校正应一起理解，不能只比较默认学习率。
\paragraph{从目标函数到计算图。}
计算图把一次预测拆成乘法、加法和激活函数等小步骤，再用箭头记录依赖。激活是某层算出的中间表示；激活函数给加权和加入非线性，例如ReLU（修正线性单元）把负数置零、正数原样保留：$\phi(a)=\max(0,a)$。没有非线性时，多层线性变换的复合仍是一个线性变换；若含偏置，则复合后仍只是一个仿射变换，即一次加权和加偏置。因此增加层数本身不能替代激活函数带来的非线性表达。反向传播只负责求导，优化器才负责用导数更新权重，二者不能混称为同一个算法。
训练目标常写成$J(\theta)=\frac1n\sum_i\ell_i(\theta)$。神经网络把它实现为多层复合函数，计算过程可以表示为有向无环图（依赖箭头不形成回路）。前向传播计算预测值，反向传播按照拓扑逆序（先处理依赖者、再处理其输入）计算每个节点对损失的梯度。自动微分系统保存中间激活，因此显存消耗和计算图结构密切相关。

\paragraph{一个两层网络的完整反向传播。}
取单样本$x\in\R^d$，隐藏宽度$m$，标量标签$y$。令
\begin{equation}
 a=W_1x+b_1\in\R^m,\quad h=\phi(a)\in\R^m,\quad
 \hat y=w_2^{\mathsf T}h+b_2,\quad \ell=\tfrac12(\hat y-y)^2,
\end{equation}
其中$W_1\in\R^{m\times d}$，$b_1,w_2\in\R^m$，$b_2\in\R$，$\phi$逐元素作用。先求输出误差$\delta_2=\partial\ell/\partial\hat y=\hat y-y$，然后逐层应用链式法则：
\begin{align}
 \nabla_{w_2}\ell&=\delta_2h,&\quad \partial\ell/\partial b_2&=\delta_2,\\
 \delta_1=\nabla_a\ell&=(w_2\delta_2)\odot\phi'(a),&
 \nabla_{W_1}\ell&=\delta_1x^{\mathsf T},\qquad \nabla_{b_1}\ell=\delta_1.
\end{align}
这里$\delta_1\in\R^m$，外积$\delta_1x^{\mathsf T}$恰为$m\times d$，各项梯度与对应参数形状一致。若同一个中间变量流向多个分支，反向传播必须把各分支贡献相加；这正是计算图反向累积的作用。Rectified Linear Unit (ReLU)在零点不可微，实际实现选择约定的次梯度，通常取零。

小批量损失取平均时，上述梯度也按样本平均；若误用求和，更新尺度会随批量线性增大。可在小型平滑测试网络上用中心差分$[J(\theta+\epsilon v)-J(\theta-\epsilon v)]/(2\epsilon)$核对$\nabla J^{\mathsf T}v$，但应避开不可微点并固定随机屏蔽，以免把随机性误判为梯度错误。

\paragraph{沿一条计算链手算。}
取一个输入、一个隐藏单元，$x=2,y=1,W_1=1,b_1=0,w_2=1,b_2=0$，激活用ReLU。前向依次得到$a=2,h=2,\hat y=2,\ell=1/2$。反向先得$\delta_2=1$，所以输出权重梯度为$2$、输出偏置梯度为$1$；再得$\delta_1=1$，输入权重梯度为$2$、隐藏偏置梯度为$1$。符号$\odot$表示对应位置相乘，不作求和。用步长$0.1$同时更新后，$W_1=w_2=0.8$、$b_1=b_2=-0.1$，新预测为$0.8(0.8\times2-0.1)-0.1=1.1$，新损失为$0.005$。所有梯度必须用更新前的同一套参数算完，不能改完后一层权重才计算前一层的梯度。

\paragraph{一阶优化。}
一阶方法主要使用梯度；二阶方法还使用描述梯度变化的曲率。随机梯度下降（SGD）每步用随机样本或小批量估计梯度；一次更新叫一个训练步，完整遍历训练集一次叫一个训练轮次（epoch）。批量为$B$、样本数为$n$时，通常每轮约有$\lceil n/B\rceil$步，最后一个小批量可能较小。
\begin{figure}[htbp]
\centering
\begin{tikzpicture}[x=1cm,y=0.8cm]
 \draw[->] (-0.2,0)--(6.0,0) node[right] {$\theta_1$};
 \draw[->] (0,-0.2)--(0,3.6) node[above] {$\theta_2$};
 \draw[gray!45] (0.9,0.5) ellipse (2.6cm and 1.2cm);
 \draw[gray!45] (1.7,1.0) ellipse (1.8cm and 0.8cm);
 \draw[gray!45] (2.4,1.45) ellipse (1.0cm and 0.42cm);
 \fill (4.9,3.0) circle (1.5pt) node[above right,font=\scriptsize]{初始点};
 \fill (4.0,2.5) circle (1.5pt); \fill (3.25,2.0) circle (1.5pt);
 \fill (2.55,1.55) circle (1.5pt); \fill (2.4,1.45) circle (1.8pt) node[below right,font=\scriptsize]{极小值附近};
 \draw[idi arrow] (4.9,3.0)--(4.0,2.5); \draw[idi arrow] (4.0,2.5)--(3.25,2.0);
 \draw[idi arrow] (3.25,2.0)--(2.55,1.55); \draw[idi arrow] (2.55,1.55)--(2.4,1.45);
\end{tikzpicture}
\caption{梯度更新沿目标函数的局部下降方向移动；等高线的狭长形状会造成不同参数方向上的尺度差异。}
\label{fig:optimization-descent-path}
\end{figure}
图\ref{fig:optimization-descent-path}用等高线和更新箭头对照目标的局部形状与参数轨迹，说明不同方向的尺度会影响步长选择。梯度下降、随机梯度下降和小批量梯度下降的差别在于梯度估计的精度和更新频率。若小批量梯度$\hat g_t$满足$\mathbb{E}[\hat g_t\mid\theta_t]=\nabla J(\theta_t)$，它虽然带有噪声，却能降低每次更新的计算成本。这里鞍点指梯度为零、任意足够小邻域内都既有更高也有更低函数值的点，例如$u^2-v^2$在原点沿两个坐标方向分别上升和下降。噪声若能扰动到下降方向，可能帮助离开鞍点附近，但无偏性本身并不保证这种逃离。

带动量的更新从$t=0$开始，取$v_0=0$、$0\leq\beta<1$：
\begin{align}
 v_{t+1}&=\beta v_t+\hat g_t,\\
 \theta_{t+1}&=\theta_t-\eta v_{t+1}.
\end{align}
Adam（自适应矩估计）维护近期梯度的平均方向，以及平方大小的平均值；指数移动平均给较早的观测以几何衰减的权重。下面改用$t=1,2,\ldots$编号Adam更新，$\hat g_t$在当前$\theta_t$处计算，初始参数为$\theta_1$，矩状态从$m_0=s_0=0$开始：
\begin{align}
 m_t&=\beta_1m_{t-1}+(1-\beta_1)\hat g_t,\\
 s_t&=\beta_2s_{t-1}+(1-\beta_2)\hat g_t\odot\hat g_t,\\
 \theta_{t+1}&=\theta_t-\eta\frac{m_t/(1-\beta_1^t)}{\sqrt{s_t/(1-\beta_2^t)}+\epsilon}.
\end{align}
其中$m_0=s_0=\mathbf0$，$0\leq\beta_1,\beta_2<1$，分式、平方根和加法均按元素进行；符号$\odot$也表示对应坐标相乘，且$\epsilon>0$为数值稳定常数。自适应学习率不等同于自动选择超参数，学习率、权重衰减和批量大小仍需验证。

\paragraph{用一维抛物线看步长。}
对$J(\theta)=\theta^2/2$，梯度为$\theta$，一步更新为$\theta^+=(1-\eta)\theta$。从$\theta=1$出发，$\eta=0.5$把损失从$0.5$降到$0.125$；$\eta=2$使参数变成$-1$，损失不变；$\eta=3$使参数变成$-2$，损失升到$2$。下降方向正确，并不意味着步子多大都行。

\paragraph{单步下降需要怎样的学习率。}
下面的光滑条件限制“斜率变化有多快”，常数$L$越大，允许可靠使用的步长通常越小。
假设$J$可微且梯度为$L$-Lipschitz，其中$L>0$，即$\|\nabla J(u)-\nabla J(v)\|_2\leq L\|u-v\|_2$。沿线段积分，令$\Delta=v-u$，则
\begin{align}
 J(u+\Delta)-J(u)
 &=\int_0^1\nabla J(u+s\Delta)^{\mathsf T}\Delta\,ds\\
 &\leq\nabla J(u)^{\mathsf T}\Delta+\frac L2\|\Delta\|_2^2.\label{eq:optimization-descent-lemma}
\end{align}
上界来自柯西不等式：梯度差与$\Delta$的内积至多为$Ls\|\Delta\|_2^2$，对$s\in[0,1]$积分便得到$L/2$。在式\eqref{eq:optimization-descent-lemma}中代入$\Delta=-\eta g$、$g=\nabla J(u)$，得到
\begin{equation}
 J(u-\eta g)\leq J(u)-\eta(1-L\eta/2)\|g\|_2^2.
\end{equation}
所以$0<\eta<2/L$且$g\ne0$时可以保证下降；常取$\eta\leq1/L$留出余量。若$J$下有界于$J_{\inf}$，取固定$\eta=1/L$，对正整数$T$步求和，得到
\begin{equation}
 \min_{0\leq t<T}\|\nabla J(\theta_t)\|_2^2
 \leq\frac{2L(J(\theta_0)-J_{\inf})}{T}.
\end{equation}
这是到近似驻点的保证，不是非凸情形下到全局最优的保证。ReLU网络并不处处满足这里的光滑假设，因此该证明是理解步长的基准，而不能未经检查直接套用到全部网络。

\paragraph{随机梯度下为何不能要求每步下降。}
对当前参数条件化，假设$\E[\hat g_t\mid\theta_t]=g_t$、$\E[\|\hat g_t-g_t\|^2\mid\theta_t]\leq\sigma^2/B$。因为交叉项的条件期望为零，$\E[\|\hat g_t\|^2\mid\theta_t]\leq\|g_t\|^2+\sigma^2/B$。将随机更新代入式\eqref{eq:optimization-descent-lemma}并取条件期望，得到
\begin{equation}
 \E[J(\theta_{t+1})\mid\theta_t]
 \leq J(\theta_t)-\eta(1-L\eta/2)\|g_t\|^2
       +\frac{L\eta^2\sigma^2}{2B}.
\end{equation}
噪声增加了一个正项：接近驻点时下降项很小，固定步长可能在邻域内波动。减小步长或增大批量可以减小这一项，但各有时间和计算代价。无偏性依赖抽样方式；非均匀采样若不做相应权重修正，实际优化的目标可能已经改变。

\paragraph{动量与Adam校正项的来源。}
从$v_0=0$展开动量得$v_{t+1}=\sum_{k=0}^t\beta^{t-k}\hat g_k$。若梯度恒定为$g$，则$v_{t+1}\to g/(1-\beta)$，说明本章未乘$(1-\beta)$的动量约定会改变有效更新尺度，不能与归一化约定共用学习率而不加说明。

同理，$m_t=(1-\beta_1)\sum_{k=1}^t\beta_1^{t-k}\hat g_k$。若各步梯度具有同一均值$\mu$，则$\E[m_t]=(1-\beta_1^t)\mu$，所以除以$1-\beta_1^t$可消除零初始化导致的偏低；$s_t$的二阶矩校正同理。训练中梯度分布持续变化，校正并不使整个比值成为某个无偏梯度，也不自动给出收敛保证。$s_t$估计的是逐坐标二阶原点矩，而非中心化方差或Hessian。

\paragraph{二阶信息与曲率。}
二阶泰勒展开为
\begin{equation}
 J(\theta+\Delta)\approx J(\theta)+g^{\mathsf T}\Delta+\frac12\Delta^{\mathsf T}H\Delta.
\end{equation}
若 Hessian 正定，Newton 步为$\Delta=-H^{-1}g$。它在接近极小点时可能快速收敛，但大模型中显式存储和求逆 Hessian 代价过高。拟牛顿法、自然梯度和共轭梯度通过近似曲率获得折中。

\paragraph{Newton步和条件数。}
对二次近似$q(\Delta)=J+g^{\mathsf T}\Delta+\tfrac12\Delta^{\mathsf T}H\Delta$求导，得到$\nabla_\Delta q=g+H\Delta$；若$H$正定，唯一最小点满足$H\Delta=-g$，实现上应解线性方程而非显式形成逆矩阵。若$H$不定，二次近似可能无下界，Newton方向也不一定下降；加阻尼$H+\lambda I$使其正定，或使用信赖域限制步幅，是处理这一问题的思路。

对严格凸二次函数$J(\theta)=\tfrac12\theta^{\mathsf T}H\theta-b^{\mathsf T}\theta$，其中$\theta,b\in\R^d$且$H\in\R^{d\times d}$对称正定，记误差$e_t=\theta_t-\theta_*$。梯度下降给出$e_{t+1}=(I-\eta H)e_t$。在$H$的第$j$个特征方向上，误差乘以$1-\eta\lambda_j$，所以所有方向稳定要求$0<\eta<2/\lambda_{\max}$。当$\kappa=\lambda_{\max}/\lambda_{\min}$很大时，为陡峭方向选的小步长会使平缓方向收敛很慢。这就是标准化与预条件可能加速训练的几何原因。

\section{尺度、信号传播与正则化}
\paragraph{学习率与尺度。}
学习率过大导致损失震荡或发散，过小则训练缓慢。预热（warm-up）可在训练初期逐步增大学习率，余弦退火则使学习率平滑下降。不同层的梯度尺度可能相差很大，归一化、残差连接和合理初始化帮助维持稳定的信号传播。

\paragraph{初始化与信号传播。}
若网络层的权重方差过大，前向激活和反向梯度会指数放大；若过小，则会逐层衰减。Xavier初始化试图保持线性层输入输出方差相近，He初始化适用于 ReLU 激活。ReLU为$\max(0,a)$，计算简单且正区间梯度稳定，但可能产生永久不激活单元。

\paragraph{正则化技术。}
Dropout在训练时随机屏蔽部分激活，相当于在许多子网络上进行参数共享。Batch Normalization对小批量激活进行标准化，Layer Normalization则在单个样本的特征维度上归一化，后者更适合序列模型。早停把验证集性能作为停止信号，实质上限制了有效训练时间。

\paragraph{先区分几种尺度操作。}
初始化只设置训练开始时的权重，输入标准化使用训练数据拟合的尺度，网络归一化作用于内部激活，权重衰减直接限制参数大小。它们操作的对象不同，不能用其中一种代替其他所有操作。比如一层有$100$个输入，He初始化给权重方差$2/100=0.02$，若用零均值正态分布抽样，应传入标准差$\sqrt{0.02}\approx0.141$，不能把方差直接当标准差。

\paragraph{Dropout的期望保持与噪声代价。}
若某激活为$4$、保留概率为$1/2$，训练时各以一半概率输出$0$或$8$，平均仍为$4$；推理时输出$4$。若保留概率记为$q$，丢弃概率就是$1-q$，不同接口可能使用相反的参数名称。
训练时采用保留概率$q\in(0,1]$，令独立屏蔽变量$m_j\sim\operatorname{Bernoulli}(q)$，倒置Dropout输出$\tilde h_j=m_jh_j/q$。条件于激活$h_j$，
\begin{equation}
 \E[\tilde h_j\mid h_j]=h_j,\qquad
 \Var(\tilde h_j\mid h_j)=h_j^2(1-q)/q.
\end{equation}
因此推理时可以直接使用$h_j$而无需再缩放。但对非线性网络，先求平均再计算输出，一般不等于先计算输出再求平均，即一般不能写成
\[
 \E[f(\tilde h)]=f(\E\tilde h).
\]
所以不能将推理网络称为所有子网络输出的精确平均。保留概率过小会显著增加噪声，是否改善泛化仍须验证。

\paragraph{归一化到底沿哪个维度进行。}
以单个通道的一批激活$a_1,\ldots,a_B$为例，Batch Normalization计算
\begin{equation}
 \mu_B=B^{-1}\sum_i a_i,\quad
 v_B=B^{-1}\sum_i(a_i-\mu_B)^2,\quad
 z_i=\gamma\frac{a_i-\mu_B}{\sqrt{v_B+\epsilon}}+\beta.
\end{equation}
标准化部分的批均值为零，批方差为$v_B/(v_B+\epsilon)$，并非严格等于一；$\gamma,\beta$恢复可学习的尺度和平移。训练输出依赖同批其他样本，推理通常改用训练阶段累计统计量，因而必须切换模式并防止测试数据更新这些统计量。Layer Normalization对一个样本的指定特征轴求均值与方差，不依赖批内其他样本；但选择哪些轴仍取决于张量结构。归一化不是保证泛化的定理，也不能替代输入单位审计。

\section{训练诊断与实验解释}
\paragraph{训练故障诊断。}
先在固定的极小数据集上关闭随机增强和Dropout，检查一个足够有表达力的模型能否显著降低训练损失；若不能，优先核查标签对应、梯度和更新实现，而不是立即搜大型结构。这只是排错试验：拟合少量样本既不证明泛化，也不排除测试泄漏。比较训练和验证损失时，还要统一损失定义，区分是否包含正则项，以及是否启用了随机屏蔽。
训练损失不下降时，优先检查标签、损失实现、输入范围、学习率和梯度是否为零。训练损失下降但验证损失上升，通常意味着过拟合或训练测试分布不一致。训练与验证都很差，可能是欠拟合、标签噪声、特征不充分或优化尚未稳定。梯度范数、激活均值和预测熵是有用的中间诊断量。

\paragraph{算法框。}
\begin{algorithm}
\caption{小批量梯度训练的基本流程}
\begin{algorithmic}[1]
\Require 非空训练集$\mathcal{D}$与验证集$\mathcal V$，初始参数$\theta$，学习率$\eta>0$，步数$T\geq1$
\State 设最佳验证损失为$+\infty$，最佳模型状态为空
\For{$t=1,\ldots,T$}
  \State 切换训练模式，从$\mathcal{D}$抽取小批量$B_t$，清空上一轮梯度
  \State 计算$J_t=|B_t|^{-1}\sum_{i\in B_t}\ell_i(\theta)$
  \State 通过反向传播计算$g_t=\nabla_\theta J_t$
  \State 更新$\theta\leftarrow\theta-\eta g_t$
  \State 切换评价模式，关闭梯度记录，在$\mathcal V$计算未加正则项的平均损失
  \If{验证损失小于已记录的最佳值}
    \State 更新最佳值，复制保存参数及归一化运行统计量
  \EndIf
\EndFor
\State \Return 保存的最佳模型状态
\end{algorithmic}
\end{algorithm}

算法中小批量必须非空；训练损失、梯度、更新后参数或验证损失出现非有限值时应中止并报告数值失败，而不是返回空状态或把旧检查点冒充本次成功结果。正常结束后，推理前须加载所保存的最佳状态并切换评价模式。

本算法为清楚起见每步验证；实际可预先约定每轮或每隔若干步验证一次，但不能用测试表现选择频率或检查点。保存必须是状态的独立副本，否则后续更新可能把所谓“最佳参数”覆盖。清空梯度适用于不打算跨批累积的情形；需要累积时必须明确累计多少批及其归一化方式。

\paragraph{损失景观与可辨识性。}
参数空间中的不同点可能实现同一个函数。例如隐藏单元交换顺序不会改变网络输出，因而神经网络损失景观存在大量等价参数化。优化器并不是在寻找唯一“真实参数”，而是在这些等价表示中找到一个训练和泛化都可接受的解。观察训练损失下降时，还应检查参数范数、梯度范数和验证性能是否同步改善。

\paragraph{批量大小与梯度噪声。}
固定当前参数，设均匀抽取的单样本梯度为$g_i$，总体均值为$\bar g$，协方差为$\Sigma$。独立有放回抽取$B$个梯度时，
\begin{equation}
 \operatorname{Cov}\left(B^{-1}\sum_{b=1}^Bg_{i_b}\right)
 =B^{-2}\sum_{b=1}^B\operatorname{Cov}(g_{i_b})=\Sigma/B.
\end{equation}
若从$n>1$个固定样本无放回抽取，且$\Sigma=n^{-1}\sum_i(g_i-\bar g)(g_i-\bar g)^{\mathsf T}$，则结果为$(n-B)\Sigma/[B(n-1)]$，$B=n$时方差为零。这些推导都依赖当前参数下的抽样规则，相关窗口或有偏抽样需要重新计算。增大批量会降低估计噪声并提高并行效率，却可能减少单位样本带来的更新次数。批量大小改变时，学习率、正则化强度和训练步数也应重新校准；不能只替换 batch size 再比较结果。

\paragraph{动量、Adam 与权重衰减。}
动量保留历史方向，能减弱狭长谷底中的来回振荡。Adam 对不同坐标使用自适应尺度，适合稀疏梯度和初期训练，但其自适应归一化不等同于解耦权重衰减。实际实现中应区分把$L_2$项加到梯度里的正则化和直接衰减参数的 AdamW 式更新，因为二者在自适应优化器下并不完全相同。

对普通Stochastic Gradient Descent (SGD)，把$\lambda\|\theta\|^2/2$加入目标产生更新$\theta^+=(1-\eta\lambda)\theta-\eta g$，恰好等于直接衰减。若优化器当前的坐标预条件矩阵为对角阵$D_t$，把正则项加入梯度则得到$\theta^+=\theta-\eta D_tg-\eta\lambda D_t\theta$；解耦衰减则为$\theta^+=(1-\eta\lambda)\theta-\eta D_tg$。前者的收缩强度随坐标尺度变化，后者不随$D_t$变化；在Adam中正则项还会进入历史矩估计，差异更进一步。应同时记录衰减公式及哪些参数被排除，如偏置和归一化参数。

\paragraph{反向传播的内存占用。}
反向传播需要用到前向计算中的输入、激活和部分中间结果，训练时通常要把它们保存在显存里。激活检查点只保存其中一部分，其余在反向传播时重新计算，代价是多花计算时间。混合精度利用低精度运算提高吞吐，常配合高精度参数副本；是否需要损失缩放（loss scaling），还取决于所用数值格式。部署预测通常只做前向计算，因此训练时占用的显存不能直接当作推理需求。

\paragraph{初始化的方差推导。}
令预激活$a=\sum_{j=1}^{n_{in}}w_jx_j$，偏置初始为零。假设权重彼此独立、均值为零、方差为$\sigma_w^2$，并独立于输入；输入各坐标具有相同二阶矩$q_x=\E[x_j^2]$。交叉项因$\E[w_jw_k]=0$消失，故
\begin{equation}
 \E[a]=0,\qquad \Var(a)=\E[a^2]
 =\sum_j\E[w_j^2]\E[x_j^2]=n_{in}\sigma_w^2q_x.
\end{equation}
只有输入均值为零时$q_x$才等于$\Var(x_j)$。若$a$关于零对称且$h=\max(0,a)$，则对称性给出$\E[h^2]=\tfrac12\E[a^2]$。为逐层保持激活二阶矩，要求$n_{in}\sigma_w^2/2=1$，于是得到He方差$\sigma_w^2=2/n_{in}$。这里保持的是二阶矩，不是ReLU输出的方差，因为$\E[h]$通常大于零。

对近似线性的激活，前向保持尺度要求$n_{in}\sigma_w^2\approx1$；反向传播在类似独立近似下要求$n_{out}\sigma_w^2\approx1$。当输入输出宽度不同，不能同时精确满足，Xavier取折中值$\sigma_w^2=2/(n_{in}+n_{out})$。训练后权重与激活会相关，因此这是初始化阶段的尺度分析，不是整个训练期的保证。

多层网络的反向梯度含有Jacobian乘积$J_1^{\mathsf T}\cdots J_L^{\mathsf T}$，其范数上界由$\prod_l\|J_l\|_2$控制，解释了逐层放大或衰减的风险。残差映射$h_{l+1}=h_l+F_l(h_l)$的Jacobian为$I+DF_l$，提供直接传播通路，但并不保证所有奇异值都接近一。初始化、残差和归一化需要共同检验，不能各自被当作梯度稳定性的充分条件。

\paragraph{训练曲线的诊断矩阵。}
若训练和验证损失都高，先检查标签、输入尺度和模型是否欠拟合；若训练损失低而验证损失高，检查正则化、数据切分和分布偏移；若损失下降但任务指标不变，检查指标阈值、标签定义和类别不平衡；若梯度为 NaN，检查学习率、混合精度和数值范围。每种曲线形态对应不同假设，不能用统一的“再训练久一点”处理。

\section{可复现训练设置}
检查点不只是参数文件。若要接着训练，应保存优化器的动量或Adam矩估计、学习率调度位置、训练步数以及随机数状态；只有预测时才通常不需要这些优化器状态。随机种子确定伪随机序列的起点，但不同硬件、并行算子和软件版本仍可能带来差异，因此“用了同一个种子”不等于任意平台逐位相同。

一次训练应固定数据版本、切分文件、模型配置、随机种子、优化器版本和停止规则。保存最佳验证检查点时要记录选择指标和选择时刻，否则重新加载的模型可能不是报告结果对应的模型。多次运行应报告均值和波动，尤其是小数据和元学习任务。

反向传播、自动微分和训练循环可借助下面两份官方教程对照实现，阅读时把代码步骤对应回本章的梯度推导。
\section*{辅助学习材料}
\begingroup
\emergencystretch=3em
\begin{itemize}
\item \textbf{Introduction to PyTorch}；作者/平台：PyTorch 官方教程；英文；视频与教程。依次阅读 Tensors、Autograd、Building Models 和 Training Models，帮助理解计算图、反向传播、优化器与复现设置。\href{https://docs.pytorch.org/tutorials/beginner/introyt/introyt1\_tutorial.html}{在线阅读}。
\item \textbf{Training a Classifier}；作者/平台：PyTorch 官方教程；英文；代码教程。重点看数据加载、损失、优化器、训练/测试循环和保存模型，帮助把本章的 SGD、批训练和验证步骤落到可运行流程。\href{https://docs.pytorch.org/tutorials/beginner/blitz/cifar10\_tutorial.html}{在线阅读}。
\item \textbf{深度学习 | 反向传播详解}；知乎；中文解说。沿计算图手算链式法则，再核对自动微分；标题与主题据必应索引，原页受限。\href{https://zhuanlan.zhihu.com/p/115571464}{在线阅读}。
\item \textbf{“反向传播算法”过程及公式推导（索引标题节选）}；CSDN；中文博客。阅读三层网络的前向与反向图示，区分梯度计算和参数更新；已读取公开页面的定义与图示说明，完整标题未核实。\href{https://blog.csdn.net/ft\_sunshine/article/details/90221691}{在线阅读}。
\item \textbf{(选修)To Learn More - 反向传播(Backpropagation)}；李宏毅；bilibili；中文课程。选看第14集，对照本章计算图与链式法则。2026年10月4日官方公开元数据显示整套课程累计播放约299.5万，并非本分集播放量；已核对目录与计数，未逐帧观看。\href{https://www.bilibili.com/video/BV1Wv411h7kN/?p=14}{观看分集}。
\end{itemize}
\endgroup
\section{自测与参考答案}
\begin{enumerate}
\item 两个单样本梯度为$2,6$，批平均梯度是多少？若程序误用求和呢？\textbf{答：}平均为$4$，求和为$8$；相同学习率下后者步幅加倍。梯度累积也要核对最终归一化。
\item Adam第一步梯度为$2$，$\beta_1=0.9,\beta_2=0.999$，校正后两矩是什么？\textbf{答：}$m_1=0.2,s_1=0.004$，除以$0.1,0.001$后分别为$2,4$；因此更新量为$\eta\,2/(2+\epsilon)$，约为$\eta$，不是$2\eta$。
\item 为什么不能只关掉梯度计算就当作进入推理模式？\textbf{答：}梯度记录与网络模式是两回事；Dropout还需停止屏蔽，Batch Normalization通常还需切换到训练累计统计量。评价时既不更新参数，也不让测试数据更新运行统计量。
\end{enumerate}

\section{本章小结}
优化是把目标函数转化为参数更新的过程，训练工程则负责让这个过程数值稳定、可诊断、可复现。学习率、批量、初始化、归一化和精度共同决定优化轨迹；它们不能被当作互相独立的旋钮。

从推导看，反向传播回答的是梯度怎样计算，下降引理回答的是在光滑假设下步长何时有效，随机梯度分析回答的是噪声为何留下波动，而初始化分析回答的是多层信号为何会改变尺度。这些结论各有条件：求到驻点不等于全局最优，优化器运行稳定也不等于学到了可迁移的规律。

训练损失已经稳定下降，换一台设备后却可能预测不准。下一章讨论这种泛化问题，分别检查模型复杂度、正则化和数据分布变化。本章让参数更新尽量正确有效；下一章要检验更新得到的模型在未见数据上是否仍然有用。

\chapter{统计学习理论与模型选择}

\begin{quote}\small
\textit{本章摘要。}本章为模型比较提供统计学习视角。先说明假设空间和容量如何影响泛化界，再讨论结构风险、交叉验证、校准、不平衡数据和置信区间。最后把理论约束转成数据泄漏检查，使模型选择从“哪个分数最高”转向“哪个结论在何种不确定性下成立”。
\end{quote}

上一章说明训练误差、抽样误差和分布偏移不是一回事，但现实实验还多了一层选择：研究者会试很多窗口、特征和超参数，再报告其中最好的结果。这时需要控制的不再是某个预先固定函数的误差，而是整个选择过程可能带来的乐观偏差。本章以这个区别为出发点，把统计学习的容量约束与实际验证规则连接起来。

试过许多窗口长度以后再挑最高分，相当于又用数据作了一次选择。若仍用同一批数据报告最终成绩，就可能把偶然的好运当成稳定的提升。理论中的复杂度代价说明这种问题为什么会发生，交叉验证帮助安排有限样本的用途，校准和成本分析则帮助把分数变成告警规则。每一种方法都有适用条件，不能替代对新设备的独立检验。

阅读本章前，应能区分训练平均损失与总体风险，并理解一次数据划分为何不能证明所有工况下的表现。第一次阅读可先掌握有限候选集合的联合界和嵌套验证；无限类复杂度推导用于回答连续参数模型如何控制选择偏差，不是实施交叉验证的前置操作。

\section{假设空间与泛化约束}

% auxiliary-resource-guide
本章末的“辅助学习材料”列出与核心方法对应的讲解和实践入口，可在阅读相关推导后，按所列小节对照学习。

\begin{figure}[htbp]
\centering
\begin{tikzpicture}[x=1cm,y=0.75cm]
 \draw[->] (0,0)--(7.2,0) node[right] {模型复杂度};
 \draw[->] (0,0)--(0,4.0) node[above] {误差};
 \draw[blue!70!black, thick] plot[smooth] coordinates {(0.5,3.5)(1.5,2.6)(2.5,1.8)(3.5,1.25)(4.5,0.95)(5.5,0.8)(6.8,0.72)};
 \draw[red!70!black, thick] plot[smooth] coordinates {(0.5,1.0)(1.5,1.25)(2.5,1.7)(3.5,2.2)(4.5,2.9)(5.5,3.35)(6.8,3.7)};
 \draw[black!65, thick, dashed] plot[smooth] coordinates {(0.5,3.7)(1.5,2.9)(2.5,2.2)(3.5,1.75)(4.5,1.55)(5.5,1.65)(6.8,2.05)};
 \node[blue!70!black, font=\scriptsize] at (5.4,0.55) {训练误差};
 \node[red!70!black, font=\scriptsize] at (5.7,3.7) {复杂度代价（示意）};
 \node[font=\scriptsize] at (3.9,1.95) {真实风险};
 \draw[dashed] (4.35,0)--(4.35,1.5);
 \node[font=\scriptsize, anchor=north] at (4.35,-0.08) {折中点};
\end{tikzpicture}
\caption{模型复杂度同时影响拟合误差与估计方差；模型选择的目标是控制真实风险，而不是单独追求训练误差。}
\label{fig:statistical-complexity-tradeoff}
\end{figure}
图\ref{fig:statistical-complexity-tradeoff}是定性示意，不是风险分解的数值相加，也不保证真实风险总呈U形或只有一个最低点。它提醒读者，选模应控制未见样本风险，而非仅追求更低训练误差；理论复杂度上界也不等于实测的估计方差。
\paragraph{假设空间。}
学习算法从假设空间$\mathcal{H}$中选择函数。有限训练集只能约束$\mathcal{H}$在观测点上的行为，因此假设空间越丰富，插值训练数据越容易，但泛化分析也越困难。模型容量不能简单等同于参数数量，还与参数化方式、数据几何和优化偏置有关。

\paragraph{读概率界前先确定谁在随机变化。}
假设空间$\mathcal H$是一组允许的预测函数，而不是数据表。以标量温度$x$为输入，三个预先固定的阈值分类器$f_c(x)=\mathbb 1\{x>c\}$、$c\in\{60,70,80\}$就组成含三个函数的空间。若阈值可以取任意实数，则空间无限，尽管每个函数只有一个参数。

以下$\Pr$指重新抽取训练集时的概率，$\epsilon>0$是允许的风险误差，$\delta$是允许理论保证失败的概率。$\sup$表示上确界，有限集合时就是最大值；$|\mathcal H|$是函数个数，$\log$是自然对数。固定函数的界像保证一个预先选定候选；同时界像保证整张候选名单，随后再看数据选人也仍在名单内。它不是说每个样本的预测误差都小于$\epsilon$。

\paragraph{统一收敛的直观证明。}
对固定函数$f$，若损失$\ell(f(X),Y)\in[0,1]$，Hoeffding不等式给出
\begin{equation}
 \Pr\left(|R(f)-\hat R(f)|>\epsilon\right)\leq2\exp(-2n\epsilon^2).
\end{equation}
若$\mathcal{H}$是有限集合，对所有$f\in\mathcal{H}$使用联合界可得
\begin{equation}
 \Pr\left(\sup_{f\in\mathcal{H}}|R(f)-\hat R(f)|>\epsilon\right)
 \leq2|\mathcal{H}|\exp(-2n\epsilon^2).\label{eq:statistical-union-bound}
\end{equation}
这说明样本量需要随假设空间复杂度增加。对无限空间，可用 Vapnik--Chervonenkis (VC) 维、Rademacher复杂度或覆盖数替代$|\mathcal{H}|$。

\paragraph{从概率尾界到样本量。}
以上$n$个样本必须独立同分布，$\mathcal H$应预先固定且非空；平方损失无界时不能不加条件地使用式\eqref{eq:statistical-union-bound}。为看到指数项的来源，令$Z_i=\ell(f(X_i),Y_i)\in[0,1]$。有界变量的指数矩满足$\E\exp(t(\E Z_i-Z_i))\leq\exp(t^2/8)$；独立性使指数矩相乘。Markov不等式给出
\begin{equation}
 \Pr(R-\hat R>\epsilon)
 \leq\exp(-tn\epsilon+nt^2/8).
\end{equation}
取$t=4\epsilon$得到单侧$\exp(-2n\epsilon^2)$，对另一侧重复并相加即得双侧界。联合界不要求各函数的误差独立，因此适用于在同一数据上比较的多个模型。

令式\eqref{eq:statistical-union-bound}右端为$\delta\in(0,1)$，解得以至少$1-\delta$的概率，对所有$f\in\mathcal H$同时成立
\begin{equation}
 |R(f)-\hat R(f)|\leq b_n,
 \qquad b_n=\sqrt{\frac{\log(2|\mathcal H|/\delta)}{2n}}.\label{eq:statistical-uniform-radius}
\end{equation}
若希望误差不超过$\epsilon$，充分条件是$n\geq\log(2|\mathcal H|/\delta)/(2\epsilon^2)$。这是一条保守的充分保证，不是对具体任务所需数据量的精确预测。

设$\hat f$最小化经验风险，$f_{\mathcal H}^*$最小化类内真实风险，则在式\eqref{eq:statistical-uniform-radius}对所有函数同时成立的事件上
\begin{align}
 R(\hat f)&\leq\hat R(\hat f)+b_n
 \leq\hat R(f_{\mathcal H}^*)+b_n
 \leq R(f_{\mathcal H}^*)+2b_n.
\end{align}
第一步和第三步都需要同时界；只给预先固定$f$的界无法直接代入数据依赖的$\hat f$。若优化仅达到$\eta$精度，最后再加$\eta$。相对所有可测函数中最优风险的差还包含类外近似误差，故增大$n$不能弥补错误的函数类。

\paragraph{代入数字检验界的含义。}
若$|\mathcal H|=10$、$n=1000$、$\delta=0.05$，则$b_n=\sqrt{\log(400)/2000}\approx0.055$。经验错误率为$0.08$的所选函数，在同时界成立时真实错误率不超过约$0.135$；相对类内最优函数的保证则需$2b_n\approx0.110$，不能把这两个误差半径混用。若要求$b_n\leq0.05$，充分样本数为$\lceil\log(400)/(2\times0.05^2)\rceil=1199$。这些数值较保守，说明“理论能证明”与“实验已经达到”是两种证据。

无限类段落中的VC维可先用一维阈值理解：一个点可以任意标为零或一，但两个由小到大排列的点不能得到“一、零”的标记，因此该类的VC维为一。Rademacher随机符号则相当于独立抛硬币生成的正负记号；复杂度询问函数类能在多大程度上追随这种无规律波动。读者可先掌握这两个对象，再阅读它们的上界推导。

\paragraph{结构风险最小化。}
经验风险最小化只看训练误差，结构风险最小化在一系列嵌套假设空间中选择复杂度。实际的正则化、树深度限制、网络宽度和数据增强都可视为控制有效复杂度的方法。理论上更强的模型仍可能泛化良好，因为优化算法偏爱某些简单或平滑的解。

\paragraph{复杂度惩罚如何随候选类分配。}
对预先指定的每个有限类$\mathcal H_m$，选择正权重$\pi_m$。权重之和为一，并据此分配失败概率：
\begin{equation*}
 \sum_m\pi_m=1,\qquad \delta_m=\delta\pi_m.
\end{equation*}
再次使用联合界，得到所有类同时有效的偏差界
\begin{equation}
 b_m=\sqrt{\frac{\log(2|\mathcal H_m|/(\delta\pi_m))}{2n}}.
\end{equation}
若在所有$m,f\in\mathcal H_m$中最小化$\hat R(f)+b_m$，则所选模型$\tilde f$满足
\begin{equation}
 R(\tilde f)\leq\inf_m\inf_{f\in\mathcal H_m}\{R(f)+2b_m\}.
\end{equation}
推导是先用$R(\tilde f)\leq\hat R(\tilde f)+b_{\tilde m}$，再用最小化性质与每个候选比较，最后把候选的经验风险换成真实风险。复杂类要靠真实的拟合改善抵消较大的惩罚，而不是仅靠训练误差获胜。

\paragraph{无限类需要什么替代量。}
连续参数类通常有无限多个函数，不能把参数个数代入$\log|\mathcal H|$。二分类的增长函数$\Pi_{\mathcal H}(n)$计数在$n$个点上最多能实现多少种二元标记；若VC维为$d_{\rm VC}$，则$n\geq d_{\rm VC}\geq1$时有$\Pi_{\mathcal H}(n)\leq(en/d_{\rm VC})^{d_{\rm VC}}$。它衡量可区分的标记，而非参数存储量。

另一个视角是随机符号拟合能力。令损失类$\mathcal G=\{(x,y)\mapsto\ell(f(x),y):f\in\mathcal H\}$，独立符号$\epsilon_i$等概率为$\pm1$，定义
\begin{equation}
 \mathfrak R_n(\mathcal G)=\E_{D,\epsilon}\sup_{g\in\mathcal G}
 \frac1n\sum_i\epsilon_i g(X_i,Y_i).
\end{equation}
这里另记$Z_i=(X_i,Y_i)$为一条样本，而非前面Hoeffding推导中的标量损失。引入独立同分布的复制样本$D'=\{Z_i'\}_{i=1}^n$后，用$R(g)=\E_{D'}\hat R_{D'}(g)$及上确界的凸性，得到
\begin{align}
 \E_D\sup_g(R(g)-\hat R_D(g))
 &\leq\E_{D,D'}\sup_g\frac1n\sum_i(g(Z_i')-g(Z_i))\\
 &\leq2\mathfrak R_n(\mathcal G).
\end{align}
第二步利用样本对可交换性插入随机符号，再将两个上确界分开。这解释了“能拟合纯随机波动”为什么意味着更大的乐观风险。若损失在$[0,1]$内，替换一个样本使上确界至多改变$1/n$，有界差分集中进一步给出以至少$1-\delta$的概率，$\sup_g(R-\hat R)\leq2\mathfrak R_n(\mathcal G)+\sqrt{\log(1/\delta)/(2n)}$。这些结论仍依赖独立抽样和适当可测性，并非对重叠窗口的自动保证。

\section{模型选择、校准与不平衡}
模型选择偏差要求把调参与最终评价分离\cite{cawley2010selection}。深度网络校准研究则表明，分类正确率与概率可信度需要分别检验，温度缩放是一种后处理基线\cite{guo2017calibration}。
\paragraph{模型选择与交叉验证。}
若数据近似独立同分布，$K$折交叉验证将训练集分为$K$份，轮流使用一份验证、其余训练。时间序列必须使用滚动窗口或前向验证，避免未来信息泄漏。超参数搜索也应嵌套在训练数据内部；否则测试集会逐渐变成隐形训练集。

\paragraph{交叉验证估计的是哪个算法。}
设折索引为$I_1,\ldots,I_K$，超参数为$h$，$\hat f_h^{(-k)}$表示只用第$k$折以外的数据完成预处理和训练后的模型。加权交叉验证风险为
\begin{equation}
 \widehat R_{\rm CV}(h)=\frac1n\sum_{k=1}^K\sum_{i\in I_k}
 \ell(\hat f_h^{(-k)}(x_i),y_i).
\end{equation}
每个损失的标签没有参与产生该预测，但不同折的训练集重叠，因此这些损失不是独立试验。选$\hat h=\argmin_h\widehat R_{\rm CV}(h)$后，最小分数有选择偏差：对带噪估计$\hat r_h$，总有$\E\min_h\hat r_h\leq\min_h\E\hat r_h$。搜索越广，越容易挑中有利噪声。

嵌套验证在外层训练集内部完成全部内层选择，再对外层留出样本预测，评估的是“选择加训练”这一完整算法。最后可用所有开发数据重新选择并拟合，但外层分数并非这个最终参数向量的无偏测量。设备切分估计跨设备能力，前向切分估计未来能力；两者对应不同目标，不能由普通随机$K$折替代。

\paragraph{一轮嵌套选择到底留下了什么。}
例如有六台开发设备，外层先留下设备5、6；只在设备1至4内部比较两个窗口长度，每个内层训练折都重新拟合填补和标准化。选好长度后，用设备1至4重新训练，最后才预测设备5、6。外层换组后，内层可能选到不同长度，这正是完整选择算法的表现，而不是执行错误。开发结束后另有独立测试设备时，应冻结最终流程再测试。超参数是训练前需选择的设置（如窗口长度），模型参数则是给定设置后从训练数据估计的数值（如回归系数）。

\paragraph{校准与决策。}
分类器输出的分数不一定是概率。校准要求在所有预测为$p$的样本中，正例频率接近$p$。可靠性图和 Expected Calibration Error 可用于检查校准。若预测要驱动检修、放行或停机决策，还应定义效用函数$U(a,y)$，最终选择
\begin{equation}
 a^*(x)=\argmax_a\mathbb{E}_{p(y\mid x,\mathcal{D})}[U(a,y)].
\end{equation}
这说明“最准确的分类器”不一定产生“最优的行动”。

\paragraph{不平衡与置信区间。}
少数类任务应报告每类召回率、宏平均$F_1$和 Area Under the Precision--Recall Curve (PR-AUC)。重采样、类别权重和阈值调整针对的是不同环节，不能混为一谈。性能估计还需要不确定性，可使用自助法在测试样本上重采样得到指标分布。跨设备评估时，设备级重采样比样本级重采样更能反映真实部署波动。

\paragraph{先读懂混淆矩阵，再读少数类指标。}
把故障记为正类，TP、FP、FN分别为正确告警、错误告警和漏报的数量。召回率为$\mathrm{TP}/(\mathrm{TP}+\mathrm{FN})$，精确率为$\mathrm{TP}/(\mathrm{TP}+\mathrm{FP})$，$F_1=2\mathrm{TP}/(2\mathrm{TP}+\mathrm{FP}+\mathrm{FN})$。例如100台设备中有10台故障，告警12台且其中8台故障，则召回率为$0.8$、精确率为$2/3$、$F_1=16/22\approx0.727$。一直报正常也有$90\%$准确率，却漏掉全部故障。若没有真实正例，召回率分母为零；若没有正类预测，精确率分母为零，此时原比值无定义，应说明报告约定，不能默认为性能完美。宏平均是先分别计算各类指标再等权平均，而不是先合并计数。PR曲线扫描阈值形成精确率与召回率的对应关系；面积还应声明插值或平均精确率等计算约定。

\paragraph{类别权重为何改变概率含义。}
令真实条件正例概率为$p$，预测为$q$，正负类权重为$a,b>0$。固定输入处的加权交叉熵为$L(q)=-ap\log q-b(1-p)\log(1-q)$。求导并令其为零，得到
\begin{equation}
 -\frac{ap}{q}+\frac{b(1-p)}{1-q}=0,
 \qquad q^*=\frac{ap}{ap+b(1-p)}.
\end{equation}
所以加权训练的理想输出一般不再是原始分布的$p$。若模型恰好达到该最优且权重已知，可反解$p=bq/[a(1-q)+bq]$；有限样本和错设模型下还需独立校准。重采样同样可能改变类别先验；它不应与只在预测后调整阈值混为一谈。

例如$p=0.1,a=9,b=1$时，理想加权输出为$q^*=0.5$，并不表示真实故障率升到了50\%；代入反解仍得到$p=0.1$。

精确率也直接受基率影响。若正例比例为$\pi$、真正率为$u$、假正率为$v$，Bayes公式给出
\begin{equation}
 \operatorname{Precision}=\frac{\pi u}{\pi u+(1-\pi)v}.
\end{equation}
即使排序和条件错误率不变，部署正例更稀少时精确率也会下降。跨工况比较精确率--召回率（Precision--Recall，PR）曲线应同时报告基率。

\paragraph{区间与配对比较的抽样单位。}
若独立同分布测试单位上的损失$L_i$具有有限方差，令$\bar L=n^{-1}\sum_iL_i$、$s_L^2=(n-1)^{-1}\sum_i(L_i-\bar L)^2$（$n>1$）。$s_L$描述单个单位的损失波动，均值的标准误估计$s_L/\sqrt n$描述重新抽取测试集时平均损失的波动，两者不能混用。大样本近似区间为$\bar L\pm1.96s_L/\sqrt n$，其中1.96是标准正态分布的97.5\%分位数，左右尾各留2.5\%。例如$n=100,\bar L=4,s_L=2$时，标准误为0.2，区间为$[3.608,4.392]$。这只针对固定模型和既定测试分布，不包括训练随机性；稀有事件或小样本时正态近似可能不良。

比较模型A、B时应先在同一测试单位上计算$d_i=L_{A,i}-L_{B,i}$，再估计$\bar d$的区间，因为
\begin{equation}
 \Var(L_A-L_B)=\Var(L_A)+\Var(L_B)-2\operatorname{Cov}(L_A,L_B).
\end{equation}
忽略配对会丢失共同设备难度的信息。对$G$台近似独立设备，自助法每次有放回抽取$G$个设备编号，保留所抽设备的全部窗口并重算指标；重复$B$次后取经验分位数作为近似区间。每台等权与每窗口等权是不同的目标，必须在每次重采样中维持原定义。设备数很少、设备之间相关或许多重采样没有正例时，区间本身也不可靠；增加重采样次数不会增加独立设备信息。

\paragraph{置信区间的概率究竟修饰什么。}
设固定模型的总体平均损失为$\mu$，由随机测试数据$D$算出的区间为$[L(D),U(D)]$。95\%置信区间的目标是$P_D\{L(D)\leq\mu\leq U(D)\}\approx0.95$：重复采集测试集并构造区间，约95\%的区间覆盖固定的$\mu$。一次数据已经到手后，不能仅凭这一定义把它读成“$\mu$有95\%概率位于当前两端点之间”。后续贝叶斯章节的可信区间则在给定数据和先验后，对未知参数的后验概率作陈述。预测区间又针对一个未来观测，不是总体均值；三者的对象和随机性不同。

自助法不是制造新设备。例如原有设备编号为$(A,B,C)$，一次有放回抽样可能为$(C,C,A)$，必须把$C$的全部记录按两次抽中计入，而非去重；另一次可能为$(B,A,B)$。每次重算相同指标，所得波动近似原实验的抽样波动。只重采样测试设备、不重新训练时，估计的是固定模型的测试不确定性，不包含重训和选模波动。三个设备仅用于演示操作，不足以支撑可靠的95\%区间。

\section{从统计结论到实验设置}
\paragraph{数据泄漏清单。}
常见泄漏包括：用全数据计算标准化参数，把同一设备的相邻窗口分到训练和测试，利用故障发生后的变量预测故障发生前的状态，以及根据测试集反复选择模型。泄漏往往带来异常漂亮的结果，因此每个实验都应记录切分单位、时间边界、预处理拟合范围和标签生成时刻。

\paragraph{从固定函数到数据依赖的模型选择。}
Hoeffding 界直接控制固定函数的经验风险偏差，但实际算法会根据数据选择函数。联合界的作用是同时控制整个假设空间，使“先看数据再选模型”仍有可分析的上界。代价是模型空间越复杂，需要更多样本或更强的结构约束。

\paragraph{容量不是参数数量的同义词。}
参数数量相同，网络能表达的函数也未必相同。激活函数、连接方式和参数共享决定函数类，优化路径则影响训练最后选到其中哪些函数。例如卷积在不同位置使用同一组权重，注意力允许相距较远的输入直接交互。分析泛化时，因此还要看函数类、参数范数、分类间隔和算法稳定性，不能仅按参数个数给模型排队。

\paragraph{交叉验证的工程细节。}
预处理、特征选择和超参数搜索都必须在每个训练折内部拟合。若先对全体数据做标准化或选择相关特征，再执行交叉验证，验证折的信息已经泄漏。时间序列使用前向验证时，验证窗口必须位于训练窗口之后，必要时在二者之间设置间隔以避免标签窗口重叠。

\paragraph{不平衡、阈值与区间估计。}
ROC-AUC 衡量排序能力，PR-AUC 更关注少数正类下的精确率和召回率。阈值改变会改变混淆矩阵，却不改变排序本身。若最终动作依赖阈值，应报告阈值扫描、成本曲线和置信区间。自助法重采样时，设备级任务应按设备重采样；否则同一设备的多个窗口会让区间过窄。

\paragraph{校准的后处理。}
温度缩放在校准集上用一个正标量调整 logits（归一化前的类别得分），保持同一样本内部的类别次序，但不保证改善独立测试集的校准；Platt scaling 用逻辑函数映射分数，isotonic regression（保序回归）拟合非递减映射，二者对数据量和函数形状的要求不同。校准模型不能使用测试集反复调参，且分布变化后需要重新检查。对安全告警而言，概率校准和拒识机制常比少数百分点的准确率更重要。

\paragraph{温度缩放的目标与边界。}
先冻结基模型及其预处理，再在校准集上保存logits并只拟合温度。部署时沿用相同变换；若用校准标签继续更新基模型，这批数据便不再是该模型的独立校准集，需要重新安排划分。
对$C$类logits $z_i\in\R^C$，在独立校准集上令$T>0$，定义$q_{ic}(T)=\exp(z_{ic}/T)/\sum_j\exp(z_{ij}/T)$，最小化
\begin{equation}
 L(T)=\sum_i\left[-z_{i,y_i}/T+\log\sum_c\exp(z_{ic}/T)\right].
\end{equation}
取逆温度$\beta=1/T$，则
\begin{align}
 L'(\beta)&=\sum_i\left(\sum_cq_{ic}z_{ic}-z_{i,y_i}\right),\\
 L''(\beta)&=\sum_i\Var_{c\sim q_i}(z_{ic})\geq0.
\end{align}
求导时使用$\partial_\beta q_{ic}=q_{ic}(z_{ic}-\sum_jq_{ij}z_{ij})$，故二阶导数正是得分的加权方差。因此这是一个一维凸优化问题（约束$\beta>0$）；下确界也可能只在$\beta\to0$或$\beta\to\infty$时逼近，不能保证存在有限最优温度。正温度保持每个样本内部的类别次序和最大类别；多分类中不同样本的最大概率排序则未必保持，不能据此保证所有排序指标不变。它只能整体调整尖锐程度，无法修复所有类别和工况上的系统偏差。

二分类可靠性图把样本分入$B_m$，计算箱内预测均值$\bar q_m$和事件比例$\bar y_m$；常用误差为$\sum_m(|B_m|/N)|\bar q_m-\bar y_m|$。这是依赖分箱的经验统计量，不是证明条件校准的定理。箱数太少会掩盖局部偏差，太多则产生高方差，应同时报告样本量与可靠性图。

\paragraph{自测：样本数怎样改变保证。}
固定候选集合和失败概率，若独立样本数从$n$增至$4n$，统一界半径$b_n$如何变化？把每条旧样本复制四次是否等价？

\textit{参考解。}式\eqref{eq:statistical-uniform-radius}给出$b_{4n}=b_n/2$；经验风险最小化相对类内最优的$2b_n$项也减半。复制不产生新的独立观测，不能把复制后的行数代入这一保证。

\section{从统计量到决策}
\paragraph{效用与损失只差一个决策方向。}
效用$U(a,y)$越大越好，损失$C(a,y)$越小越好，取$U=-C$便得到同一动作。若不检修时故障损失为100、无故障损失为零；检修固定花费8且在这个简化例子中消除该次故障损失，则故障概率为$0.1$时不检修的期望损失为10，检修为8，故选检修。这里必须写明检修效果；若只检测而不消除风险，需要把后续动作和残余损失也列入成本，不能直接沿用阈值。

模型输出只是决策系统的一个中间量。若动作集合为$\mathcal A$，效用函数为$U(a,y)$，最优动作是期望效用最大的动作。不同部门可能面对不同成本，因此同一个概率模型可以配合不同阈值，而不必为每种成本重新训练。这个分离有助于审计模型能力和业务策略各自的变化。

统计学习中的复杂度界与实验估计可结合官方交叉验证文档和重采样教程复习，注意它们服务于不同的不确定性来源。
\section*{辅助学习材料}
\begingroup
\emergencystretch=3em
\begin{itemize}
\item \textbf{Cross-validation: evaluating estimator performance}；作者/平台：scikit-learn Developers；英文；官方文档。阅读 K-fold、group、stratified 与 nested selection 的说明，帮助落实本章的模型选择和泛化估计。\href{https://scikit-learn.org/stable/modules/cross\_validation.html}{在线阅读}。
\item \textbf{Prediction Intervals for Gradient Boosting Regression}；作者/平台：scikit-learn Developers；英文；官方示例。阅读分位数回归、区间覆盖率与调参部分，区分预测区间和参数置信区间；该示例不提供分布无关的覆盖保证。\href{https://scikit-learn.org/stable/auto\_examples/ensemble/plot\_gradient\_boosting\_quantile.html}{在线阅读}。
\item \textbf{【机器学习】什么是交叉验证？为什么以及如何进行交叉验证？}；知乎；中文解说。区分模型选择与最终泛化估计；标题与主题据必应索引，原页受限。\href{https://zhuanlan.zhihu.com/p/81738187251}{在线阅读}。
\item \textbf{交叉验证方法汇总【附代码】（索引标题节选）}；CSDN；中文博客。比较留一法、K折与分层划分，另行检查设备分组要求；标题与主题据必应索引，原页受限，完整标题未核实。\href{https://blog.csdn.net/WHYbeHERE/article/details/108192957}{在线阅读}。
\item \textbf{2022 - 再探宝可梦、数码宝贝分类器 — 浅谈机器学习原理}；李宏毅；bilibili；中文课程。选看第23集，结合本章假设空间与泛化条件讨论分类。2026年10月4日官方公开元数据显示整套课程累计播放约299.5万，并非本分集播放量；已核对目录与计数，未逐帧观看。\href{https://www.bilibili.com/video/BV1Wv411h7kN/?p=23}{观看分集}。
\end{itemize}
\endgroup
\section{本章小结}
统计学习提供的是关于数据、模型复杂度和未见样本风险的推理框架。它不能保证某个模型在所有工业分布上有效，却能迫使实验者明确假设空间、验证规则、泄漏风险、不确定性和决策成本。

统一界考虑看过数据以后再选模型的影响，嵌套验证把挑选模型和报告成绩分开，配对比较与设备级区间估计则显示结果会怎样波动。比较平均误差前，要先核对任务、抽样单位和尝试方案的预算是否一致；选择告警阈值时，还要看校准和成本。下一章回到原始记录：如果不知道记录何时可用、标签何时截止、样本来自哪台设备，再严谨的统计计算也无法回答原先的维护问题。


\part{贝叶斯机器学习：用概率表达知识与不确定性}
前一部分已经规定了输入、标签、时间边界和验收动作，但点预测仍没有说明“这次判断有多可靠”。本部分保持贯穿案例不变，只改变答案的表达方式：先用贝叶斯更新描述有限证据，再比较参数化与非参数化模型如何表达函数不确定性。读者应看到，概率不是给结果加一个置信度装饰，而是改变训练目标、预测输出和维护决策的依据；下一部分将在这些边界上选择更强的表示模型。
\chapter{贝叶斯机器学习}
\begin{quote}\small
\textit{本章摘要。}本章将点预测扩展为对未知量及其依赖关系的分布表达。先建立贝叶斯更新和后验预测，再用概率图模型描述设备隐状态、观测和相邻标签之间的关系。Hidden Markov Model (HMM)、最大熵马尔科夫模型、Conditional Random Field (CRF) 和马尔科夫随机场分别揭示生成建模、条件建模与局部依赖的不同取舍。在此基础上讨论参数后验的近似计算，最后以校准、敏感性分析和维护决策检验概率是否可信。
\end{quote}
前一部分已经固定数据可用时间、预测目标与验证规则。现在假设这些前提都成立，仍然存在一个未解决的问题：少量故障能支持多精确的风险估计，多个解释相近的参数应如何共同参与预测？本章不改变任务，而是把单一答案扩展成对可能答案及其依赖关系的分布。

阅读时可以把未知量分成三类：设备长期的故障倾向属于参数，轴承此刻是否磨损属于隐状态，下一次传感器读数还包含新的随机波动。对状态求和和对参数积分处理的是不同的未知量，决定是否停机还要另外比较动作成本。先从能手算的共轭模型开始，再看链式图模型怎样减少枚举，最后讨论近似计算的误差。使用贝叶斯方法以后，仍要检查假设是否合适、概率是否校准，不能省去验证。

第一次阅读应先掌握条件概率、有限求和与期望。可先读本节的两点先验例子，再读“从共轭模型到近似推断”中的Beta更新，建立参数更新与预测的区别；随后回到图模型，学习怎样对多个隐状态求和。矩阵配方和近似梯度是计算层的展开，不必与全部图模型名称同时记忆。

\section{贝叶斯更新与后验预测}

% auxiliary-resource-guide
本章末的“辅助学习材料”列出与核心方法对应的讲解和实践入口，可在阅读相关推导后，按所列小节对照学习。

\paragraph{从参数点到参数分布。}
贝叶斯建模把未知参数视为随机变量，是概率机器学习的基本表述之一\cite{murphy2022pml,bishop2006pattern}。给定先验$p(\bm{\theta})$和似然$p(\mathcal{D}\mid\bm{\theta})$，后验为
\begin{equation}
 p(\bm{\theta}\mid\mathcal{D})=\frac{p(\mathcal{D}\mid\bm{\theta})p(\bm{\theta})}{p(\mathcal{D})},
 \qquad p(\mathcal{D})=\int p(\mathcal{D}\mid\bm{\theta})p(\bm{\theta})\,d\bm{\theta}.
\end{equation}
先验表达观测数据之前的结构知识，似然描述数据如何由参数生成，后验则是二者结合后的知识状态。

\paragraph{先验、似然和后验不是三个同义概率。}
对事件$A,B$且$P(B)>0$，条件概率$P(A\mid B)=P(A\cap B)/P(B)$意为只在$B$已发生的情形中计算$A$的比例，通常不等于$P(B\mid A)$。设设备有两种可能的故障率$\theta\in\{0.1,0.5\}$，观测前分别赋予概率$0.8,0.2$。一次试验观察到故障$Y=1$后，两种解释的未归一化权重为$0.8\times0.1=0.08$和$0.2\times0.5=0.10$；证据为$P(Y=1)=0.18$，后验概率分别为$4/9$和$5/9$。数据把较高故障率的权重提高了，却没有把另一种可能性删除。

在这个例子中，先验与后验都是对$\theta$归一化；似然则固定已发生的$Y=1$，把$P(Y=1\mid\theta)$当作$\theta$的函数。两个似然值之和$0.1+0.5$不需要等于一。连续参数时，把求和改成积分，$p(\theta)$表示密度而非单点概率；参数区间的概率是密度曲线下的面积，密度值本身可以大于一。符号$\propto$表示两边只差一个不依赖正在推断的变量的正比例常数。

若下一次试验给定$\theta$后与本次独立，其故障概率为$(4/9)\times0.1+(5/9)\times0.5=29/90$。这是保留两种解释后加权预测，不是先宣布$\theta=0.5$再报概率$0.5$。贝叶斯分析中的随机参数表达对未知参数的知识，不要求设备的真实故障率在每次抽样时重新变化。

\paragraph{后验预测。}
\begin{figure}[htbp]
\centering
\begin{tikzpicture}[node distance=0.8cm and 0.75cm]
 \node[idi input, minimum width=2.1cm] (prior) {先验知识\\$p(\bm\theta)$};
 \node[idi input, right=of prior, minimum width=2.1cm] (data) {观测数据\\$\mathcal D$};
 \node[idi process, below=0.85cm of $(prior)!0.5!(data)$, minimum width=2.6cm] (like) {似然\\$p(\mathcal D\mid\bm\theta)$};
 \node[idi state, below=0.85cm of like, minimum width=2.6cm] (post) {后验\\$p(\bm\theta\mid\mathcal D)$};
 \node[idi output, below=0.85cm of post, minimum width=2.6cm] (pred) {后验预测\\$p(y_*\mid x_*,\mathcal D)$};
 \draw[idi arrow] (prior)--(like); \draw[idi arrow] (data)--(like);
 \draw[idi arrow] (like)--(post); \draw[idi arrow] (post)--(pred);
\end{tikzpicture}
\caption{贝叶斯更新把先验与观测结合为后验，再对参数不确定性积分得到预测分布。}
\label{fig:bayes-predictive-flow}
\end{figure}
图\ref{fig:bayes-predictive-flow}区分后验更新与预测两个环节：先用数据修正参数分布，再以整个后验而非单个参数点形成新观测的预测。
对新输入$x_*$，预测分布需要对参数不确定性积分：
\begin{equation}
 p(y_*\mid x_*,\mathcal{D})=\int p(y_*\mid x_*,\bm{\theta})p(\bm{\theta}\mid\mathcal{D})\,d\bm{\theta}.\label{eq:bayes-predictive-integral}
\end{equation}
式\eqref{eq:bayes-predictive-integral}中积分涉及的不确定性包含数据噪声和参数不确定性。前者称为偶然不确定性（aleatoric uncertainty），后者称为认知不确定性（epistemic uncertainty）。当设备出现训练数据未覆盖的工况时，后者尤其值得关注。

\paragraph{更新与预测各用了哪条概率规则。}
联合分布可分别写成$p(\theta,\mathcal D)=p(\mathcal D\mid\theta)p(\theta)$和$p(\theta\mid\mathcal D)p(\mathcal D)$；两式相等并除以正且有限的$p(\mathcal D)$即得Bayes公式。连续观测使用的是密度，不要求单个精确观测有正概率。若数据分批到达且给定$\theta$后条件独立，则
\begin{equation}
 p(\theta\mid\mathcal D_1,\mathcal D_2)
 \propto p(\mathcal D_2\mid\theta)p(\theta\mid\mathcal D_1).
\end{equation}
同一条记录只能更新一次；相关窗口若被误当独立批次，会使后验虚假收窄。

后验预测来自全概率公式，再使用$Y_*\perp\mathcal D\mid(x_*,\theta)$。固定设计回归把输入视作已知条件；如果$x_*$的生成机制本身依赖待推断参数，则还应对它的证据更新，不能无条件沿用$p(\theta\mid\mathcal D)$。这一细节说明概率公式总是在特定条件结构下成立。

\section{概率图模型：从依赖关系到分布分解}
单个故障率只能回答总体上有多大风险，却不能解释设备如何从正常运行进入退化状态。振动和温度往往共同受某个不可直接观测的状态影响，相邻时刻的状态又存在持续性。概率图模型以节点表示随机变量，以边和局部因子表达依赖结构，将一个庞大的联合分布拆成可以解释和计算的部分\cite{bishop2006pattern,murphy2022pml}。

有向无环图对应贝叶斯网络。若变量集合为$X_1,\ldots,X_m$，节点$X_j$的父节点集合记为$\operatorname{pa}(j)$，则
\begin{equation}
 p(X_1,\ldots,X_m)=\prod_{j=1}^{m}p(X_j\mid X_{\operatorname{pa}(j)}).
\end{equation}
没有父节点的变量使用边缘分布。这个分解规定了条件独立假设，而不是仅用箭头描绘相关性。例如，在$X\leftarrow Z\rightarrow Y$中，给定共同原因$Z$后，模型假定$X$与$Y$独立。若$Z$表示设备工况，该假设意味着知道工况后，两个传感器不再提供额外的相互依赖信息。真实设备仍可能存在共享噪声，此时需要修改观测模型。

条件化也可能引入依赖。在碰撞结构$X\rightarrow Z\leftarrow Y$中，路径原本被$Z$阻断，而观测$Z$或其后代可能使$X$与$Y$相关。因而不能简单地认为“给定更多变量总会减少依赖”。一般有向图中的这类关系由 d-separation 判断。图中的方向也不自动具有因果含义，因果解释还需要额外的领域假设或干预证据。

无向图适合表达对称的相容关系，例如相邻表面区域倾向于具有相近状态。因子图则把变量节点与因子节点分开，统一表示$p(X)\propto\prod_a\psi_a(X_a)$。它强调计算所依赖的局部分解，可以表示有向模型，也可以表示无向模型。树上的求和积消息传递能够精确计算边缘概率，最大积用于求最优联合配置；含环图的复杂度则取决于连接结构，不能只按节点数判断。

这里需要分清两类未知量。设备在某一时刻的隐状态属于随机变量推断，转移矩阵和观测分布的参数属于参数学习。概率图模型可以用最大似然估计参数，也可以为参数指定先验并计算后验。因此，“贝叶斯网络”描述图的分解形式，“贝叶斯学习”描述对参数不确定性的处理方式，两者不是同义词。

\paragraph{用条件表理解独立，再读图。}
条件独立$X\perp Y\mid Z$意味着对每个有正概率的$z$，$p(x,y\mid z)=p(x\mid z)p(y\mid z)$。例如正常工况下两个传感器各以$0.1$概率告警，高载荷下各以$0.8$概率告警，并假设给定工况后独立；混合工况时它们仍可能一起告警，不能把条件独立删去条件后使用。

碰撞结构可用两个独立二元原因说明。令$X,Y$各以$1/2$概率为一，$Z=1$表示至少一个原因为一。只知$Z=1$时，$(X,Y)$的三种可能为$(1,0),(0,1),(1,1)$且等概率，故$P(Y=1\mid Z=1)=2/3$；再知$X=0$后，$Y=1$的概率变为一。独立原因在选择了结果以后产生依赖。d-separation的链式读法是：链和共同原因路径在中间节点被条件化后阻断，碰撞路径则相反；一般图须检查所有路径，不能只看一条边。

消息公式中的$N(i)$是变量$i$相邻的因子集合，$N(a)$是因子$a$涉及的变量集合；$a\setminus i$表示该因子中除$i$外的变量。消息是尚未归一化的局部证据汇总，不是新观测，也不一定是概率。

\paragraph{分解与消息的代数来源。}
按有向无环图的拓扑顺序排列变量，概率链式法则先给出$\prod_jp(X_j\mid X_1,\ldots,X_{j-1})$；只有加入“给定父节点后独立于其余先前变量”的假设，才能删去多余条件得到图分解。边因此代表被允许的依赖，不是从链式法则自动推出来的物理事实。

对树形因子图，边缘化要计算$\sum_{X\setminus X_i}\prod_a\psi_a(X_a)$。利用乘法对加法的分配律，把树的一个分支先求和，就得到消息
\begin{align}
 m_{i\to a}(x_i)&=\prod_{b\in N(i)\setminus a}m_{b\to i}(x_i),\\
 m_{a\to i}(x_i)&=\sum_{x_{a\setminus i}}\psi_a(x_a)
                  \prod_{j\in N(a)\setminus i}m_{j\to a}(x_j).
\end{align}
从叶子向内再向外传递后，$p(x_i)\propto\prod_{a\in N(i)}m_{a\to i}(x_i)$。有限离散状态下这些都是有限和；连续变量则需积分或近似。树上各分支只相交于分隔变量，所以每条证据恰好计算一次；含环时重复传递不再对应这样的独立分支消元。把求和换成最大值并保存取值索引可求最优配置，但不能由最大消息恢复边缘概率。

\section{隐马尔科夫模型与状态推断}
HMM 的前向后向、最优路径与参数估计可参照 Rabiner 的经典教程\cite{rabiner1989hmm}。下面将这些运算分别对应到在线状态识别、离线诊断和历史数据训练，避免混淆可用信息。
设$x_{1:T}$为一段设备观测，$z_t\in\{1,\ldots,K\}$为时刻$t$的隐状态。隐马尔科夫模型（HMM）假定状态只依赖上一时刻，且观测在给定当前状态后不再依赖其他状态和观测：
\begin{align}
 p(z_t\mid z_{1:t-1})&=p(z_t\mid z_{t-1}),\\
 p(x_t\mid z_{1:T},x_{1:t-1})&=p(x_t\mid z_t).
\end{align}
记初始概率为$\pi_k$，转移概率为$A_{jk}=p(z_t=k\mid z_{t-1}=j)$，观测密度或概率为$b_k(x)=p(x\mid z_t=k)$。其中$\pi\in\R^K$、$A\in\R^{K\times K}$，行索引$j$是出发状态，列索引$k$是到达状态；离散观测的$b_k$为概率，连续观测时为密度。联合分布分解为
\begin{equation}
 p(x_{1:T},z_{1:T})=\pi_{z_1}b_{z_1}(x_1)
 \prod_{t=2}^{T}A_{z_{t-1},z_t}b_{z_t}(x_t).\label{eq:bayes-hmm-joint}
\end{equation}
\begin{figure}[htbp]
\centering
\begin{tikzpicture}[>=Latex,
 state/.style={circle,draw,minimum size=10mm,fill=white},
 obs/.style={circle,draw,minimum size=10mm,fill=blue!12},
 every node/.style={font=\small}]
 \foreach \t in {1,2,3,4}{
  \node[state] (z\t) at (2.4*\t,0) {$z_{\t}$};
  \node[obs] (x\t) at (2.4*\t,-1.6) {$x_{\t}$};
  \draw[->,thick] (z\t) -- (x\t);
 }
 \foreach \a/\b in {1/2,2/3,3/4}{\draw[->,thick] (z\a) -- (z\b);}
 \node[above=3mm of z1] {$\pi$};
 \node at (6,0.65) {状态转移：$p(z_t\mid z_{t-1})$};
 \node at (6,-2.45) {观测生成：$p(x_t\mid z_t)$};
\end{tikzpicture}
\caption{HMM 的有向图结构（以四个时刻为例）。白色圆为隐状态，蓝色圆为已观测变量；横向箭头表示状态转移，向下箭头表示观测生成。箭头表示概率分解关系，不自动代表因果关系。}
\label{fig:bayes-hmm-structure}
\end{figure}
图\ref{fig:bayes-hmm-structure}可以直接读成式\eqref{eq:bayes-hmm-joint}的联合分布：每个节点贡献一个给定父节点的条件分布，首个状态贡献初始分布。给定$z_t$后，$x_t$与其他时刻的变量条件独立；但只给定观测而未固定隐状态时，各时刻观测通常仍然相关。图中的向下箭头是模型定义的生成方向，推断时则可以从观测反过来更新隐状态的概率。

每行转移概率之和为一。多通道观测可以用联合高斯分布或混合分布描述，HMM 并不要求同一时刻的各传感器通道独立。真正受限制的是跨时刻的依赖：一旦给定状态序列，观测之间便被假定条件独立。

以“正常、退化、故障”三个状态为例，转移矩阵可以表达状态持续性。禁止某条转移必须有物理依据，例如只有在明确排除维修过程的运行片段内，才可能限制故障状态不能直接返回正常。无监督学习得到的三个隐状态也不会自动具有这三个语义，需要通过标注、特征分布或专家检查建立对应关系。

\paragraph{过滤、平滑与最优路径。}
在线维护需要$p(z_t\mid x_{1:t})$，称为过滤；事后分析使用$p(z_t\mid x_{1:T})$，称为平滑。后者利用未来观测，不能直接作为实时告警的性能。若需要解释整段运行过程，还可以寻找联合概率最大的状态路径。这三个问题的输出不同，即使使用同一个 HMM，也不能混用结果。

直接枚举所有状态路径需要考虑$K^T$种配置。前向算法利用链结构复用公共计算，令$\alpha_t(k)=p(x_{1:t},z_t=k)$，则
\begin{align}
 \alpha_1(k)&=\pi_kb_k(x_1),\\
 \alpha_t(k)&=b_k(x_t)\sum_{j=1}^{K}\alpha_{t-1}(j)A_{jk},\\
 p(z_t=k\mid x_{1:t})&=\frac{\alpha_t(k)}{\sum_j\alpha_t(j)}.
\end{align}
最后一时刻的前向量求和即为$p(x_{1:T})$。每步对$K^2$个状态对计算，总复杂度为$O(TK^2)$。长序列中概率连乘会下溢，应逐步缩放或在对数域使用 log-sum-exp，并保留缩放因子以恢复序列似然。

后向量$\beta_t(j)=p(x_{t+1:T}\mid z_t=j)$满足
\begin{equation}
 \beta_T(j)=1,\qquad
 \beta_t(j)=\sum_{k=1}^{K}A_{jk}b_k(x_{t+1})\beta_{t+1}(k).
\end{equation}
于是平滑概率为$\gamma_t(k)=\alpha_t(k)\beta_t(k)/p(x_{1:T})$。前向量汇总过去证据，后向量汇总未来证据，这也给出了链式消息传递的具体形式。

Viterbi 算法将前向递推中的求和替换为最大值。在对数域中，
\begin{equation}
 \delta_t(k)=\log b_k(x_t)+\max_j\{\delta_{t-1}(j)+\log A_{jk}\},
 \qquad \delta_1(k)=\log\pi_k+\log b_k(x_1).
\end{equation}
保存每步最大值对应的前驱，便可从终点回溯路径。逐时刻选择边缘概率最大的状态与 Viterbi 路径通常不同，前者优化逐点分类，后者优化整条路径的联合概率。在存在禁止转移时，逐点选择甚至可能产生不合法路径。

\paragraph{两状态、两时刻的过滤手算。}
设状态为正常$N$、退化$D$，$\pi=(0.8,0.2)$，按$N,D$排列的转移矩阵为
\begin{equation}
 A=\begin{bmatrix}0.9&0.1\\0.2&0.8\end{bmatrix},\qquad
 P(\text{高振动}\mid N)=0.1,\quad P(\text{高振动}\mid D)=0.7.
\end{equation}
首次看到高振动时，$\alpha_1=(0.08,0.14)$，所以退化过滤概率为$0.14/0.22=7/11$。第二次仍为高振动，先转移再乘观测概率，得到$\alpha_2(N)=(0.08\times0.9+0.14\times0.2)\times0.1=0.0100$，$\alpha_2(D)=(0.08\times0.1+0.14\times0.8)\times0.7=0.0840$。新的退化概率为$0.084/0.094\approx0.894$。

若事后回看第一个时刻，$\beta_1(N)=0.9\times0.1+0.1\times0.7=0.16$，$\beta_1(D)=0.2\times0.1+0.8\times0.7=0.58$，所以退化平滑概率为$0.14\times0.58/0.094\approx0.864$。它不同于当时可算的$7/11$，因为第二次高振动已经提供了未来证据。此处$\alpha_t,\beta_t,\gamma_t$各有$K$个分量，$A$为$K\times K$矩阵，$\xi_t$也为$K\times K$相邻状态概率表。

\paragraph{前向后向递推的推导与稳定实现。}
固定当前状态$k$，对前驱$j$边缘化：
\begin{align}
 p(x_{1:t},z_t=k)
 &=\sum_jp(x_t\mid z_t=k)p(z_t=k\mid z_{t-1}=j)
             p(x_{1:t-1},z_{t-1}=j)\\
 &=b_k(x_t)\sum_jA_{jk}\alpha_{t-1}(j).
\end{align}
条件独立允许把观测因子提到求和外。相似地，从$z_{t+1}$边缘化得到后向递推；给定$z_t$后过去和未来的观测分离，因此$\alpha_t(k)\beta_t(k)$就是$p(x_{1:T},z_t=k)$。求和表示保留所有可能路径，Viterbi的最大值则丢弃每个终点下非最优前缀；链结构保证最优前缀可扩展为最优整段路径。

数值上令$h_{t-1}(j)=p(z_{t-1}=j\mid x_{1:t-1})$，计算
\begin{equation}
 u_t(k)=b_k(x_t)\sum_jh_{t-1}(j)A_{jk},\quad
 c_t=\sum_ku_t(k),\quad h_t(k)=u_t(k)/c_t.
\end{equation}
第一步用$u_1(k)=\pi_kb_k(x_1)$。于是$c_t=p(x_t\mid x_{1:t-1})$且$\log p(x_{1:T})=\sum_t\log c_t$；缩放后的后向递推对应除以$c_{t+1}$。精确运算下$c_t=0$表示当前证据不在模型支持内；浮点计算中的零也可能来自观测密度或乘积下溢，应先用对数域区分二者，再检查零转移与观测支持。对数域的$\log\sum_j e^{a_j}=m+\log\sum_je^{a_j-m}$（$m=\max_ja_j$有限）提供稳定计算；若所有$a_j=-\infty$，结果直接取$-\infty$，不能计算$-\infty-(-\infty)$。

\paragraph{从隐状态推断到参数学习。}
隐状态未标注时，Baum--Welch 算法以 Expectation Maximization (EM) 交替估计状态的后验占用和模型参数。E 步计算$\gamma_t(j)$以及相邻状态的后验概率
\begin{equation}
 \xi_t(j,k)=\frac{\alpha_t(j)A_{jk}b_k(x_{t+1})\beta_{t+1}(k)}{p(x_{1:T})}.
\end{equation}
对一条序列，M 步的转移更新为
\begin{equation}
 A_{jk}^{\mathrm{new}}=
 \frac{\sum_{t=1}^{T-1}\xi_t(j,k)}{\sum_{t=1}^{T-1}\gamma_t(j)}.
\end{equation}
初始分布更新为$\pi_k^{\mathrm{new}}=\gamma_1(k)$，观测分布则按$\gamma_t(k)$进行加权估计。多条独立运行序列需要分别推断后再汇总计数，不能把不同设备的尾部与头部连接成一次状态转移。EM 依赖初始化，只保证在精确 E、M 步下不降低训练似然，不能保证找到全局最优参数。

\paragraph{EM为什么得到期望计数比。}
取$q(z)=p(z\mid x,\theta^{\rm old})$，由Jensen不等式，
\begin{equation}
 \log p_\theta(x)\geq\E_q\log p_\theta(x,z)-\E_q\log q(z).
\end{equation}
在旧参数处等号成立；M步保持$q$固定并提高右端，便不降低观测似然。转移相关部分是$\sum_{jk}N_{jk}\log A_{jk}$，其中$N_{jk}=\sum_{t<T}\xi_t(j,k)$。对每行加入约束乘子$\lambda_j(\sum_kA_{jk}-1)$，一阶条件$N_{jk}/A_{jk}+\lambda_j=0$与行和为一给出$A_{jk}=N_{jk}/\sum_kN_{jk}$。又因$\sum_k\xi_t(j,k)=\gamma_t(j)$，便得到前述更新。

若$x_t\in\R^d$且$b_k=\mathcal N(\mu_k,\Sigma_k)$，高斯对数密度的加权一阶条件给出
\begin{equation}
 \mu_k^{\rm new}=\frac{\sum_t\gamma_t(k)x_t}{N_k},\qquad
 \Sigma_k^{\rm new}=\frac{\sum_t\gamma_t(k)(x_t-\mu_k^{\rm new})(x_t-\mu_k^{\rm new})^{\mathsf T}}{N_k},
 \quad N_k=\sum_t\gamma_t(k).
\end{equation}
这里$\mu_k\in\R^d$、$\Sigma_k\in\R^{d\times d}$。若状态从未被占用，分母为零，参数不可由本轮数据识别；协方差也可能退化，需预先规定先验、协方差结构或数值下界，而非把无定义更新当作成功训练。约束或近似M步需要实际检查目标是否改善。

若状态转移已经观测到，可以为$A_{j,:}$设置 Dirichlet 先验，其后验参数等于先验参数加转移计数。状态未知时，转移计数本身也不确定。EM 的期望计数可用于点估计，但将其代入 Dirichlet 分布并不自动得到精确的参数后验，完整贝叶斯学习需要联合处理状态和参数的不确定性。

标准齐次 HMM 的自转移还隐含持续时间假设。若$A_{kk}<1$，进入状态$k$后停留$d$步的概率为
\begin{equation}
 p(D=d\mid k)=A_{kk}^{d-1}(1-A_{kk}),\qquad d=1,2,\ldots.
\end{equation}
这是一种无记忆的几何分布。若退化持续时间明显集中在某个范围，增加状态数未必是合适修复，隐半马尔科夫模型可以显式描述持续时间。动态贝叶斯网络则允许每个时刻包含多个相互作用的隐变量。线性高斯状态空间模型把离散状态换成连续状态，在相应假设下可用卡尔曼滤波精确计算过滤分布，后续序列章将讨论其动态含义。

\section{从局部条件模型到条件随机场}
最大熵马尔科夫模型与条件随机场分别代表局部条件归一化和全局路径归一化\cite{mccallum2000memm,lafferty2001crf}。比较二者的关键是概率质量如何在状态转移间分配，而不是是否使用神经网络提取特征。
HMM 同时学习观测如何产生和状态如何演化。当目标仅是识别状态，而输入包含频谱、载荷变化和多个重叠窗口特征时，准确建立$p(x_t\mid z_t)$可能比分类本身更困难。最大熵马尔科夫模型（Maximum Entropy Markov Model, MEMM）直接对状态转移建模，以对数线性分类器给出
\begin{equation}
 p_\theta(z_t=k\mid z_{t-1}=j,x)=
 \frac{\exp\{\theta^{\mathsf T}f(j,k,x,t)\}}
 {\sum_{k'}\exp\{\theta^{\mathsf T}f(j,k',x,t)\}}.
\end{equation}
其中$f$为特征向量，$x$表示当前任务允许使用的观测范围。最大熵的名称来自特征期望约束下的条件熵最大化，其解具有上述指数族形式。令$z_0$为固定起始状态，序列条件概率为各步条件概率之积。对于完整标注序列，训练可转化为以真实前驱状态为条件的局部分类；预测时前驱状态未知，仍需进行链上推断。

\begin{figure}[htbp]
\centering
\begin{tikzpicture}[>=Latex,
 state/.style={circle,draw,minimum size=10mm},
 every node/.style={font=\small}]
 \node[draw,rounded corners,fill=blue!12,minimum width=9cm,minimum height=9mm] (x) at (5,1.8) {已给定的可用观测 $x$（不为其建立生成分布）};
 \node[state,fill=gray!15] (z0) at (0,0) {$z_0$};
 \foreach \t in {1,2,3,4}{\node[state] (z\t) at (2.4*\t,0) {$z_{\t}$};}
 \foreach \a/\b in {0/1,1/2,2/3,3/4}{\draw[->,thick] (z\a) -- (z\b);}
 \foreach \t in {1,2,3,4}{\draw[->,blue!65!black] (x.south -| z\t.north) -- (z\t.north);}
 \node[align=center] at (4.8,-1.15) {每一步固定前驱 $z_{t-1}$，只对后继 $z_t$ 归一化\\$\displaystyle \sum_k p_\theta(z_t=k\mid z_{t-1},x)=1$};
\end{tikzpicture}
\caption{MEMM 的条件有向结构。灰色节点为固定起始状态，观测同时参与各步局部分类；与 HMM 相比，模型不再使用由状态指向观测的生成边。}
\label{fig:bayes-memm-structure}
\end{figure}
图\ref{fig:bayes-memm-structure}中，每个待预测标签的输入来自前驱标签与已给定观测，因此整条路径概率是这些局部条件概率的乘积。顶部的$x$可以包含任务允许的多个时刻和重叠特征，并不意味着在线任务能够读取未来。横向箭头保留了标签顺序依赖，却没有取消每一步各自归一化的限制；这正是后文标签偏置的结构来源。

\paragraph{最大熵为何导出softmax。}
把前驱与可用输入合称上下文$c$，令经验上下文权重为$\hat p(c)$，特征$\phi(c,k)\in\R^p$。在归一化约束$\sum_kq(k\mid c)=1$及特征期望约束$\sum_{c,k}\hat p(c)q(k\mid c)\phi(c,k)=\bar\phi$下，最大化条件熵$-\sum_{c,k}\hat p(c)q(k\mid c)\log q(k\mid c)$。对$q(k\mid c)$求拉格朗日一阶条件得到
\begin{equation}
 -\log q(k\mid c)-1+\theta^{\mathsf T}\phi(c,k)+\lambda(c)=0.
\end{equation}
指数化后，归一化消去$\lambda(c)$，即得$q(k\mid c)=\exp(\theta^{\mathsf T}\phi(c,k))/\sum_{k'}\exp(\theta^{\mathsf T}\phi(c,k'))$。这要求约束可行；特征完全分离时有限最优权重未必存在，通常以正则化松弛精确矩约束。名称中的熵指这个构造原则，不表示模型输出必须高熵。

MEMM 可以使用相关、重叠的输入特征，不需要描述它们的联合分布，但每个前驱状态都单独归一化，这带来了标签偏置。若某个状态只有一个允许的后继，该转移的条件概率恒为一，无论对应特征得分多低，都不能借助该步降低路径概率。更一般地，后继较少或转移分布较集中的状态可能获得结构性优势。问题来自局部归一化对路径评分的约束，并非简单的类别数量不平衡。

线性链条件随机场（CRF）保留相邻标签依赖，却把归一化放到整条路径上\cite{murphy2022pml}。定义路径得分
\begin{align}
 F_\theta(z,x)&=\sum_{t=1}^{T}\theta^{\mathsf T}f(z_{t-1},z_t,x,t),\\
 p_\theta(z\mid x)&=\frac{\exp F_\theta(z,x)}{Z_\theta(x)},\\
 Z_\theta(x)&=\sum_{z'_{1:T}}\exp F_\theta(z',x).
\end{align}
\begin{figure}[htbp]
\centering
\begin{tikzpicture}[
 state/.style={circle,draw,minimum size=9mm},
 factor/.style={rectangle,draw,fill=orange!25,minimum size=7mm,inner sep=1pt},
 every node/.style={font=\small}]
 \node[draw,rounded corners,fill=blue!12,minimum width=10cm,minimum height=9mm] (x) at (5,1.8) {给定观测 $x$：各局部势函数共享的条件};
 \node[state,fill=gray!15] (z0) at (0,0) {$z_0$};
 \foreach \t in {1,2,3,4}{
  \node[state] (z\t) at (2.5*\t,0) {$z_{\t}$};
  \node[factor] (f\t) at (2.5*\t-1.25,0) {$\psi_{\t}$};
  \draw[thick] (f\t) -- (z\t);
  \draw[dashed,blue!65!black] (x.south -| f\t.north) -- (f\t.north);
 }
 \foreach \a/\b in {0/1,1/2,2/3,3/4}{\draw[thick] (z\a) -- (f\b);}
 \node[align=center] at (5,-1.2) {$\psi_t=\exp\{\theta^{\mathsf T}f(z_{t-1},z_t,x,t)\}$\\局部势相乘，再用同一个 $Z_\theta(x)$ 对整条路径归一化};
\end{tikzpicture}
\caption{线性链 CRF 的因子图。圆形表示标签，方形表示相邻标签的势函数；实线无方向，表示变量参与因子。蓝色虚线仅表示势函数依赖已给定观测，不是标签图中的随机变量边或生成箭头。}
\label{fig:bayes-crf-structure}
\end{figure}
图\ref{fig:bayes-crf-structure}把条件分布写成可逐项检查的乘积：每个方形读取相邻两个标签及观测，输出一个非负相容度。固定$z_0$后，第一个因子实际只涉及一个待预测标签。去掉因子节点、将共享因子的标签直接相连，就得到标签的无向链。与 MEMM 不同，$\psi_t$不要求对后继标签求和为一，所有因子共同决定路径得分，最后才由一个全局配分函数归一化。

配分函数$Z_\theta(x)$比较所有允许路径的得分。一个只有唯一后继的状态仍然可以获得较低的路径得分，因此消除了 MEMM 局部归一化造成的上述限制。全局归一化不意味着自动解决标注噪声、分布偏移或特征不足。

对一条标注序列$(x,z)$，带二次正则的负对数条件似然为
\begin{equation}
 \mathcal L(\theta)=-F_\theta(z,x)+\log Z_\theta(x)
 +\frac{\lambda}{2}\|\theta\|_2^2.
\end{equation}
其梯度等于模型期望的特征计数减去真实路径的特征计数，再加$\lambda\theta$。这个形式解释了训练为何需要边缘概率：只找出当前最优路径不足以计算配分函数的梯度。在线性链、有限状态和相邻标签因子下，前向后向算法可在$O(TK^2)$时间内计算配分函数和边缘概率，Viterbi 则计算最优路径。

\paragraph{两条路径怎样体现全局归一化。}
假设只有两条合法路径，得分分别为0和$\log3$，则未归一化权重为1和3，$Z=4$，路径概率为$1/4$和$3/4$。若观测证据使第二条路径增加$\log2$，它的权重变成6，概率成为$6/7$。即使该路径某一步只有一个合法后继，该步得分仍能改变整条路径的权重；MEMM在那个唯一后继处的局部概率却恒为一。这个小例子只解释归一化的区别，并不是说CRF必然取得更高测试准确率。

特征向量$f(j,k,x,t)\in\R^p$可以把“从正常到退化”“当前振动超过阈值”等指示量列成$p$项，$\theta\in\R^p$为其权重，内积得到一个标量得分。指数化把任何实得分变成正权重，最后归一化才得到概率。后文神经编码器只是在学习怎样产生这些局部得分，不改变概率必须归一化这一要求。

\paragraph{配分函数、边缘概率与梯度的完整联系。}
令$\psi_t(j,k)=\exp\{\theta^{\mathsf T}f(j,k,x,t)\}$。固定起点$z_0$，设置$a_0(z_0)=1$、其余为零，递推$a_t(k)=\sum_ja_{t-1}(j)\psi_t(j,k)$，则$Z=\sum_ka_T(k)$。后向量$b_T(k)=1$、$b_{t-1}(j)=\sum_k\psi_t(j,k)b_t(k)$给出边缘
\begin{equation}
 p_\theta(z_{t-1}=j,z_t=k\mid x)
 =\frac{a_{t-1}(j)\psi_t(j,k)b_t(k)}{Z}.
\end{equation}
这些$a,b$是未归一化消息，不是HMM的观测似然参数。定义总特征
\begin{equation}
 \Phi(z,x)=\sum_tf(z_{t-1},z_t,x,t)\in\R^p.
\end{equation}
直接对有限和求导得
\begin{align}
 \nabla_\theta\log Z
 &=\frac{\sum_z e^{\theta^{\mathsf T}\Phi(z,x)}\Phi(z,x)}{Z}
 =\E_{p_\theta(z\mid x)}\Phi(z,x),\\
 \nabla_\theta\mathcal L
 &=-\Phi(z^{\rm obs},x)+\sum_{t,j,k}p_\theta(j,k\mid x)f(j,k,x,t)+\lambda\theta.
\end{align}
再求导得到$\nabla^2\mathcal L=\operatorname{Cov}_{p_\theta}(\Phi)+\lambda I_p$。因此固定特征的线性CRF目标凸，$\lambda>0$时严格凸；若特征由待训练神经网络产生，这一参数凸性便不再成立。禁止路径可赋零势，但至少必须存在一条合法路径，否则$Z=0$使条件分布无定义。

神经网络可以提供 CRF 的局部观测得分。例如，令$s_t(k;x)$为编码器输出，$B_{jk}$为可学习转移得分，则
\begin{equation}
 F_\theta(z,x)=\sum_{t=1}^{T}s_t(z_t;x)+\sum_{t=2}^{T}B_{z_{t-1},z_t}.
\end{equation}
与 HMM 的$A_{jk}$不同，$B_{jk}$是未归一化得分，不要求非负或行和为一。Bidirectional Long Short-Term Memory (BiLSTM)--CRF 和 Transformer--CRF 将特征提取交给神经网络，将标签关系交给结构化输出层。双向编码器适合完整记录的离线标注；在线预警必须限制输入和推断范围，或者明确等待未来观测所引入的延迟。即使局部特征只读取过去，对整段序列进行平滑仍可能利用后续证据。

\section{马尔科夫随机场与空间一致性}
图像中的局部相互作用可用 Gibbs 分布和马尔科夫随机场描述，Geman 与 Geman 将这一建模方式用于图像恢复\cite{geman1984mrf}。本节据此解释邻域势函数与能量最小化，但不把平滑先验当作真实边界的替代。
时间链只是一种邻接关系。在缺陷图像或设备网络中，每个变量可能有多个邻居，马尔科夫随机场（Markov Random Field, MRF）用无向图描述这种依赖。对严格正的分布，在相应马尔科夫性质下可以按图中的团分解为
\begin{equation}
 p(z)=\frac{1}{Z}\prod_{C\in\mathcal C}\psi_C(z_C)
 =\frac{1}{Z}\exp\{-E(z)\}.
\end{equation}
\begin{figure}[htbp]
\centering
\begin{tikzpicture}[
 site/.style={circle,draw,minimum size=11mm,fill=white},
 every node/.style={font=\small}]
 \foreach \r in {1,2,3}{
  \foreach \c in {1,2,3}{\node[site] (v\r\c) at (1.8*\c,-1.5*\r) {$z_{\r\c}$};}
 }
 \foreach \r in {1,2,3}{\foreach \a/\b in {1/2,2/3}{\draw[thick] (v\r\a) -- (v\r\b);}}
 \foreach \c in {1,2,3}{\foreach \a/\b in {1/2,2/3}{\draw[thick] (v\a\c) -- (v\b\c);}}
 \node[site,fill=orange!30] at (v22) {$z_{22}$};
 \foreach \n in {12,21,23,32}{\node[site,fill=green!15] at (v\n) {$z_{\n}$};}
 \node[align=left,text width=5.2cm,anchor=west] at (6.3,-3) {橙色：当前标签 $z_{22}$\\[3pt]绿色：四邻域 $N(22)$\\[3pt]白色：其余标签\\[6pt]给定绿色邻居后，中心标签与其余标签条件独立。};
\end{tikzpicture}
\caption{四邻域网格上的成对 MRF。每个节点对应一个像素或区域标签，每条无向边对应一个二元势；不存在箭头，并不表示各标签独立，而是表示相互作用不采用有向生成分解。}
\label{fig:bayes-mrf-structure}
\end{figure}
图\ref{fig:bayes-mrf-structure}展示的是一种具体的网格结构，并非所有 MRF 都必须是网格。此处$z_{rc}$的两个下标分别表示行和列。局部马尔科夫性质可写为
\begin{equation}
 p(z_{22}\mid z_{\setminus 22})
 =p(z_{22}\mid z_{12},z_{21},z_{23},z_{32}),
\end{equation}
其中$z_{\setminus 22}$表示除中心外的全部标签。对于图示的一元和二元势分解，固定其他标签后，与$z_{22}$无关的因子在条件概率归一化时相消，因而
\begin{equation}
 p(z_{22}=k\mid z_{\setminus 22})
 =\frac{\psi_{22}(k)\prod_{j\in N(22)}\psi_{22,j}(k,z_j)}
 {\sum_{k'}\psi_{22}(k')\prod_{j\in N(22)}\psi_{22,j}(k',z_j)}.
\end{equation}
这也解释了 Gibbs 采样为何只需读取当前节点的邻居。不过，局部更新只访问邻域，并不意味着联合分布能够分解成相互独立的节点分布。若这些势还依赖给定图像$x$，图中的标签邻接关系保持不变，得到的就是后文讨论的网格 CRF。

团$C$内的节点两两相连，势函数$\psi_C$表达局部相容程度，并不是归一化概率。能量$E(z)=-\sum_C\log\psi_C(z_C)$越低，配置概率越高。配分函数$Z$对全部配置求和，保证总概率为一。若使用零势函数表达硬约束，则不能直接套用依赖严格正性的等价结论。

在给定图像$x$的二元缺陷标注中，可以设置
\begin{equation}
 E(z;x)=\sum_i U_i(z_i;x)+\lambda\sum_{(i,j)\in\mathcal E}
 w_{ij}(x)\mathbb 1\{z_i\ne z_j\},\qquad \lambda,w_{ij}(x)\geq0.
\end{equation}
一元项$U_i$衡量局部像素证据，二元 Potts 项惩罚邻居标签不一致。相邻区域越相似，$w_{ij}$可以越大；跨越明显边界时降低权重，可减轻过度平滑。若直接将分类概率转换成一元负对数得分，这构成条件模型的一种参数化，并不自动成为生成似然。

此时$p(z\mid x)\propto\exp\{-E(z;x)\}$也是一个 CRF。MRF 强调无向马尔科夫依赖，CRF 强调给定观测后的结构化条件分布，两者不是互斥的模型名单。线性链 CRF 是其中计算特别方便的一类，二维网格 CRF 通常不能直接沿用链式动态规划。

能量中的平滑项可以移除孤立误报，也可能抹掉真实细裂纹，因此应依据缺陷宽度和边界误差选择强度。一般含环图的精确推断可能十分昂贵，变量消元的复杂度随树宽增长。置信传播在树上精确，在含环图上通常只是近似且不保证收敛。均值场以较简单的分布逼近联合依赖，Gibbs 采样则逐变量更新。对于二元、次模的成对能量，图割可求全局最小能量，但这一结论不能推广到任意多类别势函数或边缘概率计算。

\begin{table}[htbp]
\centering
\caption{链式与无向图模型的建模对象和计算边界。}
\label{tab:bayes-graph-models}
\small
\begin{tabular}{p{1.3cm}p{3.1cm}p{3.3cm}p{4.2cm}}
\hline
模型 & 建模对象 & 归一化方式 & 关键限制 \\
\hline
HMM & 联合分布$p(x,z)$ & 转移与观测局部归一化 & 观测假设和状态持续时间需检验 \\
MEMM & 条件分布$p(z\mid x)$ & 每个前驱状态局部归一化 & 标签偏置影响路径评分 \\
链式 CRF & 条件分布$p(z\mid x)$ & 所有标签路径全局归一化 & 训练需计算模型期望 \\
MRF & 无向图上的联合分布 & 全部配置全局归一化 & 一般图中配分函数和推断困难 \\
\hline
\end{tabular}
\end{table}

表\ref{tab:bayes-graph-models}按建模对象、归一化方式和推断边界比较图模型；应先决定需要联合生成还是条件预测，再判断链式动态规划是否适用。
\paragraph{MRF的局部更新与均值场。}
对固定其他标签$z_{-i}$，不包含$i$的能量项在条件概率的分子分母中相消，因此
\begin{equation}
 p(z_i=k\mid z_{-i},x)\propto
 \exp\left\{-U_i(k;x)-\lambda\sum_{j:(i,j)\in\mathcal E}w_{ij}(x)\mathbb 1\{k\ne z_j\}\right\}.
\end{equation}
逐点从此分布抽样就是Gibbs更新；若改为取最小局部能量，则变为坐标下降，只能保证能量不增而非全局最优。有限状态、适当遍历性和非周期性下Gibbs链以目标分布为平稳分布，强耦合时仍可能混合极慢。

若用乘积分布$q(z)=\prod_iq_i(z_i)$近似，最小化$\KL(q\|p)$等价于最小化以下目标：
\begin{equation}
 \E_qE+\sum_i\sum_kq_i(k)\log q_i(k).
\end{equation}
固定其他因子并对$q_i$加归一化乘子求导，可得
\begin{equation}
 q_i(k)\propto\exp\{-\E_{q_{-i}}E(k,z_{-i};x)\}
 \propto\exp\{-U_i(k;x)-\lambda\sum_jw_{ij}(1-q_j(k))\}.
\end{equation}
这以邻居概率代替确定标签，但忽略联合相关性；逐坐标精确更新改善变分目标，却不保证找到全局最佳近似。二元成对能量可用图割的关键条件是$V_{ij}(0,0)+V_{ij}(1,1)\leq V_{ij}(0,1)+V_{ij}(1,0)$，非负Potts项满足它。图割解决最大概率配置，不提供配分函数或后验区间。

\paragraph{两个像素的能量可以全部列完。}
设两个相邻标签$z_1,z_2\in\{0,1\}$没有一元偏好，平滑惩罚为$\lambda=\log3$、边权为一。同标签$(0,0),(1,1)$的能量为零、权重为1，不同标签$(0,1),(1,0)$的能量为$\log3$、权重为$1/3$。配分函数为$8/3$，因而$P(0,0)=P(1,1)=3/8$、$P(0,1)=P(1,0)=1/8$。未观察邻居时，像素1的边缘概率为$P(z_1=0)=3/8+1/8=1/2$；已知$z_2=0$时，条件概率为$P(z_1=0\mid z_2=0)=(3/8)/(3/8+1/8)=3/4$。两种配置并列最优，并不使单个像素的边缘判断自动偏向零。在邻居为零时，Gibbs更新以该概率抽零，以$1/4$抽一；若总选择当前条件概率最高的标签，则变成逐坐标寻找局部最优配置，已经不是从后验抽样。配置概率最高与某像素的边缘概率，是两个不同问题。

\paragraph{把图结构带回维护任务。}
假设要从一次运行记录中识别正常、退化与故障阶段，可以先用独立分类器建立逐点基线。若预测频繁跳变，HMM 用转移持续性与观测模型共同解释状态；若观测特征复杂而状态标注充分，MEMM 和 CRF 可直接学习条件边界。若任务转为表面缺陷的空间连通区域，则链结构不再适用，需要结合图像邻接关系定义条件随机场。模型变化由输出依赖决定，而非由名称的新旧决定。

这些比较必须使用相同的可用观测、标签定义和设备划分。逐点准确率之外，还应检查状态切换延迟、阶段持续时间误差和关键故障段漏检。离线平滑结果只能与离线结果比较，在线过滤应另报实际延迟。若由隐状态概率推导未来故障风险，还需要定义故障事件与状态的关系，并验证预测时间窗，不能把“当前处于退化状态”的概率直接解释成“未来二十四小时故障”的概率。

上述动态规划在参数给定时可以精确计算链式隐状态分布，但这不意味着参数已经确定。小样本设备的转移矩阵、观测噪声和条件特征权重仍有估计误差。接下来回到参数层，讨论哪些模型可以解析更新，哪些模型需要近似后验。

\section{从共轭模型到近似推断}
\paragraph{共轭模型与近似推断。}
共轭是指给定某种似然后，所选先验族在更新下保持封闭：后验仍属于同一分布族，只更新其参数。这不是说先验和似然必须具有相同分布名称。例如，伯努利观测的成功概率$p$取Beta先验$p\sim\mathrm{Beta}(a,b)$，观察到$k$次成功、$n-k$次失败后有
\begin{equation}
 p\mid\mathcal{D}\sim\mathrm{Beta}(a+k,b+n-k).
\end{equation}
后验均值为
\begin{equation}
 \mathbb{E}[p\mid\mathcal D]=\frac{a+k}{a+b+n}.\label{eq:bayes-beta-mean}
\end{equation}
式\eqref{eq:bayes-beta-mean}将先验均值与样本比例按先验强度和观测数量加权。$a+b$越大，先验影响越强。用于故障率时，这个更新假定各次观测在给定$p$后独立且共享同一事件概率，不能直接把大量重叠时间窗视为独立试验。

\paragraph{Beta更新的归一化与预测。}
对$a,b>0$，先验密度为$p^{a-1}(1-p)^{b-1}/B(a,b)$，其中$B(a,b)=\int_0^1u^{a-1}(1-u)^{b-1}\,du$。给定独立伯努利序列，似然正比于$p^k(1-p)^{n-k}$，相乘后指数为$a+k-1$和$b+n-k-1$，用$B(a+k,b+n-k)$归一化便得到后验。若观测对象只是计数，似然还有$\binom nk$，它不依赖$p$，不改变后验。

记$A=a+k$、$B=b+n-k$（此处大写$B$为参数，不是Beta函数），由归一化积分之比得到
\begin{equation}
 \E[p\mid\mathcal D]=\frac{A}{A+B},\quad
 \Var(p\mid\mathcal D)=\frac{AB}{(A+B)^2(A+B+1)},\quad
 P(Y_*=1\mid\mathcal D)=\E[p\mid\mathcal D].
\end{equation}
其中$\E p=\frac{a+b}{a+b+n}\frac a{a+b}+\frac n{a+b+n}\frac kn$（$n>0$），显式给出收缩权重。未来$m$次试验的成功数$K_*$满足
\begin{equation}
 P(K_*=r\mid\mathcal D)=\binom mr
 \frac{B(A+r,B+m-r)}{B(A,B)},\quad 0\leq r\leq m.
\end{equation}
这是先对二项分布的未知$p$积分所得，不等同于把后验均值代入二项分布：积分后未来试验通过共享$p$相关，从而保留参数不确定性。

\paragraph{用最简单的Beta后验区分三类区间。}
均匀先验$\operatorname{Beta}(1,1)$配合一次成功，得到$\operatorname{Beta}(2,1)$，后验密度为$2p$、累积分布为$F(p)=p^2$。后验均值为$2/3$；95\%等尾可信区间的端点满足$F(l)=0.025$、$F(u)=0.975$，所以$[l,u]=[\sqrt{0.025},\sqrt{0.975}]\approx[0.158,0.987]$。给定此次数据和先验，参数$p$位于此区间的后验概率是$0.95$。它不是频率学派“重复采样时95\%的区间覆盖固定参数”的定义，也不声称一次新结果会落在这个实数区间内：新伯努利结果只有0或1，其成功概率为$2/3$。

若预测未来两次都成功，完整积分为$\E[p^2\mid\mathcal D]=\int_0^1p^2(2p)\,dp=1/2$，而把均值代入会给出$(2/3)^2=4/9$。两次未来结果给定$p$后独立，但在$p$未知并被积分掉后相关；“条件独立”不能随意删掉条件。可信区间也不是所有模型下都具有同样的重复抽样覆盖率，仍需后续覆盖检查。

共轭关系也适用于连续参数。考虑$y=X\beta+\varepsilon$，其中$\varepsilon\sim\mathcal N(0,\sigma^2I)$，先验$\beta\sim\mathcal N(m_0,S_0)$。当$\sigma^2$已知时，贝叶斯线性回归的后验仍为高斯分布，其精度矩阵和均值为
\begin{align}
 S_n^{-1}&=S_0^{-1}+\sigma^{-2}X^{\mathsf T}X,\\
 m_n&=S_n\left(S_0^{-1}m_0+\sigma^{-2}X^{\mathsf T}y\right).
\end{align}
先验精度与数据精度相加，说明收缩先验如何稳定共线特征下的小样本回归。对新输入$x_*$，预测均值为$x_*^{\mathsf T}m_n$，预测方差为$\sigma^2+x_*^{\mathsf T}S_nx_*$，分别包含观测噪声和参数不确定性。

\paragraph{高斯后验的配方步骤。}
令$X\in\R^{n\times d}$、$\beta,m_0\in\R^d$、$S_0\in\R^{d\times d}$正定，$y\in\R^n$，且$\sigma^2>0$已知。负对数联合密度中依赖$\beta$的部分为
\begin{align}
 &\frac12(\beta-m_0)^{\mathsf T}S_0^{-1}(\beta-m_0)
   +\frac1{2\sigma^2}(y-X\beta)^{\mathsf T}(y-X\beta)\\
 &=\frac12\beta^{\mathsf T}\Lambda\beta-\beta^{\mathsf T}h+\text{常数}\\
 &=\frac12(\beta-\Lambda^{-1}h)^{\mathsf T}\Lambda
   (\beta-\Lambda^{-1}h)+\text{常数},
\end{align}
其中$\Lambda=S_0^{-1}+\sigma^{-2}X^{\mathsf T}X$、$h=S_0^{-1}m_0+\sigma^{-2}X^{\mathsf T}y$。识别高斯指数可得$S_n=\Lambda^{-1}$、$m_n=S_nh$。正定先验使$X$秩不足时后验仍适定；实际通过线性方程求解，不必显式求逆。

新观测为$Y_*=x_*^{\mathsf T}\beta+\varepsilon_*$，其中$x_*\in\R^d$且新噪声独立，因此线性高斯和仍为高斯，均值$x_*^{\mathsf T}m_n$、方差$x_*^{\mathsf T}S_nx_*+\sigma^2$。潜在均值区间只含第一项，观测预测区间包含两项。若噪声尺度未知并为其指定共轭先验，积分后常得到Student分布，而非直接沿用已知$\sigma^2$的高斯结论。

复杂深度模型通常没有解析后验，可采用最大后验估计、Laplace近似、变分推断或Markov Chain Monte Carlo (MCMC)。变分推断选择易处理的$q_{\phi}(\bm{\theta})$逼近后验，并最小化
\begin{equation}
 \KL\left(q_{\phi}(\bm{\theta})\,\|\,p(\bm{\theta}\mid\mathcal{D})\right).
\end{equation}
通过证据下界（Evidence Lower Bound，ELBO）可写成
\begin{equation}
 \log p(\mathcal{D})\geq \mathbb{E}_{q_{\phi}}[\log p(\mathcal{D}\mid\bm{\theta})]-\KL(q_{\phi}(\bm{\theta})\|p(\bm{\theta})).\label{eq:bayes-elbo}
\end{equation}
式\eqref{eq:bayes-elbo}右端的第一项鼓励解释数据，第二项控制近似后验偏离先验的程度。

这一优化目标来自恒等式
\begin{equation}
 \KL(q_\phi(\theta)\|p(\theta\mid\mathcal D))
 =\log p(\mathcal D)-\operatorname{ELBO}(\phi).
\end{equation}
证据对$\phi$为常数，因此最大化 ELBO 等价于最小化所选方向的 KL 散度。均值场近似$q(\theta)=\prod_jq_j(\theta_j)$计算方便，却忽略参数相关性，可能低估后验方差。结构化变分分布保留部分依赖，代价是更复杂的计算。

若变分变量$z\sim\mathcal N(\mu_\phi,\sigma_\phi^2)$，可写成$z=\mu_\phi+\sigma_\phi\epsilon$，其中$\epsilon\sim\mathcal N(0,1)$与$\phi$无关。重参数化把随机性移到外部噪声，使 ELBO 的梯度能够通过可微变换估计。离散隐状态则未必具有这种可微表示，应优先检查是否能像 HMM 那样精确求和，再考虑松弛或得分函数估计。

\paragraph{ELBO恒等式与可计算梯度。}
设$q$在联合密度有支持的区域内，且所需积分有限。把Bayes公式代入KL定义，逐项展开为
\begin{align}
 \KL(q\|p(\theta\mid\mathcal D))
 &=\E_q\log q-\E_q\log p(\mathcal D,\theta)+\log p(\mathcal D),\\
 \operatorname{ELBO}(q)
 &=\E_q\log p(\mathcal D\mid\theta)+\E_q\log p(\theta)-\E_q\log q.
\end{align}
KL非负给出下界，等号只在$q$等于真实后验（几乎处处）时成立。固定模型时，提高ELBO会缩小这一KL；跨模型比较时，各自的对数证据也在变化，下界差距还混有近似误差，不能仅据下界大小判断后验近似的准确程度。

对均值场$q(\theta)=\prod_jq_j(\theta_j)$，固定其他因子，把ELBO视为$q_j$的函数并加归一化乘子，可得
\begin{equation}
 \log q_j^*(\theta_j)=\E_{q_{-j}}\log p(\mathcal D,\theta)+\text{归一化常数}.
\end{equation}
这给出坐标上升更新的来源；期望或归一化未必能解析算出。若$\theta=\mu+s\odot\epsilon$、$\epsilon\sim\mathcal N(0,I_d)$、$s_j>0$，在可交换微分与积分的可积条件下，
\begin{equation}
 \nabla_\mu\E_\epsilon h(\mu+s\odot\epsilon)
 =\E_\epsilon\nabla_\theta h(\theta),\qquad
 \partial_{s_j}\E_\epsilon h(\theta)=\E_\epsilon[\partial_{\theta_j}h(\theta)\epsilon_j].
\end{equation}
正态$q$的熵为$\frac12\sum_j\log(2\pi e s_j^2)$，其对$s_j$的导数为$1/s_j$，与联合对数密度梯度合成完整ELBO梯度。可用$s_j=\exp(\rho_j)$保证正值。采样梯度只近似当前变分目标的梯度，不能消除变分族的结构偏差。

\paragraph{把近似推断看成三个不同问题。}
MAP问哪个参数点的后验密度最高；变分推断问预先选定的简化分布族中，哪个分布最接近后验；MCMC则生成一串相关的参数样本，用平均近似积分。更多变分抽样可以减小梯度噪声，却不能修复分布族遗漏的相关性；更多MCMC迭代在满足混合条件时改善积分精度，却不改变数据不足造成的后验宽度。

KL散度$\KL(q\|p)=\E_q\log(q/p)$不是对称距离：若$q$把质量放在$p=0$处，代价会发散。ELBO中的“下界”是对对数证据的下界，不是对预测准确率的下界。重参数化中的$\odot$表示逐分量相乘；$\mu,s,\epsilon\in\R^d$，先抽固定分布噪声再变换，便可在保持噪声不变时求$\mu,s$的导数。

\paragraph{贝叶斯更新的几何直观。}
贝叶斯公式不仅是概率比例式，也可以理解为在函数空间中重新分配权重。先验把概率质量放在一组可能的参数上，似然根据参数对观测数据的解释能力重新加权，后验再进行归一化。数据量较小时，先验仍然有明显影响；数据量增加后，似然通常占据主导，但这依赖模型假设没有严重错配。

\begin{figure}[htbp]
\centering
\begin{tikzpicture}[>=Latex, node distance=1.8cm, every node/.style={font=\small}]
 \node[draw, rounded corners, fill=blue!8, minimum width=2.6cm, minimum height=1cm] (prior) {先验 $p(\theta)$};
 \node[draw, rounded corners, fill=orange!12, right=of prior, minimum width=2.8cm, minimum height=1cm] (like) {似然 $p(\mathcal D\mid\theta)$};
 \node[draw, rounded corners, fill=green!12, right=of like, minimum width=2.8cm, minimum height=1cm] (post) {后验 $p(\theta\mid\mathcal D)$};
 \draw[->, thick] (prior) -- node[above]{结合} (like);
 \draw[->, thick] (like) -- node[above]{归一化} (post);
 \draw[->, dashed] ($(prior.south)!0.5!(like.south)$) .. controls +(0,-1.0) and +(0,-1.0) .. node[below]{模型假设} ($(like.south)!0.5!(post.south)$);
\end{tikzpicture}
\caption{贝叶斯更新把先验结构与数据证据结合为后验。}
\label{fig:bayes-update-components}
\end{figure}
图\ref{fig:bayes-update-components}将先验与数据证据汇入后验，强调后验结论同时依赖结构假设和观测；检验预测可信度时不能只检查数据拟合。

\section{模型比较与不确定性计算}
\paragraph{从 Maximum A Posteriori (MAP) 到正则化。}
最大后验估计最大化
\begin{equation}
 \hat\theta_{MAP}=\argmax_\theta \log p(\mathcal D\mid\theta)+\log p(\theta).
\end{equation}
以高斯先验$p(\theta)\propto\exp(-\lambda\|\theta\|_2^2/2)$为例，负对数后验等价于经验损失加上$\lambda\|\theta\|_2^2/2$。因此正则化并不是一个与概率建模无关的工程技巧，而是在特定先验下的推断结果。需要注意的是，正则化系数、噪声方差和先验方差之间存在耦合，不能脱离似然尺度解释$\lambda$。

\paragraph{边际似然与模型比较。}
边际似然
\begin{equation}
 p(\mathcal D)=\int p(\mathcal D\mid\theta)p(\theta)\,d\theta
\end{equation}
同时惩罚拟合不足和模型过度灵活。它与只比较训练集似然不同，因为参数空间中大量不能解释数据的区域会降低平均证据。实际计算边际似然通常困难，可使用 ELBO、Laplace 近似或信息准则近似。比较时应使用同一观测数据与同一观测测度，并指定归一化的先验。不同先验可以比较，但比较的是包含各自先验的完整模型，而非只比较似然结构；任意常数未定的不适当先验通常不能直接用于证据比。

\paragraph{Laplace 近似。}
在后验众数$\hat\theta$附近，对负对数后验$U(\theta)=-\log p(\theta\mid\mathcal D)$做二阶 Taylor 展开：
\begin{equation}
 U(\theta)\approx U(\hat\theta)+\frac12(\theta-\hat\theta)^{\mathsf T}H(\theta-\hat\theta),
 \qquad H=\nabla^2U(\hat\theta).
\end{equation}
于是后验近似为$\mathcal N(\hat\theta,H^{-1})$。该近似在后验近似单峰且局部二次时有效；深度网络存在对称性、多峰和平坦方向时，单个高斯可能严重低估不确定性。

\paragraph{从局部二次式到归一化常数。}
若$\hat\theta\in\R^d$为内部驻点且$H$正定，线性项因梯度为零而消失。指数化二阶式，并用高斯积分
\begin{equation}
 \int\exp\{-\tfrac12u^{\mathsf T}Hu\}\,du=(2\pi)^{d/2}|H|^{-1/2}
\end{equation}
归一化，便得到协方差$H^{-1}$。若从未归一化密度$\tilde p(\theta)=p(\mathcal D\mid\theta)p(\theta)$出发，同样推得
\begin{equation}
 \log p(\mathcal D)\approx\log\tilde p(\hat\theta)
 +\frac d2\log(2\pi)-\frac12\log|H|.
\end{equation}
它同时考虑峰高与峰附近体积，而不是只看最大似然。边界众数、奇异Hessian、多峰或不适当先验会使这一步失效；加入数值阻尼可改善线性代数，却不是证明原后验本来就是该高斯。

\paragraph{MCMC 的接受率直观。}
Metropolis--Hastings 从当前状态$\theta$提出候选$\theta'$，以
\begin{equation}
 \alpha=\min\left\{1,\frac{p(\theta'\mid\mathcal D)q(\theta\mid\theta')}{p(\theta\mid\mathcal D)q(\theta'\mid\theta)}\right\}
\end{equation}
的概率接受。候选分布步长过小会导致样本高度相关，步长过大则接受率过低。实际诊断不能只看链是否“运行完成”，还应检查 burn-in、有效样本量、自相关和多条链之间的混合。

\begin{table}[htbp]
\centering
\caption{贝叶斯推断方法的工程取舍}
\label{tab:bayes-inference-methods}
\scriptsize
\setlength{\tabcolsep}{3pt}
\resizebox{\linewidth}{!}{%
\begin{tabular}{p{2.5cm}p{3.6cm}p{3.6cm}p{4.0cm}}
\hline
方法 & 近似对象 & 优势 & 主要风险与适用场景 \\
\hline
MAP & 后验众数 & 成本低、易扩展、适合在线更新 & 忽略参数不确定性，适合作为基线 \\
Laplace & 众数附近的高斯曲率 & 能快速产生局部区间 & 多峰、对称和非二次后验会导致低估 \\
变分推断 & 可处理的$q_\phi(\theta)$ & 可用随机梯度扩展到大模型 & 均值场假设常低估相关性和方差 \\
MCMC & 后验样本链 & 理论上逼近目标后验 & 计算昂贵，必须检查混合和有效样本量 \\
\hline
\end{tabular}}
\end{table}

表\ref{tab:bayes-inference-methods}区分局部近似、变分优化与采样的计算取舍；选用方法时应把其近似误差和诊断要求与后验形状共同考虑。
\paragraph{接受率为何保持目标分布。}
记目标密度为$\pi(\theta)$。对不同状态$\theta,\theta'$，Metropolis--Hastings (MH)的双向概率流满足
\begin{align}
 \pi(\theta)q(\theta'\mid\theta)\alpha(\theta,\theta')
 &=\min\{\pi(\theta)q(\theta'\mid\theta),
          \pi(\theta')q(\theta\mid\theta')\}\\
 &=\pi(\theta')q(\theta\mid\theta')\alpha(\theta',\theta).
\end{align}
这就是细致平衡；拒绝时停留原位使转移核的总概率为一。对当前状态积分可验证$\pi$不变，且后验的未知证据常数在比值中消去。平稳性不等于有限步已收敛；还需要能到达目标支持、非周期性及相应遍历条件。对称提议可约去$q$，非对称提议遗漏该比值则会改变目标分布。

\paragraph{一次接受、一次拒绝怎样写进样本链。}
假设对称提议在两个参数值之间移动，未归一化后验权重分别为2和6。从权重2的点提出权重6的点，接受率为一；反向提议的接受率为$1/3$。实际抽$u\sim\operatorname{Uniform}(0,1)$，$u\leq1/3$才接受，否则下一样本仍记当前点。若只保存接受的不同点而丢掉停留，会改变样本在各区域的频率，不再近似目标后验。未知证据常数相消，所以这里不必先算出后验归一化分母。

burn-in指弃去受初始化影响的早期部分，并非预先固定丢掉10\%就保证收敛。多链指标$\hat R$比较链内与链间波动，接近一是必要诊断线索而不是正确性的证明，尤其不能排除所有链一起漏掉某个模式。

\section{校准、诊断与工业决策}
\paragraph{不确定性的分解与校准。}
预测方差可用全概率方差公式分解：
\begin{equation}
 \Var(y_*\mid x_*,\mathcal D)=\mathbb E_{p(\theta\mid\mathcal D)}[\Var(y_*\mid x_*,\theta)]
 +\Var_{p(\theta\mid\mathcal D)}[\mathbb E(y_*\mid x_*,\theta)].\label{eq:bayes-predictive-variance}
\end{equation}
式\eqref{eq:bayes-predictive-variance}右端的第一项主要对应观测噪声，第二项对应参数不确定性。模型给出的区间还必须经过覆盖率检查：预测百分之九十区间不应只在训练分布内覆盖百分之六十的真实值。对于工业决策，区间宽度与错误成本同样重要。

\paragraph{全方差公式的推导与估计。}
省略给定的$x_*,\mathcal D$，记$m_\theta=\E[Y_*\mid\theta]$、$m=\E_\theta m_\theta$，则$Y_*-m=(Y_*-m_\theta)+(m_\theta-m)$。平方后取期望，交叉项为零，因为$\E[Y_*-m_\theta\mid\theta]=0$，从而得到全方差公式。这个分解是在指定模型内的，不包含模型类别错设所遗漏的不确定性。

若有$S$个参数样本，并能计算条件均值$m_s$及方差$v_s$，令$\bar m=S^{-1}\sum_sm_s$，则经验混合分布的精确方差为
\begin{equation}
 \hat v=\frac1S\sum_sv_s+\frac1S\sum_s(m_s-\bar m)^2.
\end{equation}
若目的是无偏估计独立后验样本的均值间方差，第二项可改用$S-1$分母，但这时不再是该有限混合的精确方差。回归区间需从完整预测混合求分位数，不能仅对$m_s$取分位数后声称覆盖有噪声观测。对于重尾或多峰预测，$\bar m\pm1.96\sqrt{\hat v}$也不一定是95\%区间。

\paragraph{参数抽样与观测抽样要分两层。}
设两个等权参数样本对应的条件预测分别为
\begin{equation*}
 Y_*\mid\theta^{(1)}\sim\mathcal N(0,1),\qquad
 Y_*\mid\theta^{(2)}\sim\mathcal N(2,1).
\end{equation*}
这里省略共同给定的输入和数据。混合均值为1，平均条件方差为1，条件均值的方差也为1，总预测方差为2。仅保存两个条件均值$(0,2)$会遗漏观测噪声。

生成预测样本时应先选择参数样本，再从对应的条件观测分布抽$Y_*$；生成参数可信区间则只对参数样本取分位数。把抽样次数由100增至10000，只会使这些分位数估计更稳定，并不等于增加了9900次真实设备试验。

\paragraph{贝叶斯方法的常见误用。}
使用贝叶斯模型时，可以从结果倒查计算过程。先检查报告的概率与实际发生频率是否相符，再检查陌生工况下模型是否仍然过度自信。后验方差只反映当前模型表达的未知量，模型没有考虑到的故障机制可能根本不会使它增大。先验也应写明来源和尺度；如果近似算法尚未稳定，计算误差甚至可能大于想要估计的参数不确定性。

\paragraph{预测分布的蒙特卡洛近似。}
当后验积分无法解析计算时，可从后验得到$S$个样本$\theta^{(s)}$，使用
\begin{equation}
 p(y_*\mid x_*,\mathcal D)\approx\frac1S\sum_{s=1}^{S}p(y_*\mid x_*,\theta^{(s)}).
\end{equation}
注意先平均概率再取最大类别；平均 logits 或平均类别标签一般不等价。对回归，可由每个参数样本产生一条预测，再计算总方差和分位数区间。$S$太小会造成蒙特卡洛噪声，样本相关时有效样本量还会进一步下降。

\paragraph{MCMC 诊断不能省略。}
链的轨迹图可以检查是否进入稳定区域，自相关函数衡量样本冗余，有效样本量近似为
\begin{equation}
 ESS\approx\frac{S}{1+2\sum_{k\geq1}\rho_k}.
\end{equation}
多条不同初始状态的链还应比较$\hat R$等收敛指标。若链只在后验的一个模式附近游走，表面上的稳定并不代表全局混合良好。工业高维模型通常更偏向近似推断，但这不意味着可以跳过近似质量和预测覆盖率检查。

\paragraph{Effective Sample Size (ESS)是针对哪个估计量。}
对平稳链上的标量$h(\theta^{(s)})$，设其方差为$v_h$、滞后相关为$\rho_k$，展开样本均值的协方差和得
\begin{equation}
 \Var\left(\frac1S\sum_sh(\theta^{(s)})\right)
 =\frac{v_h}{S}\left[1+2\sum_{k=1}^{S-1}(1-k/S)\rho_k\right].
\end{equation}
当相关可求和且链足够长时，匹配独立样本均值的方差便得到前述ESS。不同参数、尾部概率和维护成本函数有不同的自相关，因此不存在适用于所有结论的单一ESS。短链直接相加所有经验相关会很不稳定，需截断或稳定的谱估计；尚未进入平稳区或未跨模式时，这个局部诊断不能证明准确性。

\paragraph{先验敏感性分析。}
先验选择应通过敏感性分析验证。对故障率、噪声尺度或参数收缩强度使用几组合理先验，比较后验预测和决策是否稳定。若结果对先验极其敏感，说明数据不足、模型可识别性弱，或先验承担了过多结论。此时应增加数据、改进测量或明确报告先验依赖，而不是只选择产生最好结果的先验。

\paragraph{校准、覆盖率与风险。}
预测区间覆盖率和区间宽度需要同时观察。过宽的区间容易达到覆盖率，却可能没有决策价值；过窄的区间看似精确，却会造成过度自信。对分类概率，可按置信度分箱比较预测概率和经验事件率；对回归区间，可按设备、工况和预测提前量分层检查。校准应在独立验证数据上完成，并在部署漂移后持续监测。

\paragraph{小样本故障率的后验决策。}
某设备在新工况下运行$20$次，观察到$2$次异常。若采用$\operatorname{Beta}(1,9)$先验，则后验为$\operatorname{Beta}(3,27)$，后验均值为$0.1$。直接使用样本比例得到的估计也是$0.1$，但二者对未来风险的表达不同：后验还包含小样本造成的宽区间。若把“安排人工复检”的成本记为$C_I$，把一次漏检损失记为$C_M$，则当后验异常概率满足
\begin{equation}
 p(\text{异常}\mid\mathcal D)C_M>C_I
\end{equation}
时才应触发复检。这里假定复检及随后的处理能消除所计漏检损失，且只产生固定成本$C_I$。例如$C_M=100,C_I=12$时，后验概率$0.1$对应不复检损失10，小于12，故不复检；若复检不完美，应重新列出残余损失。这样，概率输出通过成本函数转化为动作，而不是凭经验固定一个看似通用的阈值。

该例还说明先验不是隐藏的常数。若先验改为$\operatorname{Beta}(1,1)$，后验均值变为$3/22$；两种结果的差异提示当前数据不足以完全决定风险。工程报告应同时给出先验、后验和决策阈值，使不同团队可以复核结论。

\paragraph{回归预测区间的覆盖实验。}
假设模型对每台设备给出正态观测预测分布，其均值为$\mu_i$、标准差为$\sigma_i$，则名义95\%预测区间为$[\mu_i-1.96\sigma_i,\mu_i+1.96\sigma_i]$。在验证集上统计真实值落入区间的比例，并分别按设备、环境温度和预测提前量分组。若整体覆盖率为$0.94$，但高温工况只有$0.71$，则总体数字掩盖了局部失效。可以在独立校准集上估计分位数修正因子，再在另一段时间上复验，避免把校准集上的偶然波动误认为改进。

区间宽度也应进入报告。两个模型可能具有相同覆盖率，但一个区间宽得无法指导维护，另一个区间更窄且只在少数工况失效。因而应联合绘制覆盖率、平均区间宽度和漏覆盖成本，并说明校准是否使用了未来数据。

\paragraph{自测：把数据更新与模拟分开。}
先验为$\operatorname{Beta}(2,2)$，两次条件独立试验一成功一失败。求后验及下一次成功概率；随后模拟100次预测结果，是否应把模拟成功数再加进后验参数？

\textit{参考解。}后验为$\operatorname{Beta}(3,3)$，预测成功概率为$1/2$。模拟结果来自当前模型，不是新观测，不能再次更新为设备证据；增加模拟次数只改善数值估计。未来两次均成功的概率为$\E[p^2]=3\times4/(6\times7)=2/7$，不等于把$1/2$平方所得的$1/4$。

\paragraph{后验预测检查的操作步骤。}
用于拟合后验的真实数据也参与了这项检查，因此它不是独立测试集上的泛化评价。复制数据只是模型内的模拟结果，不能作为新设备证据再次更新同一后验。
后验预测检查不是把训练数据重新拟合一次，而是从后验参数样本生成复制数据$\tilde y$，比较$\tilde y$与真实数据的统计量。例如可比较均值、方差、极值数量、故障持续时间和相邻残差相关性。若模型能匹配均值却无法产生真实的长尾，说明其似然假设不适合风险评估。工业案例中，检查统计量应优先选择与维护动作相关的量，而不是只选择容易展示的总体误差。

本章的后验、图模型和近似推断可分别参考下面的课程材料与官方概率编程教程，前者建立概念，后者展示可计算实现。
\section*{辅助学习材料}
\begingroup
\emergencystretch=3em
\begin{itemize}
\item \textbf{Probabilistic Graphical Models}；作者/平台：Stanford University CS228；英文；课程讲义与课程资源。重点学习 Bayesian networks、Markov networks、因子分解和推断主题，帮助理解本章的 HMM、MRF、消息传递与条件独立。\href{https://ermongroup.github.io/cs228/}{课程主页}。
\item \textbf{Bayesian Regression - Introduction (Part 1)}；作者/平台：Pyro 社区；英文；官方教程。阅读模型与先验、变分推断和后验预测部分，将点估计与参数不确定性作对照。\href{https://pyro.ai/examples/bayesian\_regression.html}{在线阅读}。
\item \textbf{【深度学习数学基础 23】统计推断与贝叶斯框架}；知乎；中文解说。先写先验、似然与后验，再复算本章一维更新例子；标题与主题据必应索引，原页受限。\href{https://zhuanlan.zhihu.com/p/1910026999415673453}{在线阅读}。
\item \textbf{贝叶斯推断概览}；CSDN；中文博客。比较参数后验与后验预测，注意噪声假设；标题与主题据必应索引，原页受限。\href{https://blog.csdn.net/qq\_41282102/article/details/84928605}{在线阅读}。
\item \textbf{基于GMM-HMM的语音处理全过程讲解-通俗易懂}；上传者：\_静静老师；bilibili；中文技术讲解。以语音为例对照本章隐马尔可夫模型，区分隐状态转移与观测分布，再比较工业状态估计的条件假设。2026年10月4日公开元数据播放4,404次；本次候选中选择主题直接相关的讲解，不宣称达到十万播放或全站最高；已核对标题与计数，未逐帧观看。\href{https://www.bilibili.com/video/BV1Rv4y1c7SA/}{观看视频}。
\end{itemize}
\endgroup
\paragraph{本章小结。}
贝叶斯更新表达参数在当前证据下仍有多少不确定性，概率图模型则明确状态、观测和标签之间保留哪些依赖。HMM 建模状态与观测的联合分布，MEMM 和 CRF 直接建模条件状态分布，MRF 将邻接关系推广到一般无向图。结构决定可以使用哪些推断算法，参数学习决定这些算法依据什么概率进行计算。即使推断在数学上精确，错误的观测假设或时间边界仍会使维护决策失效。后验预测检查、覆盖率与先验敏感性因此构成模型使用前的检验，而不是公式推导后的附加说明。

贯穿这些方法的共同步骤是先写联合或条件分布，再利用结构进行求和、积分或近似，最后在实际成本下解释预测。解析解提供清楚的参照，动态规划节省状态枚举，变分与采样处理无法解析的参数积分；各自的适用边界都来自前提，而非算法名称。小样本复检案例中的阈值还隐含复检能消除所计漏检损失的简化假设，若检查不完美或有后续成本，应扩充动作损失再取后验期望。

下一章继续使用概率，但不只给有限个系数设分布，还直接考虑可能的函数。线性高斯后验中用过的配方，将再次用于高斯过程的条件分布；核函数则说明两个输入有多相似、相应的函数值应有多大相关性。读者可以据此比较参数化与非参数化模型的灵活程度、先验假设和计算开销。

\chapter{参数化与非参数化模型}
\begin{quote}\small
\textit{本章摘要。}本章比较两种表达函数复杂度的方式：参数化模型用有限维参数概括规律，非参数化模型则通过核、邻域或函数分布保留更灵活的形状。高斯过程和核方法贯穿全文，用来说明先验结构、计算代价和小样本性能之间的权衡，并与上一章的贝叶斯推断自然衔接。
\end{quote}
上一章用参数后验解释有限证据下的未知规律，但一个分布再精确，也受它所依附的模型结构限制。本章进一步问：一定要先选有限个回归系数再描述函数吗？能否直接规定不同输入处函数值如何共同变化，再由观测决定局部形状？参数化与非参数化的区别，正落在这个结构选择上。

假设要根据转速预测温升，可以先找转速相近的历史记录取平均，这就是近邻方法的直观起点。核方法用相似程度决定记录之间怎样相互影响，高斯过程进一步描述可能的函数及其概率。方法越灵活，越需要考虑样本是否足够、计算是否承受得起。还要分清一个预测数值和一整个预测分布：两个方法给出相同的平均温升，并不意味着它们对预测有同样的把握。

本章默认读者理解样本平均、协方差以及上一章的后验预测。可先用近邻例子理解“让附近样本参与预测”，再用单观测GP例子核对均值与方差；函数空间和对偶推导用于解释算法形式，不是动手计算这两个例子的先修门槛。

\section{两个概念}

% auxiliary-resource-guide
本章末的“辅助学习材料”列出与核心方法对应的讲解和实践入口，可在阅读相关推导后，按所列小节对照学习。

参数化模型用固定维度参数描述分布，例如线性回归$y=\bm{w}^{\mathsf T}x+b+\varepsilon$。无论样本数如何增加，参数维度仍由输入维度和模型设计决定。非参数化模型不预先固定同样意义上的有限参数维度，而允许模型复杂度随数据增长，例如$k$近邻、高斯过程和Dirichlet过程模型。

\paragraph{近邻回归：由局部常数假设得到估计。}
设$x_i\in\R^d$、$Y_i=f^*(x_i)+\varepsilon_i$，给定输入后噪声独立、均值为零且方差为$\sigma^2$。固定查询$x$，按预先规定距离选出$k$个近邻集合$N_k(x)$。若假设该邻域内函数近似常数$c$，最小化$\sum_{i\in N_k(x)}(Y_i-c)^2$的一阶条件为$2\sum_i(c-Y_i)=0$，所以
\begin{equation}
 \hat f_k(x)=\frac1k\sum_{i\in N_k(x)}Y_i.
\end{equation}
给定训练输入，条件偏差为$k^{-1}\sum_{i\in N_k(x)}f^*(x_i)-f^*(x)$，条件方差为$\sigma^2/k$。若$f^*$为$L_f$-Lipschitz函数，近邻最远距离为$r_k(x)$，则
\begin{equation}
 |\operatorname{Bias}(x)|\leq L_fr_k(x),\qquad
 \E[(\hat f_k(x)-f^*(x))^2\mid X]\leq L_f^2r_k(x)^2+\sigma^2/k.
\end{equation}
增大$k$降低噪声方差，却扩大邻域偏差。若局部密度为正且规则，典型半径量级为$(k/n)^{1/d}$，平衡两项得到$k$约为$n^{2/(d+2)}$、均方误差量级约为$n^{-2/(d+2)}$；这是基于光滑与密度条件的局部量级论证，而非对任何数据的保证。维度越高，邻域收缩越慢，这就是距离方法的维数困难。分类近邻则平均类别指示以估计局部类别概率。

一致性通常要求$k\to\infty$且$k/n\to0$，并附加分布与正则条件；固定$k$一般保留不可消失的方差。距离尺度、无关特征和跨工况混合都可能破坏局部相似假设。朴素查询需约$O(nd)$距离计算，快速索引在高维也可能失效。

\paragraph{近邻手算与“非参数”的含义。}
历史输入为$(0,1,3)$，对应温升为$(2,4,8)$，查询$x=1.2$，按绝对距离找两个邻居，距离依次为$1.2,0.2,1.8$，因此预测为$(4+2)/2=3$；三个邻居都用则为$14/3$。$k$是控制平滑程度的超参数，储存的训练样本仍参与预测，所以“非参数”绝不是“没有需要设置或保存的数值”。多维距离还受单位影响：转速与温度相加求距离前，必须明确尺度或领域距离，不能把数值大的单位误当成更重要的变量。

Lipschitz条件$|f^*(x)-f^*(x')|\leq L_f\|x-x'\|$表示输入靠近时真实均值不会任意跳变。偏差上界来自对邻居逐个应用该条件再取平均，方差$\sigma^2/k$则来自$k$个独立噪声平均。若邻居来自同一条高度相关轨迹，后一个计算需加入协方差；不能仅增加窗口数量就宣称方差按$1/k$下降。

\paragraph{高斯过程示例。}
\begin{figure}[htbp]
\centering
\begin{tikzpicture}[x=1.15cm,y=0.85cm]
 \filldraw[draw=blue!40,fill=blue!12]
 (0.3,1.65) .. controls (1.2,1.65) and (1.7,1.1) .. (2.5,0.72)
 .. controls (2.8,0.62) and (3.2,0.1) .. (3.5,0.02)
 .. controls (3.8,0.02) and (4.2,0.55) .. (4.5,0.67)
 .. controls (4.9,0.8) and (5.2,0.58) .. (5.5,0.42)
 .. controls (6.2,0.6) and (6.8,1.65) .. (7.7,1.65)
 -- (7.7,-1.55)
 .. controls (6.8,-1.55) and (6.2,-0.65) .. (5.5,0.08)
 .. controls (5.2,0.22) and (4.9,0.42) .. (4.5,0.33)
 .. controls (4.2,0.2) and (3.8,-0.38) .. (3.5,-0.32)
 .. controls (3.2,-0.25) and (2.8,0.25) .. (2.5,0.38)
 .. controls (1.7,-0.2) and (1.2,-1.55) .. (0.3,-1.55) -- cycle;
 \draw[->] (0,-1.85)--(8,-1.85) node[right] {$x$};
 \draw[->] (0,-1.85)--(0,2.1) node[above] {$f(x)$};
 \draw[gray,dashed] (2.5,-1.85)--(2.5,1.05);
 \draw[gray,dashed] (5.5,-1.85)--(5.5,1.05);
 \draw[blue!70!black,thick]
 (0.3,0.05) .. controls (1.2,0.05) and (1.7,0.45) .. (2.5,0.55)
 .. controls (2.8,0.435) and (3.2,-0.075) .. (3.5,-0.15)
 .. controls (3.8,-0.18) and (4.2,0.375) .. (4.5,0.5)
 .. controls (4.9,0.61) and (5.2,0.4) .. (5.5,0.25)
 .. controls (6.2,-0.025) and (6.8,0.05) .. (7.7,0.05);
 \foreach \x/\y in {2.5/0.55,3.5/-0.15,4.5/0.5,5.5/0.25}
   \fill[black] (\x,\y) circle (2pt);
 \draw[<->,blue!70!black] (0.8,-1.42)--(0.8,1.52);
 \draw[<->,blue!70!black] (7.2,-1.42)--(7.2,1.52);
 \node[font=\small] at (1.05,2.05) {远处区间宽};
 \node[font=\small] at (6.95,2.05) {远处区间宽};
 \node[font=\small] at (4,1.55) {观测附近区间窄};
 \node[font=\scriptsize] at (4,-1.55) {观测输入范围};
 \draw[blue!70!black,thick] (0.6,-2.4)--(1.1,-2.4);
 \node[font=\scriptsize,anchor=west] at (1.2,-2.4) {后验均值};
 \fill[black] (3.5,-2.4) circle (2pt);
 \node[font=\scriptsize,anchor=west] at (3.7,-2.4) {观测点};
 \filldraw[draw=blue!40,fill=blue!12] (5.4,-2.52) rectangle (5.85,-2.28);
 \node[font=\scriptsize,anchor=west] at (5.95,-2.4) {逐点后验区间};
\end{tikzpicture}
\caption{高斯过程不确定性示意：观测附近的后验区间较窄，远离观测后逐渐变宽。}
\label{fig:kernel-gp-prediction}
\end{figure}
高斯过程（Gaussian process, GP）直接在函数空间上定义先验：任意有限输入集合对应的函数值服从联合高斯分布，记为$f\sim\mathcal{GP}(m,k)$。给定训练输入$X$和观测$y=f(X)+\varepsilon$，其中$\varepsilon\sim\mathcal{N}(0,\sigma^2I)$，新点$x_*$的后验均值为
\begin{equation}
 \mu_* = m(x_*)+k_*^{\mathsf T}(K+\sigma^2I)^{-1}(y-m(X)),\label{eq:kernel-gp-mean}
\end{equation}
后验方差为
\begin{equation}
 \sigma_*^2=k(x_*,x_*)-k_*^{\mathsf T}(K+\sigma^2I)^{-1}k_*.\label{eq:kernel-gp-variance}
\end{equation}
这里$X$的$n$行分别为训练输入的转置，$y\in\R^n$为观测列向量，$m$为先验均值函数，$m(X)=[m(x_1),\ldots,m(x_n)]^{\mathsf T}$，$K_{ij}=k(x_i,x_j)$，$k_*=[k(x_1,x_*),\ldots,k(x_n,x_*)]^{\mathsf T}$。噪声与潜在函数独立。式\eqref{eq:kernel-gp-variance}针对潜在函数值$f(x_*)$；若预测含独立同方差噪声的新观测，还需加上$\sigma^2$。核函数$k$表达“相似输入应有相似函数值”的先验。远离训练样本时，预测方差通常增大，这提供了自然的不确定性信号。

图\ref{fig:kernel-gp-prediction}是相关性随输入距离衰减时的定性示意，并非某组数据的数值拟合结果。黑点集中在两条虚线之间，蓝线表示后验均值，浅蓝色区域表示潜在函数值的逐点后验区间：观测附近较窄，两侧远离观测后逐渐变宽。区间宽度取决于后验标准差，而不是均值曲线的高低。例如，逐点95\%后验区间为$\mu_*\pm1.96\sigma_*$；当$k_*\to0$时，式\eqref{eq:kernel-gp-variance}中的方差缩减项趋于零，区间宽度回到先验所决定的水平，而非无限增大。这里的图示针对相关性随距离衰减的核，不能推广到所有核。
\paragraph{GP后验不是另一个需要记忆的公式。}
设$C=K+\sigma^2I_n$，训练中心化观测$r=y-m(X)$，新点中心化函数值$g_*=f(x_*)-m(x_*)$。GP定义和独立高斯噪声给出
\begin{equation}
 \begin{bmatrix}r\\g_*\end{bmatrix}\sim\mathcal N\left(0,
 \begin{bmatrix}C&k_*\\k_*^{\mathsf T}&k_{**}\end{bmatrix}\right),
 \qquad k_{**}=k(x_*,x_*).
\end{equation}
构造残差$u=g_*-k_*^{\mathsf T}C^{-1}r$，直接计算$\operatorname{Cov}(u,r)=k_*^{\mathsf T}-k_*^{\mathsf T}C^{-1}C=0$。联合高斯中不相关意味着独立，且
\begin{equation}
 \Var(u)=k_{**}-k_*^{\mathsf T}C^{-1}k_*.
\end{equation}
因此给定$r$时$u$仍为同一零均值高斯，恢复先验均值便得到式\eqref{eq:kernel-gp-mean}与式\eqref{eq:kernel-gp-variance}。这个证明同时说明减去的项为何非负：$C$正定时$k_*^{\mathsf T}C^{-1}k_*\geq0$；整个条件方差非负则由联合协方差的半正定性保证。

对$m$个新点，$K_{X*}\in\R^{n\times m}$、$K_{**}\in\R^{m\times m}$，后验协方差为$K_{**}-K_{X*}^{\mathsf T}C^{-1}K_{X*}$。各点之间通常相关，不能把多个边际区间当作整条曲线的同时置信带。若核超参数固定，此协方差只依赖输入而不依赖观测值；异常的$y$不会自动增大方差。

$\sigma^2>0$使半正定$K$对应的$C$正定；无噪声且有重复输入时需处理奇异性。只有当远处$k_*$趋近零时，后验才回到先验均值和先验方差；周期核等结构可能在距离很远时仍高度相关。因此“远离数据更不确定”是常见核下的行为，而非GP定义本身的普适结论。

\paragraph{只有一个观测时，GP公式就是分数运算。}
取零均值先验，一个训练观测$y=2$，训练点与新点的先验方差都为1，交叉协方差$k_*=0.5$，噪声方差$\sigma^2=0.25$。此时$C=1.25$，后验均值为$0.5\times2/1.25=0.8$，潜在函数方差为$1-0.5^2/1.25=0.8$；新观测方差还要加噪声，得到$1.05$。输入与训练点完全不相关时$k_*=0$，预测就回到均值0、函数方差1的先验，而不是断言温升必为零。

在训练点本身预测，交叉协方差变为1，后验均值为$1.6$而非观测值2，函数方差为$0.2$而非零：因为观测可能带噪声，不必完全通过观测点。区间半宽取标准差而非方差，例如前一个新点的逐点95\%可信区间为$0.8\pm1.96\sqrt{0.8}$。这依赖当前核和噪声假设，也不是整条函数同时位于区间带内的95\%保证。

\section{核方法、函数空间与计算边界}
支持向量机展示了核函数如何把线性判别推广到非线性空间，稀疏高斯过程则以诱导变量降低函数后验近似的计算成本\cite{cortes1995svm,titsias2009sparsegp}。核化与贝叶斯函数建模可以共享工具，但优化目标和不确定性含义不同。
\paragraph{联系而非对立。}
参数化和非参数化并不是“现代”和“传统”的划分。神经网络是有限参数模型，但宽网络和核方法之间存在理论联系；贝叶斯神经网络在有限参数空间上保留后验；核方法则把参数学习转化为函数空间中的推断。真正的选择依据是数据量、计算资源、先验结构、预测任务和对不确定性的要求。

\begin{remark}
 当数据量很少而风险可解释性重要时，结构清晰的参数化概率模型往往更合适；当关系复杂且数据充足时，深度参数化模型可以提供更强的表示能力。两者都必须通过验证数据和任务指标检验。
\end{remark}

\paragraph{参数模型中的先验结构。}
参数化并不意味着没有先验。线性模型假设输入与输出的关系可由有限维方向表达；多项式模型增加曲率；状态空间模型把时间依赖写入转移方程。参数越少通常越容易估计，但若结构假设错误，偏差会成为主要误差来源。选择参数化模型就是选择哪些变化被认为值得显式表达。

\paragraph{核方法的函数空间视角。}
核函数$k(x,x')$定义了输入相似性。核岭回归的预测可以写成训练样本核函数的线性组合：
\begin{equation}
 f(x)=\sum_{i=1}^{n}\alpha_i k(x_i,x).
\end{equation}
核技巧把高维特征映射隐含在内积中，但训练通常需要处理$n\times n$核矩阵，样本规模增大时计算和存储成为瓶颈。稀疏近似、随机特征和诱导点（inducing points）可以降低这一代价。下文的Gram矩阵就是有限输入两两核值组成的矩阵$K$；“核矩阵”与“Gram矩阵”在此指同一对象。

\paragraph{什么相似性可以成为核。}
合法实值核必须对称，并对任意有限输入和$a\in\R^n$满足$\sum_{ij}a_ia_jk(x_i,x_j)\geq0$，即Gram矩阵半正定。若$k(x,x')=\phi(x)^{\mathsf T}\phi(x')$，该二次式等于$\|\sum_ia_i\phi(x_i)\|^2$，因此自动满足条件；任意手写距离的负值却不一定是合法核。

平方指数核$k(x,x')=s_f^2\exp[-\frac12\sum_j(x_j-x'_j)^2/\ell_j^2]$用$\ell_j>0$控制各维相关长度，$s_f^2>0$控制函数幅度。长度很大意味着该维变化难以使函数改变，长度很小则允许快速波动。例如一维中两点相距$\ell$时，核值为$s_f^2e^{-1/2}$；相距$2\ell$时为$s_f^2e^{-2}$，长度尺度不是超过后相关性突然归零的硬阈值。核的非负加权和仍为核，可组合平滑、线性或周期结构，但组合之前应确定各分量对应什么物理变化。

\paragraph{核岭回归的表示与系数求解。}
在核对应的再生核Hilbert空间$\mathcal H_k$中，再生性质为
\begin{equation}
 f(x_i)=\langle f,k(x_i,\cdot)\rangle.
\end{equation}
考虑$\lambda>0$时的目标
\begin{equation}
 J(f)=\frac1{2n}\sum_i(y_i-f(x_i))^2+\frac\lambda2\|f\|_{\mathcal H_k}^2.
\end{equation}
把$f$正交分解为训练核函数张成空间中的$f_\parallel$与$f_\perp$。再生性质给出$f_\perp(x_i)=0$，故去掉$f_\perp$不改变训练拟合，却将范数平方减少$\|f_\perp\|^2$。最优解必可写为$f(\cdot)=\sum_i\alpha_i k(x_i,\cdot)$，这就是此处表示结论的推导。

代入后$J(\alpha)=\|y-K\alpha\|^2/(2n)+\lambda\alpha^{\mathsf T}K\alpha/2$，梯度条件为$K[(K+n\lambda I)\alpha-y]=0$。可选稳定解
\begin{equation}
 \alpha=(K+n\lambda I_n)^{-1}y.\label{eq:kernel-ridge-coefficients}
\end{equation}
若$K$奇异，不能在梯度式中随意约去$K$；但式\eqref{eq:kernel-ridge-coefficients}的解仍有效，系数的零空间差异表示同一个函数。正则化保证函数解唯一。零均值GP在$\sigma^2=n\lambda$时与核岭回归均值相同，但核岭目标本身不定义观测噪声或预测方差，不能直接把两者的不确定性等同。

\paragraph{再生核空间先读作带长度规则的函数集合。}
$\mathcal H_k$的元素是函数，内积$\langle f,g\rangle$衡量函数之间的关系，范数$\|f\|_{\mathcal H_k}$衡量核规定的函数复杂度，而不是简单把有限训练点的预测平方求和。$k(x_i,\cdot)$表示固定第一个输入、把第二个输入留给查询的函数；再生性质使“在$x_i$求值”转成一次内积。$f_\perp$与所有训练核函数正交，所以它在训练点为零，却可能在别处改变预测；正则化删掉的是没有训练拟合收益的额外部分。

继续使用单观测$y=2,K=[1]$，取$n\lambda=0.25$，核岭系数为$\alpha=2/1.25=1.6$，新点核值为0.5时预测为0.8，与前面的GP均值一致。$\alpha\in\R^n$每个分量对应一条训练记录，既不是类别概率，也不是输入的$d$维回归系数。核岭计算到这里就结束；还要报告方差，则必须另有概率或区间构造依据。

\paragraph{从有限参数高斯先验得到核。}
令$f(x)=\phi(x)^{\mathsf T}w$，$\phi(x)\in\R^p$、$w\sim\mathcal N(0,S)$。任意有限函数值均为高斯线性变换，协方差为$k(x,x')=\phi(x)^{\mathsf T}S\phi(x')$。令$\Phi\in\R^{n\times p}$，则参数后验给出的均值与函数空间表达相等：
\begin{equation}
 (S^{-1}+\sigma^{-2}\Phi^{\mathsf T}\Phi)^{-1}\sigma^{-2}\Phi^{\mathsf T}
 =S\Phi^{\mathsf T}(\Phi S\Phi^{\mathsf T}+\sigma^2I)^{-1}.
\end{equation}
在$S$正定时，左右两边同乘左侧精度矩阵即可验证恒等式。它说明有限秩核GP也可有固定有限参数表示；所谓非参数性主要针对可随样本扩展的无限秩函数类，而不是“出现核就绝无有限参数”。

\paragraph{核支持向量机的对偶来源。}
对$y_i\in\{-1,1\}$，软间隔分类在特征空间中求
\begin{equation}
 \min_{w,b,\xi}\ \tfrac12\|w\|^2+C\sum_i\xi_i,
 \quad y_i(\langle w,\phi(x_i)\rangle+b)\geq1-\xi_i,
 \quad\xi_i\geq0.
\end{equation}
给两组不等式加入乘子$\alpha_i,\mu_i\geq0$。此处$\alpha\in\R^n$是SVM约束乘子，不是前面核岭回归中可取任意实数的展开系数；两个模型只是沿用同一常见符号。对$w,b,\xi_i$的驻点条件分别为$w=\sum_i\alpha_iy_i\phi(x_i)$、$\sum_i\alpha_iy_i=0$、$C-\alpha_i-\mu_i=0$。代回拉格朗日函数得
\begin{equation}
 \max_{0\leq\alpha_i\leq C,\ \sum_i\alpha_iy_i=0}
 \sum_i\alpha_i-\frac12\sum_{ij}\alpha_i\alpha_jy_iy_jk(x_i,x_j).
\end{equation}
凸性和可行条件使对偶可用于求解；预测得分为$\sum_i\alpha_iy_ik(x_i,x)+b$。若$0<\alpha_i<C$，互补松弛给出对应样本位于间隔边界，可用于计算$b$。若没有$0<\alpha_i<C$的样本，不能用该等式硬求$b$；应由$\alpha_i=0$时$y_if(x_i)\geq1$、$\alpha_i=C$时$y_if(x_i)\leq1$所给出的上下界选择可行$b$，并检查数值容差。非零系数对应支持向量。间隔得分不是故障概率，仍须单独校准；核矩阵半正定则是凸优化性质的重要前提。

\paragraph{支持向量机中的松弛变量与成本参数。}
$\xi_i$允许第$i$个点违反单位间隔约束。取$C>0$时，最优解不会无故增大松弛量，故$\xi_i=\max\{0,1-y_if(x_i)\}$。在这个最小松弛约定下，$\xi_i=0$表示满足间隔，$0<\xi_i<1$表示分类正确但离边界太近，$\xi_i=1$表示位于判别边界，$\xi_i>1$才对应得分落在错误一侧。任意尚未优化的可行松弛量可能更大，不能直接用它判断分类对错。增大$C$会更重地处罚违反间隔；这与核岭中增大$\lambda$加强平滑不是同一个方向的调节。乘子$\alpha_i$是约束在最优解处的权重，互补松弛说明未触及约束的点可以有零权重，这就是只保留支持向量仍能计算判别函数的原因。这里$f(x)>0$选择正类，不代表$f(x)$是$[0,1]$内的概率。

\paragraph{高斯过程的边际似然。}
高斯过程不仅给出预测均值，还能通过边际似然选择核参数和噪声尺度：
\begin{equation}
 \log p(y\mid X)=-\frac12(y-m(X))^{\mathsf T}K_y^{-1}(y-m(X))-\frac12\log|K_y|-\frac n2\log(2\pi).\label{eq:kernel-marginal-likelihood}
\end{equation}
其中$K_y=K+\sigma^2I$，$|K_y|$表示行列式；$m\equiv0$时得到零均值先验的简式。
二次型衡量残差相对协方差的大小，行列式项来自高斯密度的归一化；两项共同形成拟合与分布扩散程度的折中，不能把第二项单独当作通用的复杂度排序。边际化所体现的这种 Occam 效应与上一章的模型证据相连，矩阵分解成本通常为$O(n^3)$。

\paragraph{边际似然的积分、导数与数值实现。}
由$f_X\sim\mathcal N(m_X,K)$及$y=f_X+\varepsilon$，独立高斯相加后$y\sim\mathcal N(m_X,C)$，$C=K+\sigma^2I$；把这个$n$维高斯密度取对数就得到式\eqref{eq:kernel-marginal-likelihood}。这一步已经积分掉函数值，并不是将最优函数值代入似然。

设超参数为$\eta$，$r=y-m_X$、$a=C^{-1}r$。利用$d\log|C|=\operatorname{tr}(C^{-1}dC)$和$dC^{-1}=-C^{-1}(dC)C^{-1}$，逐项微分得到
\begin{equation}
 \frac{\partial\log p(y\mid X)}{\partial\eta}
 =\frac12\operatorname{tr}\left[(aa^{\mathsf T}-C^{-1})\frac{\partial C}{\partial\eta}\right]
  +\left(\frac{\partial m_X}{\partial\eta}\right)^{\mathsf T}a.
\end{equation}
核长度、幅度和噪声可据此优化，正尺度通常取对数参数化。目标不保证凸，应检查多起点和边界。将超参数估计后固定属于经验贝叶斯，预测区间未包含超参数不确定性；充分贝叶斯做法还要对其后验积分。

数值上分解$C=LL^{\mathsf T}$，解$Lv=r$、$L^{\mathsf T}a=v$，可计算$r^{\mathsf T}C^{-1}r=v^{\mathsf T}v$及$\log|C|=2\sum_i\log L_{ii}$。预测时解$Lu=k_*$，方差为$k_{**}-u^{\mathsf T}u$，无需存储逆矩阵。浮点误差可能出现极小负方差，需结合残差和条件数诊断；不能把显著负值简单截断后忽略错误。训练分解为$O(n^3)$时间、$O(n^2)$存储，单点方差通常需$O(n^2)$求解。

训练完成后，应保存训练侧预处理参数、训练输入、均值与核配置、噪声尺度、Cholesky因子$L$及解向量$a$。新点只应用已保存的预处理并计算交叉核，再求均值与方差；不能在每批查询上重新拟合标准化。若加入新观测或更换核参数，则需更新相应分解与解向量，不能混用旧缓存。

\paragraph{行列式项不是独立的万能复杂度标尺。}
若$C=U\operatorname{diag}(c_j)U^{\mathsf T}$，中心化边际负对数似然为$\frac12\sum_j[(u_j^{\mathsf T}r)^2/c_j+\log c_j]+\frac n2\log(2\pi)$。增大某方向方差会减小残差惩罚，却增加归一化代价，两项共同决定折中。不能脱离数据与核谱，把行列式项简单解释成参数越多就惩罚越大。

\paragraph{非参数模型的计算边界。}
非参数模型可以表达复杂结构，却不代表复杂度无限且没有代价。kNN 的查询成本随样本数增加；核方法受核矩阵限制；高斯过程的后验更新需要处理越来越大的协方差矩阵。实际部署常用近似推断、原型压缩、局部模型或分层采样。

\paragraph{诱导变量如何降低代价。}
选$m\ll n$个诱导输入$Z$，令$u=f(Z)\in\R^m$。为简化取零均值，$K_{uu}\in\R^{m\times m}$可逆，$K_{Xu}\in\R^{n\times m}$。高斯条件公式给出
\begin{equation}
 p(f_X\mid u)=\mathcal N(K_{Xu}K_{uu}^{-1}u,\ K-Q),\qquad
 Q=K_{Xu}K_{uu}^{-1}K_{uX}.
\end{equation}
这里$K-Q$是诱导变量未解释的条件协方差，不能只因为使用低秩表示就断言它为零。取近似联合$q(f_X,u)=p(f_X\mid u)q(u)$，ELBO化为$\E_q\log p(y\mid f_X)-\KL(q(u)\|p(u))$。对高斯噪声，先在$p(f_X\mid u)$下取期望，得到
\begin{equation}
 \E[\log p(y\mid f_X)\mid u]
 =\log\mathcal N(y\mid K_{Xu}K_{uu}^{-1}u,\sigma^2I)
 -\frac{\operatorname{tr}(K-Q)}{2\sigma^2}.
\end{equation}
再对自由$q(u)$优化，利用$\sup_q\{\E_q\log h(u)-\KL(q\|p)\}=\log\int h(u)p(u)\,du$，可得
\begin{equation}
 \mathcal L=\log\mathcal N(y\mid0,Q+\sigma^2I)
 -\frac{\operatorname{tr}(K-Q)}{2\sigma^2}.
\end{equation}
迹修正惩罚遗漏的函数变异，因此这不是随意用$Q$替换$K$的精确模型。借助低秩恒等式主要成本可降至$O(nm^2+m^3)$、存储$O(nm)$；诱导点过少或位置不当会遗漏局部故障结构。$m=n$、诱导输入等于训练输入且矩阵适定时，$Q=K$恢复精确训练证据。

\paragraph{诱导点不是额外测量点。}
$Z$是人为选择的代表输入位置，$u=f(Z)$仍未知，需要由$q(u)$描述；不是在这些位置真的测到了新温升。$m$表示诱导点个数，不再表示均值函数$m(x)$。矩阵$Q$表达经这些代表位置传递的相关性，$\operatorname{tr}(K-Q)$为未解释的各训练点条件方差之和，迹就是矩阵对角线元素相加。因而更多计算上的代表点可以改善近似，却不能增加真实观测证据。上一章的ELBO在这里仍是对数证据下界，固定近似结构的后验区间也仍需独立检查覆盖。

\paragraph{随机特征的另一路近似。}
对适当的平稳正定核，写$k(x,x')=s_f^2\E_\omega\cos(\omega^{\mathsf T}(x-x'))$，其中$s_f^2=k(x,x)>0$。独立采样$\omega_j$和$b_j\sim\operatorname{Uniform}(0,2\pi)$，设$\phi_j(x)=s_f\sqrt{2/m}\cos(\omega_j^{\mathsf T}x+b_j)$。积化和差给出$2\E_b[\cos(u+b)\cos(v+b)]=\cos(u-v)$，因为含$2b$的余弦在整周期上的平均为零；再对频率和$m$项取期望，得到$\E[\phi(x)^{\mathsf T}\phi(x')]=k(x,x')$。这把核问题变成$m$维线性问题，但有限$m$引入随机近似误差；适用的频率分布由核决定，不能对任意相似性随意采样。

\paragraph{参数化与非参数化的混合。}
这些方法也可以组合使用。例如先用神经网络把振动窗口压成低维特征，再用 GP 预测并给出模型内的不确定性。这里若把编码器固定，GP 的区间通常不会自动包含编码器本身的不确定性。另一些任务只需在冻结特征上重新训练线性头，或用状态空间先验约束序列预测。组合前要说明每部分负责什么、哪些误差仍未计入，不能因为接上概率模型就认为整个系统已经校准。

\paragraph{自测：相同均值不等于相同区间。}
零均值GP只有一个观测$y=3$，$K=[2]$，噪声方差为1；查询点先验方差为2，交叉核值为1。求预测均值、潜在函数方差与新观测方差，并说明核岭何时给出相同均值。

\textit{参考解。}$C=[3]$，均值为$1\times3/3=1$，函数方差为$2-1/3=5/3$，新观测方差为$8/3$。同核的核岭取$n\lambda=1$时均值相同，但其目标本身不提供这两个方差。两点先验核矩阵为$\left[\begin{smallmatrix}2&1\\1&2\end{smallmatrix}\right]$，特征值为1和3，故本题协方差设置合法。

\paragraph{模型选择案例。}
小样本温度预测可以从线性回归、核岭回归和高斯过程开始，比较留出设备上的误差、区间覆盖率、训练时间和增量更新成本。若数据规模扩大且输入包含图像或长序列，再考虑深度模型。这个顺序能够先验证任务是否可预测，再决定是否需要更复杂的表示学习。

下列材料可用于对照近邻、支持向量机与高斯过程回归的实现；普通近邻方法不需要先构造核矩阵。
\section*{辅助学习材料}
\begingroup
\emergencystretch=3em
\begin{itemize}
\item \textbf{Support Vector Machines}；作者/平台：scikit-learn Developers；英文；官方文档。重点阅读 linear/nonlinear kernels、C、epsilon 和 dual formulation 小节，帮助理解本章的 RKHS、核岭和 SVM 对偶。\href{https://scikit-learn.org/stable/modules/svm.html}{在线阅读}。
\item \textbf{Gaussian Processes}；作者/平台：scikit-learn Developers；英文；官方文档。重点阅读 GP regression、kernels、noise level 和 predictive uncertainty，帮助理解核函数、协方差矩阵与不确定性。\href{https://scikit-learn.org/stable/modules/gaussian\_process.html}{在线阅读}。
\item \textbf{高斯过程回归：推导，实现和理解}；知乎；中文解说。对照核矩阵与条件高斯分布，检查预测均值和方差；标题与主题据必应索引，原页受限。\href{https://zhuanlan.zhihu.com/p/104601803}{在线阅读}。
\item \textbf{高斯过程回归 (Gaussian Processes Regression, GPR)简介}；CSDN；中文博客。先复算少量观测点的预测，再区分潜在函数与含噪观测的区间；标题与主题据必应索引，原页受限。\href{https://blog.csdn.net/HelloWorldTM/article/details/126980872}{在线阅读}。
\item \textbf{03-KNN算法}；黑马程序员；bilibili；中文课程分集。选看《黑马程序员3天快速入门python机器学习》第21集，对照本章近邻方法的距离、邻居数和非参数性质，比较标准化前后的邻域；该分集不替代高斯过程推导。2026年10月4日公开元数据所列整套课程累计播放1,793,918次（不是本分集计数）；已核对目录与计数，未逐帧观看。\href{https://www.bilibili.com/video/BV1nt411r7tj/?p=21}{观看分集}。
\end{itemize}
\endgroup
\paragraph{本章小结。}
参数化模型以明确结构换取估计效率和解释性，非参数化模型以更灵活的函数空间换取计算成本。二者不是互斥阵营，关键是让模型复杂度、数据量、先验知识和不确定性需求相互匹配。

从局部常数的近邻估计，到再生核空间的正则解，再到高斯联合分布的条件化，本章展示了三种共享样本信息的机制。核规定哪些输入相关，正则或先验规定允许多大变化，近似方法规定在预算内保留哪些变化；任何一个环节都可能成为偏差来源。相同预测均值不意味着相同概率解释，较宽函数空间也不意味着在未覆盖工况下必然诚实地表达不确定性。

到这里，我们已经从参数分布走到了函数分布。下一章把读数按时间排好：平均温度相同，持续升温和持续降温却可能需要不同处理。序列网络、状态空间模型和时间核各有办法利用这种顺序。比较时继续检查它们如何表达依赖、需要多少计算，以及给出的不确定性究竟针对潜在规律还是未来读数。


\part{深度学习模型}
本部分从维护流程中尚未解决的问题出发。统计特征遗漏演化过程时，需要序列表示；风险分数不能指出缺陷位置时，需要视觉定位；检测结果不足以确定复检步骤时，需要检索有效规程；单路证据含糊时，需要跨模态对齐；真实故障覆盖不足时，才讨论生成数据。后续章节沿用第\ref{chap:pipeline}章的信息可用性原则，但分别定义预测、定位、问答和合成的验收对象，不能用一种任务的高分证明另一种任务也有效。
\chapter{序列建模模型}
\begin{quote}\small
\textit{本章摘要。}本章围绕一个核心问题展开：如何把历史信息压缩成对未来有用的状态。先介绍 Recurrent Neural Network (RNN)、Long Short-Term Memory (LSTM) 和 Gated Recurrent Unit (GRU) 的递归记忆，再解释注意力与 Transformer 如何改变信息访问路径，随后由状态空间滤波转入连续系统的离散化与 Mamba 的选择性记忆，讨论并行扫描、长序列计算、多步预测和异常告警。模型比较最终要检查时间边界、提前量、非平稳性与资源消耗。
\end{quote}
窗口统计量可以描述振动强度，却可能抹去冲击出现的顺序。本章由此研究如何表示截至$t$的历史信息。故障概率、剩余寿命和未来轨迹是三种不同输出，后文的分钟级轨迹预测用于说明多步误差，不能与二十四小时故障概率直接比较。递归状态、注意力和状态空间模型的取舍，取决于输出需要保留哪些历史证据。
前两章已经区分参数化表示、函数先验和预测不确定性，本章进一步追问：当观测按时间到达时，哪些信息应随状态一起保留下来？一个温度均值相同的窗口，可能对应平稳运行，也可能对应持续升温后骤降；只把窗口当作无序样本，就丢掉了区分这两种过程的依据。因此，序列建模不是在回归器前多加一个时间下标，而是规定历史如何进入当前判断，以及新的观测怎样修正已有判断。

读公式时，先看模型怎样保存历史。递归模型把过去压进一个状态，注意力可以回头读取历史记录，卷积和状态空间模型则按各自的结构传递信息。再看这个状态是否还带有协方差：一组确定数值和一个状态分布，能表达的不确定性不同。最后核对输出到底是什么。预测未来温度曲线、预测故障概率和给当前窗口打异常分数，需要不同的标签，也对应不同的维护动作。
\section{序列与状态}

% auxiliary-resource-guide
本章末的“辅助学习材料”列出与核心方法对应的讲解和实践入口，可在阅读相关推导后，按所列小节对照学习。

图\ref{fig:sequence-recurrent-state}中，纵向箭头把当前输入写入状态，横向箭头将历史传到下一时刻，最终预测因而依赖状态保留下来的信息。
\begin{figure}[htbp]
\centering
\begin{tikzpicture}[node distance=0.55cm and 0.6cm]
 \node[idi input, minimum width=1.45cm] (x1) {$x_{t-2}$};
 \node[idi input, right=of x1, minimum width=1.45cm] (x2) {$x_{t-1}$};
 \node[idi input, right=of x2, minimum width=1.45cm] (x3) {$x_t$};
 \node[idi state, above=0.75cm of x1, minimum width=1.45cm] (h1) {$h_{t-2}$};
 \node[idi state, above=0.75cm of x2, minimum width=1.45cm] (h2) {$h_{t-1}$};
 \node[idi state, above=0.75cm of x3, minimum width=1.45cm] (h3) {$h_t$};
 \draw[idi arrow] (x1)--(h1); \draw[idi arrow] (x2)--(h2); \draw[idi arrow] (x3)--(h3);
 \draw[idi arrow] (h1)--(h2); \draw[idi arrow] (h2)--(h3);
 \node[above=0.25cm of h2,font=\scriptsize] {状态沿时间向前传递并压缩历史};
 \node[idi output, right=0.7cm of h3, minimum width=1.45cm] (y) {$\hat y_{t+1}$};
 \draw[idi arrow] (h3)--(y);
\end{tikzpicture}
\caption{递归序列模型把历史压缩为状态；状态维度和更新机制决定哪些信息能够持续保留。}
\label{fig:sequence-recurrent-state}
\end{figure}
LSTM 的门控记忆用于改善长时依赖中的信息与梯度传播\cite{hochreiter1997lstm}。与之相对，线性高斯状态估计中的 Kalman 滤波显式传播状态均值和协方差\cite{kalman1960filter}；二者不能因为都有隐状态就被视为相同的推断方法。
序列数据$(x_1,\ldots,x_T)$具有顺序依赖。状态空间模型用隐状态$h_t$压缩历史信息：
\begin{equation}
 h_t=g(h_{t-1},x_t),\qquad p(y_t\mid x_{\leq t})=p(y_t\mid h_t).
\end{equation}
状态的核心作用不是记住所有历史，而是在给定任务下保留预测未来所需的信息。

初读时只需先掌握向量、矩阵乘法和损失函数：沿着“输入窗口—历史状态—任务头—标签”走通一次预测，再读梯度与概率推导。链式法则解释模型怎样学会保留历史，高斯条件分布解释新观测怎样修正不确定性；二者分别服务于训练与状态估计，不必先掌握滤波才能理解RNN。

\paragraph{先把一条序列送进模型。}
这里的隐状态不是另一个传感器，而是网络内部保存的一组数。设每分钟读取温度、转速、振动三个量，则$x_t\in\R^3$；把连续四分钟按行排放得到$4\times3$的输入，八台设备构成一批时是$8\times4\times3$。批次维只是同时计算的样本数，不应在不同设备之间传递状态。若状态宽度取$m=2$，每读一分钟就把三个观测与两个旧状态数一起变成两个新状态数。初始$h_0$可设为零，或学习一个共同初值，但不是把每台设备的最后状态混在一起。

状态之后还需要任务头，即把表示变成业务输出的最后一层。例如$\hat y_{t+1}=w^{\mathsf T}h_t+b$输出一个温度，$\sigma(w^{\mathsf T}h_t+b)$输出指定未来窗口的故障概率，矩阵投影则可输出多个提前量。前述$p(y_t\mid x_{\leq t})=p(y_t\mid h_t)$是模型采用的压缩假设，不是任意有限维状态都能无损保存历史的定理。训练用标签计算误差并调整共享权重；推理通常固定权重，只随观测更新状态。

\paragraph{递归状态与门控记忆。}
最基本的循环神经网络（Recurrent Neural Network，RNN）为$h_t=\tanh(W_hh_{t-1}+W_xx_t+b)$；这里$W_x$把当前输入映射到状态空间，$W_h$把旧状态映射到新状态，$b$是偏置。反向传播时，梯度包含多个$W_h$的乘积，可能出现梯度消失或爆炸。LSTM通过门控缓解这一问题：
\begin{align}
 i_t&=\sigma(W_ix_t+U_ih_{t-1}),& f_t&=\sigma(W_fx_t+U_fh_{t-1}),\label{eq:sequence-lstm-input-forget}\\
 o_t&=\sigma(W_ox_t+U_oh_{t-1}),& \tilde c_t&=\tanh(W_cx_t+U_ch_{t-1}),\label{eq:sequence-lstm-output-candidate}\\
 c_t&=f_t\odot c_{t-1}+i_t\odot\tilde c_t,& h_t&=o_t\odot\tanh(c_t).
\end{align}
这里$\sigma(a)=1/(1+e^{-a})$把数压到零与一之间，$\tanh$把数压到$(-1,1)$，$\odot$表示对应分量分别相乘，不是矩阵乘法。输入门$i_t$决定写多少候选内容，遗忘门$f_t$决定留多少旧内容，输出门$o_t$决定公开多少记忆；细胞状态$c_t$是存储通道，$h_t$是供下一步和任务头使用的表示。

可先手算一个分量：若$c_{t-1}=2,f_t=0.9,i_t=0.2,\tilde c_t=0.5,o_t=1/2$，则$c_t=0.9\times2+0.2\times0.5=1.9$，$h_t=\tanh(1.9)/2\approx0.478$。候选值受限不代表累积记忆也限于$[-1,1]$。这些门值是网络根据输入算出的，不是人为判定“重要”后再开门。细胞状态提供了更稳定的信息通道。GRU用更少的门实现类似的更新：常见写法把更新门$z_t$和重置门$r_t$分别用于混合旧状态与候选状态、控制候选状态读取旧状态；不同资料的门符号和先后约定可能互换，比较公式时应先核对定义。

门控的必要性可以从训练梯度理解。
对简单RNN，隐藏状态对早期状态的导数包含
\begin{equation}
 \frac{\partial h_T}{\partial h_t^{\mathsf T}}=J_TJ_{T-1}\cdots J_{t+1},\qquad
 J_s=\diag\left(1-\tanh^2(a_s)\right)W_h.
\end{equation}
这里$a_s=W_hh_{s-1}+W_xx_s+b$，Jacobian（雅可比矩阵）的第$(i,j)$项为输出第$i$分量对输入第$j$分量的偏导；矩阵因子按式中次序排列不可交换，作用于早期扰动列向量时则由最右侧因子先起作用；反向传播列梯度时使用这一 Jacobian 的转置。
若乘积中的最大奇异值长期小于一，梯度趋向零；若大于一，梯度可能爆炸。梯度裁剪在$g\ne0$时把$g$替换为$g\min(1,\tau/\|g\|)$，其中阈值$\tau>0$；$g=0$时直接保留零，避免除零。它可以控制爆炸但不能修复长期消失。LSTM的细胞状态通过门$f_t$保留一条近似加法路径，使信息保存时间由门值控制。

非线性导数$1-\tanh^2(a_s)$在饱和区域接近零，因此仅控制循环权重的尺度仍不足以保证长期梯度。对于长流，可以每$K$步截断计算图以控制训练内存，前向状态仍可携带更早的信息，但早期输入不能通过被截断的路径直接贡献梯度。窗口长度与截断长度因而是不同的实验变量。实现截断时，应保留状态数值但切断它与上一计算图的梯度连接；切断梯度不等于把状态清零。若下一批已换设备或不再时间连续，则应按样本边界重置，不能沿用上一批状态。

\paragraph{从单步导数到跨时间反向传播。}
设$x_t\in\R^p$、$h_t\in\R^m$，于是$W_x\in\R^{m\times p}$、$W_h\in\R^{m\times m}$、$b\in\R^m$。令总损失为$\mathcal L=\sum_{t=1}^T\ell_t(h_t)$，用列向量$g_t=\nabla_{h_t}\mathcal L$表示包含所有后续损失的梯度。链式法则逐步给出
\begin{align}
 g_T&=\nabla_{h_T}\ell_T,\\
 g_t&=\nabla_{h_t}\ell_t+J_{t+1}^{\mathsf T}g_{t+1},\\
 g_t&=\sum_{s=t}^T(J_sJ_{s-1}\cdots J_{t+1})^{\mathsf T}\nabla_{h_s}\ell_s,
\end{align}
其中$s=t$时空乘积为$I$。再令$\delta_t=\diag(1-h_t\odot h_t)g_t$，共享权重的梯度不是某一步的梯度，而是所有出现位置的累加：
\begin{equation}
 \nabla_{W_h}\mathcal L=\sum_{t=1}^T\delta_t h_{t-1}^{\mathsf T},\qquad
 \nabla_{W_x}\mathcal L=\sum_{t=1}^T\delta_t x_t^{\mathsf T}.
\end{equation}
由矩阵范数的次乘性，$\|J_s\cdots J_{t+1}\|_2\leq\prod_{r=t+1}^s\|J_r\|_2$。若每项至多为$\rho<1$，远期损失传回的上界按$\rho^{s-t}$衰减；若上界大于一，仅意味着可能放大，不能据此断言实际梯度必然爆炸，因为奇异方向会随时间改变。

在LSTM中，$x_t\in\R^p$，各门、$c_t,h_t\in\R^m$，所以各$W$为$m\times p$、各$U$为$m\times m$；式\eqref{eq:sequence-lstm-input-forget}和式\eqref{eq:sequence-lstm-output-candidate}省略了可训练偏置。固定门值时，沿细胞状态直接路径有
\begin{equation}
 \left.\frac{\partial c_T}{\partial c_t^{\mathsf T}}\right|_{\text{直接路径}}
 =\diag\left(\prod_{s=t+1}^T f_s\right).
\end{equation}
乘积按元素计算。若某维遗忘门近似恒为$f\in(0,1)$，经过$k$步保留$f^k$，其衰减到初值$e^{-1}$所需步数为$-1/\log f$。例如$f=0.9$时这一时间尺度约为$9.49$步，$f=0.99$时约为$99.50$步；它是衰减到约$36.8\%$的时间，不是到零的时间。这解释了门值与记忆时间尺度的关系，但不是完整Jacobian：门还通过$h_{t-1}$依赖早期状态，其他路径仍会放大或衰减梯度。门控改善可训练性，不保证无限记忆。

\paragraph{GRU 的参数化。}
GRU使用更新门和重置门
\begin{align}
 z_t&=\sigma(W_zx_t+U_zh_{t-1}),\\
 r_t&=\sigma(W_rx_t+U_rh_{t-1}),\\
 \tilde h_t&=\tanh(W_hx_t+U_h(r_t\odot h_{t-1})),\\
 h_t&=(1-z_t)\odot h_{t-1}+z_t\odot\tilde h_t.
\end{align}
更新门接近零时保留旧状态，接近一时写入候选状态。它比LSTM参数更少，却牺牲了对细粒度记忆通道的独立控制。

\section{注意力、Transformer 与长依赖}
图\ref{fig:sequence-attention-query}区分历史键用于匹配、历史值用于输出的两条路径；当前查询通过归一化权重决定读取哪些历史内容。
\begin{figure}[htbp]
\centering
\begin{tikzpicture}[node distance=0.8cm and 0.7cm]
 \node[idi input, minimum width=1.8cm] (q) {当前查询\\$q_i$};
 \node[idi state, right=of q, minimum width=2.0cm] (score) {与历史键\\计算相似度};
 \node[idi process, right=of score, minimum width=2.0cm] (weight) {softmax 权重\\$a_{ij}$};
 \node[idi output, right=of weight, minimum width=2.0cm] (out) {加权值\\$o_i$};
 \draw[idi arrow] (q)--(score); \draw[idi arrow] (score)--(weight); \draw[idi arrow] (weight)--(out);
 \node[idi input, below=0.75cm of score, minimum width=2.0cm] (memory) {历史键和值\\$K,V$};
 \draw[idi arrow] (memory)--(score); \draw[idi arrow] (memory)--(out);
\end{tikzpicture}
\caption{注意力把“压缩全部历史”改为“按查询重新读取历史”，代价是保存和计算位置间的相互作用。}
\label{fig:sequence-attention-query}
\end{figure}
递归模型把全部历史压入有限维状态。当预测需要重新访问某次冲击或某段工况时，这种压缩可能成为瓶颈。注意力改为按当前查询读取历史表示，其缩放点积形式为\cite{vaswani2017attention}
\begin{equation}
 \operatorname{Attention}(Q,K,V)=\softmax\left(\frac{QK^{\mathsf T}}{\sqrt{d_k}}\right)V.
\end{equation}
直接访问缩短了远距离位置之间的信息路径，但标准全局注意力的计算量随序列长度呈二次增长。Transformer 将多头注意力、位置编码、前馈网络、残差连接和归一化组合起来。因此，从 GRU 转向 Transformer 是以保存和访问历史的成本换取更灵活的信息选择，而非无条件提高预测能力。

\paragraph{注意力的逐位置解释。}
对于第$i$个查询，注意力权重为
\begin{equation}
 a_{ij}=\frac{\exp(q_i^{\mathsf T}k_j/\sqrt{d_k})}{\sum_{r=1}^{T}\exp(q_i^{\mathsf T}k_r/\sqrt{d_k})},
 \qquad o_i=\sum_j a_{ij}v_j.
\end{equation}
因此每个输出是值向量的加权平均。查询表示“我现在需要什么”，键表示“我具有什么索引”，值表示“被取出的内容”。多头注意力使用多组投影，让不同头分别学习局部邻接、长距离指代或语义匹配。

\paragraph{一次可以手算的读取。}
把一个采样点或一个时间片转换成向量后，该向量也称为一个token，即模型处理的一个序列单元。对当前查询，假设两个可见历史位置的缩放得分为$0,\log3$，则softmax先取指数得到$1,3$，再除以总和得到$1/4,3/4$。若对应的标量值为$10,20$，输出就是$17.5$，而不是直接复制得分最大的$20$。键用于算权重，值用于形成输出，所以两者可以有不同维度。若第二个位置实际来自未来，必须在归一化前屏蔽它，此时权重成为$1,0$，输出为$10$。

注意力输出仍是表示，并不直接等于预测。每个头完成一次读取，多头结果按特征维拼接并投影回模型宽度；残差把原表示加回来，前馈网络再对每个位置应用同一组非线性变换。归一化控制特征尺度。这些步骤使模型既能交换位置间信息，也能加工当前位置的特征；头的具体分工由训练形成，不保证某个头必然学到某种语义。

\paragraph{缩放、归一化与访问成本的推导。}
采用逐行存放位置表示的约定。各矩阵形状为
\begin{equation}
 X\in\R^{T\times d},\quad W_Q,W_K\in\R^{d\times d_k},\quad W_V\in\R^{d\times d_v},
\end{equation}
因此$Q,K\in\R^{T\times d_k}$、$V\in\R^{T\times d_v}$。为理解缩放，暂设$q_i,k_j$的坐标彼此独立、均值为零、方差为一，则
\begin{equation}
 \E(q_i^{\mathsf T}k_j)=0,\qquad
 \Var(q_i^{\mathsf T}k_j)=\sum_{r=1}^{d_k}\Var(q_{ir}k_{jr})=d_k.
\end{equation}
因此除以$\sqrt{d_k}$使得分方差约为一，避免维度增加时softmax过早饱和。这是初始化尺度的解释，不是对训练后相关特征的严格分布假设。记缩放得分为$s_{ij}$，由商法则得到
\begin{equation}
 \frac{\partial a_{ij}}{\partial s_{ir}}=a_{ij}(\mathbf1_{j=r}-a_{ir}).
\end{equation}
当一个权重已接近一时，其余权重及大量导数接近零，模型便难以调整访问对象。softmax使$a_{ij}\geq0$且$\sum_j a_{ij}=1$，单头输出因而位于值向量的凸包内；但后续投影、残差和前馈变换不受这个凸包约束。

计算$QK^{\mathsf T}$需$O(T^2d_k)$次标量运算，计算$AV$需$O(T^2d_v)$；显式保存分数需$O(T^2)$空间。若无位置编码、无遮罩，令$P$为位置置换矩阵，则$Q'=PQ,K'=PK,V'=PV$，有$A'=PAP^{\mathsf T}$及$O'=PO$。这叫置换等变，而不是输出不随置换变化；网络本身无法从内容之外识别先后顺序。因果遮罩和位置编码正是为打破这种不足而加入的结构。

\paragraph{位置编码与因果遮罩。}
没有位置编码和遮罩的自注意力具有前述置换等变性：调换输入行会相应调换输出行，而不是让输出数表完全不变。要表达明确的先后或时间距离，需加入位置信息。正弦位置编码使用不同频率的正余弦函数；可学习位置嵌入直接为每个位置分配向量。因果遮罩令$j>i$的注意力得分为$-\infty$，从而保证第$i$个位置只能使用$1$到$i$的信息。解码时的 Key--Value (KV) cache 保存已经计算的键和值，使每次生成新 token 不必重复计算全部历史。

\section{时间序列预测与状态估计}
单步预测学习$p(x_{t+1}\mid x_{\leq t})$，多步预测可递归调用模型，也可一次输出整个未来窗口。递归策略容易累积误差，直接多步预测避免把自身输出反复作为输入，但参数量取决于输出头与共享结构。概率预测可以输出分位数，使用 pinball loss
\begin{equation}
 \ell_\tau(y,\hat y)=\max\{\tau(y-\hat y),(\tau-1)(y-\hat y)\}
\end{equation}
学习条件分位数，从而形成预测区间。

\paragraph{未来窗口与分位数的具体含义。}
若输入过去四分钟的三个传感器，直接预测未来两分钟的温度，则单样本输入为$4\times3$，标签和点预测都是长度为二的向量。若每个提前量输出$0.1,0.5,0.9$三个分位数，输出形状就是$2\times3$，不是六个不同传感器。训练时将预测与实际未来温度比较；推理时未来标签尚不存在，只能输出估计。递归模型第一步把当前观测作为输入，第二步起需要反馈预测；若还使用未来转速等外生变量，必须确认部署时已有计划值，否则也需要预测它们。

以$\tau=0.9$、真实值$y=10$为例，预测$8$的损失为$0.9\times2=1.8$，预测$12$的损失为$0.1\times2=0.2$。因此高分位数更重罚低估，学的是分布靠上的位置，而不是给均值简单加一个固定安全余量。实际覆盖率要在留出数据中统计真实值落入区间的比例。

\paragraph{多步误差为何会积累。}
为隔离递归效应，设真实向量状态满足$x_{t+1}=f_*(x_t)+\eta_{t+1}$，预测器递归使用$\hat x_{t+h}=f_\theta(\hat x_{t+h-1})$，初值$\hat x_t=x_t$。假设在所考察区域，$f_\theta$为$L$-Lipschitz，即$\|f_\theta(a)-f_\theta(b)\|\leq L\|a-b\|$，且$\|f_\theta(x)-f_*(x)\|\leq\delta$、$\|\eta_t\|\leq\epsilon$。插入并减去$f_\theta(x_{t+h-1})$，可得
\begin{align}
 e_h&=\|\hat x_{t+h}-x_{t+h}\|\\
 &\leq L e_{h-1}+\delta+\epsilon\\
 &\leq(\delta+\epsilon)\sum_{j=0}^{h-1}L^j.
\end{align}
当$L\ne1$时，上界为$(\delta+\epsilon)(L^h-1)/(L-1)$，当$L=1$时为$h(\delta+\epsilon)$。即便一步拟合误差小，迭代的敏感性仍可能放大误差；$L<1$时该上界则有界。这是有界噪声条件下的示意推导，高斯噪声不能直接使用统一的确定性$\epsilon$。直接多输出方法避免反馈自身误差，却不消除未来不可预测性，也不必然需要更多参数，参数量取决于共享骨干和输出头设计。

\paragraph{分位数损失为何学习分位数。}
固定历史信息$\mathcal F_t$，设未来标量$Y$的条件分布函数为$F$，候选预测为$a$，分位水平$\tau\in(0,1)$。当分布连续且允许交换微分与积分时，条件期望损失为
\begin{align}
 R(a)&=(1-\tau)\int_{-\infty}^{a}(a-y)\,dF(y)
       +\tau\int_a^\infty(y-a)\,dF(y),\\
 R'(a)&=(1-\tau)F(a)-\tau(1-F(a))=F(a)-\tau.
\end{align}
故最优解满足$F(a)=\tau$；有原子质量时改用次梯度，条件为$F(a^-)\leq\tau\leq F(a)$。上下分位数$\tau=\alpha/2$与$1-\alpha/2$由此形成名义覆盖率$1-\alpha$的区间。有限样本、模型失配和分布漂移会破坏实际覆盖率，不同分位数独立训练还可能交叉。每个提前量分别覆盖并不等于整条未来轨迹同时覆盖，后者需要联合分布或额外的同时覆盖校准。

\paragraph{工艺参数异常预警。}
可以把设备传感器流切成长度为$T$的窗口，输入过去$T$步，预测未来$H$步的正常轨迹。异常分数可以是预测误差、概率负对数似然或隐状态与正常原型的距离。阈值应在验证集上按告警提前量和误报预算选择，而不能直接取训练误差的某个分位数。

\paragraph{状态空间模型与滤波。}
序列模型可以拆成状态转移和观测生成两部分：
\begin{align}
 p(h_t\mid h_{t-1},u_t)&=\mathcal N(Ah_{t-1}+Bu_t,Q),\\
 p(x_t\mid h_t)&=\mathcal N(Ch_t,R).
\end{align}
这里$u_t$可以表示控制量，$Q$和$R$分别表示状态演化噪声与观测噪声。卡尔曼滤波先用转移方程预测状态，再用新观测修正预测。深度 RNN 可以看作学习非线性转移函数，但如果忽略观测噪声和状态不确定性，就无法直接给出滤波意义下的置信区间。

从这一统一视角看，现代时间序列模型是在不同位置放置“记忆”和“预测分布”。DeepAR 用自回归 RNN 输出下一时刻分布的参数，适合许多设备共享统计规律的场景；Neural Basis Expansion Analysis for Interpretable Time Series Forecasting (N-BEATS) 用堆叠的前馈块逐步拟合趋势和季节性，能够通过基函数解释预测来源，但需要固定长度窗口。Temporal Convolutional Network (TCN) 用因果空洞卷积并行处理整个窗口，感受野由卷积核和空洞率显式控制，适合高吞吐批处理；若故障信号必须在每个新采样到达时更新，带显式状态的 GRU 或状态空间模型通常更自然。

Transformer 系列则把时间点或时间片看作 token。Informer 使用稀疏概率注意力和序列蒸馏，减少长序列的二次成本；Autoformer 先将序列分解为趋势和季节项，再用自相关机制寻找周期模式；FEDformer 在频域或低秩结构中进行全局建模；PatchTST 把连续时间片切成 patch，以更少 token 保留局部形状，并在规则采样的长预测任务中形成强基线。PatchTST 对缺失时间戳和事件驱动数据并不天然稳健，因此应显式加入时间间隔、缺失掩码和设备身份，而不能直接把图像 Vision Transformer (ViT) 的 patch 机制照搬到所有传感器流。

当序列长度使注意力矩阵无法存储时，结构化状态空间模型提供了另一种计算路径。Structured State Space Sequence Model (S4) 将连续时间状态方程离散化并用结构化参数高效计算长卷积，兼顾长记忆与高效计算；选择性状态空间模型 Mamba 进一步让输入控制状态更新、遗忘和输出，使模型可以对不同事件选择不同的记忆时间尺度。它们适合超长振动流、日志流和在线扫描，但状态压缩也可能丢失需要精确回看的单次事件。比较 Transformer、TCN、S4 和 Mamba 时，应同时报告短暂冲击的召回、长周期趋势误差、流式延迟、显存和窗口外推能力。

\begin{table}[htbp]
\centering
\caption{序列模型路线与工业选择}
\label{tab:sequence-model-selection}
\scriptsize
\setlength{\tabcolsep}{3pt}
\resizebox{\linewidth}{!}{%
\begin{tabular}{p{2.4cm}p{3.8cm}p{3.6cm}p{4.1cm}}
\hline
模型路线 & 记忆机制 & 计算与部署特征 & 更适合的工业问题 \\
\hline
LSTM/GRU & 显式递归状态与门控 & 顺序计算，流式更新简单 & 在线状态跟踪；不规则到达须显式编码时间间隔 \\
TCN & 因果空洞卷积与固定感受野 & 可并行训练，窗口外推有限 & 高频批处理、固定提前量预测 \\
Transformer & 全局注意力与位置编码 & 长度平方级注意力，缓存可降低解码重复 & 长距离依赖、数据量较大、离线分析 \\
PatchTST/\allowbreak Informer & patch 或稀疏注意力 & 减少 token，依赖规则采样假设 & 长窗口、多变量、规则传感器流 \\
S4/Mamba & 结构化或选择性状态空间 & 结构化卷积或选择性扫描，支持流式更新 & 长振动流、日志流、低延迟连续监控 \\
\hline
\end{tabular}}
\end{table}

表\ref{tab:sequence-model-selection}按记忆机制和部署成本比较序列模型，选型时应将流式更新、固定窗口与长历史访问的要求对应到具体工业任务。
\paragraph{多头注意力的张量形状。}
设批大小为$B$、序列长度为$T$、模型维度为$d$、头数为$H$，每头维度为$d_h=d/H$，此处假设$d$可被$H$整除；本段$H$表示头数，不是预测提前量。输入形状为$B\times T\times d$，投影后重排为$B\times H\times T\times d_h$，注意力分数形状为$B\times H\times T\times T$。最后一个$T\times T$矩阵解释了长序列的显存瓶颈：即使参数量不大，二次注意力也可能无法处理高频传感器流。

\begin{figure}[htbp]
\centering
\begin{tikzpicture}[>=Latex, node distance=0.65cm, every node/.style={font=\small}]
 \node[draw, rounded corners, fill=blue!10, minimum width=2.2cm, minimum height=0.75cm] (x) {$x_{1:T}$};
 \node[draw, rounded corners, fill=orange!12, right=of x, minimum width=2.4cm, minimum height=0.75cm] (qkv) {$Q,K,V$ 投影};
 \node[draw, rounded corners, fill=green!12, right=of qkv, minimum width=2.5cm, minimum height=0.75cm] (att) {多头注意力};
 \node[draw, rounded corners, fill=yellow!18, right=of att, minimum width=2.2cm, minimum height=0.75cm] (ffn) {前馈网络};
 \node[draw, rounded corners, fill=red!10, right=of ffn, minimum width=2.0cm, minimum height=0.75cm] (y) {$h_{1:T}$};
 \draw[->, thick] (x)--(qkv); \draw[->, thick] (qkv)--(att); \draw[->, thick] (att)--(ffn); \draw[->, thick] (ffn)--(y);
 \draw[->, dashed] (x.north) .. controls +(0,0.85) and +(0,0.85) .. node[above,font=\scriptsize]{注意力子层残差} ($(att.east)!0.5!(ffn.west)$);
 \node[below=0.2cm of ffn,font=\scriptsize] {归一化与前馈残差略};
\end{tikzpicture}
\caption{Transformer 编码块的主要信息流。}
\label{fig:sequence-transformer-flow}
\end{figure}

图\ref{fig:sequence-transformer-flow}沿投影、注意力和前馈变换展示编码路径，虚线则标出绕过注意力子层的残差通路；二者共同形成各位置的输出表示。
\paragraph{位置编码的选择。}
绝对位置编码把位置直接加到 token 表示中，简单但长度外推能力有限；相对位置编码直接描述 token 间距离，更适合长度变化；旋转位置编码通过复数相位或二维旋转改变查询和键，使点积携带相对位置信息。选择位置编码时要同时考虑训练长度、部署长度和缓存机制，不能只比较短序列上的损失。

\paragraph{时间序列中的非平稳性。}
趋势、季节性、工况切换和传感器老化会使$p(x_{t+1}\mid x_{\le t})$随时间变化。差分、归一化、分解和时间特征可以减少非平稳性，但也可能删除故障前的有用趋势。滚动验证比随机交叉验证更适合检验未来性能；模型更新还应区分概念漂移与测量系统改变。

\section{训练、预测与工业告警}
\paragraph{序列模型实验矩阵。}
至少应比较：持久性或季节性基线、线性模型、树模型、RNN/\allowbreak TCN 和 Transformer。固定预测提前量后，再改变窗口长度、模型容量和缺失处理。评价不仅包括 MAE 或交叉熵，还应包含延迟、显存、吞吐、异常事件召回率和预测区间覆盖率。只有当复杂模型在关键约束下仍有稳定收益，才说明其结构提供了有效归纳偏置。

以轴承剩余寿命（Remaining Useful Life，RUL）预警为例，首先把原始振动、温度、转速和负载按照设备时钟对齐，并保留每一路的采样时间、缺失标记和传感器版本。随后用设备级而非窗口级划分数据：同一轴承的相邻窗口必须进入同一个集合，测试设备和测试时间段完全留出。输入可以由高频振动的统计量、频谱 patch 和低频工况变量组成，目标则根据维护流程选择连续 RUL、未来$H$分钟内的失效概率或两者的多任务输出。

训练阶段先比较持久性、线性外推、随机森林和 GRU，再引入 TCN、PatchTST 或 Mamba。所有模型使用相同的预测提前量、相同的训练设备和相同的标准化统计量；对 RUL 回归可以采用 Huber 损失，对失效概率采用加权交叉熵，对预测区间增加分位数损失。若失效事件很少，不能只随机抽取正常窗口，否则模型会在类别先验上获得虚假的高准确率；应按设备和事件构造训练批次，并报告事件级召回与每台设备的误报数。

上线时，模型每分钟更新一次状态，告警状态机可以要求连续三个窗口超过阈值才进入高等级告警；“至少提前20分钟”是依据事后确认的故障时刻验收的目标，不能作为在线状态机已知的触发条件。输入缺失超过五分钟时，系统暂停风险结论并生成数据质量告警；恢复后先执行状态重置或冷启动窗口，避免把长时间缺失后的填充值当作真实趋势。监控面板同时展示预测轨迹、置信区间、模型版本、输入覆盖率和人工确认结果，形成从模型输出到维护动作的可追溯的处理记录。

\paragraph{滤波、平滑与预测。}
滤波只使用截至$t$的观测估计$h_t$，平滑则允许使用未来观测修正过去状态。在线故障预警只能使用滤波信息，离线诊断可以使用平滑信息。若把平滑后的状态用于训练在线预警器，模型可能隐式读取未来，导致离线指标偏高。实验报告应明确状态估计使用的信息边界。

\paragraph{卡尔曼滤波的两步更新。}
在线线性高斯状态空间模型中，预测步骤为
\begin{align}
 \hat h_{t\mid t-1}&=A\hat h_{t-1\mid t-1}+Bu_t,\\
 P_{t\mid t-1}&=AP_{t-1\mid t-1}A^{\mathsf T}+Q.
\end{align}
观测到$x_t$后，计算创新$r_t=x_t-C\hat h_{t\mid t-1}$（新观测减去预测观测的残差，并非技术创新）和增益
\begin{equation}
 K_t=P_{t\mid t-1}C^{\mathsf T}(CP_{t\mid t-1}C^{\mathsf T}+R)^{-1},
\end{equation}
再更新$\hat h_{t\mid t}=\hat h_{t\mid t-1}+K_tr_t$。创新序列可以用于检测模型失配：若其长期偏离零或自相关显著，说明转移、噪声或传感器模型需要重新估计。

\paragraph{把滤波更新算成一个温度例子。}
设当前温度的预测均值为$20$，预测方差$P^-=4$，传感器噪声方差$R=1$，且$C=1$。读到$22$后，创新是$22-20=2$，增益为$4/(4+1)=0.8$，修正均值为$21.6$，修正方差为$(1-0.8)4=0.8$。若传感器更不可靠，$R=16$，增益降为$0.2$，均值只移到$20.4$。方差的单位是温度单位的平方；它衡量估计不确定性，不是另一条温度曲线。滤波中的“更新”是在固定系统参数下修正状态分布，与神经网络训练时更新权重不是同一操作。

\paragraph{由高斯条件化推导滤波增益。}
设$h_t\in\R^m,x_t\in\R^p,u_t\in\R^q$，则$A\in\R^{m\times m},B\in\R^{m\times q},C\in\R^{p\times m}$；$Q,P\in\R^{m\times m}$，$R\in\R^{p\times p}$。假设转移噪声与观测噪声均为零均值、彼此独立、跨时间独立，并与先前状态独立。记预测均值和协方差为$m^-,P^-$。线性变换的均值与协方差规则先给出前述预测步骤，再给出联合分布
\begin{equation}
 \begin{bmatrix}h_t\\x_t\end{bmatrix}\Bigm|x_{<t}
 \sim\mathcal N\left(
 \begin{bmatrix}m^-\\Cm^-\end{bmatrix},
 \begin{bmatrix}P^-&P^-C^{\mathsf T}\\CP^-&CP^-C^{\mathsf T}+R\end{bmatrix}\right).
\end{equation}
为看清条件化的来源，令$e=h_t-m^-$、$r=x_t-Cm^-$、$S=CP^-C^{\mathsf T}+R$，并写出修正后的误差$e^+=e-Kr$。要求它与已经利用的创新不相关，有
\begin{equation}
 \operatorname{Cov}(e-Kr,r)=P^-C^{\mathsf T}-KS=0,
 \qquad K=P^-C^{\mathsf T}S^{-1}.
\end{equation}
联合高斯使不相关等价于独立，因此$\E(h_t\mid x_t,x_{<t})=m^-+Kr$，且后验协方差为
\begin{align}
 P^+&=P^--P^-C^{\mathsf T}S^{-1}CP^-\\
 &=(I-KC)P^-(I-KC)^{\mathsf T}+KRK^{\mathsf T}.\label{eq:sequence-joseph-covariance}
\end{align}
式\eqref{eq:sequence-joseph-covariance}称为Joseph形式，更适合在有限精度计算中保持对称性和半正定性。这里要求$S$可逆，实际实现应解线性方程而非显式求逆。若$p=m=1,C=1$，则$K=P^-/(P^-+R)$：观测越噪，修正越小；状态预测越不确定，越相信新观测。在线性高斯且参数正确时，$r^{\mathsf T}S^{-1}r$服从自由度为$p$的卡方分布；用它设阈值仍需检查噪声、参数估计与重复告警假设，不能直接把名义尾概率当成事件级误报率。

\subsection{从固定状态空间到 Mamba 的选择性记忆}
\label{sec:sequence-mamba}
前面的状态空间模型（State Space Model，SSM）用概率分布描述状态与观测，并通过卡尔曼滤波估计不确定性。本节沿用状态转移的语言，但把 SSM 用作可训练的确定性序列层：它输出表示而非状态后验，不自动携带协方差或预测区间。先说明固定线性系统为何能形成长记忆，再解释 Mamba 如何根据输入选择写入、保留和读取的信息。以下讨论原始 Mamba 的选择性 SSM 与模块结构，不把后续 Mamba-2 的算法细节混入其中。

\paragraph{连续动力学怎样成为离散记忆。}
以$\tau$表示连续时间，$t$仍表示离散位置。设状态$h(\tau)\in\R^m$、输入$u(\tau)\in\R^q$、层输出$v(\tau)\in\R^r$，连续系统为
\begin{equation}
 \frac{d h(\tau)}{d\tau}=A_c h(\tau)+B_c u(\tau),\qquad
 v(\tau)=Ch(\tau)+Du(\tau).
 \label{eq:sequence-ssm-continuous}
\end{equation}
其中$A_c\in\R^{m\times m}$、$B_c\in\R^{m\times q}$、$C\in\R^{r\times m}$、$D\in\R^{r\times q}$均不随位置变化，$D$为输入到输出的直接通路。这里$u_t$是由观测$x_t$提取的输入特征，不限于前文滤波模型中的控制量；$v_t$是中间表示，不与任务标签$y_t$混用。式\eqref{eq:sequence-ssm-continuous}是线性时不变（Linear Time-Invariant，LTI）系统。

在长度为$\Delta>0$的每个区间内保持输入不变，即零阶保持，积分因子法给出
\begin{align}
 h_t&=\bar A h_{t-1}+\bar B u_t,\qquad v_t=Ch_t+Du_t,
 \label{eq:sequence-ssm-discrete}\\
 \bar A&=\exp(\Delta A_c),\qquad
 \bar B=\int_0^\Delta\exp(sA_c)B_c\,ds.
 \label{eq:sequence-ssm-zoh}
\end{align}
只有$A_c$可逆时，才可把式\eqref{eq:sequence-ssm-zoh}的积分写为$A_c^{-1}(\bar A-I)B_c$；积分形式始终有效。反复代入式\eqref{eq:sequence-ssm-discrete}可得
\begin{equation}
 v_t=C\bar A^t h_0+\sum_{j=1}^t C\bar A^{t-j}\bar B u_j+Du_t.
 \label{eq:sequence-ssm-convolution}
\end{equation}
式\eqref{eq:sequence-ssm-convolution}表明，固定线性系统既可递推，也可使用只依赖时间差的卷积核$G_k=C\bar A^k\bar B$计算。S4通过结构化参数高效表示和计算长核；卷积实现的复杂度取决于具体算法，不能将其一概说成严格线性。若$A_c$的特征值实部为负，则$\bar A$的特征值模小于一，长期记忆衰减，但非正规矩阵仍可能短期放大扰动。固定核对同一时间差使用相同权重，不能仅因某次冲击内容重要就改变记忆规则，这正是选择性更新要解决的限制。

\paragraph{输入依赖的步长、写入与读取。}
为明确维度，先看一个标量输入通道$u_t\in\R$，其状态$h_t\in\R^m$。令$\xi_t\in\R^d$为当前位置经过因果局部处理后的特征，定义
\begin{align}
 \Delta_t&=\operatorname{softplus}(w_\Delta^{\mathsf T}\xi_t+b_\Delta)>0,
 \label{eq:sequence-mamba-selection}\\
 B_t&=W_B\xi_t\in\R^m,\qquad C_t=W_C\xi_t\in\R^m,\qquad
 A_c=\diag(-\exp(a_1),\ldots,-\exp(a_m)).\nonumber
\end{align}
这里$w_\Delta\in\R^d$、$W_B,W_C\in\R^{m\times d}$以及$a_i$是训练得到、推理时共享的参数；$\operatorname{softplus}(s)=\log(1+e^s)$。式\eqref{eq:sequence-mamba-selection}是便于推导的单通道参数化，实际模块可用低秩投影产生各通道步长，并在通道间共享部分$B_t,C_t$。输入依赖的是$\Delta_t,B_t,C_t$，不是所有参数；尤其$A_c$和直接通路系数$D$不必随输入变化。

逐步冻结这些系数，精确零阶保持对应
\begin{align}
 \bar A_t&=\exp(\Delta_t A_c),\qquad
 \bar B_t^{\mathrm{ZOH}}=\int_0^{\Delta_t}\exp(sA_c)B_t\,ds,
 \label{eq:sequence-mamba-zoh}\\
 h_t&=\bar A_t h_{t-1}+\bar B_t u_t,\qquad
 v_t=C_t^{\mathsf T}h_t+Du_t.
 \label{eq:sequence-mamba-recurrence}
\end{align}
式\eqref{eq:sequence-mamba-recurrence}中的$\bar B_t$必须与所选离散化一致。原始 Mamba 的官方选择性扫描实现对输入项采用$\bar B_t=\Delta_t B_t$，同时保留$\bar A_t=\exp(\Delta_t A_c)$；它并非直接使用式\eqref{eq:sequence-mamba-zoh}的精确输入积分。两者在$\|\Delta_t A_c\|$较小时一阶相符，不能将简化输入项写成精确零阶保持的恒等式。

对第$i$个状态分量，保留系数为$\exp(-\Delta_t\exp(a_i))$。较小的$\Delta_t$使旧状态较少衰减，较大的$\Delta_t$使更新更快；$B_t$控制写入方向和强度，$C_t$控制当前读出。步长因此是一种学习到的记忆时间尺度，不天然等于传感器的物理采样间隔。对于不规则采样，还需显式提供实际时间间隔或约束连续时间模型。选择也不是“自动丢弃噪声”的保证：是否保留短时冲击取决于训练目标与数据，必须通过事件级验证检查。

尽管式\eqref{eq:sequence-mamba-recurrence}对给定系数下的旧状态仍是线性的，系数本身依赖输入，整个输入到输出映射已不是 LTI 系统。历史输入的权重同时依赖其后的选择，不能再用式\eqref{eq:sequence-ssm-convolution}中那个固定的时间差卷积核替代。

\paragraph{先手算两步，再理解扫描。}
取单维$h_0=0$，第一步$F_1=0.5,b_1=2$，第二步$F_2=0.8,b_2=1$。依次递推得到$h_1=2,h_2=0.8\times2+1=2.6$。合并成一步则是$h_2=(0.8\times0.5)h_0+(0.8\times2+1)=0.4h_0+2.6$。第三步既可先与第二步合并，也可后接到前两步的结果上，答案相同；交换前两步却得到$0.5\times1+2=2.5$，说明顺序不能交换。这里的$b_t$是已经计算好的写入量，不是神经网络所有层都共用的固定偏置。

离散化也可以手算：取$A_c=-1,B_t=1,u_t=2,h_{t-1}=0,\Delta_t=\log2$，则保留系数为$1/2$。精确积分的输入系数为$1-e^{-\Delta_t}=1/2$，得到$h_t=1$；简化输入系数为$\Delta_t$，得到$h_t=2\log2\approx1.386$。这不是计算错误，而是两种离散化在步长不小时确实不同。学习参数可以适应所用实现，但不能因此把两种公式称为相等。

\paragraph{为何能递推推理，也能并行训练。}
记$F_t=\bar A_t$、$b_t=\bar B_t u_t$，每步状态更新都是仿射映射$h\mapsto F_t h+b_t$。定义按时间先后合并两个映射的运算
\begin{equation}
 (F_2,b_2)\star(F_1,b_1)
   =(F_2F_1,\ F_2b_1+b_2).
 \label{eq:sequence-mamba-scan}
\end{equation}
函数复合的结合性使式\eqref{eq:sequence-mamba-scan}满足结合律，单位元为$(I,0)$；它通常不满足交换律，不能打乱时间顺序。对前缀$(F_t,b_t)\star\cdots\star(F_1,b_1)$做并行扫描，即可得到所有$h_t$。由于选择参数由输入特征而非尚未求得的$h_{t-1}$生成，训练时可先并行计算各位置系数，再合并前缀；这与一般非线性 RNN 的顺序状态依赖不同。理想并行前缀算法具有$O(\log T)$深度，但总工作量仍随$T$增长；实际硬件上的分块、访存和同步决定实际延迟，不能据此宣称运行时间只随$\log T$增长。

\paragraph{从选择性 SSM 到 Mamba 模块。}
原始 Mamba 不只是一个状态递推式。输入先经线性投影分成内容分支与门控分支；内容分支经过短核因果深度卷积和 SiLU 激活，再进入选择性 SSM。这里深度卷积指每个通道分别沿时间使用短核，SiLU为$\operatorname{SiLU}(a)=a\sigma(a)$，两者分别负责局部混合与非线性变换。扫描结果与另一分支的 SiLU 输出逐元素相乘，经输出投影后接入残差路径；堆叠模块配合归一化构成深层网络。因果卷积只读取当前及过去的位置，承担局部混合，SSM承担长历史压缩，通道投影混合不同传感器特征。流式部署必须同时维护 SSM 状态和短卷积缓存。输出还须接任务头，才能成为未来轨迹、故障概率或分位数预测；选择性状态本身不是概率置信度。

\paragraph{与 Transformer 的计算和记忆取舍。}
设层内通道数为$d$，每通道状态维度为$m$。对角转移下，选择性扫描核心的总工作量为$O(Tdm)$，流式状态占用为$O(dm)$，与已处理历史长度无关；固定短卷积核还需要与核宽有关的缓存。密集通道投影通常另需$O(Td^2)$工作，不能把整个模块的成本仅写成$O(T)$而省略宽度与状态大小。训练仍需存储或重算随序列长度增长的激活，硬件感知扫描通过融合、分块与重计算减少访存，不意味着训练显存为常数。

相较之下，标准全局注意力的交互计算为$O(T^2d)$，自回归 KV 缓存随历史增长到$O(Td)$，每个新位置还需读取历史键和值。FlashAttention 等精确算法可避免显式存储完整注意力矩阵，但不消除全局交互的二次算术量。Mamba 在长流与受限缓存下有结构性优势，Transformer 则能按当前查询重新访问具体历史位置；压缩进有限状态的证据一旦丢失，Mamba 无法像注意力那样直接回读原记录。精确事件检索、短序列、小批量或缺少优化扫描内核时，不能预设 Mamba 更准确或更快；混合注意力与 SSM 也是可比较的选择。

\paragraph{工业实例：保留冲击还是保留工况。}
沿用轴承案例，可将因果振动片段特征、当前温度、转速、负载、缺失掩码与实际时间间隔组成$x_t$，经投影和局部卷积形成$\xi_t$。稳定工况中，一部分状态维度可学习较慢衰减以积累磨损趋势；负载切换时，选择机制可调整写入，避免把正常升温直接解释为退化；稀疏冲击是否进入长期记忆则由故障目标约束。这只是机制解释，不是已经观察到的实验结果，更不能把某个门值当作物理故障原因。

在相同设备划分、时间边界、预测提前量和参数预算下，应比较 GRU、Transformer 与 Mamba，并消融输入依赖的$\Delta_t,B_t,C_t$，检查收益究竟来自选择机制还是局部卷积与额外容量。除轨迹误差外，报告短冲击召回、事件级误报、流式延迟和峰值显存。每台设备独立维护状态；跨设备拼接、长缺失后的恢复和模型更新必须按预定规则重置或预热，训练时也要模拟这些边界。这样，式\eqref{eq:sequence-mamba-recurrence}中的长记忆才与上线时真正可用的历史一致。

\paragraph{长序列的计算策略。}
Transformer 的二次注意力使高频工业流难以直接处理。常见策略包括下采样、分块、局部注意力、稀疏注意力、线性注意力和层级编码器。下采样降低成本却可能删除短暂故障；局部注意力保留邻域细节却限制远距离依赖；层级模型通过低频状态连接长时间尺度。选择策略应由故障持续时间和提前量决定。

\paragraph{教师强制与部署解码。}
教师强制规定的是输入使用真实历史，而不是限定损失种类；连续值任务也可用均方误差或概率损失。离散词元任务中，teacher forcing 下的交叉熵衡量“给定真实历史时的下一步预测”，并不等于自由运行时的轨迹质量。评估多步预测时应让模型使用自己的历史输出，记录误差随预测步长的增长。对连续值预测，还应区分递归策略中的误差累积和直接多输出策略中的共享表示误差。

\paragraph{非平稳序列的分层验证。}
若设备老化导致均值、方差或周期变化，随机切分会把未来状态混入训练。时间滚动验证可以模拟上线过程：用较早区间训练，在后续区间验证，再向前滚动。若存在新设备部署，还应增加设备留出维度。结果应按时间、设备、工况和传感器版本分层，以区分概念漂移和测量系统变化。

\paragraph{异常分数的选择。}
预测残差适合发现可预测模式的偏离，负对数似然还会考虑预测方差，隐状态距离适合表示空间中已有正常原型的偏离。三种分数的阈值含义不同：残差阈值依赖量纲，似然阈值依赖概率校准，距离阈值依赖表示几何。可以组合多个分数，但必须在验证集上固定组合规则，不能根据测试事件逐个调阈值。

\paragraph{序列模型的资源消耗。}
报告模型时应记录每步延迟、端到端窗口延迟、显存、参数量和吞吐。一个离线速度快的模型，如果等待完整窗口才能输出，可能不满足提前预警要求。流式 RNN 可以逐步更新状态，Transformer 则可能需要缓存和窗口重算。计算预算是序列建模问题的一部分，而非部署阶段才补充的细节。

\paragraph{滚动预测的完整设置。}
设传感器以一分钟为间隔采样，使用过去$T=120$步预测未来$H=30$步。样本构造应令输入窗口为$[t-T+1,t]$，标签窗口为$[t+1,t+H]$，并保证同一故障事件不会同时出现在训练和测试窗口。若窗口高度重叠，随机切分会使相邻片段泄漏到不同集合；更稳妥的做法是按时间块或设备划分，并在块之间留出间隔。

评价时分别计算$h=1,5,15,30$的误差，观察模型是否只在短提前量有效。若采用递归预测，必须在每一步输入模型自己的预测；若采用直接多输出头，则应说明每个提前量是否共享参数。除 MAE 外，还应报告故障前提前量、告警持续时间和每台设备每天误报数，因为这些量直接决定维护团队的工作量。

\paragraph{从预测到异常告警。}
轨迹预测误差只有在真实观测到达后才能计算，不能用未来误差触发当前预警。设$\hat x_{t\mid t-h}$及$\hat\sigma_{t\mid t-h}$是在$t-h$时刻对$t$的预测和标准差。对一个传感器通道，时刻$t$可计算的标准化残差为
\begin{equation}
 s_t=\max_{1\leq h\leq H}\frac{|x_t-\hat x_{t\mid t-h}|}{\hat\sigma_{t\mid t-h}+\epsilon}.
\end{equation}
其中$\epsilon>0$须与标准差使用相同单位。即使单个提前量的标准化残差近似正态，取绝对值再对多个提前量取最大值后，也不能直接套用标准正态阈值；应按实际组合规则在验证集上定阈值。这个分数检测当前观测相对于历史预测的偏离，不等于未来故障概率。若要提前预警，还需验证偏离与后续故障的关系，或直接训练未来事件概率输出头。
为了避免一个瞬时噪声触发告警，可以要求连续$m$个窗口超过阈值，或在事件级别合并相邻告警。阈值和$m$只能在验证数据上根据漏报成本、误报预算和最小提前量确定。测试集只用于一次性报告，不能在看到故障时间后反复调整规则。

传感器掉线时，要先区分设备真的异常，还是根本没有可靠读数。直接填零可能在曲线上制造一次突变。可以用缺失掩码注明哪些值没测到；若沿用上次读数，还应限制最多沿用多久。长时间缺失应触发数据质量告警，并按预定规则暂停判断或转人工处理。这样才不会把通信故障误报成设备故障。

\paragraph{在线部署检查表。}
上线前至少应验证初始状态、窗口边界、时区和时间戳顺序；上线后监控输入缺失率、窗口延迟、预测分布、异常分数和告警确认结果。模型更新时保留旧版本作为影子模型，比较同一批输入上的分数变化。这样可以区分真实工况变化、数据管线故障和模型版本引入的行为改变。

\paragraph{序列模型的基线与消融。}
持久性预测$\hat x_{t+h}=x_t$是许多连续传感器任务的强基线；具有周期性的过程还应加入季节性复制和线性外推。若深度模型不能稳定超过这些基线，应先检查窗口构造、目标定义和数据泄漏，而不是立即增加层数。消融实验可以依次去掉时间特征、缺失掩码、注意力、残差或多步损失，以判断性能来自真正的结构还是额外的参数量。

\paragraph{不规则采样与时间间隔。}
许多工业事件并非严格等间隔。直接把事件编号当作时间步，会把十秒间隔和十分钟间隔视为同样的动态变化。可以显式输入$\Delta t$，或使用连续时间状态转移；但填补缺失时间点也会改变噪声结构。实验应分别比较规则重采样、时间间隔特征和事件驱动建模，并报告重采样造成的延迟和短时故障损失。

\paragraph{多变量序列的因果边界。}
多个传感器之间可能存在同步采集、控制回路和滞后影响。相关性注意力权重不能直接解释为因果关系，因为共同工况变量会同时改变多个通道。若预测目标是控制动作，还应区分可用的外生变量、动作之后才产生的变量和部署时不可获得的维修记录。特征可用性应按时间戳审计，并在数据字典中固定定义。

\paragraph{概率预测的评价。}
分位数预测使用 pinball loss，完整分布预测还可以使用连续排序概率得分。评价时既要检查尖锐性，也要检查校准：过宽的预测分布可能轻易覆盖真实值，却没有调度价值。对多步预测，应分别评估每个提前量，并观察区间是否随着预测跨度合理变宽。若预测分布出现不合理的收缩，通常意味着模型忽略了过程噪声或训练目标过度追求均值误差。

以下非论文材料覆盖递归状态、Transformer 与 SSM/Mamba：课程提供整体框架，图解解释信息流，官方教程和项目文档用于核对实现。
\section*{辅助学习材料}
\begingroup
\emergencystretch=3em
\begin{itemize}
\item \textbf{CS224N: Natural Language Processing with Deep Learning}；作者/平台：Stanford University；英文；课程讲义、作业与视频。重点看 self-attention、Transformers 和 evaluation，帮助理解本章注意力、位置编码与评价边界；课程内容随学年更新。\href{https://web.stanford.edu/class/cs224n/}{课程主页}。
\item \textbf{NLP From Scratch: Classifying Names with a Character-Level RNN}；作者/平台：Sean Robertson / PyTorch；英文；代码教程。阅读字符编码、循环网络、训练与混淆矩阵部分，观察隐藏状态如何逐字符更新以及如何完成序列分类。\href{https://docs.pytorch.org/tutorials/intermediate/char\_rnn\_classification\_tutorial.html}{在线阅读}。
\item \textbf{The Illustrated Transformer}；作者/平台：Jay Alammar / 作者博客；英文；图解教程。对照本章查询、键、值、多头注意力和位置编码，沿编码器与解码器的信息流理解 Transformer；注意区分解码器因果遮罩与编码器的双向访问。\href{https://jalammar.github.io/illustrated-transformer/}{在线阅读}。
\item \textbf{A Visual Guide to Mamba and State Space Models}；作者/平台：Maarten Grootendorst / 作者博客；英文；图解教程。配合第\ref{sec:sequence-mamba}节，从连续与离散 SSM、递推与卷积读到输入依赖选择和扫描；并行扫描依赖结合律而非交换律，时间顺序不能改变。\href{https://www.maartengrootendorst.com/blog/mamba/}{在线阅读}。
\item \textbf{Mamba}；作者/平台：state-spaces / GitHub 官方项目文档；英文；实现与使用说明，非论文阅读条目。重点看 Usage 下的 Selective SSM、Mamba Block，沿文档链接核对扫描接口与模块实现；区分原始 Mamba 与后续版本，部署前检查平台和加速内核要求。\href{https://github.com/state-spaces/mamba\#usage}{使用文档}。
\end{itemize}
\endgroup
\paragraph{自测与参考答案。}
在$t$时刻，能否用刚生成的$t+5$预测与其真实值之差触发告警？若把训练计算图每20步切断，是否必须每20步清零状态？参考答案：未来真实值尚未到达，不能计算该残差；只能使用当前观测与过去对当前的预测比较，或使用另行训练的未来事件概率。截断只切断梯度连接，连续序列可保留状态数值；换设备或独立序列才按边界规则重置。

\paragraph{本章小结。}
窗口统计不足以保留演化顺序时，可以引入递归状态；状态压缩丢失关键历史时，可以改用注意力；历史访问成本过高时，再比较分块、卷积和状态空间。Mamba 用输入依赖的步长、写入与读取改变固定 SSM 的记忆规则，并以结合性的仿射复合支持扫描；它降低长流缓存需求，却不能消除有限状态的信息损失。选定表示后，还必须区分未来轨迹、故障概率和已到达观测的异常分数。梯度乘积说明长记忆为何难以训练，高斯条件化说明新观测为何应按不确定性加权，分位数风险与误差递推则说明多步输出不能只用一步均方误差评价。这些推导共同指向一个原则：表示能访问的信息、训练所优化的风险和上线时真正可计算的量必须一致。

振动曲线可以提示设备发生了变化，但通常不能指出表面裂纹在哪里。下一章转向图像，用卷积、区域匹配和像素分割寻找缺陷位置。数据仍需独立测试，也仍要检查预测时能否取得；不过标签将变为图像类别、检测框或像素标记，因此不能拿定位成绩直接代替未来故障预测的成绩。

\chapter{视觉模型}
\begin{quote}\small
\textit{本章摘要。}本章从图像的局部空间结构出发，说明卷积为何形成层级表示、残差为何改善深层优化，再逐步扩展到检测、分割、Vision Transformer 和自监督学习。最后将评价从单张图像分数推进到跨相机、跨批次、拒识和数据追溯，强调视觉模型必须服务于可验收的工业决策。
\end{quote}
风险预警告诉维护人员何时需要关注设备，却不能指出缺陷位于哪里。本章转向巡检图像和表面质检：先判断是否存在缺陷，再定位区域，最后估计边界与尺寸。这些任务使用图像级、框级或像素级标签，与未来故障标签不同。拆检后才能获得的照片可以支持离线诊断，却不能直接作为拆检前预警的输入。
上一章中，时间顺序决定同一组读数代表稳定运行还是异常演化；本章中，空间关系决定相同的颜色和纹理代表正常表面还是裂纹。把图像展平成无结构向量虽然在数学上可行，却放弃了邻近像素之间的规律和不同位置共享的成像机制。卷积把这种共享写进参数化，注意力则让远距离区域按内容交互，两者是在不同数据和资源条件下组织空间证据的方式。

识别“这里有裂纹”和指出裂纹的准确位置，对图像细节的要求不同。下采样能减少计算、扩大观察范围，但一条很细的裂纹也可能随之消失。即使网络能比较图像中相距很远的部件，也无法恢复输入时已经丢掉的细节。因此先确定最小要检出的缺陷有多大、需要输出类别还是位置，再选择分辨率和网络。下面从卷积和感受野推到检测框、像素损失与校准，逐步说明这些选择怎样影响结果。
\section{视觉表示：从卷积骨干到视觉 Transformer}

% auxiliary-resource-guide
本章末的“辅助学习材料”列出与核心方法对应的讲解和实践入口，可在阅读相关推导后，按所列小节对照学习。

卷积神经网络（Convolutional Neural Network，CNN）把小窗口内的加权运算应用于整幅图像；骨干是负责提取共享特征的主体，任务头则把这些特征变成类别、框或掩码。通道表示同一位置上的不同数值：输入彩色图常有红、绿、蓝三个通道，隐藏层的通道则是学习得到的不同特征，不再对应颜色。张量可以先理解成带多个索引的数表。本章公式通常按高、宽、通道排列，实际软件也常按批次、通道、高、宽排列；轴的顺序必须核对。

Residual Network (ResNet) 通过残差连接改变深层网络的优化路径，视觉Transformer（Vision Transformer，ViT）将图像块（patch）编码为向量单元（token）后使用 Transformer\cite{he2016resnet,dosovitskiy2021vit}。下文从这两种机制比较局部归纳偏置、全局交互和预训练需求。
\begin{figure}[htbp]
\centering
\begin{tikzpicture}[node distance=0.65cm and 0.65cm]
 \node[idi input, minimum width=1.7cm] (image) {工业图像\\$H\times W$};
 \node[idi process, right=of image, minimum width=2.0cm] (backbone) {视觉骨干\\CNN / ViT};
 \node[idi state, right=of backbone, minimum width=2.0cm] (features) {多尺度特征\\$F_1,\ldots,F_s$};
 \node[idi output, right=of features, minimum width=1.8cm] (class) {分类\\图像级};
 \node[idi output, below=0.75cm of class, minimum width=1.8cm] (detect) {检测\\框级};
 \node[idi output, above=0.75cm of class, minimum width=1.8cm] (segment) {分割\\像素级};
 \draw[idi arrow] (image)--(backbone); \draw[idi arrow] (backbone)--(features);
 \draw[idi arrow] (features)--(class); \draw[idi arrow] (features)--(detect); \draw[idi arrow] (features)--(segment);
\end{tikzpicture}
\caption{视觉系统先由卷积或 Transformer 骨干提取表示，再按标签粒度输出分类、检测或分割结果。}
\label{fig:vision-output-granularity}
\end{figure}
图\ref{fig:vision-output-granularity}将共享骨干与三种输出粒度分开：分类使用图像级标签，检测和分割还要求保留位置与边界。图像是具有二维邻域结构的信号。卷积层以局部窗口提取平移共享的特征：
\begin{equation}
 z_{i,j,c}=\sum_{u,v,c'}W_{u,v,c',c}x_{i+u,j+v,c'}+b_c.
\end{equation}
共享权重降低参数量，也把局部性作为归纳偏置。堆叠卷积、非线性和下采样后，感受野逐步扩大，表示从边缘、纹理发展到部件和对象。

本章先修只需向量与矩阵乘法、非线性激活和用标签计算损失。先手算一个卷积窗口，再追踪它怎样变成类别、框或掩码；后面的梯度推导用于解释训练难点，不是理解输出格式的前提。若只有“有无缺陷”的图像标签，不能直接套用需要框或像素标注的训练损失。

\paragraph{一个窗口怎样变成一个输出。}
取单通道窗口$\left[\begin{smallmatrix}1&2\\3&4\end{smallmatrix}\right]$与核$\left[\begin{smallmatrix}1&0\\0&-1\end{smallmatrix}\right]$，偏置为零，则对应输出为$1\times1+2\times0+3\times0-4\times1=-3$。核移到右边后仍使用这四个权重，只更换被读取的像素。若随后使用ReLU，即$\max(0,z)$，这一响应变为零；多个核可学习不同方向的响应。核权重由训练损失调整，示例核仅用于理解计算，不是要求人工为每种裂纹设计模板。

步幅是相邻窗口起点移动的像素数，填充是在图像边缘补值，空洞率是核内相邻采样点的间隔。例如$5\times5\times3$输入经过$3\times3$核、步幅一、无填充、四个输出通道，得到$3\times3\times4$，含$3\times3\times3\times4+4=112$个参数。每个输出通道都汇总三个输入通道，并不是分别对红绿蓝得到三个互不交流的答案。

\paragraph{卷积的索引、参数与等变性。}
设输入$X\in\R^{H\times W\times C_{\rm in}}$，方形核大小$k$，输出通道数$C_{\rm out}$，步幅$s$、填充$p$、空洞率$d_a$。采用从零开始的空间索引，有
\begin{equation}
 Z_{i,j,c}=b_c+\sum_{u,v=0}^{k-1}\sum_{a=1}^{C_{\rm in}}
 K_{u,v,a,c}X_{si-p+d_au,\,sj-p+d_av,\,a},
\end{equation}
越界输入按指定规则填充。这个深度学习中惯称的卷积实际采用互相关索引；若使用数学卷积，只需翻转核。沿高度方向，为使最下端采样不超过填充后的输入，需满足$s i+d_a(k-1)\leq H+2p-1$，因此
\begin{equation}
 H_{\rm out}=\left\lfloor\frac{H+2p-d_a(k-1)-1}{s}\right\rfloor+1,
\end{equation}
宽度同理。参数量为$k^2C_{\rm in}C_{\rm out}+C_{\rm out}$，与图像面积无关；不共享的局部连接则须为每个输出位置另存一组权重。

在无限网格或忽略边界、$s=1$时，令$(T_\delta X)_{i,j}=X_{i-\delta_1,j-\delta_2}$，直接代入求和式得$K*(T_\delta X)=T_\delta(K*X)$。这就是平移等变性，而不是平移后输出保持不变。步幅$s>1$时通常仅保留对$s$整数倍平移的对应关系，边界填充还会引入额外差异。图像分类可以通过汇聚追求不变性，定位任务却应保留位置随输入一起移动的关系。

深度可分离卷积先对每个输入通道施加一个空间核，再用逐点矩阵混合通道。忽略偏置，每个输出位置的乘加数由$k^2C_{\rm in}C_{\rm out}$降为$k^2C_{\rm in}+C_{\rm in}C_{\rm out}$，比值为$1/C_{\rm out}+1/k^2$。节省来自结构约束：对固定输入通道，通向各输出通道的空间模板被限制为同一个模板的不同倍数，因此它不等价于任意标准卷积，且实际延迟还受内存访问影响。

卷积骨干的演化可以理解为三个连续问题：如何让网络更深而仍然可优化，如何在有限计算量下提高特征复用，如何让不同尺度的目标都得到足够的空间分辨率。ResNet通过残差连接$h_{l+1}=h_l+F_l(h_l)$改善深层网络的优化；DenseNet进一步把前面各层的特征拼接到当前层，使梯度和低层纹理可以直接复用，但显存访问成本更高。MobileNet 使用深度可分离卷积，将标准卷积拆成逐通道空间卷积和$1\times1$逐点卷积，计算量近似从$k^2C_{in}C_{out}$降为$k^2C_{in}+C_{in}C_{out}$，适合边缘质检；EfficientNet 则联合缩放深度、宽度和输入分辨率，强调在给定算力下平衡三者，而不是单独堆叠层数。

对工业图像而言，骨干网络并不只是分类器的前端。ResNet-18/34 通常适合小数据和低延迟基线，ResNet-50 适合迁移学习的性能--成本折中，ConvNeXt 把大卷积核、归一化和训练策略现代化，在保留卷积局部归纳偏置的同时接近 Transformer 的表示能力。选择骨干时应同时记录参数量、特征图步幅、显存、冻结层数和输入分辨率；同一个“准确率更高”的模型，如果把最小可检缺陷下采样掉，就不适合作为生产模型。

ViT 把图像划分为$P\times P$的 patch，展平后通过线性投影得到 token：$z_i=x_iW_E+e_i^{pos}$，再用标准 Transformer 建模全局关系。其优势是任意两个区域可以在少量层数内直接交互，代价是 patch 数量为$N=HW/P^2$时注意力开销近似为$O(N^2d)$，并且从零训练往往需要更多数据。Swin Transformer 通过窗口注意力和交错移位把复杂度降为局部窗口规模的平方，并通过层级下采样适配检测和分割；当缺陷同时依赖局部纹理和较大范围工件结构时，Swin 或 ConvNeXt 往往比原始 ViT 更自然。

卷积与 Transformer 不是简单的替代关系。卷积提供局部性、平移等变性和较强的数据效率，ViT 提供长距离交互和可扩展的全局建模。工业迁移实验应固定预训练数据、输入分辨率和微调预算，至少比较从零训练、冻结骨干线性探测、全量微调和参数高效微调四种设置，避免把预训练规模或训练预算误认为模型结构的收益。

\begin{table}[htbp]
\centering
\caption{视觉模型族谱与工业选型依据}
\label{tab:vision-model-families}
\scriptsize
\setlength{\tabcolsep}{3pt}
\resizebox{\linewidth}{!}{%
\begin{tabular}{p{2.4cm}p{3.7cm}p{3.5cm}p{4.0cm}}
\hline
模型族 & 核心机制 & 主要优势 & 典型工业选择依据 \\
\hline
ResNet & 残差块与层级卷积 & 优化稳定、预训练模型与工具较完善 & 小中型数据集、稳定基线、低风险迁移 \\
DenseNet & 跨层特征拼接 & 特征复用充分、参数相对紧凑 & 纹理细节重要但显存带宽充足 \\
MobileNet/\allowbreak EfficientNet & 深度可分离卷积或复合缩放 & 参数和延迟较低 & 边缘设备、相机侧实时推理 \\
ConvNeXt & 现代化大卷积核骨干 & 局部先验与表达能力平衡 & 需要卷积归纳偏置的高性能分类或检测 \\
ViT/Swin & patch 注意力或窗口注意力 & 全局关系与大规模预训练迁移 & 长距离结构、充足预训练或多尺度输出 \\
\hline
\end{tabular}
}
\end{table}

表\ref{tab:vision-model-families}把结构机制对应到工业资源条件，比较时应区分残差优化、特征复用、边缘计算和长距离交互各自解决的问题。
\paragraph{表示学习的评价。}
分类常用top-1准确率，检测常用平均精度，分割常用交并比。工业视觉还应报告跨设备、跨光照、跨批次和少见缺陷的结果。视觉模型的显著性图可以辅助分析，但可视化不等于因果解释，仍需结合遮挡实验、反事实扰动或专家标注验证。

\paragraph{卷积的空间归纳偏置。}
卷积核是局部加权模板。平滑核抑制高频噪声，边缘核响应局部变化，多个卷积核共同形成一组可学习滤波器。步幅大于一会降低空间分辨率，空洞卷积在不显著增加参数的情况下扩大感受野。池化提供一定平移鲁棒性，却也可能丢失小缺陷位置，因此检测任务常使用特征金字塔融合仍保留在不同层级的多尺度信息，而不是凭空恢复已被所有分支丢弃的像素细节。

\paragraph{残差优化与特征复用。}
普通深层网络学习$H(x)$，残差网络学习$F(x)=H(x)-x$。若恒等映射已经足够好，优化器只需令$F$接近零。对残差块$h_{l+1}=h_l+F_l(h_l)$，梯度包含直接项$I$，因此比连续相乘的纯变换更容易沿深层传播。残差连接也成为Transformer和扩散模型的基础组件。

\paragraph{残差梯度不再只有单一路径。}
设各层$h_l\in\R^d$，且残差分支$F_l:\R^d\to\R^d$可微。令$J_l=\partial F_l/\partial h_l^{\mathsf T}$，则
\begin{align}
 \frac{\partial h_{l+1}}{\partial h_l^{\mathsf T}}&=I+J_l,\\
 \nabla_{h_l}\mathcal L&=(I+J_l)^{\mathsf T}\nabla_{h_{l+1}}\mathcal L.
\end{align}
连续两个块的前向Jacobian为$(I+J_{l+1})(I+J_l)=I+J_l+J_{l+1}+J_{l+1}J_l$，包含不经过任何残差分支的恒等项、经过一个分支的项和经过两个分支的项。相比普通网络只有$J_{l+1}J_l$，残差结构允许梯度走较短路径。若$\|J_l\|_2\leq\epsilon<1$，则对任意向量$v$有
\begin{equation}
 (1-\epsilon)\|v\|_2\leq\|(I+J_l)v\|_2\leq(1+\epsilon)\|v\|_2.
\end{equation}
小残差分支使单块接近恒等，但多块的上下界仍相乘，所以残差不是绝不消失或爆炸的定理。改变通道或空间尺寸时，要把捷径替换为投影$S_lh_l$，导数变为$S_l+J_l$，此时也不再有字面上的恒等路径。

\paragraph{相加、拼接与图像块不要混淆。}
若残差输入和分支输出都是$8\times8\times16$，逐元素相加后仍是$8\times8\times16$；按通道拼接则成为$8\times8\times32$。因此ResNet的相加与DenseNet、U-Net常用的拼接不是同一个操作。若空间大小不同，应先下采样或上采样到一致大小；若相加的通道数不同，还需投影。上采样是把低分辨率数表变大，不会自动找回被删除的真实细节。

一幅$32\times32\times3$图像以$8\times8$划块，得到$4\times4=16$个块，每块展平为$192$个数。投影到宽度$64$后得到$16\times64$表示，单头注意力得分为$16\times16$。分类可把这些块表示汇聚成一个$64$维向量，再输出类别得分；定位则需要保留每个块的空间对应关系。若加入分类token，位置数变为$17$，注意力尺寸也随之变化。窗口注意力减少单层读取范围，但跨窗口联系要靠后续层建立，不能同时假定它既完全局部又单层全局。

\paragraph{图像块与窗口注意力的成本。}
假设$H,W$均可被$P$整除，令第$i$个展平图像块为列向量$x_i\in\R^{P^2C}$，投影矩阵$E\in\R^{d\times P^2C}$，则$z_i=Ex_i+e_i^{\rm pos}\in\R^d$；前文$x_iW_E$采用的是其转置后的行向量记法。共有$N=HW/P^2$个块，$QK^{\mathsf T}$需要$O(N^2d)$计算。若每个窗口包含$M\times M$个token，且token网格的高宽均可被$M$整除，共有$N/M^2$个窗口；不能整除时需填充并屏蔽无效位置，成本按填充后位置数计算。可整分时局部注意力开销为
\begin{equation}
 O\left(\frac{N}{M^2}(M^2)^2d\right)=O(NM^2d).
\end{equation}
固定$M$时对$N$线性增长，但单层不能直接连接不同窗口，移位窗口用后续层交换边界信息。另一方面，线性patch投影若$d<P^2C$，其秩最多为$d$，必有非零输入变化被映射为零；这说明不可假定所有细节都保留在token中，却不意味着小于patch的缺陷必然不可见，是否丢失还取决于投影和训练。

\section{空间输出：检测、分割与多尺度特征}
卷积神经网络（Convolutional Neural Network，CNN）可以用于不同粒度的视觉预测。U-Net 的同尺度跳跃连接用于恢复像素级定位，Feature Pyramid Network (FPN，特征金字塔网络) 通过自顶向下路径和横向连接融合多尺度语义，Mask Region-based Convolutional Neural Network (Mask R-CNN) 则在实例检测中增加掩码分支\cite{ronneberger2015unet,lin2017fpn,he2017mask}。这些方法解决的是不同粒度的空间输出，不能仅按网络深度排序。
从图像分类转向定位时，高层语义与空间细节之间的矛盾变得直接。高分辨率特征保留小缺陷边界，低分辨率特征提供较大范围的上下文。特征金字塔通过自顶向下路径与横向连接融合不同尺度，因此应先根据缺陷在输入图像中的像素尺寸选择输出层级，再选择检测头。对于 ViT，patch 过大同样可能丢失细缺陷，卷积 stem（图像进入主干前的浅层卷积模块）、重叠 patch 和层级表示可以提供不同折中。小、中、大目标的分层召回有助于区分表示分辨率不足与后续匹配错误。

检测和分割的差异不只是输出张量的形状，而是模型必须建立从候选区域到空间决策的完整机制。两阶段检测器 Faster R-CNN 先由区域建议网络产生候选框，再由感兴趣区域（Region of Interest，RoI）的池化或 RoIAlign 提取每个候选区域的特征并分别分类、回归；它通常具有较高定位质量，但候选区域和多阶段计算会增加延迟。单阶段 You Only Look Once (YOLO)、RetinaNet 和 Fully Convolutional One-Stage Object Detection (FCOS) 直接在密集特征图上预测类别与位置，速度更快，却必须处理大量背景位置和正负样本分配。RetinaNet 的 focal loss
\begin{equation}
 \mathcal{L}_{\mathrm{focal}}=-\alpha_t(1-p_t)^\gamma\log p_t
\end{equation}
其中$t$表示该样本的真实类别，$p_t$是模型对真实类别的预测概率，$\alpha_t$是该类别的预设权重，$\gamma\geq0$控制对易分类样本的降权。二分类时，$p_t$按真实标签在$p$与$1-p$之间选取。
这一损失通过降低易分类背景的权重缓解类别不平衡；FCOS 不依赖预设锚框（anchor，即预先放置的参考矩形），直接回归位置到边界的距离，减少 anchor 尺寸和比例的超参数。工业小缺陷任务应比较 anchor-based、anchor-free 和不同输入分辨率，而不能只更换 backbone。

\paragraph{先分清训练时的框与推理时的框。}
锚框（anchor）是在特征位置预置的参考矩形，不是真实标注。假设一幅图有$100$个候选、两类缺陷，采用包含背景类的简化检测头，类别得分为$100\times3$，类别无关的框偏移为$100\times4$。训练先把候选与标注匹配，得到“这是背景还是哪类缺陷”及正样本的坐标目标，再计算损失。推理没有真实框，只能解码偏移、筛选低分框和消除重复；最终框数因图像而异。不同检测器的输出格式可能不同，例如独立的目标存在分数或类别专属回归头，不能把这个示例当作所有YOLO版本的接口。

非极大值抑制（Non-Maximum Suppression，NMS）通常逐类执行：先保留最高分框，再去掉与它交并比（IoU，即交集面积除以并集面积）超过阈值的其他框，重复直到候选用完。它解决“同一缺陷报了多个框”，不是证明保留框正确；紧邻的两个真实缺陷也可能被误合并。RoIAlign则是在特征图上对候选区域进行插值采样，变成固定大小的数表，供后续分类和回归使用，并不是把候选位置当成已知正确答案。

\paragraph{检测框、匹配与定位损失。}
检测器为每个候选框输出类别概率和边界框参数。两个框$A,B$的交并比为
\begin{equation}
 \operatorname{IoU}(A,B)=\frac{|A\cap B|}{|A\cup B|}.
\end{equation}
训练通常由分类损失和定位损失组成。推理时非极大值抑制保留高分框并删除与其重叠过大的框。小缺陷检测需要特别关注输入分辨率、正负样本比例和标注框的一致性。

\paragraph{从几何框到可训练的回归目标。}
设参考框中心及宽高为$(a_x,a_y,a_w,a_h)$，真实框为$(g_x,g_y,g_w,g_h)$，宽高均严格为正且使用同一像素坐标系。直接回归像素位移会使大框产生更大尺度的梯度，因此定义无量纲目标
\begin{equation}
 t_x^*=\frac{g_x-a_x}{a_w},\quad t_y^*=\frac{g_y-a_y}{a_h},\quad
 t_w^*=\log\frac{g_w}{a_w},\quad t_h^*=\log\frac{g_h}{a_h}.
\end{equation}
网络输出$t\in\R^4$后，以$b_x=a_x+a_wt_x,b_y=a_y+a_ht_y,b_w=a_we^{t_w},b_h=a_he^{t_h}$解码。中心相对化降低尺度依赖，对数宽高保证解码尺寸为正；但$\partial b_w/\partial t_w=b_w$，极端输出仍可能带来不稳定。

对匹配到真实框的正样本，常用残差$e=t-t^*$的分段损失
\begin{equation}
 \rho_\beta(e)=\begin{cases}e^2/(2\beta),&|e|<\beta,\\|e|-\beta/2,&|e|\geq\beta,\end{cases}
 \qquad \beta>0.
\end{equation}
在接点$e=\pm\beta$，两段函数值均为$\beta/2$，两侧导数也分别相等，所以此处分段不妨碍一阶求导。其导数在小误差处为$e/\beta$，大误差处为$\operatorname{sign}(e)$，兼顾局部平滑与异常残差的有界影响。完整损失可写为
\begin{equation}
 \mathcal L=\frac1{N_{\rm cls}}\sum_i\ell_{\rm cls}(p_i,y_i)
 +\frac{\lambda}{\max(1,N_+)}\sum_{i:y_i>0}\sum_{j\in\{x,y,w,h\}}\rho_\beta(t_{ij}-t_{ij}^*).
\end{equation}
这里$p_i$是候选$i$的类别概率向量，$y_i=0$表示背景、$y_i>0$表示缺陷类别；$N_{\rm cls}>0$为参与分类计分的候选数，$N_+$为正样本数，$\lambda\geq0$控制定位项权重。背景样本没有真实框回归目标；分类与定位的归一化方式必须报告，否则相同$\lambda$可能代表不同梯度比例。

IoU直接比较几何重叠。设两个边长为$a$的正方形仅水平错位$\delta\in[0,a)$，则交集为$a(a-\delta)$，并集为$a(a+\delta)$，所以$\operatorname{IoU}=(a-\delta)/(a+\delta)$，在$\delta=0$右侧的导数为$-2/a$。这说明相同一像素偏移对小框更加苛刻。不相交框的普通IoU为零，在保持不相交的区域内不给出靠近方向，因此实际定位还会结合参数回归或加入外接框、中心距离等项。

\paragraph{焦点损失究竟改变了什么梯度。}
对正类样本设logit为$z$、$p=\sigma(z)$，则$\ell=-\alpha(1-p)^\gamma\log p$。不能只把交叉熵梯度乘以$(1-p)^\gamma$，因为这个权重本身依赖参数。乘积法则给出
\begin{equation}
 \frac{d\ell}{dz}=\alpha\gamma p(1-p)^\gamma\log p-\alpha(1-p)^{\gamma+1}.
\end{equation}
令$p=1-\varepsilon$并利用$\log(1-\varepsilon)\approx-\varepsilon$，得梯度约为$-\alpha(\gamma+1)\varepsilon^{\gamma+1}$；普通交叉熵则约为$-\alpha\varepsilon$。所以易样本梯度以更高阶衰减，困难样本占据更大的相对权重。但标注错误往往也是困难样本，焦点损失不具备自动甄别错误标签的能力，也不能保证输出概率仍然校准。

\paragraph{语义分割、实例分割与 U-Net。}
像素级交叉熵适合类别互斥场景，Dice损失
\begin{equation}
 \mathcal{L}_{Dice}=1-\frac{2\sum_i p_iy_i+\epsilon}{\sum_i p_i+\sum_i y_i+\epsilon}\label{eq:vision-soft-dice}
\end{equation}
对小目标更友好。U-Net 用逐级下采样的编码器提取语义上下文，再用逐级上采样的解码器恢复空间分辨率，并通过同尺度跳跃连接把编码器中的细粒度特征送入解码器。跳跃连接缓解了下采样造成的边界丢失，但也会把纹理噪声带入输出，因此需要结合标注质量、通道宽度和后处理验证。U-Net++通过重新设计跳跃路径减少编码器与解码器之间的语义鸿沟，DeepLab则以空洞空间金字塔池化扩大多尺度上下文；它们解决的是不同的空间信息瓶颈，不能仅凭网络名称比较优劣。实例分割还要进一步区分相邻目标，Mask R-CNN 在检测框分支之外增加 RoI 级掩码分支。边界质量、缺陷面积误差和漏检区域的连通性可能比平均 IoU 更接近工业验收标准。对关键缺陷，可以使用多个模型或随机增强估计预测波动。

\paragraph{把像素输出还原成掩码。}
语义分割为每个像素判类别，同类的两条裂纹可以用同一个标签；实例分割还要给不同裂纹分配不同实例身份。以一张$32\times32$图像的二值分割为例，模型输出$32\times32$个logit，即尚未转成概率的实数得分，经sigmoid成为同形状概率图，训练与同形状零一标签比较。推理时才按验证集确定的阈值变成硬掩码。多类互斥分割通常输出$32\times32\times K$，在每个像素的$K$个类别上归一化；不能在整幅图的全部像素上做一次softmax。

四个像素的真实掩码为$(1,1,0,0)$，预测硬掩码为$(1,0,1,0)$，则交集为一、并集为三，IoU为$1/3$，Dice为$2/(2+2)=1/2$。训练中的软Dice用概率代替硬预测，以便求导；阈值化通常不参与这条训练梯度。U-Net在解码端把低分辨率特征放大，与同尺度编码特征拼接，再经卷积输出像素得分，拼接提供的细节仍必须由数据教会网络如何使用。

\paragraph{像素似然、Dice与稀疏前景。}
设一幅图含$n$个像素，二值标签$y_i\in\{0,1\}$，模型输出$p_i=\sigma(z_i)$。在给定图像后用逐像素Bernoulli分布作为训练分解，有
\begin{align}
 p(y\mid X)&=\prod_{i=1}^n p_i^{y_i}(1-p_i)^{1-y_i},\\
 \mathcal L_{\rm BCE}&=-\frac1n\sum_i[y_i\log p_i+(1-y_i)\log(1-p_i)],\\
 \frac{\partial\mathcal L_{\rm BCE}}{\partial z_i}&=\frac{p_i-y_i}{n}.
\end{align}
这个似然分解不表示真实裂纹像素相互独立，而是输出头采用的近似。前景极少时，大量背景项会主导平均损失。Dice从集合重叠出发：硬掩码的$D=2|A\cap B|/(|A|+|B|)$，将预测指示函数换成概率便得到式\eqref{eq:vision-soft-dice}中的软Dice。其中$\epsilon>0$是固定的小平滑常数，用于避免空掩码时分母为零，求导时不随$p_i$变化。记$N=2\sum_i p_iy_i+\epsilon$、$D_0=\sum_i p_i+\sum_i y_i+\epsilon$，则
\begin{equation}
 \frac{\partial\mathcal L_{\rm Dice}}{\partial p_j}
 =\frac{N-2y_jD_0}{D_0^2},\qquad
 \frac{\partial\mathcal L_{\rm Dice}}{\partial z_j}
 =\frac{N-2y_jD_0}{D_0^2}p_j(1-p_j).
\end{equation}
一个像素的梯度依赖全图预测总量，体现其区域级归一化，而不是各像素独立计分。对无前景图，$N=\epsilon$，当预测面积较大时惩罚梯度可能很小；可结合Binary Cross Entropy (BCE)并明确空掩码的评价约定。Dice也不是概率校准目标，因此不应把优化后的$p_i$未经验证就解释为可靠风险概率。

硬掩码下，令$J=\operatorname{IoU}$，由$|A|+|B|=|A\cup B|+|A\cap B|$得到$D=2J/(1+J)$。同一对掩码的两个指标单调对应，但先对图像平均再作非线性变换不等价；软Dice与阈值后二值IoU更不能直接替换。边界和拓扑验收正是补充这些面积指标遗漏的信息。

\paragraph{数据增强与标签语义。}
裁剪、翻转、颜色扰动和噪声增强应保持标签语义不变。若翻转会改变设备方向，或颜色变化会掩盖真实缺陷，增强反而制造错误先验。几何增强必须同步变换图像、框和掩码；离散类别掩码缩放通常使用最近邻插值，不能把类别编号插值成新类别。裁剪若删除了全部缺陷，应重新检查图像级标签，并按约定裁切或删除对应框。增强策略应由物理过程和相机成像条件约束，并通过消融实验验证。

\paragraph{工业检测案例。}
对表面缺陷数据，可先建立“人工特征加树模型”和“小型卷积网络”两个基线，再比较预训练视觉编码器。实验应按生产批次或相机划分，报告少数缺陷的召回率和每千张图的误报数。IoU 对边界移动的敏感性、正负样本不平衡的来源以及不改变缺陷标签的增强边界，都应作为验收规则的一部分明确记录。

\paragraph{感受野与有效感受野。}
理论感受野描述某个深层单元能够访问的输入区域。若连续两层卷积核大小为$k_1,k_2$、步幅为$s_1,s_2$，感受野可递推为
\begin{equation}
 r_2=r_1+(k_2-1)j_1,\qquad j_2=j_1s_2,
\end{equation}
其中$j$表示相邻深层单元在原图上的间隔。更一般地，令$r_0=j_0=1$，第$l$层核大小、步幅与空洞率为$k_l,s_l,d_l$。该层最左与最右采样点间隔为$d_l(k_l-1)$个上一层步距，每个采样点又覆盖$r_{l-1}$个输入位置，故
\begin{equation}
 r_l=r_{l-1}+d_l(k_l-1)j_{l-1},\qquad j_l=s_lj_{l-1}.
\end{equation}
例如三层$3\times3$卷积的步幅依次为$1,2,1$、空洞率均为一，则各层的感受野与步距为
\begin{align*}
 (r_1,j_1)&=(3,1),\\
 (r_2,j_2)&=(5,2),\\
 (r_3,j_3)&=(9,2).
\end{align*}
第三层相邻位置在原图相隔两像素，各自理论上覆盖九像素宽的区域。这里计算的是内点的覆盖范围，不保证边界处拥有同样多的真实像素，也不保证这个范围内每个像素都获得显著权重。理论感受野很大不代表所有像素贡献相同，实际梯度通常集中在更小的有效区域。小缺陷任务需要检查下采样是否在进入高层前已经抹去目标。

\paragraph{等变性、池化与空间细节。}
理想卷积具有平移等变性：输入平移后，输出特征也相应平移。但边界填充、步幅和池化会破坏严格等变性。池化提供一定不变性，却可能让两个相邻缺陷无法区分。工业视觉中应根据缺陷尺寸选择下采样比例，并在增强实验中单独测试平移、旋转、尺度和光照变化。

\paragraph{归一化层与批次统计。}
Batch Normalization 使用批内均值和方差，训练与推理阶段的统计量来源不同；小批量或设备分布高度异质时，估计会很噪。Layer Normalization 按样本内特征归一化，更适合序列和 Transformer。Group Normalization 在视觉小批量训练中常更稳定。归一化不是无害的数值操作，它改变了网络的参数化和对批次组成的依赖。

\paragraph{检测匹配与正样本分配。}
检测训练首先要决定预测框与真实框如何匹配。基于交并比（Intersection over Union，IoU）的正负样本分配简单，但小目标和稀疏目标可能没有足够正样本；一对多匹配会导致重复预测，需配合非极大值抑制或集合匹配。定位损失的尺度也影响分类与回归梯度的平衡，应报告损失权重和正样本定义。

\paragraph{分割边界与拓扑。}
两个掩码具有相同 IoU，不代表它们对工业决策同样有用。细长裂纹、孔洞和连通结构需要边界距离、连通性和骨架指标。前文说Dice对小目标更友好，是指区域归一化有助于减轻大量背景像素的支配，不保证所有小目标都受益；无前景图的梯度问题见前述推导，边界位置仍需单独评价。边界损失可以补充几何信息。训练标签若边缘本身含糊，模型输出精细边界并不一定是更真实的结果。

\begin{figure}[htbp]
\centering
\begin{tikzpicture}[>=Latex, every node/.style={font=\small}]
 \draw[fill=blue!8, draw=blue!60] (0,0) rectangle (2.4,2.4);
 \foreach \x in {0.4,0.8,1.2,1.6,2.0} {\draw[gray!35] (\x,0)--(\x,2.4);}
 \foreach \y in {0.4,0.8,1.2,1.6,2.0} {\draw[gray!35] (0,\y)--(2.4,\y);}
 \draw[fill=orange!55, draw=orange!80!black] (0.8,0.8) rectangle (1.6,1.6);
 \node[anchor=north] at (1.2,-0.15) {局部卷积窗口};
 \draw[->, thick] (2.8,1.2)--(4.0,1.2);
 \draw[fill=green!10, draw=green!60!black] (4.4,0.2) rectangle (6.8,2.2);
 \draw[fill=green!45, draw=green!70!black] (5.1,0.9) rectangle (6.1,1.5);
 \node[anchor=north] at (5.6,0.05) {高层特征图};
 \node[above] at (5.6,2.35) {更大感受野、较低空间分辨率};
\end{tikzpicture}
\caption{卷积网络从局部窗口逐步形成层级表示。}
\label{fig:vision-receptive-field}
\end{figure}

图\ref{fig:vision-receptive-field}对照输入局部窗口与高层特征，突出扩大感受野同时降低空间分辨率的代价；细小缺陷是否被保留须在下采样前后检查。
\section{预训练、跨域泛化与工业验收}
标注稀缺时，监督微调并不是唯一选择。Contrastive Language--Image Pretraining (CLIP) 使用大规模图文对比学习把图像和文本映射到共享空间，给定“有裂纹”“无裂纹”或更细的缺陷描述即可进行零样本分类；但工业术语、相机域和细粒度缺陷通常与通用图文语料不同，因此 CLIP 只能作为迁移起点，必须通过领域文本模板、少量标注或 Low-Rank Adaptation (LoRA) 适配校准。Bootstrapping Language--Image Pretraining (BLIP)-2、LLaVA 等视觉语言模型把视觉编码器与语言模型连接起来，适合生成检验报告和解释性问答，但它们的文字流畅性不能替代像素级定位，报告中的缺陷位置仍应由检测器或分割器提供可验证证据。

\paragraph{迁移学习究竟训练了哪一部分。}
预训练是先在已有数据上学习参数，迁移则将这些参数用于当前质检任务。冻结骨干的线性探测只训练最后的分类矩阵；全量微调同时更新骨干与任务头；低秩适配只训练少量附加参数，第八章再解释其分解。自监督不使用人工类别标签，而从图像自身构造训练目标；它仍需训练，不等于完全没有监督信号。零样本分类是没有使用当前任务的标注示例训练分类头，不等于模型从未见过任何图像或语义知识。

例如冻结骨干每图输出$128$维向量，三类分类头用$3\times128$矩阵和三个偏置得到得分；训练只改变这$387$个数。上线时固定它们，对新图前向计算。训练时的随机增强通常不在推理时任意重复，Batch Normalization通常改用保存的运行统计量；因此“冻结权重”与“设为推理模式”是两个独立设置，必须同时明确。

\paragraph{自监督与视觉表征。}
当缺陷标签稀缺时，可以先用无标签图像学习表示。对比学习通过不同增强视图构造正样本，遮挡重建则要求模型恢复被隐藏的 patch。增强必须与工业语义一致：如果随机裁剪会删除缺陷，模型可能学会“正常背景”；如果两种视图来自不同设备，正样本关系也可能并不成立。预训练表示仍需在跨设备测试上验证，而不能只看线性探测分数。

\paragraph{校准与拒识。}
这里的校准针对分类置信度，不是相机内外参数的几何标定，也不是调整检测框坐标。校准检查预测概率与实际发生频率是否一致：在大量被报为$0.8$缺陷概率的样本中，缺陷比例应接近$80\%$，而非要求每张图都答对。拒识则是不自动判定，转交后续复核。高置信度错误在质检系统中尤其危险。可以用温度缩放或保序回归修正分类概率，再根据风险设置拒识区域。拒识不是把困难样本删除，而是将其转交人工或更精细的模型。报告覆盖率--风险曲线可以说明系统在自动处理多少样本时保持怎样的错误水平。

\paragraph{任务边界与模型选择。}
分类回答“图像中是否存在某类状态”，检测还要回答“状态位于哪里”，分割则要求逐像素或逐区域描述。任务标签越细，标注成本越高，评价也越依赖边界和拓扑。若业务只需判断是否需要复检，像素级分割未必是最经济的方案；若需要估计裂纹长度和扩展方向，分类分数就不足以支持决策。

\paragraph{检测指标与匹配链。}
检测评价先按置信度排序预测框，再依据 IoU 和类别匹配真实框，形成精确率--召回率曲线。平均精度依赖匹配阈值、类别平均方式和置信度排序。小目标的一个像素偏移可能显著改变 IoU，因此应同时报告按目标尺寸分层的结果。非极大值抑制的阈值也会影响最终告警数量，不能只记录网络训练损失。

\paragraph{正负样本与困难样本。}
工业缺陷通常是少数类，检测器会面对大量容易分类的背景位置。焦点损失通过降低易样本权重，把梯度集中到困难样本；类别权重则直接改变不同类别的损失比例。两者都可能提高少数类召回，却也可能放大错误标注和难例噪声。验收时应同时查看每类召回、误报空间分布和人工复核负担。

\paragraph{分割后处理与几何指标。}
像素预测之后常需要连通域分析、孔洞填充、形态学操作或面积过滤。后处理可以改善视觉结果，却也可能删除细小真实缺陷。应把后处理规则纳入验证规则，并在不同缺陷尺寸、边界模糊程度和相机分辨率下检查稳定性。对裂纹等细长对象，骨架长度、断裂数和最大偏差可能比单一 Dice 更有解释力。

\paragraph{域偏移与测试时监控。}
光照、相机、镜头污染、材料批次和背景工装都会改变输入分布。输入统计监控可以发现均值、对比度和频谱变化，但不能直接证明标签关系改变。部署后应结合少量延迟标签检查分层性能，区分数据质量告警、需要校准的协变量偏移和需要重新标注的概念偏移。

\paragraph{解释工具的证据边界。}
梯度显著性、遮挡图和类激活图可以帮助定位模型关注的区域，但高亮处不一定就是导致故障判断的真实原因。例如模型可能看的是裂纹旁边的批次标记。可以分别遮住缺陷、替换背景或只保留缺陷区域，观察输出如何变化，并请专家核对。不过这些图像修改也可能产生不自然的输入，结果应结合跨相机测试解释，不能仅凭一张热力图作因果判断。

\paragraph{视觉系统验收。}
验收应覆盖正常样本、已知缺陷、少见缺陷、遮挡、模糊、光照变化、相机切换和未知输入。指标包括缺陷召回、每千张误报、定位偏差、处理延迟、拒识率和人工复核时间。若模型在实验室图像上表现优秀，却无法承受生产线速度和污染条件，就不能视为完成部署。

\paragraph{分层实验设置。}
设生产线上有四台相机、六个批次和三类缺陷。训练集可以使用前三台相机的早期批次，验证集使用相同相机的后期批次，测试集则完整留出第四台相机。这样的划分同时检验时间变化和设备迁移。若把所有图像随机打散，来自同一工件的近重复帧可能进入不同集合，检测结果会明显偏乐观。

实验应先建立三个层次的基线：多数类或人工规则基线、冻结预训练编码器的线性头、端到端微调模型。每次只改变一个因素，例如输入分辨率、增强策略、正负样本分配或后处理阈值。除 Mean Average Precision (mAP) 外，报告缺陷级召回、每千张正常图的误报、最小可检缺陷尺寸和单张图延迟。对生产线而言，漏掉一个关键缺陷的代价通常不能由平均准确率代表。

\paragraph{分割结果验收。}
对裂纹掩码，先计算 Dice 和 IoU，再计算边界到真实边界的平均距离、最大距离和骨架断裂数。若模型通过扩大预测区域提高召回，IoU 可能下降但安全风险降低；若模型只保留粗主干，Dice 可能尚可，却无法估计裂纹长度。因此验收规则应先由业务定义允许的漏检和过检，再选择指标和阈值。

标注不确定时，可以由多名标注者独立勾画并保留分歧区域。模型在分歧区域给出低置信度并不一定是失败，反而可能提示需要复核。训练、评价和部署中的掩码格式、像素坐标系、缩放方式必须保持一致，否则边界误差会被错误归因于模型。

\section{数据反馈与部署决策}
少样本适配不应只比较准确率。新相机只有少量标注时，可以冻结 ResNet、ConvNeXt、Swin 或 ViT 的大部分编码器，仅更新归一化统计量、分类头或低秩适配参数；也可以使用 CLIP 的文本原型进行零样本筛查，再把高不确定样本交给人工。验证时必须保留未参与适配的批次，并比较全量微调、参数高效微调和直接拒识三种方案。若新域置信度整体下降，拒识率应作为业务容量约束，而不是被隐藏的失败样本比例。

\paragraph{从评价数字到一批工件的处理。}
精确率是报为缺陷的结果中有多少正确，召回率是真实缺陷中有多少被找到。例如十个真实缺陷中找对八个，另误报两个，则精确率和召回率均为$0.8$；如果改成一百个真实缺陷而仍只找对八个，精确率可不变，召回率却降为$0.08$。检测计算这些数之前还必须完成类别和空间匹配。平均精度（AP）汇总不同置信度阈值下的精确率--召回率关系，mAP再按指定类别或IoU阈值平均，所以必须一起说明评价协议。

拒识后不能只报告留下样本的成绩。例如一百张图中自动处理八十张，其中两张错误，则自动覆盖率为$80\%$，已处理样本风险为$2/80=2.5\%$；其余二十张仍消耗人工容量，不可从系统总账中消失。主动学习可以优先标注这些新颖或困难样本，但固定测试集不能拿来选样、微调或挑阈值。

\paragraph{数据处理记录。}
每张图像应关联相机编号、镜头、曝光、生产批次、工件编号、标注版本和采集时间。增强后的图像保留原图引用，预测结果记录模型版本和预处理参数。发生误报时，只有根据这些记录回溯，才能判断问题来自光学条件、标注规则、数据切分、模型还是后处理。可追溯性是视觉智能进入生产系统的必要条件。

\paragraph{主动学习与标注预算。}
当人工标注昂贵时，可优先选择高不确定性、代表新域或具有高业务价值的图像。只按模型置信度采样可能反复选择相似噪声，因而需要结合特征多样性、设备覆盖和缺陷稀有度。每轮主动学习都应保留固定测试集，记录新增标注数量、类别分布、性能变化和人工时间，才能判断单位标注成本带来的真实收益。

以下材料按“残差与图像块表示—检测—像素分割”对应本章内容；卷积基础只保留一个复习入口。视频名称与链接已对照作者课程目录，未逐帧观看，播放量未能可靠核实，故不作热度排序。
\section*{辅助学习材料}
\begingroup
\emergencystretch=3em
\begin{itemize}
\item \textbf{Convolutional Neural Networks (CNNs / ConvNets)}（Stanford CS231n，英文笔记）。仅在需要复习时阅读卷积层、输出尺寸与参数共享，回到本章“卷积的索引、参数与等变性”检查步幅和填充。\href{https://cs231n.github.io/convolutional-networks/}{在线阅读}。
\item \textbf{ResNet网络讲解}（中文视频）。对照“残差梯度不再只有单一路径”，区分恒等捷径与改变尺寸的投影捷径，解释残差块为何比单纯加深网络更易优化。\href{https://www.bilibili.com/video/BV1T7411T7wa}{观看讲解}；中文博客\href{https://zhuanlan.zhihu.com/p/42706477}{《详解残差网络》}可作图文补充，其标题和残差块主题已由必应索引核对，但知乎原页访问受限，未核读全文。
\item \textbf{Vision Transformer网络讲解}（中文视频）。重点看图像如何切成 patch、投影为 token 并加入位置信息，再对照“图像块与窗口注意力的成本”计算分辨率变化带来的开销；这是 ViT 的图像建模讲解，不是泛自注意力入门。\href{https://www.bilibili.com/video/BV1Jh411Y7WQ}{观看讲解}。
\item \textbf{区域卷积神经网络（R-CNN）系列}（《动手学深度学习》，中文教程）。先读其中的 Fast R-CNN 小节与 RoI pooling 示例，将共享整图特征、候选区域、分类和框回归对应到“检测框、匹配与定位损失”。Fast R-CNN 使用选择性搜索；随后对照 Faster R-CNN 小节，辨明后者才引入区域提议网络（RPN），不可混称。\href{https://zh.d2l.ai/chapter\_computer-vision/rcnn.html}{在线阅读}。
\item \textbf{YOLO系列网络讲解(V1\textasciitilde V3)}（中文视频）。沿单阶段预测到多尺度检测的变化，比较网格、锚框与类别／位置输出，对照本章“空间输出”和“检测匹配与正样本分配”；不要把早期版本的机制直接套给所有新版 YOLO。\href{https://www.bilibili.com/video/BV1yi4y1g7ro}{观看讲解}。
\item \textbf{U-Net网络讲解}（中文视频）。沿编码器、解码器与同尺度特征拼接追踪分辨率，回到“语义分割、实例分割与 U-Net”检查细裂纹边界为何需要低层特征；并与 ResNet 的相加捷径区分。\href{https://www.bilibili.com/video/BV1Vq4y127fB/}{观看讲解}。上述四个模型视频均可从\href{https://github.com/WZMIAOMIAO/deep-learning-for-image-processing}{作者课程目录}查到配套代码。
\item \textbf{TorchVision Object Detection Finetuning Tutorial}（PyTorch 官方，英文实践）。跟随自定义数据集、预训练 Mask R-CNN 微调与评价流程，核对框、标签和掩码格式，对应本章“预训练、跨域泛化与工业验收”；它是检测／实例分割实践，不替代前面的 Fast R-CNN 原理解说。\href{https://docs.pytorch.org/tutorials/intermediate/torchvision\_tutorial.html}{在线阅读}。
\end{itemize}
\endgroup
\paragraph{自测与参考答案。}
参考框的中心及宽高为$(10,10,4,2)$，真实框为$(12,9,8,2)$，求回归目标并解码核对。参考答案：$t^*=(1/2,-1/2,\log2,0)$；解码中心为$(10+4/2,10-2/2)=(12,9)$，宽高为$(4e^{\log2},2e^0)=(8,2)$。四个目标均无量纲，不能把对数宽度$\log2$直接当作像素宽度。若候选被匹配为背景，则不应给它随意指定这组定位目标。

\paragraph{本章小结。}
视觉表示的选择应从输出要求出发：分类需要区分状态，检测需要定位区域，分割还要恢复边界。卷积和残差提供层级特征，注意力扩展区域之间的信息交互，预训练则改变这些表示所需的标注预算。最终方案必须同时满足缺陷尺度、跨域可靠性、推理资源和人工复核容量的约束。卷积索引与感受野决定证据是否仍在特征中，残差路径决定这些特征能否有效训练，框回归、焦点损失和Dice则分别改变位置误差、难例与区域重叠的权重。任何一个损失下降都只回答其中一部分问题，不能代替空间匹配、跨相机验证和部署后的复核。

找到裂纹后，还得查这台设备允许多大缺陷、该按哪一版规程复检。图像本身不能回答这些问题，下一章因此讨论语言模型怎样检索和使用手册、工单。生成报告时应保留裂纹坐标、尺寸尺度、拍摄时间和模型版本，不能只用一句“发现异常”替代可复查的检测结果。

\chapter{语言大模型}
\begin{quote}\small
\textit{本章摘要。}本章沿着“文本如何变成可预测的序列”这条主线介绍语言大模型。先讲 token、表示和 Transformer，再解释预训练、缩放、监督微调、参数高效适配和解码，最后讨论检索增强、工具调用、评测与风险。这样可以把模型能力、外部知识和系统行为区分开来。
\end{quote}
在贯穿案例中，本章不让语言模型直接替代故障预测器，而是处理设备手册、维修记录和告警上下文这类离散证据。它改变的是“如何检索、组织和调用知识”，而不是把事后记录伪装成预测时刻可用的信息；任何维修建议都必须回到原始证据、权限和人工确认。
上一章把缺陷落实为空间位置和边界，但维护判断还需要回答“这个型号允许多大偏差”“哪一版条款适用于当前设备”。这些知识并不天然存在于像素或传感器轨迹中，而分散在带条件、例外和有效期的文档里。本章研究的不是让语言流畅性覆盖证据缺口，而是把文本概率模型、外部知识访问和受约束的系统行为分开建模，再考察它们怎样协同。

可以把本章的方法分开理解：预训练学习文本怎样接续，微调和偏好学习调整回答方式，检索则为这次回答找来相关资料。三者不能互相替代。模型即使把句子续得很顺，也可能记错规程；回答即使受人欢迎，也可能漏掉适用条件；放入整本手册，也不保证它读到了关键条款。下面通过似然梯度、低秩更新和偏好概率解释训练，再检查生成答案能否追溯到当前有效的原文。
\section{Token、表示与语言建模}

% auxiliary-resource-guide
本章末的“辅助学习材料”列出与核心方法对应的讲解和实践入口，可在阅读相关推导后，按所列小节对照学习。

本章的自注意力结构以 Transformer 为基础\cite{vaswani2017attention}。Bidirectional Encoder Representations from Transformers (BERT) 的双向掩码预训练与自回归语言模型采用不同的信息可见性，因而需要分别理解表示与生成目标\cite{devlin2019bert,brown2020gpt3}。
分词器（tokenizer）先按固定规则把文本转换为词元（token）编号。词元可以是字、词的一部分、标点或字节片段，不一定是完整词。为便于手算，假设一个教学词表把“温度升高”切成“温度”“升高”两个词元，编号为$3,7$；真实切分取决于所用分词器。编号只用于查表，$7$不是比$3$“语义更大”。嵌入就是从可训练矩阵中取出相应行，把离散编号变成稠密向量。

若一批有两段文本，补齐后各四个词元，则编号张量为$2\times4$；嵌入宽度取八，查表后是$2\times4\times8$；词表有十项，输出得分为$2\times4\times10$。最后一维表示十个候选词元，不是十个独立答案。logit是softmax之前的实数打分，可正可负；softmax把一组得分转换成和为一的概率。语言模型学习
\begin{equation}
 p_\theta(x_{1:T})=\prod_{t=1}^{T}p_\theta(x_t\mid x_{<t}),\label{eq:llm-chain-probability}
\end{equation}
训练目标是最小化负对数似然。Transformer中的 token 表示由词嵌入、位置编码和多层注意力变换得到。

初读可先复习条件概率、对数和矩阵乘法，再按“词元编号—嵌入—词表得分—下一个词元”追踪数据。最大似然的推导解释为什么用交叉熵训练；LoRA和DPO的推导分别解释如何限制参数更新、如何利用偏好，不是开始使用检索问答前必须先实现的三个训练阶段。

\paragraph{一段文本怎样产生训练标签。}
沿用两个词元的例子，输入为“起始符、温度、升高”，目标为“温度、升高、结束符”。第一位置只能看起始符，第二位置能看起始符与“温度”，第三位置能看此前全部输入。因此一次前向计算虽有三个目标，每个目标仍只能依据其前缀。若三个正确目标的概率分别为$1/2,1/4,1/2$，整段含终止符的概率为$1/16$，平均负对数似然为$\log16/3$，困惑度为$16^{1/3}\approx2.52$。计算使用自然对数。

生成时没有“升高”这个真实标签可以提前喂入。模型先选一个词元，将其加入输入，再预测下一个，直至结束符或长度上限。训练时输入真实前缀称为教师强制；生成时使用自己生成的前缀，错误因此可能影响后续内容。两者都调用同一个参数化网络，区别在于输入从哪里来以及是否更新参数。

\paragraph{链式分解如何变成可计算的训练目标。}
设词表大小为$V_{\rm voc}$，$x_t\in\{1,\ldots,V_{\rm voc}\}$，嵌入矩阵$E\in\R^{V_{\rm voc}\times d}$。由联合概率的乘法公式，反复展开$p(x_{1:T})=p(x_{1:T-1})p(x_T\mid x_{<T})$即可得到式\eqref{eq:llm-chain-probability}；这不要求token独立。以起始符（BOS）提供第一个位置的条件，以结束符（EOS）学习何时停止，才构成包括长度的序列分布。设因果网络对前缀$x_{<t}$产生$h_t\in\R^d$，输出矩阵$U\in\R^{V_{\rm voc}\times d}$，则
\begin{align}
 z_t&=Uh_t+b,\qquad p_{t,j}=\frac{e^{z_{t,j}}}{\sum_{k=1}^{V_{\rm voc}}e^{z_{t,k}}},\\
 \ell_t&=-\log p_{t,x_t}=-z_{t,x_t}+\log\sum_ke^{z_{t,k}},\\
 \frac{\partial\ell_t}{\partial z_{t,j}}&=p_{t,j}-\mathbf1_{j=x_t}.
\end{align}
若把这一位置的logit当作独立变量做梯度下降，真实token的logit增加，其余logit减小，真实token的概率随之上升。共享网络参数的一步更新还受其他位置梯度影响，不能据此保证每个样本的真实token概率都上升。对隐藏状态的梯度为
\begin{equation}
 \frac{\partial\ell_t}{\partial h_t}=U^{\mathsf T}(p_t-e_{x_t}),
\end{equation}
其中$p_t,e_{x_t}\in\R^{V_{\rm voc}}$，$e_{x_t}$仅在真实词元编号处为一，其余为零；$\mathbf1_{j=x_t}$是该条件成立时取一、否则取零的指示量。所得梯度为$d$维，再经网络反传。数值实现用$\log\sum_ke^{z_k}=m+\log\sum_ke^{z_k-m}$，其中$m=\max_k z_k$，避免指数溢出。

训练时可一次输入整段右移后的真实token，并用因果遮罩并行计算所有位置；第$t$个输出的目标必须是尚未作为该位置可见输入的$x_t$，否则会退化成复制任务。填充位置不计损失，指令微调还可只对回答位置计损失。注意力遮罩控制能读取哪些位置，损失掩码控制哪些输出计分，二者不能互换；不计提示损失仍可允许回答读取提示。部分模型接口内部完成标签移位，传入前须核对契约，避免手动移位后再被内部移位。归一化应按实际计分token数而非填充后的矩阵大小进行。训练并行并不改变生成时必须逐步追加未知token的事实。

固定同一个输入前缀$c$，设真实条件分布为$q(\cdot\mid c)$，则期望交叉熵分解为
\begin{equation}
 -\sum_jq_j\log p_j=-\sum_jq_j\log q_j+\KL(q\|p).
\end{equation}
第一个熵项与参数无关，最大似然因此在可表达范围内逼近训练分布，而不是直接优化事实核查。若训练分布本身含有过时或错误陈述，降低这项损失反而可能提高相应陈述的概率。这是外部证据管理不能被预训练规模替代的数学原因。

Tokenizer 在有限词表下平衡覆盖率与序列长度。词表过小会产生较长序列，词表过大则增加嵌入参数，并使部分低频 token 缺少训练。设备编号、故障码和单位可能被切成多个片段，因此语言模型是否能正确复制这些字符串，需要单独检验。标准注意力的计算量随序列长度呈二次增长，比较上下文成本时也应记录 token 数，而非只比较字符数。

自注意力由输入$H$生成查询、键和值：
\begin{equation}
 Q=HW_Q,\quad K=HW_K,\quad V=HW_V,
 \qquad A=\operatorname{softmax}\left(\frac{QK^{\mathsf T}}{\sqrt{d_k}}+M\right),
\end{equation}
其中$d_k$为键的维度。自回归模型令$M_{ij}=-\infty$（当$j>i$），其余位置为零，使当前位置不能读取后续 token。多头注意力以不同投影组合上下文，前馈网络进一步变换表示，残差连接和归一化控制深层计算。各个头可能学到不同关系，但不能预先认定某个头必然负责特定语义。

对单条长为$T$的序列，$H\in\R^{T\times d}$，$W_Q,W_K\in\R^{d\times d_k}$，$W_V\in\R^{d\times d_v}$，所以$A\in\R^{T\times T}$，注意力读取结果为$AV\in\R^{T\times d_v}$。softmax沿每一行可见位置计算。查询用于提出当前读取需求，键用于匹配位置，值是实际读出的内容。例如一行可见得分为$0,\log3$，权重便为$1/4,3/4$；不是从两个位置中硬选一个。位置编码说明内容出现在哪里，前馈网络在各位置独立使用相同参数加工表示，多头结果经拼接和投影恢复宽度$d$后才能与原输入相加。

预归一化将归一化放在子层之前，保留较直接的残差梯度路径；后归一化改变了残差输出的尺度。因而 Transformer 是一组可以组合的结构选择，而非唯一固定的网络。理解注意力方向后，才能区分双向表示学习与因果生成的任务边界。

\begin{figure}[htbp]
\centering
\begin{tikzpicture}[>=Latex, node distance=0.65cm, every node/.style={font=\small}]
 \node[draw, rounded corners, fill=blue!10, minimum width=3cm, minimum height=0.7cm] (tok) {Token 序列};
 \node[draw, rounded corners, fill=orange!12, below=of tok, minimum width=3cm, minimum height=0.7cm] (emb) {词嵌入 + 位置};
 \node[draw, rounded corners, fill=green!12, below=of emb, minimum width=3cm, minimum height=0.7cm] (att) {多头因果注意力};
 \node[draw, rounded corners, fill=yellow!20, below=of att, minimum width=3cm, minimum height=0.7cm] (ffn) {前馈网络 + 残差};
 \node[draw, rounded corners, fill=red!10, below=of ffn, minimum width=3cm, minimum height=0.7cm] (head) {词表 logits};
 \draw[->, thick] (tok)--(emb); \draw[->, thick] (emb)--(att); \draw[->, thick] (att)--(ffn); \draw[->, thick] (ffn)--(head);
\end{tikzpicture}
\caption{自回归语言模型的基本计算链条。}
\label{fig:llm-computation-chain}
\end{figure}

图\ref{fig:llm-computation-chain}从离散 token 追踪到词表 logits，说明嵌入、因果注意力和前馈变换如何共同产生下一词元的打分，而非直接检索一条既存答案。
\section{预训练、对齐与推理}
图\ref{fig:llm-decoding-loop}的反馈箭头表示选定 token 被追加到上下文；这使后续分布依赖先前生成结果，也解释了训练并行与推理逐步执行的区别。
\begin{figure}[htbp]
\centering
\begin{tikzpicture}[node distance=0.72cm and 0.7cm]
 \node[idi input, minimum width=1.7cm] (ctx) {上下文\\$x_{<t}$};
 \node[idi state, right=of ctx, minimum width=1.8cm] (net) {Transformer\\隐藏状态};
 \node[idi process, right=of net, minimum width=1.8cm] (logit) {词表 logits\\$z_t$};
 \node[idi output, right=of logit, minimum width=1.8cm] (prob) {温度 softmax\\$p(x_t)$};
 \node[idi input, below=0.75cm of prob, minimum width=1.8cm] (sample) {采样/贪心\\选定 token};
 \draw[idi arrow] (ctx)--(net); \draw[idi arrow] (net)--(logit); \draw[idi arrow] (logit)--(prob); \draw[idi arrow] (prob)--(sample);
 \draw[idi feedback] (sample.west) -- node[below,font=\scriptsize]{追加到上下文} (ctx.south |- sample.west) -- (ctx.south);
\end{tikzpicture}
\caption{自回归解码的单步迭代：模型先形成词表分布，再把选中的 token 反馈给下一步。}
\label{fig:llm-decoding-loop}
\end{figure}
指令微调和人类反馈训练的典型路线可参照 InstructGPT，直接偏好优化则把偏好学习转写为对语言模型的优化目标\cite{ouyang2022instruct,rafailov2023dpo}。LoRA 的低秩参数更新有独立的结构假设\cite{hu2021lora}，不应与对齐目标混为一谈。
语言模型的主要分歧来自预训练任务和注意力方向，而不只是参数规模。BERT 属于 encoder-only 架构，通过双向注意力和 masked language modeling 学习上下文表示，适合分类、序列标注和检索重排序；Generative Pretrained Transformer (GPT)、LLaMA 等 decoder-only 模型使用因果掩码，只根据前文预测下一个 token，天然适合连续生成和工具调用；Text-to-Text Transfer Transformer (T5) 属于 encoder--decoder 架构，把分类、摘要、问答等任务统一写成文本到文本的条件生成。工业系统应依据输出形式选择架构：需要字段分类时，可先用BERT类编码器接固定类别的分类头，并在验证集上校准；能否校准良好仍取决于数据和目标，而非架构名称；需要长文本生成时，GPT/LLaMA 类模型更直接；需要输入输出长度差异明显的转换任务时，T5 类结构可减少提示工程依赖。

编码器（encoder）把已有输入变成表示，解码器（decoder）在这里指按条件逐步产生输出的模块，并非把压缩文件原样还原。BERT可根据“温度［遮蔽］，请复检”的左右文恢复被遮蔽词，做分类时再接任务头；GPT只能根据左侧前缀接续；T5先完整编码输入，再由因果解码器读取编码表示生成输出。BERT的双向不是偷看分类答案，只要整段输入在决策时已可取得便合理；但不能把未来维修记录放入在线预警输入。MoE则主要替换部分前馈计算，它可以与这些注意力方向组合，不是第四种互斥的文本可见性。

decoder-only 模型对整段序列优化
\begin{equation}
 \mathcal{L}_{\mathrm{CLM}}=-\sum_{t=1}^{T}\log p_\theta(x_t\mid x_{<t}),
\end{equation}
而 BERT 的掩码目标只在被遮蔽位置计算损失。两者的训练信号不同，因此不能仅用参数量推断迁移效果。对于设备手册和维修记录，可以先用 encoder 模型做语义检索和事件分类，再用 decoder 模型生成带证据引用的维修建议；这种组合通常比让一个生成模型独立完成所有步骤更容易审计。

当模型容量继续增加时，稀疏混合专家（Mixture-of-Experts, MoE）用路由器为每个 token 选择少数专家：
\begin{equation}
 y=\sum_{e\in\mathcal{S}(x)}g_e(x)E_e(x),\qquad |\mathcal{S}(x)|\ll N_{\rm exp}.
\end{equation}
这里$N_{\rm exp}$是专家总数，$\mathcal S(x)$是路由器为当前词元表示$x$选中的专家索引集合，$E_e(x)$为第$e$个专家的输出，$g_e(x)$为相应路由权重；各专家输出须同维才能加权相加。专家是可训练子网络，并非人工专家；$N_{\rm exp}$也不是前文的嵌入矩阵$E$。
Switch Transformer 采用近似 top-1 路由，降低每个 token 的激活计算；更复杂的 top-2 路由能够提高容量利用，但需要处理专家负载不均和通信开销。MoE 的“总参数量”不等于“单次推理成本”，工业部署必须分别报告总存储、激活参数、路由丢弃率和跨卡通信。

\begin{table}[htbp]
\centering
\caption{语言模型架构与工业任务匹配}
\label{tab:llm-architectures}
\scriptsize
\setlength{\tabcolsep}{3pt}
\resizebox{\linewidth}{!}{%
\begin{tabular}{p{2.7cm}p{3.8cm}p{3.7cm}p{3.8cm}}
\hline
架构路线 & 训练与注意力方向 & 优势 & 工业任务和风险 \\
\hline
BERT 类 encoder-only & 双向掩码建模 & 表示稳定、分类和检索高效 & 事件分类、实体抽取、重排序；不直接生成长文本 \\
GPT/LLaMA 类 decoder-only & 因果下一 token 预测 & 生成、对话和工具调用自然 & 报告草拟、脚本生成；需防幻觉和提示注入 \\
T5 类 encoder--decoder & 条件文本到文本 & 输入输出长度灵活、任务统一 & 摘要、规范化和文本转换；部署链路更复杂 \\
MoE 语言模型 & 路由器选择少数专家 & 总容量大而激活成本可控 & 多领域共享；需监控负载和通信 \\
\hline
\end{tabular}}
\end{table}

表\ref{tab:llm-architectures}将注意力方向和训练任务与工业输出要求对应，提示读者把分类检索、连续生成和条件转换分开选型，并单独核算 MoE 的通信成本。
\paragraph{预训练目标与困惑度。}
语言模型的平均负对数似然为
\begin{equation}
 \mathcal{L}_{LM}=-\frac1T\sum_{t=1}^{T}\log p_\theta(x_t\mid x_{<t}).
\end{equation}
困惑度定义为$\operatorname{PPL}=\exp(\mathcal{L}_{LM})$，可以理解为模型在每个位置面对的有效候选数。困惑度适合比较相同 tokenization 和数据分布下的语言模型，但不能单独代表事实性、帮助性或工业任务质量。

\paragraph{缩放、数据质量与模型容量。}
模型参数量、训练 token 数和计算预算共同影响性能，但增加规模并不能弥补重复、污染或错误数据。高质量语料需要去重、过滤个人信息、识别模板化文本，并建立训练数据与评测数据的隔离。对于工业语料，还需处理机密信息、专有术语和时间有效性。

近重复文本可能降低训练损失，却不增加新的知识覆盖。若训练语料包含评测问题或答案模板，较低困惑度也不能证明迁移能力。固定计算预算下，应同时记录参数规模、有效 token 数、序列长度、训练步数和去重比例。工业语料还需要保存来源、许可和生效时间，使旧版本规程始终与适用设备绑定。

预训练提供语言分布的基础，监督微调则用示范答案塑造任务行为。上下文学习只改变当前输入条件，不更新参数。偏好优化进一步利用同一提示下的优选与非优选答案调整输出。Reinforcement Learning from Human Feedback (RLHF) 通常训练奖励模型并优化带参考策略约束的目标，Direct Preference Optimization (DPO) 则直接学习相对于参考模型的回答概率比。参考模型提供约束基准，但不保证事实性或原有能力完全保留，仍需独立评价。

\paragraph{监督微调与参数高效适配。}
全量微调更新所有参数，成本随模型规模增长。LoRA在权重更新上使用低秩分解
\begin{equation}
 W'=W+\frac{\alpha}{r}BA,\qquad B\in\mathbb{R}^{d\times r},\ A\in\mathbb{R}^{r\times k},\ r\ll\min(d,k).
\end{equation}
它只训练$A,B$，便于为不同设备、部门或任务保存轻量适配器。参数高效方法降低存储成本，但多个适配器之间可能产生干扰，需要明确版本和适用域。

对一个$d\times k$矩阵，全量更新需要$dk$个参数，LoRA 仅需$r(d+k)$个，参数节省比例为
\begin{equation}
 1-\frac{r(d+k)}{dk}.
\end{equation}
低秩约束限制更新所在的方向，秩过小可能无法表达目标域变化。因此应同时比较任务效果、可训练参数、训练显存及适配器合并后的推理成本，而不是只比较保存文件的大小。

\paragraph{区分训练目标、可训练参数与当前提示。}
监督微调（Supervised Fine-Tuning，SFT）需要“问题、示范回答”对，以示范词元为目标；DPO需要同一问题下优选与非优选回答对，以相对偏好为目标；LoRA只规定哪些参数能变化，可以用于不同训练目标。给提示附上两个示例属于上下文学习，只改变本次前向计算，不会自动把示例写进权重。将三者都称为“教模型”容易混淆数据需求和成本。

取$d=k=4,r=1,\alpha=1$，完整权重有$16$个数，LoRA的$A,B$共有八个数。令$A=(1,0,0,0)$、$B=(1,2,0,0)^{\mathsf T}$、$x=(3,4,5,6)^{\mathsf T}$，则$Ax=3$，附加输出$BAx=(3,6,0,0)^{\mathsf T}$。它只通过一个中间标量改变四维输出，说明“低秩”限制的是更新方向，而不是把原模型的全部知识压成一个数。推理时基座和已训练适配器通常都固定，不再利用真实答案反向传播。

\paragraph{低秩更新的前向与反向计算。}
对单个线性层，令输入$x\in\R^k$、输出$y\in\R^d$、冻结权重$W\in\R^{d\times k}$，并记$s=\alpha/r$。LoRA并非先构造完整更新再相乘，而可按
\begin{equation}
 v=Ax\in\R^r,\qquad y=Wx+sBv
\end{equation}
计算。若上游梯度为$g=\nabla_y\mathcal L\in\R^d$，由微分$d\mathcal L=g^{\mathsf T}dy$逐项整理得
\begin{align}
 \nabla_B\mathcal L&=s\,g(Ax)^{\mathsf T},\\
 \nabla_A\mathcal L&=s\,B^{\mathsf T}g\,x^{\mathsf T},\\
 \nabla_x\mathcal L&=W^{\mathsf T}g+sA^{\mathsf T}B^{\mathsf T}g.
\end{align}
即使$W$冻结，反向传播仍可能需要穿过它，且前向激活仍需存储，所以可训练参数减少不等于训练显存按同一比例减少。常用$A$随机初始化、$B=0$，使初始模型严格等于基座；第一步$\nabla_A=0$而$\nabla_B$一般不为零，随后$B$变化才使$A$获得梯度。若两者都初始化为零，两组梯度都为零，会阻断学习。

矩阵的秩是其线性独立列的最大数量，可理解为能独立改变输出的方向数；更新满足$\operatorname{rank}(BA)\leq r$。若理想全量更新$\Delta W_*$的奇异值为$\sigma_1\geq\cdots\geq\sigma_s$，其最佳秩$r$的Frobenius范数近似误差为$\sum_{j>r}\sigma_j^2$的平方根。这解释了低秩假设依赖任务变化是否集中在少数方向，但不意味着实际非凸训练一定找到这个近似。部署时可把$sBA$合并进$W$以消除额外分支，前提是适配器固定且数值格式兼容；动态切换适配器、量化权重与合并精度则需另外验证。

\paragraph{从偏好概率到直接偏好优化。}
设提示为$q$，完整回答为$y$，策略为$\pi_\theta(y\mid q)$，冻结参考策略为$\pi_{\rm ref}$。假设参考策略在所比较回答上概率为正，先考虑给定$q$的KL正则目标
\begin{equation}
 \max_\pi\left\{\sum_y\pi(y\mid q)r(q,y)
 -\beta\sum_y\pi(y\mid q)\log\frac{\pi(y\mid q)}{\pi_{\rm ref}(y\mid q)}\right\},\qquad\beta>0.
\end{equation}
用乘子$\lambda$约束$\sum_y\pi(y\mid q)=1$，对每个$\pi(y\mid q)$求导，得到
\begin{equation}
 r(q,y)-\beta\left(\log\frac{\pi(y\mid q)}{\pi_{\rm ref}(y\mid q)}+1\right)+\lambda=0.
\end{equation}
还需归一化常数$Z(q)=\sum_y\pi_{\rm ref}(y\mid q)e^{r(q,y)/\beta}$有限；例如回答空间有限且奖励有限时满足这一条件。归一化后，最优策略为$\pi^*(y\mid q)=\pi_{\rm ref}(y\mid q)e^{r(q,y)/\beta}/Z(q)$，即
\begin{equation}
 r(q,y)=\beta\log\frac{\pi^*(y\mid q)}{\pi_{\rm ref}(y\mid q)}+\beta\log Z(q).
\end{equation}
再假设偏好符合Bradley--Terry模型，优选回答$y^+$胜过$y^-$的概率为$\sigma(r(q,y^+)-r(q,y^-))$。同一提示的$\log Z(q)$相消，用$\pi_\theta$参数化最优策略便得到
\begin{align}
 a_\theta&=\beta\left[\log\frac{\pi_\theta(y^+\mid q)}{\pi_{\rm ref}(y^+\mid q)}-
 \log\frac{\pi_\theta(y^-\mid q)}{\pi_{\rm ref}(y^-\mid q)}\right],\\
 \mathcal L_{\rm DPO}&=-\E_{(q,y^+,y^-)}\log\sigma(a_\theta),\\
 \nabla_\theta\ell_{\rm DPO}&=-\beta\sigma(-a_\theta)
 [\nabla_\theta\log\pi_\theta(y^+\mid q)-\nabla_\theta\log\pi_\theta(y^-\mid q)].
\end{align}
回答的对数概率是其token条件对数概率之和，提示本身不计入回答似然。这个目标提高相对于参考模型的偏好差，并非单独保证优选回答的绝对概率上升。偏好噪声、回答长度偏差、策略表达能力和有限离线数据都会破坏理想推导的前提；若标注者偏好自信但错误的答案，目标本身不会自动发现错误。KL约束与偏好数据因而不能替代事实和规程验收。

\paragraph{偏好目标的一个数值检查。}
设参考模型对两个完整回答的概率分别为$0.2,0.2$，当前模型为$0.4,0.1$，其他回答占剩余概率。取$\beta=1$，则$a_\theta=\log(0.4/0.2)-\log(0.1/0.2)=\log4$，模型隐含的优选胜出概率为$\sigma(\log4)=0.8$，这对样本损失为$-\log0.8$。若当前模型等于参考模型，隐含胜率是$0.5$、损失为$\log2$。这里的胜率描述偏好模型，不是优选答案事实正确的概率；冻结参考模型是固定比较基准，也不是提供实时事实核验。

\paragraph{解码与生成控制。}
贪心解码每次选择最大概率 token，束搜索维护多个候选序列，温度采样使用
\begin{equation}
 p_{\tau_{\mathrm{dec}}}(x_t=j\mid x_{<t})
 =\frac{\exp(z_{t,j}/\tau_{\mathrm{dec}})}{\sum_{k=1}^{V_{\mathrm{voc}}}\exp(z_{t,k}/\tau_{\mathrm{dec}})}.
\end{equation}
这里$z_t\in\R^{V_{\mathrm{voc}}}$为当前词表 logits，$V_{\mathrm{voc}}$为词表大小，$j$为候选词元编号，$\tau_{\mathrm{dec}}>0$为解码温度，与序列长度$T$区分。
温度越高，分布越平坦。top-$k$保留概率最大的$k$个候选；核采样（nucleus sampling，也称top-$p$）保留按概率降序累加首次达到阈值$p$的最小前缀，再归一化。例如概率为$0.6,0.25,0.15$时，top-$2$与top-$0.8$都保留前两个，变为$0.6/0.85,0.25/0.85$；这里$p$是累计质量门槛，不是只保留单个概率超过$0.8$的词元。事实问答和工具调用通常需要低随机性与结构化输出，创意生成则可以接受更高多样性。

\paragraph{温度、序列概率与缓存的边界。}
固定当前logits，两个token的概率比满足
\begin{equation}
 \log\frac{p_\tau(j)}{p_\tau(k)}=\frac{z_j-z_k}{\tau}.\label{eq:llm-temperature-ratio}
\end{equation}
由式\eqref{eq:llm-temperature-ratio}可知，$\tau\to0^+$时概率集中到最大logit集合，$\tau\to\infty$时在有限词表上趋于均匀。令$\mathcal H(\tau)=-\sum_jp_\tau(j)\log p_\tau(j)$，利用$\mathcal H=\log Z-\E_{p_\tau}[z]/\tau$及$d\E[z]/d\tau=-\Var_{p_\tau}(z)/\tau^2$，得到
\begin{equation}
 \frac{d\mathcal H}{d\tau}=\frac{\Var_{p_\tau}(z)}{\tau^3}\geq0.
\end{equation}
温度增加确实增加这个固定前缀下的分布熵，但不保证生成事实更多样或更正确，因为后续前缀也随采样改变。top-$k$或top-$p$先选择集合$S$再按$p'(j)=p(j)\mathbf1_{j\in S}/\sum_{r\in S}p(r)$归一化，已经改变原始模型分布。

贪心逐步取最大条件概率，也不等于最大化完整序列概率。例如第一步两条路径概率为$0.6,0.4$，第二步各自最优续接概率为$0.5,0.9$，两条最优完整路径概率分别为$0.30,0.36$；贪心会错过第二条。束搜索只保留有限候选，也不能保证全局最优，长度处理和终止规则还会改变比较方式。

设批大小为$B$、层数为$L$、缓存长度为$T$、键值头数为$H_{\rm kv}$、每头维度为$d_h$，每元素占$b$字节，则仅键和值缓存占用约为
\begin{equation}
 M_{\rm KV}=2BLTH_{\rm kv}d_hb.
\end{equation}
例如$B=1,L=2,T=1024,H_{\rm kv}=4,d_h=64,b=2$时，缓存为$2097152$字节，即$2$ MiB；这还没有计入模型权重、当前激活和框架开销。单个新查询仍需与$T$个历史键比较，标准全局注意力的新token成本仍随上下文线性增加，连续生成的累计注意力计算仍可能呈二次增长。缓存节省的是旧键值投影和旧位置的重复计算，并不把任意长上下文变成常数成本。不同键值头共享方案还会改变缓存与表达能力的折中。

生成时只新增一个 token，KV cache 保存已计算的历史键和值以减少重复计算，但缓存占用随上下文增长。缓存管理还需保证会话隔离与敏感内容清理。量化可以降低内存开销，却应在故障码、数值、单位和工具参数上单独验证。低温度和缓存优化只改变生成行为与资源消耗，不能补足模型缺少的有效规程，因此还需要外部检索。

\section{检索、工具与安全边界}
检索增强生成将外部检索证据与生成模型结合，Retrieval-Augmented Generation (RAG) 的原始工作给出了知识密集任务中的代表性实现\cite{lewis2020rag}。工业手册检索还需要额外约束版本、生效时间和权限，这些属于本书的应用规则，不能由生成模型自身保证。
\paragraph{检索增强生成。}
给定查询$q$，检索器返回文档$D_q$，生成器计算$p(y\mid q,D_q)$。系统误差可分为召回错误、排序错误、上下文理解错误和生成错误。评估时应分别测量检索召回率、引用是否支持答案、答案完整性和拒答质量。检索系统不能把不存在的证据变成可靠事实。

\paragraph{从一个问题走完检索链。}
假设问题指定了设备型号和当前日期。离线先把手册切为保留标题、版本和页码的片段，用文档编码器将每段映射为$d$维向量并存入索引；在线把问题映射为同维向量，在权限与生效时间允许的片段中检索候选。双编码器分别编码问题和片段，便于缓存；交叉编码器（cross-encoder）把问题与某个候选共同输入再打分，更适合少量候选重排，不能直接复用独立文档向量。

若返回两段，一段有适用条件但没有处理步骤，另一段有步骤但属于已失效版本，不能因为两段都“相似”就拼成可执行结论。应继续找完整有效证据或拒答。普通RAG查询通常既不更新生成器权重，也不更新检索器权重，只改变送进上下文的文字；文档更新后需同步更新索引，若编码器版本改变还需重建对应表示。检索片段中的词元也占上下文长度，更多片段可能挤掉关键条件。

\paragraph{检索得分、训练目标与证据边缘化。}
设查询编码器$f_\phi$与文档编码器$g_\phi$输出$\R^d$中的非零向量，归一化后得分为$s_\phi(q,d)=f_\phi(q)^{\mathsf T}g_\phi(d)$。给定一个相关文档$d^+$和候选集合$\mathcal C$，对比检索损失为
\begin{align}
 P_\phi(d\mid q,\mathcal C)&=\frac{\exp(s_\phi(q,d)/\tau_r)}{\sum_{d'\in\mathcal C}\exp(s_\phi(q,d')/\tau_r)},\\
 \ell_{\rm ret}&=-\log P_\phi(d^+\mid q,\mathcal C),\qquad \tau_r>0,\\
 \frac{\partial\ell_{\rm ret}}{\partial s_\phi(q,d)}&=\frac{P_\phi(d\mid q,\mathcal C)-\mathbf1_{d=d^+}}{\tau_r}.\label{eq:llm-retrieval-gradient}
\end{align}
若把候选得分视为独立变量，式\eqref{eq:llm-retrieval-gradient}的下降方向提高正例得分、降低负例得分；实际共享编码器还会汇总不同样本的梯度，不保证每对向量距离都按此方向变化。所得概率仅相对于候选集合归一化，不是文档支持答案的校准概率。同型号不同版本尤其适合作为需要区分的候选，却必须避免把另一份同样有效的证据错误标成负例。权限和有效期先定义可搜索集合$\mathcal D(q)$，相似度排序只在这个集合内进行。

为理解检索和生成怎样组合，可把相关文档$d$视为潜变量，写出简化的单文档边缘化模型
\begin{equation}
 p(y\mid q)=\sum_{d\in\mathcal D(q)}p_\phi(d\mid q)p_\theta(y\mid q,d).
\end{equation}
实际仅取top-$K$集合$S_K$。令$Z_K=\sum_{d\in S_K}p_\phi(d\mid q)$，截断并重新归一化的分布为
\begin{equation}
 p_K(y\mid q)=\frac1{Z_K}\sum_{d\in S_K}p_\phi(d\mid q)p_\theta(y\mid q,d).
\end{equation}
若$0<Z_K<1$，完整分布满足$p(y\mid q)=Z_Kp_K(y\mid q)+(1-Z_K)p_{\rm tail}(y\mid q)$，故两者对任意回答事件的概率差不超过$1-Z_K$。这个界解释了截断遗漏的影响，但工业向量得分通常没有给出真实全库后验，因此不能从top-$K$分数直接宣称已知$Z_K$。此外，将多个片段拼接到一个提示里与上述单文档混合并不等价，跨段落推理还需要联合证据。

令$E$表示检索上下文包含足够有效证据，$A$表示答案正确。全概率公式给出
\begin{equation}
 P(A)=P(E)P(A\mid E)+P(E^{\rm c})P(A\mid E^{\rm c}).
\end{equation}
证据缺失时模型可能凭已有知识答对，所以一般不能声称正确率必然低于检索召回；但若验收事件要求“正确且由当前有效证据支持”，则它是$E$的子事件，其概率必不超过$P(E)$。这正是分别评价召回、答案正确性和引用支持率的理由。增加$K$可能提高召回，也可能加入冲突或噪声，需要固定上下文预算作分层验证。

工业知识问答应把 RAG 拆成可独立验收的流水线。离线阶段按照设备、部件、工况和有效时间切分手册，保留章节标题、表格关系、版本号和页码；切块时不能把“警告条件”和“处理步骤”拆到互不相邻的向量中。检索阶段可以先用 Best Matching 25 (BM25) 召回精确故障码，再用向量检索补充语义相近描述，最后由 cross-encoder 重排序。生成阶段要求模型输出答案、证据片段、文档版本和不确定性状态；当证据覆盖不足时，应拒答或请求更多设备上下文。

评估集需要按问题类型分层：直接查表、跨段落组合、版本冲突、过期规程、未登录故障码和恶意提示。对每个问题分别记录检索召回、证据支持率、答案正确率、引用准确率和拒答质量。一个语言模型即使在开放问答上流畅，也可能在版本冲突时选择了旧规程；因此文档有效期和权限过滤必须作为检索条件，而不是依赖生成模型自行判断。

\paragraph{工具返回值不是模型自己算出的事实。}
例如回答设备记录查询时，模型先产生包含工具名、设备编号和时间范围的请求；外部程序检查格式与权限后才执行查询，将结果作为新的上下文返回，模型再据此组织文字。请求格式合法不等于设备编号存在，更不等于查询已成功。超时、空结果或版本冲突应作为显式状态传回，不能让模型把未返回的数据补成“合理数值”。这一过程是推理中的外部交互，不会因为调用一次数据库就自动更新模型权重。

\paragraph{工具调用与评测。}
工具调用把模型输出限制为结构化参数，再由外部程序执行查询、计算或控制动作。工业系统应采用白名单工具、参数校验、权限隔离和人工确认。评测集需要覆盖正常任务、边界输入、对抗提示、过时知识和工具失败。困惑度与任务准确率并不等价，LoRA 的参数量优势也必须结合任务迁移效果和推理开销评价。

\paragraph{幻觉、提示注入与安全边界。}
幻觉不仅是“模型说错话”，还包括引用不存在、混淆设备型号、把条件性建议表达成确定事实。系统应允许模型在证据不足时拒答，并显示证据覆盖范围。提示注入则利用文档或用户输入改变系统指令；因此检索文本必须被视为不可信数据，工具调用必须通过结构化参数、权限检查和资源范围限制。

\paragraph{大模型评测矩阵。}
工业评测应至少包含四个轴：语言能力、任务正确性、证据忠实度和系统可靠性。语言能力可测困惑度或结构化遵循；任务正确性应使用专家标注或可验证程序；证据忠实度检查答案是否被文档支持；系统可靠性则记录延迟、成本、拒答、超时和工具失败。每个轴都应按设备型号、工况、输入长度和缺失信息分层报告。

以下材料沿“因果语言建模—示范微调—低秩适配—偏好学习—外部知识”组织，不再重复通用 Transformer 入门。官方网页访问不稳定时，可从相应项目的公开文档源文件核对示例；接口版本应以实际安装环境为准。
\section*{辅助学习材料}
\begingroup
\emergencystretch=3em
\begin{itemize}
\item \textbf{Causal language modeling}（Hugging Face，英文教程）。跟随文本分块、标签构造、训练与生成，核对“Token、表示与语言建模”中的下一 token 目标及“预训练目标与困惑度”。示例是在已有模型上继续训练，不是从零预训练大模型的完整配方。\href{https://huggingface.co/docs/transformers/tasks/language\_modeling}{在线阅读}。
\item \textbf{SFT Trainer}（Hugging Face TRL，英文教程）。重点读指令数据格式、标签右移、只对回答计损失与指令微调示例，把本章“监督微调与参数高效适配”落实为训练样本与损失掩码；再比较仅回答计分和整段文本计分的差别。\href{https://huggingface.co/docs/trl/sft\_trainer}{在线阅读}。
\item \textbf{PEFT：在低资源硬件上对十亿规模模型进行参数高效微调}（Hugging Face，中文博客）。阅读 LoRA 配置、可训练参数统计和适配器保存／加载示例，对照本章低秩更新的前向与反向计算，区分冻结基座和训练适配器。文中低精度训练示例不等于 QLoRA 的完整说明，也不意味着显存与参数按同一比例减少。\href{https://huggingface.co/blog/zh/peft}{在线阅读}。
\item \textbf{ChatGPT 背后的“功臣”——RLHF 技术详解}（Hugging Face，中文图解博客）。按预训练、奖励模型与强化学习微调三步读图，识别偏好数据和参考策略各自的作用，再与本章 DPO 推导比较：DPO 直接优化相对回答概率，并非照搬奖励模型加策略优化的训练流程。\href{https://huggingface.co/blog/zh/rlhf}{在线阅读}。
\item \textbf{一文彻底搞懂大模型 - RAG（检索、增强、生成）}（CSDN，中文博客）。阅读知识库构建、向量检索、排序与上下文融合流程，对照本章“检索增强生成”，把召回错误与生成错误分开检查；图示便于理解管线，但“增强”不保证事实正确，仍须执行本章的版本、权限和引用支持率验收。\href{https://blog.csdn.net/a2875254060/article/details/142468037}{在线阅读}。
\item \textbf{吴恩达 x OpenAI 的 Prompt Engineering 课程专业翻译版}（中英双语字幕视频，Datawhale 教程收录）。选看迭代提示、文本总结与聊天机器人示例，练习把维修记录转成受约束的输出，对照本章生成控制与任务评测；提示只改变输入条件，不是 SFT，也不能替代上面的训练材料。\href{https://www.bilibili.com/video/BV1Bo4y1A7FU/}{观看课程}，\href{https://github.com/datawhalechina/llm-cookbook}{中文配套教程}提供示例。名称、字幕说明和链接已核对教程目录；视频页访问受限，未逐帧观看，未核实播放量。
\end{itemize}
\endgroup
\paragraph{自测与参考答案。}
为了保持初始输出与基座一致，能否把LoRA的$A,B$同时设为零？参考答案：不能依靠这种初始化开始通常的梯度学习，因为$Ax=0$使$\nabla_B=0$，$B=0$又使$\nabla_A=0$。可令$A$随机、$B=0$，初始更新仍为零，但$B$通常能先获得梯度。再问：只把新规程放进RAG上下文是否等于完成一次微调？答案是否定的；它改变当前条件输入，不自动更新任何模型参数。

\paragraph{本章小结。}
语言大模型的能力来自token表示、因果建模、规模化训练和对齐过程的共同作用，但这些环节优化的对象不同。最大似然逼近语料分布，LoRA限制可更新的方向，偏好优化调整相对回答倾向，温度与截断改变解码分布，检索则改变当前能够引用的证据。把这些数学对象区分开来，才能判断某项改进究竟降低了什么误差，而不是把更流畅、更受偏好和更可信混为一谈。

工业应用的关键不在于把模型接入界面，而在于建立数据版本、检索证据、工具权限、输出校验和持续评测组成的控制链。本章沿用型号、生效时间和预测时刻的信息边界，使文字答案始终能返回原始来源。对于故障码、量值和复检条件，正确复制与精确引用往往比措辞自然更重要；证据不足时的拒答也是系统输出的一部分。

文字还需要与现场记录对应。“此处裂纹”指图中的哪个区域？“升温后异响”对应哪一段传感器记录？下一章讨论怎样配对、对齐和融合这些输入。即使放进共同的表示空间，时间戳、坐标和记录质量也要保留；几路观测有矛盾时，应指出矛盾并核查，不能靠一段流畅的描述把它们掩盖过去。

\chapter{多模态模型}
\begin{quote}\small
\textit{本章摘要。}本章研究图像、文本、传感器和时间序列如何在同一任务中协同。先区分共同表示空间与真正的信息融合，再讨论对比学习、跨模态注意力、模态缺失和时间对齐，最后建立面向工业数据的评价方式。关键问题不是模态越多越好，而是每个模态是否提供了可对齐、可验证的互补信息。
\end{quote}
序列、视觉和语言模型各自提供证据后，还需要判断它们是否描述同一对象和同一阶段。本章先区分共享空间与信息融合，再处理对齐误差和模态缺失。联合诊断的类别标签与预警的未来事件标签必须分别定义；维修结论只有在形成后才能作为诊断证据，不能提前进入预警模型。

把照片和振动记录放在一起之前，先核对设备编号和采集时间。若照片来自维修后，振动来自维修前，直接配成一对就可能教错模型。共同嵌入空间用于寻找匹配记录，跨模态注意力让一种输入读取另一种输入，质量门控则在照片模糊或传感器失真时调整使用程度。这几种办法需要不同的训练数据。后文都用相机与振动联合诊断来说明：多出来的信息帮了什么忙，少一路输入时又该怎么办。

\section{不同模态的共同空间}

% auxiliary-resource-guide
本章末的“辅助学习材料”列出与核心方法对应的讲解和实践入口，可在阅读相关推导后，按所列小节对照学习。

\begin{figure}[htbp]
\centering
\begin{tikzpicture}[node distance=0.75cm and 0.9cm]
 \node[idi input, minimum width=1.8cm] (img) {图像\\$v_i$};
 \node[idi input, below=0.8cm of img, minimum width=1.8cm] (txt) {文本\\$t_i$};
 \node[idi process, right=of img, minimum width=1.8cm] (ev) {视觉编码器\\$f_v$};
 \node[idi process, right=of txt, minimum width=1.8cm] (et) {文本编码器\\$f_t$};
 \node[idi state, right=1.0cm of $(ev)!0.5!(et)$, minimum width=2.0cm] (space) {共同空间\\余弦相似度};
 \node[idi output, right=of space, minimum width=2.0cm] (loss) {对比损失\\匹配/非匹配};
 \draw[idi arrow] (img)--(ev); \draw[idi arrow] (txt)--(et);
 \draw[idi arrow] (ev)--(space); \draw[idi arrow] (et)--(space); \draw[idi arrow] (space)--(loss);
\end{tikzpicture}
\caption{双塔对比学习先把不同模态映射到共同空间，再用配对关系约束表示距离。}
\label{fig:multimodal-dual-encoder}
\end{figure}
图\ref{fig:multimodal-dual-encoder}的两路编码先独立计算，再在共同空间比较匹配关系；这种交互位置允许分别缓存表示，却未提供区域级联合推理。
模态（modality）是信息的表现形式，例如图像像素、文本词元和传感器数值；同一事件的不同模态仍可能有不同采样时间、噪声和缺失模式。编码器把一种原始输入变成特征，投影层再把不同宽度的特征映射到相同宽度；共同嵌入空间指这些输出可以用同一距离或内积比较，不要求原始输入格式相同。图文对比预训练（Contrastive Language--Image Pretraining，CLIP）通过图文对比学习获得可比较表示，BLIP-2 以轻量查询模块连接冻结的视觉与语言模型，Large Language and Vision Assistant (LLaVA) 通过视觉指令微调学习多模态交互\cite{radford2021clip,li2023blip2,liu2023llava}。BLIP-2 中的查询变换器称为 Querying Transformer (Q-Former)。三者分别强调对齐、表示连接和指令响应，并非同一损失的规模变化。
本章“对齐”指不同模态在对象、时间或表示上的对应，不同于上一章使回答符合指令与偏好的行为对齐；图文匹配良好不表示回答已符合规程。多模态模型首先要决定“对齐”发生在哪里。CLIP 采用图像编码器$f_v$和文本编码器$f_t$构成双塔结构，经过归一化后以温度参数$\tau$计算相似度：
\begin{equation}
 s_{ij}=\frac{f_v(v_i)^{\mathsf T}f_t(t_j)}{\tau}.
\end{equation}
批内对比目标让匹配图文相似度高于非匹配组合，因此 CLIP 能把“有裂纹”“轴承外圈剥落”等文本描述映射为分类原型。它的优点是图像和文本编码可以独立缓存，缺点是双塔交互浅，难以表达复杂问答和精确区域关系。

进入本章前，应能识别编码器输出的形状，理解内积、softmax和用标签计算损失；不必先训练一套视觉或语言大模型。先判断任务是寻找匹配记录、联合分类还是生成报告，再选双塔、融合头或语言连接模块。它们所需的监督分别可能是配对索引、类别标签和目标回答，不能只因输入都是图文就混用训练目标。

\paragraph{双塔的输入输出与一个两对样本的计算。}
设一批有$N=2$对图文，图像分支汇聚输出$2\times d$矩阵，文本分支也输出$2\times d$矩阵；行向量归一化是把每行除以其长度，要求非零向量。图像矩阵乘以文本矩阵的转置，再除以温度，得到$2\times2$得分表，其中对角线对应真实配对，非对角线是本批次用于比较的其他组合。设$d=2$，图文表示都依次为$(1,0),(0,1)$，取$\tau=1$，则得分表为$\left[\begin{smallmatrix}1&0\\0&1\end{smallmatrix}\right]$。每行正确配对概率为$e/(e+1)\approx0.731$，损失为$-\log0.731\approx0.313$。表示已经完全正交也不意味着softmax概率为一，因为它仍比较有限得分。

构造批次时，图文应按同一配对索引同步打乱；若单独重排文本，必须同步更新目标索引，不能继续把对角线当作真配对。训练时配对索引告诉模型哪个组合应更接近，梯度更新编码器或投影层；推理时不需要提供“正确配对”或重新计算训练损失。例如先缓存三种缺陷描述的向量，再把新图与它们比较，输出最相似的描述。所谓零样本分类只表示未用当前类别的标注训练专门分类头，不表示没有预训练；候选只有三类也不说明真实输入一定属于其中某类，需要另设拒识与校准。

BLIP-2 用冻结视觉编码器、冻结语言模型和轻量 Q-Former 连接二者。Q-Former 通过一组可学习查询从视觉特征中提取与文本任务相关的信息，再投影到语言模型输入空间；这样可以显著减少端到端训练成本，但视觉细节会受查询数量和投影瓶颈限制。LLaVA 类结构则通常把视觉编码器输出经过线性投影后拼接到语言模型上下文中，适合图像问答和报告生成。工业场景应把生成式描述与可定位模型分开验收：语言模型可以组织证据，不能凭文字流畅性证明缺陷确实位于某个像素区域。

\paragraph{连接语言模型时，向量还要经过什么步骤。}
假设视觉编码器产生$64\times128$个数，即64个区域、每区128维。Q-Former可用八个可学习查询读取这些区域，输出$8\times d_q$，再投影为$8\times d_{\rm lm}$，作为语言模型的条件输入。查询是训练得到的向量，不是八句人工问题。以LLaVA式直接投影为例，$128\times d_{\rm lm}$矩阵可把全部区域转成$64\times d_{\rm lm}$，与文字嵌入在序列维拼接。两条路线分别压缩视觉位置数与保留更多视觉位置，都会消耗语言上下文和计算预算。

训练图像描述时，语言部分仍预测目标回答的下一词元；推理时固定参数，给定图像和问题逐词生成。冻结视觉或语言权重不等于无需前向计算，也不等于从文字损失到连接模块之间无需传梯度。不同训练阶段可冻结不同部分，应报告实际设置。CLIP主要输出匹配分数，连接语言模型才获得文本生成路径；生成路径存在也不保证细裂纹的区域证据被保留。

对于传感器、图像和文本，常见做法是先分别得到$z_s,z_v,z_t$，再进行加权融合或跨模态注意力。若传感器序列的采样频率远高于图像帧率，可以用时间窗池化、事件 token 或 cross-attention 将高频状态压缩到视觉时间点；若某个模态临时失效，则训练时加入模态 dropout，并让质量门控$g_m$根据缺失率、噪声估计和传感器健康度动态调整贡献。多模态系统的关键不是把向量简单拼接，而是明确每种模态提供何种互补证据。

\begin{table}[htbp]
\centering
\caption{多模态模型的交互位置与适用任务}
\label{tab:multimodal-interaction}
\scriptsize
\setlength{\tabcolsep}{3pt}
\resizebox{\linewidth}{!}{%
\begin{tabular}{p{2.7cm}p{3.7cm}p{3.8cm}p{3.8cm}}
\hline
模型路线 & 交互机制 & 优势 & 工业风险与适用任务 \\
\hline
CLIP 双塔 & 图像--文本共享嵌入空间 & 编码可缓存、零样本分类方便 & 细粒度定位弱，适合筛查和检索 \\
BLIP-2 & Q-Former 从视觉特征提取查询 & 冻结大模型即可迁移 & 查询瓶颈可能丢失小缺陷，适合报告草拟 \\
LLaVA 类 & 视觉 token 投影到语言上下文 & 问答和多轮解释自然 & 生成内容需外部证据，不能代替定位器 \\
交叉注意力融合 & 查询模态读取键值模态 & 能表达局部、时间和条件关系 & 计算较重，适合联合诊断与根因分析 \\
\hline
\end{tabular}}
\end{table}

表\ref{tab:multimodal-interaction}对照共享嵌入、查询瓶颈和跨模态注意力，选型时须把编码成本与所需细粒度关系同时考虑。
以“相机图像加主轴振动诊断”为例，先由视觉编码器输出工件区域 token，由时序编码器输出若干转速同步的振动 token，再将维修记录编码为文本证据。模型可以采用两阶段决策：第一阶段用 CLIP 或轻量分类器筛选明显正常与明显异常样本，第二阶段仅对异常候选执行跨模态注意力和根因排序。训练标签应同时包含缺陷类别、严重度、发生时间和证据模态；如果只有最终维修结论，没有记录哪一帧图像或哪段频谱支持结论，模型即使获得较高准确率也难以审计。

联合损失可以写为
\begin{equation}
 \mathcal{L}=\lambda_{\rm cls}\mathcal{L}_{\rm cls}+\lambda_{\rm align}\mathcal{L}_{\rm InfoNCE}+\lambda_{\rm miss}\mathcal{L}_{\rm miss}+\lambda_{\rm rank}\mathcal{L}_{\rm rank},
\end{equation}
其中分类项约束业务输出，对比项约束共同空间，模态缺失项要求模型在删去某一路输入时仍保持可接受风险，排序项则约束根因候选的相对顺序。权重不能只按数值大小设置，应根据漏诊代价、人工复核容量和各模态的标签可信度进行敏感性分析。
图像、文本、音频和传感器序列具有不同的统计结构。多模态模型通常为每种模态建立编码器$z_m=f_m(x_m)$，再通过共享投影或跨模态注意力交互。对配对图文$(x_i,t_i)$，对比学习可最大化匹配样本相似度：
\begin{equation}
 \mathcal{L}_{\mathrm{NCE}}=-\frac{1}{N}\sum_i\log\frac{\exp(s(z_i^x,z_i^t)/\tau)}{\sum_j\exp(s(z_i^x,z_j^t)/\tau)}.
\end{equation}
这里$N$为配对样本数，$i,j=1,\ldots,N$，$z_i^x$与$z_i^v$均表示第$i$幅图像的单位范数表示，$s(a,b)=a^{\mathsf T}b$是不含温度的余弦相似度。因此$s_{ij}=s(z_i^v,z_j^t)/\tau$是已经缩放的对比得分，不能再次除以温度，且$\tau>0$。

\paragraph{融合方式。}
早期融合将原始特征拼接，晚期融合分别预测后再组合，中间融合在表示空间通过 cross-attention 交互。编码器--解码器架构还可以让一个模态条件化另一个模态的生成。工业场景中，时间同步、传感器缺失和模态质量差异往往比网络结构更决定系统效果。

\paragraph{拼接、加权和交叉读取的区别。}
若图像和振动分别输出三维、二维向量，拼接得到五维向量，可接$K\times5$分类矩阵预测$K$种状态；直接相加则不合法，必须先投影到同维度。晚期融合可以分别得到两个$K$类概率向量，再按权重组合，但要求类别含义和概率尺度可比较。中间交互则保留多个区域或时间片，让不同查询读取不同证据，不先压成各一个全局向量。

例如图片只显示表面颜色变化，振动给出周期冲击，两者联合可以帮助区分状态；若图片已丢掉裂纹细节，简单增加文字不能恢复真实像素。训练融合分类器还需要联合任务标签，只有成对图文并不自动提供故障类别、严重度或因果来源标签。前述联合损失是可选设计而非每个模型必备四项；缺失项和排序项应给出具体定义后才能复现实验。

\paragraph{多模态推理。}
图像问答、视觉定位和设备状态描述都可看作条件生成。应分别评估感知正确性、跨模态对齐和语言表达，避免只用语言流畅度掩盖视觉判断错误。

\paragraph{对比学习的温度参数。}
设图像和文本编码器输出归一化表示$z_i^v,z_i^t$，相似度取余弦相似度。温度$\tau$控制 softmax 的尖锐程度：
\begin{equation}
 s_{ij}=\frac{(z_i^v)^{\mathsf T}z_j^t}{\tau}.
\end{equation}
小温度会放大正负样本差异，梯度更集中于难负样本，但也可能使训练不稳定。对称的图到文、文到图损失通常比单向损失利用更多配对信息。

\paragraph{跨模态注意力。}
以文本为查询、视觉 patch 为键和值时，文本 token 可以选择性读取图像区域。为说明其计算含义，令$Q\in\mathbb R^{n_t\times d}$、$K\in\mathbb R^{n_v\times d}$、$V\in\mathbb R^{n_v\times d_v}$。矩阵$QK^{\mathsf T}$的第$(a,b)$项是第$a$个文本查询与第$b$个视觉键的内积；除以$\sqrt d$是为了在各维独立、方差为一时保持内积方差在可控范围。对每个文本位置$a$，定义
\[
\alpha_{ab}=\frac{\exp(q_a^{\mathsf T}k_b/\sqrt d)}{\sum_{r=1}^{n_v}\exp(q_a^{\mathsf T}k_r/\sqrt d)},\qquad \sum_b\alpha_{ab}=1.
\]
于是输出第$a$行为$h_a'=\sum_b\alpha_{ab}v_b$，它是视觉值向量的凸组合，而不是一个自动产生定位真值的坐标。若缺陷区域很小，$n_v$过小或下采样会使多个区域共享一个 token；若时间掩码错误，注意力仍会在错误窗口内给出看似合理的权重。实际部署应同时检查注意力之外的区域标注、遮挡实验和跨设备稳定性。
矩阵形式为
\begin{equation}
 H_{text}'=\softmax\left(\frac{Q_{text}K_{image}^{\mathsf T}}{\sqrt d}\right)V_{image}.
\end{equation}
此处$Q_{text}\in\R^{n_t\times d}$、$K_{image}\in\R^{n_v\times d}$、$V_{image}\in\R^{n_v\times d_v}$，$n_t,n_v$分别为文本与视觉 token 数。各行存放列向量的转置，softmax 沿视觉位置逐行计算，输出为$n_t\times d_v$矩阵。
这与简单拼接不同，因为每个文本位置可使用不同的视觉证据。多层交互可以支持图像描述、区域问答和视觉定位，但其对齐质量仍受配对数据噪声影响。

\paragraph{交叉注意力并没有把两个轴变成同一个轴。}
若两个文本查询读取三个视觉区域，则权重是$2\times3$矩阵，而不是$5\times5$。假设两行归一化权重分别是$(1/2,1/4,1/4)$和$(0.1,0.2,0.7)$，三个标量值为$2,4,10$，则输出为$4.5,8$；两个词元从同一张图读取了不同组合。输出行数跟查询数一致，要得到按视觉区域排列的结果，就需另设视觉查询路径或定位头，而不能直接把两个输出当成三个区域的标签。

\paragraph{模态缺失与质量加权。}
真实工业系统常出现相机被挡、传感器掉线或文本记录晚到。模型应显式接收“这一模态当前是否可用”的掩码$m_m$，并考虑
\begin{equation}
 z=\sum_m\frac{m_m\exp(a_m)}{\sum_r m_r\exp(a_r)}z_m.\label{eq:multimodal-quality-fusion}
\end{equation}
其中$m,r$遍历模态，$m_m\in\{0,1\}$表示第$m$个模态是否可用，$a_m$为有限的质量得分，各$z_m\in\R^{d_f}$须投影到相同维度，归一化权重是标量，融合结果$z\in\R^{d_f}$而非各模态维度之和。式\eqref{eq:multimodal-quality-fusion}要求$\sum_r m_r>0$；全部缺失时不计算这一比值，而应拒绝输出或调用预先验证的回退流程。训练时进行模态 dropout，能减少模型对某一模态的脆弱依赖，但必须保留真实缺失模式。

\paragraph{缺一路输入时，分母也必须改变。}
取两路质量得分$a_v=\log3,a_s=0$，完整输入时权重为$3/4,1/4$；若图像不可用，掩码变为$m_v=0,m_s=1$，权重应变为$0,1$，而不是仍把传感器乘以$1/4$。一个数值全零的图像可能是真实黑图，也可能是缺失填充值，必须用独立可用性掩码区分。模态dropout是训练时按设计暂时屏蔽某一路，教模型适应不完整输入，不是部署时随机扔掉可用观测。

共同空间的归一化同样需要边界处理：零向量不能除以自身长度，应标记为无效表示并检查输入或编码器；用小常数保护分母虽可防止除零，却不能让零向量成为单位方向。实现时应先清理缺失占位或跳过该分支，再参与融合；不能指望乘以零消除非数值（NaN），因为浮点运算中$0\times\mathrm{NaN}$仍为NaN。质量得分只用于相对加权，不自动等于该模态判断正确的校准概率。

一种可复现的缺失项是对至少保留一路的掩码模式重新前向计算，并对同一任务标签计算交叉熵，再在这些模式上平均。但若缺陷只能在图片中观察到，删图后标签在剩余输入下可能不可判定，不能强迫模型仍然自信答对；应允许不确定性或拒识，并按缺失模式分别验收。质量门控也不是裁判真相的机制，两路证据矛盾时还需检查数据来源。

\section{时间对齐、缺失鲁棒性与工业诊断}
\paragraph{时间对齐的多模态传感器。}
视觉帧、振动窗口和控制指令可能不在同一时间尺度。可以用时间编码、可学习延迟或 cross-attention 在邻近时间段内寻找对应信息。若把未来传感器值错误地对齐到当前图像，离线结果会明显偏高，部署时却失效。因此多模态数据集必须保存采集时间、传输延迟和有效时间区间。

设第$m$种模态在时间窗$W_i=[t_i^-,t_i^+]$内产生输入$x_i^{(m)}$，编码器输出$z_i^{(m)}
\in\mathbb R^d$。所谓配对，至少应满足设备标识一致、事件窗符合任务规定的重叠或滞后关系，并且作为输入的记录在预测时刻已经到达。监督标签可以在事后确认；训练时可用它计算损失，但不能把它或其派生字段提前作为模型输入。若任务要求同一事件的同步观测，而两个模态在校正允许滞后后仍无有效窗重叠，将它们强行作为正对会把时间错配当成监督噪声。对于存在未知延迟的系统，可在训练中显式搜索有限延迟集合$\Delta$，例如令
\[
\operatorname{sim}_i=\max_{\delta\in\Delta}\operatorname{sim}\bigl(z_i^{(v)},z_{i,\delta}^{(s)}\bigr),
\]
但最大化也可能选择偶然相似的错误事件，因此必须在留出设备和固定延迟范围上验证，而不能把对齐搜索结果直接解释为真实因果延迟。

\paragraph{按时钟核对一条样本。}
在10:00作出诊断时，振动窗口可以覆盖09:59:50至10:00，照片虽在09:59:58拍摄，若10:00:03才传到服务器，当前预测仍应把图像标记为不可用。10:30人工确认的缺陷类别可以作为这条样本的事后监督标签，却不是10:00已知的输入。把图像插值到最近传感器时点时，同样要检查插值是否用到了决策后的帧；按采集时间“对齐”不等于按可用时间因果合法。

将1000 Hz振动的十秒窗口与每秒一帧的图像配对时，振动原始输入有10000个采样点，不能仅因为都属于十秒就逐元素拼接。可先把振动分成十个因果一秒片段，每片编码为$d_s$维，再与对应图像区域交互；压缩也可能抹去短冲击，需比较不同片长。若模态间确有物理传播滞后，应保留有依据的滞后窗口，而非一律强求同一时刻。

\paragraph{图像与工艺曲线联合诊断。}
视觉编码器提取表面纹理，序列编码器提取温度和压力变化，再用融合层预测缺陷类型与严重度。消融实验至少包括仅视觉、仅时序、简单拼接和跨模态注意力四种设置。若融合模型只在完整数据上有效，还需报告相机缺失或曲线缺失时的退化曲线。

\paragraph{对齐、融合与协同的区别。}
对齐回答“不同模态是否指向同一个对象或事件”，融合回答“如何联合使用这些信息”，协同则进一步要求联合表示在任务上优于任何单一模态。三者不能混为一谈：图像和文本可以在共享空间中相近，却未必能完成精确定位；多个传感器可以简单拼接，却没有学到跨模态因果关系。

\begin{figure}[htbp]
\centering
\begin{tikzpicture}[>=Latex, node distance=0.8cm, every node/.style={font=\small}]
 \node[draw, rounded corners, fill=blue!10, minimum width=2.4cm] (vision) {视觉编码器};
 \node[draw, rounded corners, fill=orange!12, right=2cm of vision, minimum width=2.4cm] (text) {文本编码器};
 \node[draw, rounded corners, fill=green!12, below=1.1cm of vision, minimum width=2.4cm] (sensor) {传感器编码器};
 \node[draw, rounded corners, fill=yellow!20, right=2cm of sensor, minimum width=2.4cm] (fusion) {对齐/融合层};
 \node[draw, rounded corners, fill=red!10, right=2cm of fusion, minimum width=2.4cm] (out) {诊断或生成};
 \draw[->, thick] (vision)--(fusion); \draw[->, thick] (text)--(fusion); \draw[->, thick] (sensor)--(fusion); \draw[->, thick] (fusion)--(out);
\end{tikzpicture}
\caption{工业多模态系统中，编码、对齐、融合和任务输出是不同层次。}
\label{fig:multimodal-industrial-fusion}
\end{figure}
图\ref{fig:multimodal-industrial-fusion}将各模态编码器与融合层分开，说明原始输入先形成各自表示，再由对齐与融合环节为诊断或生成提供联合证据。

\paragraph{对称对比损失。}
图到文和文到图两个方向都应参与训练。令$S_{ij}=s_{ij}$为已经除以温度的得分矩阵，且$i,j=1,\ldots,N$，则对称损失为
\begin{equation}
 \mathcal{L}=\frac12\left[-\frac1N\sum_i\log\frac{e^{S_{ii}}}{\sum_j e^{S_{ij}}}
 -\frac1N\sum_i\log\frac{e^{S_{ii}}}{\sum_j e^{S_{ji}}}\right].
\end{equation}
第一项固定图像$i$、沿第$i$行在文本间归一化；第二项固定文本$i$、沿第$i$列在图像间归一化，两者都把对角位置作为目标。因此第二项分母使用$S_{ji}$而非$S_{ij}$。这里正样本指匹配的图文对，负样本指作为不匹配候选的图文对，与设备“正常／故障”的类别无关。若$N=1$且没有额外负样本，两方向的分子分母相同，损失恒为零，也没有对比分数梯度；因此损失小不一定表示已学到匹配。批次中的负样本并不都是真负样本：两个不同时间的图像可能描述同一设备状态，两个文本也可能共享相同故障原因。因此对比学习需要检查假负例（false negative，即实际匹配却被当作负例的组合），并在必要时使用组标签、软标签或难负样本过滤。

\paragraph{模态不平衡与缺失鲁棒性。}
某一模态信噪比更高时，模型可能忽略其他模态；训练损失下降并不表示融合真正发生。可以通过模态 dropout、分支独立监督、梯度平衡和质量门控减轻这一问题。测试时应绘制从完整输入到逐模态缺失的性能曲线，并区分随机缺失、连续掉线、延迟到达和系统性遮挡。

\paragraph{时间对齐与事件粒度。}
图像帧、振动窗口和控制指令的时间粒度不同。对齐时要明确事件发生时间、采集时间、传输时间和标签确认时间。窗口重叠会增加样本条目数，但不等于增加同样多的独立信息，还可能把同一故障事件复制到训练和测试两侧。应以事件或设备为单位划分，并在报告中说明窗口长度、步长和允许的预测提前量。

\paragraph{一个能区分互补信息与模型规模的评价。}
先固定同一批测试事件、同一分类标签与决策时间，在相同调参预算下比较仅图像、仅振动、拼接与交叉注意力，再按照片模糊、传感器掉线和两路冲突分别统计。假设一百个事件中图像答对八十个、振动答对七十五个，这不足以推出融合能答对九十个：两路错误可能高度重合，也可能互补，须逐事件比较预测。还可打乱图像与振动的配对关系作对照；若性能不变，可能没有利用对应关系，但也可能任务本来不需要它，不能仅据此断定实现错误。

用于报告生成时，至少把“识别对了什么”“证据对应哪帧、哪段信号”“文字是否忠于证据”分别计分。正确类别加上编造的位置不算定位正确；只会在完整输入下生成流畅报告，也不能代替缺失模式下的拒识验证。所有示意数值用于解释评价方法，不代表本书已完成这些模型的实测比较。

\paragraph{多模态评价。}
多模态系统应分别报告单模态基线、简单拼接、注意力融合和缺失模态结果。分类任务可检查跨模态遮挡后的鲁棒性；定位任务需要区域级准确率；生成任务需要事实一致性和证据覆盖率。最重要的消融不是把某一层删掉，而是回答“哪一种模态信息在什么工况下真正改变了决策”。

多模态学习的共同嵌入、对比目标与跨模态注意力可用下列材料互补复习，分别侧重表示学习和工程实现。
\section*{辅助学习材料}
\begingroup
\emergencystretch=3em
\begin{itemize}
\item \textbf{Contrastive Representation Learning}；作者/平台：Lilian Weng；英文；技术教程。重点看 contrastive loss、InfoNCE、temperature 和 negative sampling，帮助理解本章的 CLIP 双塔、对齐空间与模态缺失风险。\href{https://lilianweng.github.io/posts/2021-05-31-contrastive/}{在线阅读}。
\item \textbf{Natural language image search with a Dual Encoder}；作者/平台：Khalid Salama / Keras；英文；代码教程。阅读双编码器、投影层、对比目标和文本检索图像部分，理解共享表示空间。示例使用 Keras 2，运行前需核对依赖版本。\href{https://keras.io/examples/vision/nl\_image\_search/}{在线阅读}。
\item \textbf{多模态学习综述 (MultiModal Learning)}；知乎；中文技术解说。区分表示、对齐与融合，核对不同模态的数据配对方式；标题与主题据必应索引，原页受限。\href{https://zhuanlan.zhihu.com/p/582878508}{在线阅读}。
\item \textbf{一文彻底搞懂多模态：模态表示、多模态融合、跨模态对齐（索引标题节选）}；CSDN；中文博客。比较早期与后期融合，另行检查缺失模态下的评价；标题与主题据必应索引，原页受限，完整标题未核实。\href{https://blog.csdn.net/star\_nwe/article/details/143416379}{在线阅读}。
\item \textbf{多模态 之 CLIP ｜ 1、CLIP 网络结构精讲}；上传者：Enzo\_Mi；bilibili；中文技术讲解。对照本章图文对齐，画出图像编码器、文本编码器和相似度计算，比较对比学习与融合；不要把图文匹配直接等同于生成能力。2026年10月4日公开元数据播放4,447次；优先选择结构讲解而非成品展示。视频落地页显示验证码，未绕过；标题、上传者和计数据官方公开元数据核对，未逐帧观看。\href{https://www.bilibili.com/video/BV1BeZFBGEhe/}{观看视频}。
\end{itemize}
\endgroup
\paragraph{自测与参考答案。}
沿用10:00诊断例：09:59:58拍摄、10:00:03才到达的照片，与10:30确认的缺陷类别，哪一个可用于10:00预测，哪一个可用于事后训练监督？参考答案：照片尚未到达，不能作为该时刻输入，应标记缺失；确认类别可作为监督标签，但不能回填成当时已知的特征。若图像缺失而传感器可用，质量权重须重新归一化为$0,1$；若两路均缺失，分母为零，应进入拒识或已验证的回退流程。

\paragraph{本章小结。}
图像、振动和文字的采集频率不同，质量也不同，还可能缺失。联合使用时，要说明每一路怎样编码、按什么时间和设备配对、相互矛盾时听谁的，以及少一路输入时还能完成什么任务。最后在独立数据上分别测试完整输入和缺失输入，才能判断复杂的融合结构是否真的有帮助。

\chapter{深度生成模型}
\begin{quote}\small
\textit{本章摘要。}本章从“学习数据分布而非只学习标签”出发，依次解释变分自编码器、对抗生成、可逆密度建模、扩散模型及其条件控制，再比较生成质量、覆盖性、可控性和计算代价。ELBO、噪声反演和物理约束把抽象生成目标连接到工业仿真、数据增强和异常样本分析。
\end{quote}

在贯穿案例中，生成模型不能通过制造合成故障来改变真实测试集或故障定义。它只回答数据稀缺时能否补充候选样本、恢复缺失观测或模拟退化轨迹；最终价值仍由未见设备上的真实故障召回、误报、校准和物理约束违反率决定。

本章先说明要学习怎样的联合分布或条件分布，再推导相应的训练目标。接下来检查生成样本是否遗漏某些故障类型、是否符合指定工况，以及采样误差有多大。合成样本可以用于训练，但收益要在保持独立的真实设备测试集上检验。生成的裂纹看起来逼真，也不代表现场出现过这种裂纹；真实故障很少时，模型尤其容易把已有猜测反复复制出来。因此每种方法都要说明哪些性质能够从目标函数得到，哪些仍需观测和验证。

阅读本章需要用到条件概率、期望和反向传播；遇到潜变量积分时，先把它理解为对所有可能的隐藏原因求和（连续情形为积分）。各模型共同要回答的是：没有现成输入样本时，从哪里开始随机抽样，又怎样把随机数变成具有数据结构的输出？

\section{生成建模的目标}

% auxiliary-resource-guide
本章末的“辅助学习材料”列出与核心方法对应的讲解和实践入口，可在阅读相关推导后，按所列小节对照学习。

生成模型学习数据分布$p_{\mathrm{data}}(x)$或其条件分布$p(x\mid c)$，目标是支持采样、补全或去噪；是否能直接计算似然，还取决于模型结构。判别模型关注$p(y\mid x)$或决策边界，生成模型还试图描述输入本身的结构。

\paragraph{分布不是一张平均样本。}
把一个含$D$个采样值的振动窗口写成$x\in\R^D$，随机变量$X$表示尚未取得的窗口，$x$表示它的一次具体取值。分布描述各种窗口出现的相对可能性；模型参数$\theta$则决定学习到的分布$p_\theta$。训练集只是有限观测，并不是已知的真实分布。$\E_{x\sim p_{\mathrm{data}}}$在训练中通常用一个批次的样本平均近似。对连续数据，$p_\theta(x)$是密度而非单点概率，某一区域的概率要对密度积分；似然是把已观测$x$固定后，考察不同参数对它赋予的密度。

例如只记录离散的冲击次数$X\in\{0,1,2\}$，三个值的概率为$(0.5,0.3,0.2)$。均值为$0\times0.5+1\times0.3+2\times0.2=0.7$，但每次都输出$0.7$既不合法，也没有再现分布。采样应让多次输出的频率逐渐接近这三个概率。条件$c$可以是转速或故障类别，$p(x\mid c)$描述固定条件下仍存在的变化；不是给定$c$就必须产生唯一的$x$。GAN 通常只提供采样机制而不直接提供密度，VAE 给出似然下界，Flow 才能在其结构假设下直接计算密度。

\paragraph{先把潜变量的生成方向写清楚。}
潜变量又称隐变量，是未直接观测、由模型引入的内部变量。令$z\in\R^d$，通常$d<D$；它可以组织波形的变化因素，但某一坐标不一定天然等于转速或损伤程度。生成时先从先验$p(z)$抽取$z$，再从$p_\theta(x\mid z)$抽取$x$。例如$z\sim\mathcal N(0,1)$、$x\mid z\sim\mathcal N((z,2z)^\top,\sigma_x^2I_2)$会使两个观测分量相关；潜变量是它们共享的随机来源。看见$x$后反推哪些$z$更可能产生它，才是后验推断。普通自编码器只学重构映射，没有自动规定从哪里抽取新编码；VAE 用先验和概率解码器补齐了这个生成机制。

\paragraph{变分自编码器。}
Variational Autoencoder (VAE)引入潜变量$z$，生成模型为$p_\theta(x,z)=p(z)p_\theta(x\mid z)$\cite{kingma2014vae}。后验$p_\theta(z\mid x)$通常难以计算，以$q_\phi(z\mid x)$近似。由 KL 分解可得
\begin{equation}
 \log p_\theta(x)=\mathcal{L}(x)+\KL(q_\phi(z\mid x)\|p_\theta(z\mid x)),
\end{equation}
其中证据下界（ELBO）为
\begin{equation}
 \mathcal{L}(x)=\E_{q_\phi(z\mid x)}[\log p_\theta(x\mid z)]-\KL(q_\phi(z\mid x)\|p(z)).\label{eq:generative-vae-elbo}
\end{equation}
式\eqref{eq:generative-vae-elbo}右端的第一项重构输入，第二项使潜变量分布接近先验。若$q_\phi(z\mid x)=\mathcal{N}(\mu_\phi(x),\diag(\sigma_\phi^2(x)))$，可用重参数化$z=\mu_\phi(x)+\sigma_\phi(x)\odot\epsilon$、$\epsilon\sim\mathcal{N}(0,I)$进行反向传播。

这里$\theta$是解码器参数，$\phi$是编码器参数。编码器输出$\mu_\phi(x),\sigma_\phi(x)\in\R^d$，$\sigma_j>0$；$I_d$为$d$阶单位矩阵，$\odot$表示对应元素相乘。训练时已知$x$，从$q_\phi(z\mid x)$抽样以计算重构项；生成新样本时没有输入$x$，改从$p(z)$抽样。贝叶斯公式给出$p_\theta(z\mid x)=p_\theta(x\mid z)p(z)/p_\theta(x)$，困难在分母$p_\theta(x)=\int p_\theta(x\mid z)p(z)\,dz$通常无法解析积分，故需要近似后验。

KL 散度的定义为$\KL(q\|p)=\E_q[\log q-\log p]\geq0$，衡量两分布不一致，但一般不对称，不是欧氏距离。ELBO 中的 KL 比较后验与先验，等式中的缺口则比较近似后验与真实后验，不要混为一项。最大化下界在代码中通常实现为最小化它的相反数，即负重构对数似然加 KL。

这个下界可以从 Jensen 不等式直接得到：
\begin{align}
 \log p_\theta(x)
 &=\log\int q_\phi(z\mid x)\frac{p_\theta(x,z)}{q_\phi(z\mid x)}\,dz\\
 &\geq\mathbb E_{q_\phi}[\log p_\theta(x,z)-\log q_\phi(z\mid x)].
\end{align}
将联合分布拆为$p_\theta(x\mid z)p(z)$，就得到重构项减去先验正则项。更完整地说，ELBO的缺口正是后验近似误差：
\begin{align}
\log p_\theta(x)-\mathcal L(x)
&=\KL\bigl(q_\phi(z\mid x)\|p_\theta(z\mid x)\bigr)\geq0.
\end{align}
这个等式要分清哪些参数固定：固定$\theta$时，$\log p_\theta(x)$不变，提高 ELBO 就等价于缩小近似后验与真实后验之间的 KL 缺口；联合更新$\theta,\phi$时，似然和缺口都在变化，下界上升不保证似然逐步上升、缺口逐步缩小。它也不保证重构误差与工业风险完全一致。若$x\in\mathbb R^D$采用独立高斯解码器$p_\theta(x\mid z)=\mathcal N(\mu_\theta(z),\sigma_x^2I)$，则
\[
-\log p_\theta(x\mid z)=\frac{1}{2\sigma_x^2}\|x-\mu_\theta(z)\|_2^2+\frac D2\log(2\pi\sigma_x^2),
\]
所以常见均方重构项隐含了同方差、条件独立噪声假设。若裂纹像素、高频冲击或少数故障模式更重要，应显式改变观测模型或评价权重，而不能仅增加潜变量维度。

编码器提供近似后验，解码器给出生成机制，两者通过同一个下界联合训练。对高斯后验和标准正态先验，KL 项可解析计算：
\begin{equation}
 \KL(q\|p)=\frac12\sum_{j=1}^{d}(\mu_j^2+\sigma_j^2-\log\sigma_j^2-1).
\end{equation}
可手算一个一维例子：$\mu=1,\sigma=1$时，KL 为$(1+1-0-1)/2=0.5$；若$\mu=0,\sigma=1$则为零，但这不代表重构一定准确。重参数化把随机性放入与参数无关的$\epsilon$，固定一次抽到的$\epsilon$后，$\partial z/\partial\mu=1$、$\partial z/\partial\sigma=\epsilon$，重构梯度便可沿$z$传回编码器。它并非取消随机性，而是让随机抽样后的计算路径可求导。

若解码器能绕过潜变量解释数据，或 KL 约束过强，近似后验可能接近先验而不再携带输入信息，这称为后验坍缩。KL annealing 指逐渐增大 KL 权重，free bits 则在一定阈值内不继续惩罚小 KL；它们改变训练折中，不保证潜变量具有可解释含义。仍需检查潜变量是否保留故障模式，不能只观察总损失。

\section{对抗学习与可逆密度建模}
Generative Adversarial Network (GAN) 的原始对抗目标与可逆密度模型的变量替换目标具有不同的训练信号\cite{goodfellow2014gan,dinh2017realnvp}。本节比较二者如何定义生成分布，而不把样本锐利程度等同于似然质量。
VAE 以可计算的下界换取潜变量推断的便利。若主要目标是采样，可以改用判别器提供训练信号；若需要显式密度，则可以限制变换结构以换取精确的变量替换计算。这两种选择分别引出 GAN 和 Normalizing Flow。

GAN 中$G:\R^d\to\R^D$把噪声向量变成样本，$D:\R^D\to(0,1)$输出样本来自真实训练数据的判别概率。这里$D$作为函数名表示判别器，作为维度时表示数据长度。训练一个批次时，先固定$G$，让$D$区分真实$x$和合成$G(z)$；再固定$D$的参数，让$G$产生更难被识别的样本。更新$G$仍需通过$D$对输入的导数传递梯度，“固定判别器”不是切断这条路径。采样时仅需$G$，不需要判别器。

GAN 通过生成器$G$和判别器$D$进行极小极大优化：
\begin{equation}
 \min_G\max_D\;\mathbb{E}_{x\sim p_{data}}\log D(x)+\mathbb{E}_{z\sim p(z)}\log(1-D(G(z))).
\end{equation}
在理想最优判别器下，目标与真实分布和生成分布之间的 Jensen--\allowbreak Shannon 散度相关，其英文名称为 Jensen--\allowbreak Shannon Divergence，简称 JSD。

固定生成分布$p_g$，并假设判别器可以逐点优化。对每个$x$，记
\[
 a=p_{\mathrm{data}}(x),\qquad b=p_g(x).
\]
需最大化的目标及其关于判别输出$D$的一阶导数分别为
\[
 v(D)=a\log D+b\log(1-D),\qquad
 v'(D)=\frac aD-\frac b{1-D}.
\]
当$a,b>0$时令导数为零，得到
\[
D^*(x)=\frac{p_{\mathrm{data}}(x)}{p_{\mathrm{data}}(x)+p_g(x)}.
\]
若$a+b>0$但其中一个为零，该式给出边界最优值$0$或$1$；对输出严格在$(0,1)$的判别器，只能以极限逼近。若$a=b=0$，该点不贡献目标，判别输出可任取。
为看清这里的散度，令混合分布$M=(P+Q)/2$，定义
\[
 \operatorname{JSD}(P\|Q)=\tfrac12\KL(P\|M)+\tfrac12\KL(Q\|M).
\]
它比较两分布各自偏离共同中点的程度；采用自然对数时取值在$0$与$\log2$之间，相同分布时为零。代回目标并利用 Jensen--Shannon 散度可得
\[
V(D^*,G)=-\log4+2\operatorname{JSD}(p_{\mathrm{data}}\|p_g),
\]
因此理想博弈的全局最优要求$p_g=p_{\mathrm{data}}$。这一步依赖判别器足够强、分布可比较和优化达到平衡；有限样本和有限判别器会使经验目标偏离理想散度；支撑集几乎不相交时，JSD 本身仍有定义，却可能接近饱和值而难以给出有用的生成器梯度。实际交替训练中，判别器过强可能使原始生成器目标的梯度饱和，模式崩溃则使不同噪声产生相似样本。对工业故障而言，少数典型样本逼真却遗漏罕见模式，仍不足以支持数据增强。

流模型使用可微可逆变换$z=f_\theta(x)$把数据映射到简单先验。这里$x,z\in\R^d$，$J_{f_\theta}(x)$的第$(i,j)$项为$\partial f_{\theta,i}(x)/\partial x_j$，且其行列式非零：
\begin{equation}
 \log p_X(x)=\log p_Z(f_\theta(x))+\log\bigl|\det J_{f_\theta}(x)\bigr|.
\end{equation}
推导来自微分体积变换：若$z=f_\theta(x)$，则$d z=|\det J_f(x)|d x$，而概率守恒给出$p_X(x)dx=p_Z(z)dz$，故$p_X(x)=p_Z(f(x))|\det J_f(x)|$。例如逐维仿射耦合层保持一部分坐标不变，并令另一部分$y_b=x_b\odot\exp s(x_a)+t(x_a)$，其雅可比为三角矩阵，因而$\log|\det J|=\sum_k s_k(x_a)$，同时逆变换可显式计算。可逆性和可计算的 Jacobian 行列式限制了网络设计，也使密度计算无需 VAE 的变分下界。采样时使用逆变换。精确似然只表示对模型分布的精确计算，并不保证该分布符合真实故障机制，更不保证高似然等同于低风险。

用一维变换可检查行列式因子的方向。设$z=x/2$且$z\sim\mathcal N(0,1)$，则$p_X(x)=p_Z(x/2)/2$，因此$x\sim\mathcal N(0,4)$。区间$x\in[0,2]$对应$z\in[0,1]$，两者概率相同，但前者区间宽一倍，所以密度需减半；漏掉$1/2$会使总概率变成$2$。采样方向则相反：先抽$z$，再令$x=2z$。这也说明可逆 Flow 通常保持维数，不能把有信息损失的普通压缩编码器直接当作 Flow。

表\ref{tab:generative-model-routes}比较的关键是目标与代价是否匹配。需要连续表示时检查 VAE 的压缩损失，需要快速采样时检查 GAN 的模式覆盖，需要密度时检查 Flow 的结构限制。扩散模型则把复杂生成映射分解为多步去噪，后文再推导其训练过程。

\begin{table}[htbp]
\centering
\caption{生成模型路线与工业应用边界}
\label{tab:generative-model-routes}
\scriptsize
\setlength{\tabcolsep}{3pt}
\resizebox{\linewidth}{!}{%
\begin{tabular}{p{2.5cm}p{3.7cm}p{3.7cm}p{4.0cm}}
\hline
模型路线 & 学习对象 & 主要优点 & 工业使用时的关键检查 \\
\hline
VAE & 潜变量分布与重构似然 & 潜空间连续、采样稳定 & 重构是否抹平细裂纹，潜变量是否保留工况 \\
GAN & 生成器与判别器的对抗平衡 & 样本锐利、条件生成直接 & 模式崩溃、训练不稳定、类别覆盖 \\
Flow & 可逆变换与精确似然 & 可计算密度、逆向映射清晰 & 可逆结构限制、显存和表达能力 \\
Denoising Diffusion Probabilistic Model (DDPM)/Denoising Diffusion Implicit Model (DDIM) & 逐步加噪与反向去噪 & 条件控制强、覆盖较好 & 采样步数、异常细节和条件一致性 \\
潜空间扩散 & 压缩空间中的扩散过程 & 高分辨率成本较低 & 编码器压缩损失和生成细节可信度 \\
\hline
\end{tabular}}
\end{table}

\section{扩散过程与条件控制}
\begin{figure}[htbp]
\centering
\begin{tikzpicture}[node distance=0.55cm and 0.45cm]
 \node[idi output, minimum width=1.35cm] (x0) {$x_0$};
 \node[idi process, right=of x0, minimum width=1.35cm] (x1) {$x_1$};
 \node[idi process, right=of x1, minimum width=1.35cm] (x2) {$x_2$};
 \node[idi process, right=of x2, minimum width=1.35cm] (dots) {$\cdots$};
 \node[idi input, right=of dots, minimum width=1.35cm] (xt) {$x_T$};
 \draw[idi arrow] (x0)--node[above,font=\scriptsize]{$+\epsilon$}(x1);
 \draw[idi arrow] (x1)--node[above,font=\scriptsize]{$+\epsilon$}(x2);
 \draw[idi arrow] (x2)--(dots); \draw[idi arrow] (dots)--(xt);
 \draw[idi feedback] (xt.south) .. controls +(0,-0.45) and +(0,-0.45) .. node[below,font=\scriptsize]{学习反向去噪 $\epsilon_\theta$} (x0.south);
\end{tikzpicture}
\caption{扩散模型的时间轴：前向过程逐步破坏结构，反向网络学习在每个噪声尺度恢复结构。}
\label{fig:generative-diffusion-time}
\end{figure}
图\ref{fig:generative-diffusion-time}的实线沿时间增加噪声，反馈箭头则概括学习到的反向恢复路径；训练需覆盖多个噪声尺度，而不是只拟合从末端噪声到样本的一次映射。
扩散模型通过逐步加噪和反向去噪构造生成过程\cite{ho2020ddpm}。前向过程逐步加入高斯噪声：
\begin{equation}
 q(x_t\mid x_{t-1})=\mathcal{N}(\sqrt{1-\beta_t}x_{t-1},\beta_tI).
\end{equation}
这里$t=1,\ldots,T$为扩散步，$0<\beta_t<1$为噪声日程，$x_0\sim p_{\mathrm{data}}$。记$\alpha_t=1-\beta_t$、$\bar\alpha_t=\prod_{s=1}^{t}\alpha_s$，则可直接采样
\begin{equation}
 x_t=\sqrt{\bar\alpha_t}x_0+\sqrt{1-\bar\alpha_t}\epsilon,
 \qquad \epsilon\sim\mathcal{N}(0,I).\label{eq:generative-direct-noising}
\end{equation}
所有$x_t$与$\epsilon$都在$\R^D$中；下标$t$是人工设置的加噪步，不是设备运行时间。两步递推可以检查直接采样公式的来源：$x_2=\sqrt{\alpha_2\alpha_1}x_0+\sqrt{\alpha_2\beta_1}\epsilon_1+\sqrt{\beta_2}\epsilon_2$。两个独立标准高斯噪声的加权和仍为高斯，其方差为$\alpha_2\beta_1+\beta_2=1-\alpha_1\alpha_2$；继续递推就得到$1-\bar\alpha_t$。因此训练时不必真的依次执行前面所有加噪步骤。

例如一维$x_0=2$、$\bar\alpha_t=0.64$，本次抽得$\epsilon=1$，则$x_t=0.8\times2+0.6\times1=2.2$。训练程序知道自己加入的噪声，所以不需要人工给噪声打标签；网络却只看$x_t,t$，不能偷看$x_0$。同一个$x_t$可能由多种$x_0,\epsilon$组合产生，故网络通常无法逐样本精确找回本次噪声，而是在平方损失下估计条件平均。

反向网络$\epsilon_\theta(x_t,t)$预测噪声，常用简化目标
\begin{equation}
 \mathcal{L}_{\mathrm{simple}}=\E_{t,x_0,\epsilon}\left[\|\epsilon-\epsilon_\theta(x_t,t)\|_2^2\right].
\end{equation}
此处独立采样$t\sim\operatorname{Uniform}\{1,\ldots,T\}$、$x_0\sim p_{\mathrm{data}}$和$\epsilon\sim\mathcal N(0,I)$，再由式\eqref{eq:generative-direct-noising}构造$x_t$。由条件期望的最小均方性质，若网络容量和优化均理想，则
\[
\epsilon_\theta^*(x_t,t)=\mathbb E[\epsilon\mid x_t,t].
\]
又因$x_t=\sqrt{\bar\alpha_t}x_0+\sqrt{1-\bar\alpha_t}\epsilon$，可由噪声预测得到干净样本估计
\[
\widehat x_0=\frac{x_t-\sqrt{1-\bar\alpha_t}\,\epsilon_\theta(x_t,t)}{\sqrt{\bar\alpha_t}}.
\]
当$\bar\alpha_t$很小时，这个反演会放大预测误差，说明末端高噪声阶段对模型校准和采样数值稳定性尤其敏感。

采样时从高斯噪声开始，按反向时间步逐渐去噪。直观上，模型学习的不是一次性画出样本，而是学习一系列局部的“如何去掉一点噪声”。这里还缺一个从预测到采样的步骤：一种常用 DDPM 参数化为
\[
 \mu_\theta(x_t,t)=\frac1{\sqrt{\alpha_t}}
 \left(x_t-\frac{\beta_t}{\sqrt{1-\bar\alpha_t}}\epsilon_\theta(x_t,t)\right),\qquad
 x_{t-1}=\mu_\theta(x_t,t)+\sqrt{\tilde\beta_t}\,\zeta_t,
\]
其中$t\geq2$，$\zeta_t\sim\mathcal N(0,I_D)$是本次反向步骤新抽取的噪声，不是训练时那一个$\epsilon$；$\tilde\beta_t$在下文推导。先抽$x_T\sim\mathcal N(0,I_D)$，再依次计算$t=T,\ldots,2$，最后通常以$t=1$的均值作为连续数据输出而不再加噪。起点近似成立需要日程使$\bar\alpha_T$足够小且数据尺度适当。只把一次$\widehat x_0$当成最终样本，不等于执行这条反向链；多步中的随机项用于表达仍未确定的可能性。

\paragraph{扩散模型的后验。}
由于前向过程是高斯马尔可夫链，给定$x_0,x_t$的后验仍为高斯：
\begin{equation}
 q(x_{t-1}\mid x_t,x_0)=\mathcal{N}(\tilde\mu_t(x_t,x_0),\tilde\beta_tI).
\end{equation}
对$t=2,\ldots,T$，其参数为
\begin{align}
 \tilde\mu_t(x_t,x_0)&=\frac{\sqrt{\bar\alpha_{t-1}}\beta_t}{1-\bar\alpha_t}x_0
 +\frac{\sqrt{\alpha_t}(1-\bar\alpha_{t-1})}{1-\bar\alpha_t}x_t,\label{eq:generative-posterior-mean}\\
 \tilde\beta_t&=\frac{1-\bar\alpha_{t-1}}{1-\bar\alpha_t}\beta_t.
\end{align}
其推导只需把两个高斯因子相乘。由
$q(x_t\mid x_{t-1})\propto\exp[-\|x_t-\sqrt{\alpha_t}x_{t-1}\|^2/(2\beta_t)]$
以及
$q(x_{t-1}\mid x_0)=\mathcal N(\sqrt{\bar\alpha_{t-1}}x_0,(1-\bar\alpha_{t-1})I)$，收集关于$x_{t-1}$的二次项，精度为
$\alpha_t/\beta_t+(1-\bar\alpha_{t-1})^{-1}=(1-\bar\alpha_t)/[\beta_t(1-\bar\alpha_{t-1})]$，所以方差是其倒数$\tilde\beta_t$；再收集一次项并乘以该方差，就得到式\eqref{eq:generative-posterior-mean}的均值。这说明后验同时利用当前噪声状态$x_t$和原始状态$x_0$。采样时$x_0$未知，网络必须先由$x_t$估计它或等价地估计噪声，不能把该后验公式误当作可直接使用的生成器。
约定$\bar\alpha_0=1$，则$t=1$时方差为零，条件后验退化到已知$x_0$，不作为非退化高斯参与上述 KL 计算。
对$t\geq2$，扩散模型的变分目标包含上述后验与反向模型$p_\theta(x_{t-1}\mid x_t)$之间的 KL 项。固定反向方差并采用噪声预测参数化后，这些项可写为带时间步权重的噪声预测误差；常用简化目标改变了这些权重，不能与完整变分目标直接等同。噪声预测、$x_0$预测和速度预测是不同参数化，改变梯度尺度与采样表现，但共享逐步去噪框架。

\paragraph{条件扩散与引导。}
若有条件$c$，网络写为$\epsilon_\theta(x_t,t,c)$。分类器引导使用带噪样本分类器的梯度调整生成方向；无分类器引导通过训练时随机丢弃条件，使网络同时学习条件与无条件预测，再组合为
\begin{equation}
 \hat\epsilon=\epsilon_\theta(x_t,t,\varnothing)
 +w[\epsilon_\theta(x_t,t,c)-\epsilon_\theta(x_t,t,\varnothing)].
\end{equation}
$w=0$对应无条件预测，$w=1$对应原始条件预测，$w>1$进一步加强条件方向。更强引导可能损害多样性或产生伪影，应在固定采样预算下比较条件满足率和缺陷覆盖。ControlNet 等条件分支可接收边缘、深度或几何信息，但条件输入并不自动保证生成结果满足物理约束。

\paragraph{得分匹配（Score Matching）。}
以下$p_t(x_t)=\int q(x_t\mid x_0)p_{\mathrm{data}}(x_0)\,dx_0$表示前向带噪边缘密度，不是反向模型的转移密度。
扩散模型也可从各噪声时刻的 score $\nabla_{x_t}\log p_t(x_t)$理解。去噪 score matching 估计带噪分布的对数密度梯度，而不直接计算归一化常数。高斯加噪下，噪声预测对应的 score 估计为$-\epsilon_\theta(x_t,t)/\sqrt{1-\bar\alpha_t}$。这一估计需与噪声日程共同构成反向随机过程，或对应的概率流常微分方程，才能用于分布采样；单纯沿密度梯度上升会趋向局部高密度区域，并不保证恢复目标分布的模式比例。这个观点连接了扩散模型、能量模型和随机微分方程。

\section{采样、约束与工业可信度}
DDIM 在与 DDPM 相关的训练目标下改变采样轨迹，潜空间扩散通过压缩表示降低高分辨率生成成本\cite{song2021ddim,rombach2022latent}。无分类器引导与 ControlNet 分别提供条件强度组合和结构条件分支\cite{ho2022cfg,zhang2023controlnet}，两者均不能自动保证物理可行性。
有了可采样模型，还不能直接把生成样本当成维护数据。只用正常数据训练的模型可以模仿正常模式，却没有凭据知道故障长什么样。若要按故障标签或严重度生成，应有真实故障样本、经过验证的仿真机制，或能够核查的领域约束。不同用途也要分别验收：补足少数类别要检查模式是否被覆盖，模拟退化要检查时间顺序和工况，指定位置的缺陷图像要检查位置与标注是否一致。

\paragraph{DDIM 与采样加速。}
DDPM 使用随机反向过程，需要较多时间步。DDIM 可以构造非马尔可夫的确定性或弱随机采样轨迹，在保留训练目标的同时减少采样步数。步数减少会降低延迟，却可能损害细节和条件一致性，因此应绘制质量--延迟曲线，而不是只报告单一采样配置。

\paragraph{潜空间扩散。}
直接在高维像素或原始波形空间扩散成本很高。潜空间扩散先用编码器把输入压缩到$z$，在潜空间去噪，再由解码器重构。压缩比例过大可能删除细微裂纹、冲击脉冲或高频异常，因此工业应用必须比较原空间频谱、边缘和峰值统计，而不能只看重构图像的主观质量。

\begin{figure}[htbp]
\centering
\begin{tikzpicture}[>=Latex, node distance=1.15cm, every node/.style={font=\small}]
 \node[draw, rounded corners, fill=blue!10, minimum width=2.3cm] (x) {$x_0$：真实样本};
 \node[draw, rounded corners, fill=orange!12, right=of x, minimum width=2.5cm] (noise) {\shortstack{前向加噪\\$q(x_t|x_0)$}};
 \node[draw, rounded corners, fill=gray!15, right=of noise, minimum width=2.3cm] (xt) {$x_T$：近似噪声};
 \draw[->, thick] (x)--(noise); \draw[->, thick] (noise)--(xt);
 \draw[->, thick, dashed] (xt.north) .. controls +(0,0.65) and +(0,0.65) .. node[above,font=\scriptsize]{学习反向去噪} (x.north);
 \node[draw, rounded corners, fill=green!12, below=1.1cm of noise, minimum width=3cm] (net) {$\epsilon_\theta(x_t,t,c)$};
 \draw[->] (noise.south)--(net.north);
\end{tikzpicture}
\caption{扩散模型把生成拆为前向加噪和学习反向去噪。}
\label{fig:generative-denoising-network}
\end{figure}
图\ref{fig:generative-denoising-network}把预设加噪机制与可训练的噪声预测网络分开；网络读取带噪状态、时间和条件，采样时再由这些预测逐步恢复样本。

\paragraph{物理约束生成。}
工业生成样本必须满足守恒、边界、单调性、频率范围或设备状态转移约束。可以在损失中加入残差惩罚
\begin{equation}
 \mathcal{L}=\mathcal{L}_{gen}+\lambda_{phys}\mathbb{E}[\|r(x,c)\|_2^2],
\end{equation}
也可以在采样后投影到可行域。惩罚过弱会产生不可能样本，过强则可能牺牲数据多样性；因此应分别报告真实性、多样性和约束违反率。

以非负总量约束为例，$x=(x_1,x_2)$应满足$x_1,x_2\geq0$及$x_1+x_2=1$。可取$r(x)=x_1+x_2-1$，候选$(0.8,0.5)$的平方残差为$0.09$；等距投影到该直线得到$(0.65,0.35)$，恰好也满足非负约束。一般候选应投影到非负约束与该直线的交集，而不能只校正总量。但投影改变了样本分布，多个候选还可能被压到同一边界，不能把约束满足率上升等同于分布恢复正确。物理残差具体作用于干净样本还是带噪状态也必须写明，带噪中间量通常不必满足干净数据的守恒关系。

\paragraph{记忆、隐私与生成数据。}
生成模型可能记住训练集中的罕见记录，尤其在数据量小、重复度高或模型过拟合时。近邻距离、成员推断和重复片段检测可用于筛查记忆风险。差分隐私训练可以降低单个样本的影响，但通常带来质量--隐私权衡。生成数据用于发布或跨机构共享前，应进行去标识、相似样本审计和下游泄漏测试。

\paragraph{生成模型的实验流程。}
完整实验应包含训练数据分布、采样配置、随机种子、生成数量、重复检测、专家评价、统计检验和下游任务。对故障生成尤其应使用未参与生成训练的真实故障测试集；否则增强后的模型可能只是在适应合成数据的风格。最有说服力的结论是：在不同随机种子、不同设备和不同故障类型上，真实测试性能都获得稳定改善。

评价应先检查边缘分布、相关性、频谱与缺陷几何，再检查物理约束和条件一致性，最后比较真实数据上的下游结果。Fréchet Inception Distance (FID) 或主观逼真度只覆盖部分属性。对照应至少包括仅真实数据与真实加合成数据，并匹配训练预算；重采样基线可以区分数量效应，打乱合成标签则可作为诊断对照，不能替代有效基线。

新增样本可能复制设备身份和批次标记，也可能改变故障先验比例。因此测试集须保持真实、独立且未参与生成器训练与选择，报告召回、误报和校准，而不只看训练损失。若训练指标上升而真实测试性能下降，应检查覆盖不足、生成伪影和先验变化，不能继续用更多合成数据掩盖问题。

VAE、GAN、Flow 与扩散模型的联系和差异可用下列教程按模型族复习，阅读时特别对照训练目标与采样过程。
\section*{辅助学习材料}
\begingroup
\emergencystretch=3em
\begin{itemize}
\item \textbf{What are Diffusion Models?}；作者/平台：Lilian Weng；英文；技术教程。阅读前向扩散、反向去噪、训练目标、classifier-free guidance 和 latent diffusion，核对本章的噪声预测与条件采样。\href{https://lilianweng.github.io/posts/2021-07-11-diffusion-models/}{在线阅读}。
\item \textbf{From Autoencoder to Beta-VAE}；作者/平台：Lilian Weng；英文；技术教程。阅读变分自编码器、ELBO、重参数化和 Beta-VAE 部分，理解重建误差与潜变量先验之间的权衡。\href{https://lilianweng.github.io/posts/2018-08-12-vae/}{在线阅读}。
\item \textbf{超详细的扩散模型（Diffusion Models）原理+代码}；知乎；中文解说。对照前向加噪与反向去噪，核对噪声预测目标；标题与主题据必应索引，原页受限。\href{https://zhuanlan.zhihu.com/p/624221952}{在线阅读}。
\item \textbf{扩散模型详解：从DDPM到Stable Diffusion再到DiT的技术演进}；CSDN；中文博客。比较采样过程和模型结构，生成质量不能代替真实数据上的下游评价；标题与主题据必应索引，原页受限。\href{https://blog.csdn.net/weixin\_67327688/article/details/155781843}{在线阅读}。
\item \textbf{第六节 2021 - 生成式对抗网络(GAN) (一) – 基本概念介紹}；李宏毅；bilibili；中文课程。选看第58集的基本概念讲解，对照本章对抗目标；该分集不是扩散模型教程。2026年10月4日官方公开元数据显示整套课程累计播放约299.5万，并非本分集播放量；已核对目录与计数，未逐帧观看。\href{https://www.bilibili.com/video/BV1Wv411h7kN/?p=58}{观看分集}。
\end{itemize}
\endgroup
\paragraph{本章小结。}
VAE 借助潜空间生成，GAN 通过对抗判别训练，flow 使用可逆变换计算似然，diffusion 逐步去噪得到样本。它们的计算方式不同，生成结果都仍需检查。用于工业任务时，既看样本是否逼真、是否覆盖需要的故障，也检查是否违背物理规律、是否泄露训练记录，最后用独立的真实数据检验下游收益。

\chapter{物理信息神经网络：从稀疏观测重建风场}
\label{chap:pinn}
\begin{quote}\small
\textit{本章摘要。}物理信息神经网络（Physics-Informed Neural Network，PINN）将观测误差、微分方程残差和初边值条件共同用于训练。本章以风电场稀疏多普勒激光雷达测量为例，推导从视线速度到连续风场的逆问题，解释自动微分、尺度处理、约束方式和主要变体，并讨论可辨识性、湍流闭合与独立验证。
\end{quote}


风电场中的尾流使下游风机入流速度下降，并改变剪切和载荷。扫描激光雷达能够在远距离取得风速信息，但观测沿有限波束和距离门分布，扫描过程还需要时间。目标因而不是把一个完整风速图输入网络，而是根据不完整、不同步且带有空间平均效应的观测，估计指定时空区域中的速度场。本章给出教学建模设置，不对应已完成的风电场实验，也不预设 PINN 优于传统反演或数值同化。

雷达只沿有限方向测量，却希望得到一片区域的风速，这就是本章逆问题的起点。先用观测算子写清仪器到底测到了什么，再选定重建范围和尺度，把方程、边界与初值写入训练损失。未测到的位置主要依靠这些条件推断，因此方程残差很小，并不保证风场重建准确，还要用独立观测检验。能算出结果、能从数据确定唯一结果、能验证结果是否准确，是三个不同问题。下一章的预训练可以改善初始化或共享表示，但不能补上当前场地缺失的观测与边界约束。


\section{从函数拟合到受约束逆问题}
\label{sec:pinn-formulation}

% auxiliary-resource-guide
本章末的“辅助学习材料”列出与核心方法对应的讲解和实践入口，可在阅读相关推导后，按所列小节对照学习。

不必先掌握完整的流体数值求解器才能理解训练流程，但需要区分函数值、对坐标的导数和对参数的梯度。后文先用视线投影说明观测不足，再建立风场方程与尺度约定，并在训练部分用一维微分方程手算一次更新；流体方程在这里提供约束，不充当额外的测量标签。

\paragraph{先读懂坐标、导数与残差。}
这里的网络不是把一幅风场图分类，而是把一个查询坐标映射到该处状态。三维问题中输入$(x,y,z,t)\in\R^4$，输出$(u,v,w,\pi)\in\R^4$；对$B$个查询点可分别组成$B\times4$的输入、输出数组。$x,y,z$是位置，$u,v,w$才是速度，不能因符号都只有一个字母就混用。参数$\theta\in\R^P$是网络内待学习的权重，不是物理空间坐标。

偏导数$\partial_xu$表示固定其他坐标和时间后，沿$x$方向的速度变化率；$\partial_tu$表示固定位置的时间变化率。二阶导数$\partial_{xx}u$描述变化率自身如何变化。把一个候选函数代入方程，将左右两端移到同一侧，得到的差称为残差。残差为零表示候选满足这条方程，并不直接表示它与真实解的差为零。例如$\partial_xq=0$的任意常数函数都满足方程，要确定是哪一个常数，还需要一个函数值条件。

普通监督回归通过成对样本约束函数值。PINN 进一步用已知方程约束函数在未观测位置的变化方式。设空间坐标为$\mathbf{x}\in\Omega$、时间为$t\in[0,T]$，待求状态为$\mathbf{q}(\mathbf{x},t)$，控制方程、边界条件和初始条件分别写为
\begin{align}
 \mathcal N_{\boldsymbol\eta}[\mathbf q](\mathbf x,t)&=\mathbf0,\qquad (\mathbf x,t)\in\Omega\times(0,T],\\
 \mathcal B[\mathbf q](\mathbf x,t)&=\mathbf b(\mathbf x,t),\qquad \mathbf x\in\partial\Omega,\\
 \mathbf q(\mathbf x,0)&=\mathbf q_0(\mathbf x).
\end{align}
其中$\boldsymbol\eta$表示方程中的物理参数，$\mathcal B$可以表示函数值、法向导数或通量。神经网络$\mathbf q_\theta$提供一个可微的连续函数近似，训练使其同时符合观测与方程\cite{raissi2019pinn,karniadakis2021physics}。

符号$\Omega$是重建区域，$\partial\Omega$是它的边界，$[0,T]$是观察时间窗。$\mathcal N$是把函数及其导数组合起来的规则，称为微分算子，不是额外神经网络；$\mathcal B$提取边界上的值或导数。初值规定$t=0$时整个区域的状态，边界值规定区域边缘随时间的状态。配点上的“标签为零”指方程残差为零，不是该处风速为零。

正问题在参数和初边值条件给定时求状态。逆问题则由部分观测反推状态、参数或边界输入。稀疏风场重建属于后者：即使只优化网络参数，完整风场也不是直接观测标签。若同时估计入流或有效黏度，还需要检查这些未知量是否能够被现有测量分别识别。增加待估参数会扩大解空间，不会自动提高重建可信度。

配点是用于计算方程残差的时空位置，不需要对应真实风速标签。配点越多通常越能覆盖残差，但它们没有增加独立观测信息。若目标是对区域上的残差平方积分，连续量为
\[
\mathcal J_f=\frac1{|D|}\int_D\|r(\boldsymbol\xi)\|_2^2\,d\boldsymbol\xi,
\qquad D=\Omega\times(0,T),
\]
从$\boldsymbol\xi_j\sim\rho$独立抽样并用 Monte Carlo 近似时，
\[
\widehat{\mathcal J}_f=\frac1{N_f}\sum_{j=1}^{N_f}\frac{\|r(\boldsymbol\xi_j)\|_2^2}{|D|\rho(\boldsymbol\xi_j)}
\]
在$\rho$覆盖$D$且积分有限时无偏。均匀采样时权重为常数；残差自适应采样若不除以$\rho$，就等价于改变了目标积分的空间权重。有限样本下误差还取决于残差方差，局部加点可以降低某些区域的方差，却不能创造新的观测约束。即便某个网络在全部配点上满足方程，缺少足够的初边值约束时仍可能存在多个物理解。PINN 将连续问题转化为有限样本上的优化问题，离散残差较小也不等于已证明连续解的存在、唯一性或误差上界。

本章采用坐标网络直接拟合一个时间窗内的风场。它与输入一组工况后直接输出解的神经算子不同：更换场地或边界条件后，普通 PINN 往往需要重新优化。若采用跨工况预训练与适应，则应另行报告预训练数据、适应观测及额外成本。

把逆问题写成算子组合有助于定位信息来源。记$\mathcal H$为雷达测量算子、$\mathcal N$为方程算子、$\mathcal C$为初边值约束，则训练实际上寻找
\[
\mathbf u\in\{\mathbf u:\ \mathcal H\mathbf u\approx\mathbf y,\;\mathcal N[\mathbf u]=0,\;\mathcal C[\mathbf u]=0\}.
\]
观测项缩小的是与仪器一致的解集合，物理项缩小的是满足所选方程的解集合；二者交集为空时，继续增大物理权重只会牺牲观测拟合。这个集合观点解释了为什么配点数量、网络宽度和损失权重都不能单独解决边界错误或观测几何退化。

\section{激光雷达实际观测了什么}
\label{sec:pinn-lidar}
\begin{figure}[htbp]
\centering
\begin{tikzpicture}[node distance=0.75cm and 0.7cm]
 \node[idi input, minimum width=2.0cm] (beam) {激光雷达\\视线观测 $y_m$};
 \node[idi process, right=of beam, minimum width=2.25cm] (operator) {观测算子\\投影与积分 $\mathcal H_m$};
 \node[idi state, right=of operator, minimum width=2.25cm] (field) {连续风场\\$\mathbf u_\theta(\mathbf x,t)$};
 \node[idi process, right=of field, minimum width=2.2cm] (residual) {自动微分\\方程残差};
 \node[idi output, below=0.8cm of field, minimum width=2.4cm] (loss) {数据、物理与边界损失};
 \draw[idi arrow] (beam)--(operator); \draw[idi arrow] (operator)--(field);
 \draw[idi arrow] (field)--(residual); \draw[idi arrow] (residual)--(loss);
 \draw[idi feedback] (loss) to[bend left=28] (field);
\end{tikzpicture}
\caption{PINN 的信息路径：观测通过测量算子约束场，微分方程和边界条件在未观测位置提供额外约束。}
\label{fig:pinn-constraint-path}
\end{figure}
图\ref{fig:pinn-constraint-path}区分测量算子与微分残差：前者把场映射到可观测量，后者在配点上检验物理约束，二者不能用同一种点值误差替代。
设三维速度为$\mathbf u=(u,v,w)^\top$。在以东、北、上为轴的局部坐标系中，若方位角$\alpha$从东向北计、仰角为$\beta$，则视线单位向量为
\begin{equation}
 \mathbf e(\alpha,\beta)=
 (\cos\beta\cos\alpha,\cos\beta\sin\alpha,\sin\beta)^\top.
\end{equation}
仪器原始方位角若从北起算，必须先转换。径向速度正负号也应与仪器的朝向或远离约定一致。理想点测量为$\mathbf e^\top\mathbf u$，而不是速度模长$\|\mathbf u\|_2$。

例如真实速度为$(3,4,0)^\top\,\mathrm{m/s}$，沿东向的波束只读到$3\,\mathrm{m/s}$，沿北向读到$4\,\mathrm{m/s}$，两者都不是模长$5\,\mathrm{m/s}$。仅有东向读数时，$(3,0,0)$与$(3,4,0)$都符合观测。若两条波束夹角为$\delta$且已知$w=0$，则$y_1=u$、$y_2=u\cos\delta+v\sin\delta$，故$v=(y_2-y_1\cos\delta)/\sin\delta$；$\delta$很小时，分母很小，测量噪声被放大。这就是后文秩和条件数在几何上的含义。

实际距离门对有限探测体积内的速度加权，扫描又可能包含时间平均。以固定位置雷达为例，第$m$次观测可以表示为
\begin{align}
 y_m&=\mathcal H_m[\mathbf u]+\epsilon_m,\label{eq:pinn-observation}\\
 \mathcal H_m[\mathbf u]
 &=\int\!\int K_m(s,\tau)\,
 \mathbf e_m(\tau)^\top
 \mathbf u(\mathbf x_L+s\mathbf e_m(\tau),\tau)\,ds\,d\tau,\\
 &\hspace{1cm}K_m\geq0,\qquad \int\!\int K_m(s,\tau)\,ds\,d\tau=1.
\end{align}
这里$\mathbf x_L$为雷达位置，$K_m$描述该次观测的距离与时间响应，$\epsilon_m$包含测量噪声。只有当探测体积和时间窗相对于流场变化足够小时，才可近似为门中心、时间戳处的点投影。激光雷达测风中的体积平均与湍流测量限制可参见相关综述\cite{sathe2013lidar}。运动平台还需修正位置、姿态和平台速度，本章先限定为固定平台。

在网络训练中，可用数值求积实现这一观测算子：
\begin{equation}
 \widehat y_m=\sum_{a=1}^{A_m}\omega_{ma}\,
 \mathbf e_m(t_{ma})^\top\mathbf u_\theta(\mathbf x_{ma},t_{ma}),
 \qquad \sum_{a=1}^{A_m}\omega_{ma}=1.
\end{equation}
因此数据损失应比较$\widehat y_m$和$y_m$，不能将径向速度直接当作$u$分量标签。求积节点和权重来自测量模型，不是任意增加的训练标签。忽略体积平均时，模型可能以过于平滑的局部风场解释距离门观测，也可能把未分辨的剪切错误归因于噪声。

单条波束在某个位置只提供一个标量约束。若扰动$\delta\mathbf u$满足$\mathbf e^\top\delta\mathbf u=0$，该观测无法区分扰动前后的速度。多台雷达在近似相同位置和时刻交叉观测时，可以组成
\begin{equation}
 \mathbf y=E\mathbf u+\boldsymbol\epsilon,\qquad
 E=\begin{bmatrix}\mathbf e_1^\top\\ \vdots\\ \mathbf e_J^\top\end{bmatrix}.
\end{equation}
这里$E\in\R^{J\times3}$、$\mathbf y,\boldsymbol\epsilon\in\R^J$，每一行只增加一个视线投影约束。直接恢复三维速度需要$E$满列秩且条件数可接受。更一般地，在局部线性化下，观测对未知参数或场系数$\vartheta\in\mathbb R^P$的敏感度矩阵为
\[
G=\frac{\partial(\mathcal H[\mathbf u_\vartheta])}{\partial\vartheta}.
\]
若共使用$N_d$个标量观测，则$G\in\R^{N_d\times P}$；此处$R\in\R^{N_d\times N_d}$为正定观测噪声协方差，独立等方差噪声下$R=\sigma^2I$。矩阵$G^{\mathsf T}R^{-1}G\in\R^{P\times P}$衡量按噪声大小加权后各参数方向的可观测程度。若存在非零方向$\delta\vartheta$使$G\delta\vartheta=0$，则一阶观测对该方向没有辨识能力；若$G^{\mathsf T}R^{-1}G$有很小特征值，噪声会被放大为很大的参数不确定性。物理残差和边界会增加额外行，但只有当它们对该方向提供独立敏感度时才改善秩或条件数。因而“物理项降低”不能单独证明参数可辨识。两条独立视线配合已知垂直速度等额外假设，才可能恢复两个水平分量。近乎平行的波束会放大噪声，非同时扫描也不能无条件合并为同一瞬时速度。PINN 借助方程和边界在时空中传播信息，但不能消除由测量几何和未知边界共同造成的不可辨识性。

\section{风场方程、尺度与物理假设}
\label{sec:pinn-physics}
先考虑恒密度、不可压缩流动。将压力除以流体质量密度$\rho$，记运动学压力为$\pi=p/\rho$；这里$\rho$的单位是质量除以体积，不是前文的配点采样概率密度。动量与连续性方程为
\begin{align}
 \nabla\cdot\mathbf u&=0,\\
 \partial_t\mathbf u+(\mathbf u\cdot\nabla)\mathbf u
 &=-\nabla\pi+\nu\nabla^2\mathbf u+\mathbf f.
\end{align}
这里$\nabla=(\partial_x,\partial_y,\partial_z)^\top$，散度$\nabla\cdot\mathbf u=\partial_xu+\partial_yv+\partial_zw$是一个标量，表示局部净流出程度；不可压缩并不是风速不随位置变化，而是这三项之和为零。梯度$\nabla\pi$为三维向量，指示压力变化方向。拉普拉斯$\nabla^2u=\partial_{xx}u+\partial_{yy}u+\partial_{zz}u$作用到速度每个分量上，产生三维向量。对流项$(\mathbf u\cdot\nabla)\mathbf u$表示随流体移动时因空间位置变化而看到的速度变化；即便$\partial_t\mathbf u=0$，流过弯曲或加速区域的流体仍可能有加速度。

其中$\nu$为运动黏度，$\mathbf f$为单位质量外力。重力的静水平衡部分可吸收入压力，科里奥利力、浮力或风机作用则应按目标尺度决定是否保留。忽略这些项属于模型假设，需要与场地、稳定度和观测时间尺度匹配。

对于风机尾流，不能在穿过转子的区域直接假设$\mathbf f=\mathbf0$却又期望方程解释转子动量抽取。可以用已知运行状态约束的执行器盘体力近似，也可以将重建区域限定在转子外，并提供经过验证的入口或界面条件。如果把任意体力场也交给高容量网络自由拟合，它可能吸收所有动量残差，使物理约束失去辨识作用。

实际稀疏扫描通常无法解析全部湍流尺度。如果目标是时间平均速度$\overline{\mathbf u}$，雷诺平均方程包含额外应力：
\begin{equation}
 \partial_t\overline{\mathbf u}
 +(\overline{\mathbf u}\cdot\nabla)\overline{\mathbf u}
 =-\nabla\overline\pi+\nu\nabla^2\overline{\mathbf u}
 -\nabla\cdot\overline{\mathbf u'\mathbf u'}+\overline{\mathbf f}.
\end{equation}
将瞬时速度分解为$\mathbf u=\overline{\mathbf u}+\mathbf u'$，其中$\mathbf u'$是偏离平均值的三维脉动。式中的$\overline{\mathbf u'\mathbf u'}$指外积的平均，即$3\times3$矩阵，其第$(i,j)$项是$\overline{u_i'u_j'}$，不是向量内积。即使$\overline{\mathbf u'}=0$，平方的平均仍可非零：脉动等概率取$+1,-1$时，平均为$0$而均方为$1$。所以平均速度不能确定全部脉动相关项，平均后的方程出现了新的未知量。这里的湍流闭合指用额外的模型假设表示平均流动方程中未确定的湍流应力，使方程可求解。采用涡黏度闭合时，常将黏性与湍流应力合并为$\nabla\cdot[2(\nu+\nu_t)S]$，其中$S=(\nabla\overline{\mathbf u}+\nabla\overline{\mathbf u}^{\top})/2$，各向同性湍动能项吸收入修正压力。当$\nu_t$随位置变化时，这一散度项不能简化为$(\nu+\nu_t)\nabla^2\overline{\mathbf u}$。估计$\nu_t$需要闭合假设、取值约束和额外证据，否则黏度、入流与压力可能互相补偿。

因此应先声明重建对象是瞬时场、滤波场还是平均场，再匹配方程和观测平均尺度。本章后续以已选定的恒定有效黏度模型说明算法，并将其作为需要检验的近似。真实复杂地形、热分层和非平稳尾流可能需要扩展控制方程，不能用较大的物理损失权重弥补遗漏机制。

物理损失之前先进行无量纲化。选取长度$L_0$、速度$U_0$、参考位置$\mathbf x_0$和时间$t_0$，定义
\begin{align}
 \mathbf x^*&=(\mathbf x-\mathbf x_0)/L_0,&
 t^*&=(t-t_0)U_0/L_0,\\
 \mathbf u^*&=\mathbf u/U_0,&
 \pi^*&=\pi/U_0^2.
\end{align}
压力还可减去参考值，因为不可压缩动量方程只包含其空间梯度。定义$\mathrm{Re}=U_0L_0/\nu_{\mathrm{eff}}$和$\mathbf f^*=L_0\mathbf f/U_0^2$后，残差为
\begin{align}
 r_c&=\nabla^*\cdot\mathbf u_\theta^*,\\
 \mathbf r_m&=\partial_{t^*}\mathbf u_\theta^*
 +(\mathbf u_\theta^*\cdot\nabla^*)\mathbf u_\theta^*
 +\nabla^*\pi_\theta^*
 -\mathrm{Re}^{-1}\nabla^{*2}\mathbf u_\theta^*-\mathbf f^*.
 \label{eq:pinn-momentum}
\end{align}
这使不同损失项不再因米、秒与压力的单位而任意改变相对量级。以长度尺度$L_0$和速度尺度$U_0$为例，连续性残差的自然量级是$U_0/L_0$，动量残差的自然量级是$U_0^2/L_0$；若直接平方原始残差，动量项相对连续性项会多出$(U_0)^2$的单位和尺度因素。因此应先除以各自参考尺度，或通过$\lambda_c,\lambda_m$显式定义无量纲权重。对独立观测噪声，数据残差应先除以相应噪声标准差再平方，即平方误差按方差的倒数加权；边界项则按边界测量误差和边界采样测度归一化。权重不是越大越“物理”，而是定义了不同证据之间的相对容忍度。

如果输入还被映射到$[-1,1]$，计算物理导数时必须包含额外缩放的链式法则。例如$\widetilde x=2(x-x_{\min})/(x_{\max}-x_{\min})-1$时，$\partial_x=2\partial_{\widetilde x}/(x_{\max}-x_{\min})$。自动微分只对实际计算图求导，无法修正被漏写的尺度变换。

\section{网络、损失与训练过程}
\label{sec:pinn-training}
网络以无量纲时空坐标为输入，输出三个速度分量和压力：
\begin{equation}
 (u_\theta^*,v_\theta^*,w_\theta^*,\pi_\theta^*)
 =F_\theta(x^*,y^*,z^*,t^*).
\end{equation}
自动微分对输出关于输入求一阶和二阶导数，再将它们代入式\eqref{eq:pinn-momentum}。例如第一分量的对流项为$u^*\partial_{x^*}u^*+v^*\partial_{y^*}u^*+w^*\partial_{z^*}u^*$。这与仅对网络参数反向传播不同：优化参数时还需要保留输入导数的计算图，二阶空间导数会显著增加内存与计算开销。

强形式黏性方程需要足够光滑的网络。常用$\tanh$或平滑正弦激活；分段线性的 ReLU 二阶导数几乎处处为零，不适合直接表达这里的经典二阶黏性残差。Fourier 特征可增强高频表示，但过高频率会放大导数并加重优化困难。若输入角度来自雷达元数据，其三角函数只用于观测算子，不应与风场坐标混淆。

若第$m$个径向观测的标准差为$\sigma_m$，可将数据项写成
\begin{equation}
 \mathcal L_d=\frac1{N_d}\sum_{m=1}^{N_d}
 \left(\frac{\mathcal H_m[\mathbf u_\theta]-y_m}{\sigma_m}\right)^2.\label{eq:pinn-weighted-data-loss}
\end{equation}
式中$y_m,\sigma_m$使用同一速度单位，网络虽输出$\mathbf u_\theta^*$，代入观测算子前需还原$\mathbf u_\theta=U_0\mathbf u_\theta^*$并对应到实际坐标。等价地，可同时将观测与标准差除以$U_0$，全程使用无量纲速度；不能只缩放预测而保留原始标签。例如$U_0=10\,\mathrm{m/s}$时，预测$u^*=0.3$应与$3\,\mathrm{m/s}$比较，而非与原单位标签直接相减。式\eqref{eq:pinn-weighted-data-loss}对应独立高斯噪声下的加权最小二乘。距离门或扫描之间具有相关噪声时，应考虑残差协方差及有效样本量，不能把重叠门当作独立证据重复加权。质量控制可筛除受降水、遮挡或低信噪比影响的观测，阈值应事先固定。离群值较多时可比较 Huber 损失，但稳健损失不能代替错误时间戳或仪器偏差的校准。

在$N_f$个内部配点上计算连续性和动量项，并加入边界、初值与压力规范约束：
\begin{align}
 \mathcal L_f&=\frac1{N_f}\sum_{j=1}^{N_f}
 \left[\lambda_c|r_c(\boldsymbol\xi_j)|^2
 +\lambda_m\|\mathbf r_m(\boldsymbol\xi_j)\|_2^2\right],\\
 \mathcal L&=\lambda_d\mathcal L_d+\mathcal L_f
 +\lambda_b\mathcal L_b+\lambda_i\mathcal L_i
 +\lambda_g\mathcal L_g+\lambda_\eta\mathcal R(\boldsymbol\eta).
 \label{eq:pinn-loss}
\end{align}
其中$\boldsymbol\xi_j$为无量纲时空坐标，$\mathcal L_b$和$\mathcal L_i$分别是边界、初值误差。压力可加上任意仅随时间变化的函数而不改变动量方程，因此$\mathcal L_g$可在每个采样时刻固定一个参考点的压力，或固定空间平均压力。只固定一个时空点不足以消除整个非稳态问题的压力自由度。$\mathcal R$用于约束待估物理参数；已知参数时无需这一项。

\paragraph{一个可手算的微分训练例子。}
先用无量纲的一维问题$q'(x)+q(x)=0$、$q(0)=1$说明计算顺序，区间为$[0,1]$，真实解为$e^{-x}$。暂用$q_{a,b}(x)=ax+b$代替大网络，则输入导数为$q'=a$，残差为$r(x)=a+ax+b$。取配点$x=0,1$，边界权重为$1$，损失是
\[
 \mathcal L(a,b)=\tfrac12[(a+b)^2+(2a+b)^2]+(b-1)^2.
\]
在$a=0,b=1$时，边界误差为零，但两个方程残差均为$1$，所以$\mathcal L=1$。参数梯度为$\partial_a\mathcal L=(a+b)+2(2a+b)=3$、$\partial_b\mathcal L=(a+b)+(2a+b)+2(b-1)=2$。学习率$0.1$的一步更新得到$a=-0.3,b=0.8$，残差变为$0.5,0.2$，总损失为$0.185$。这里先对坐标$x$求导以形成物理残差，再对参数$a,b$求导以学习函数，是两次不同的求导。直线模型不能在全区间精确表达$e^{-x}$，即便优化完成也有表达误差。

自动微分沿加、乘、激活等运算的链式法则计算上述导数，不是以相邻观测之差近似导数；其结果在浮点误差范围内忠实于网络计算图，而非忠实于未知真实风场。这个小例还显示，软边界可能在更新中暂时变差，必须分别监测边界项与内部项。

\paragraph{从残差到参数梯度。}
将所有配点、观测和边界条件对应的残差排列为$r(\theta)\in\R^K$，令$\theta\in\R^P$、固定权重矩阵$W\in\R^{K\times K}$，并暂时写加权平方损失为
\[
\mathcal L(\theta)=\frac12 r(\theta)^{\mathsf T}Wr(\theta),\qquad W\succeq0.
\]
若$J_r=\partial r/\partial\theta\in\mathbb R^{K\times P}$，则逐项链式法则给出
\[
\nabla_\theta\mathcal L=J_r^{\mathsf T}Wr,
\qquad
\nabla_\theta^2\mathcal L\approx J_r^{\mathsf T}WJ_r,
\]
其中第二式忽略残差对参数的二阶项，是 Gauss--Newton 近似。这里的$J_r$不仅包含$\partial F_\theta/\partial\theta$，还包含例如$\partial(\partial_xu_\theta)/\partial\theta$和$\partial(\partial_{xx}u_\theta)/\partial\theta$；因此自动微分先对输入求导，再对所得残差对参数反向传播。残差很大并不必然产生大更新：若网络在该点对参数的雅可比接近零，梯度仍可能很小，这正是饱和激活和 PINN 病态优化的一个来源。

为简化这一段推导的记号，暂用$u,f,r_m$分别表示三维速度、体力和动量残差向量。对一个内部点的$r_m=\partial_tu+(u\cdot\nabla)u+\nabla\pi-\nu\nabla^2u-f$，在所需混合偏导连续时，其任意参数$\theta_p$的导数为
\[
\frac{\partial r_m}{\partial\theta_p}=\partial_t\frac{\partial u}{\partial\theta_p}
+\left(\frac{\partial u}{\partial\theta_p}\cdot\nabla\right)u
+(u\cdot\nabla)\frac{\partial u}{\partial\theta_p}
+\nabla\frac{\partial\pi}{\partial\theta_p}
-\nu\nabla^2\frac{\partial u}{\partial\theta_p}.
\]
这说明训练不是只让输出接近标签，而是让网络参数改变后其导数结构也朝残差下降的方向变化。上述式子假定$\nu$和$f$固定。若$\nu$也待估，则还要增加$-\partial_{\theta_p}\nu\,\nabla^2u$项；若外力$f$也依赖待估参数，还要增加$-\partial_{\theta_p}f$。遗漏这些依赖会使逆参数梯度错误。

边界条件应来自物理设置，而不是为了降低损失任意补齐。对一阶时间方程，给定初值通常决定时间演化；对含空间二阶导数的黏性方程，还需要足够的空间边界信息。若入口速度、压力和出口通量同时被强行指定而彼此不相容，连续问题本身可能无解。离散训练中，这表现为数据项与边界项的残差不能同时下降；此时调权重只是在两类冲突之间选择，并没有修复问题。相反，边界条件不足会留下零空间：存在不同的$\mathbf q$具有相同观测和相同内部残差，却在未约束边界上不同。因此应通过边界类型、数量和独立观测检查唯一性，而不是以网络输出平滑作为证据。
上游测风可约束入流，地表可在适用分辨率下采用无滑移或壁面模型，开放边界则需要与流动方向和数值模型相符的条件。粗糙地表附近未被解析时，强行令全部近地网格速度为零会改变垂直剪切。初始场未知时可以用低维参数化联合估计，并明确先验来源；不能将任意零场当作真实初值。

软约束通过损失惩罚边界误差。Dirichlet 条件指直接指定边界上的函数值；对适合的这类条件，也可构造硬约束
\begin{equation}
 q_\theta(\boldsymbol\xi)=g(\boldsymbol\xi)
 +d(\boldsymbol\xi)\widetilde q_\theta(\boldsymbol\xi),
 \qquad d|_{\Gamma_D}=0.
\end{equation}
例如区间$[0,1]$上要求$q(0)=1,q(1)=0$，可取$q_\theta(x)=1-x+x(1-x)\widetilde q_\theta(x)$。代入两个端点即可检查，无论网络参数怎样更新，端点值都保持不变；内部残差仍需训练，求导时也必须对$x(1-x)$使用乘积法则。

这里$g$是满足边界值的延拓，$d$是适当光滑的消失函数。硬约束消除了该部分边界误差，却要求边界数据相容，并可能改变导数尺度。将含噪测量精确写入硬约束，会把测量误差固定在解中。三维速度也可写成光滑向量势的旋度以自动满足散度为零，但拓扑、边界构造和更高阶导数成本仍需处理。

复合损失存在梯度竞争：数据项可能快速下降，而动量项几乎不变，或者物理项主导更新使网络忽略观测。相关研究分析了 PINN 的梯度病态和训练失败模式\cite{wang2021gradient,krishnapriyan2021failure}。训练应分别跟踪各项残差及参数梯度范数，先检查尺度和方程，再考虑调整权重。自适应权重需要限幅或归一化，不能让难拟合的数据项被持续降权后掩盖模型失配。

\begin{algorithm}[htbp]
\caption{稀疏激光雷达风场重建的 PINN 训练}
\label{alg:pinn-wind}
\begin{algorithmic}[1]
\Require 训练扫描、观测算子、可用初边值信息、物理假设与验证扫描
\State 固定坐标、符号、时间戳、质量控制和无量纲尺度
\State 初始化坐标网络及有约束的待估参数
\For{每个训练阶段}
 \State 采样内部配点、边界点和初值点，保留基础覆盖
 \State 用距离门求积计算径向预测与数据损失
 \State 自动微分计算连续性、动量及边界残差
 \State 按式\eqref{eq:pinn-loss}更新参数，记录各项损失和梯度
 \State 在验证扫描上检查观测误差，并在独立配点上检查残差
\EndFor
\State 冻结模型和选择规则，在留出观测与独立仪器上验收
\end{algorithmic}
\end{algorithm}

Adam 可用于早期优化，随后可在固定配点和确定性目标上尝试 Limited-memory Broyden--Fletcher--Goldfarb--Shanno (L-BFGS)。频繁更换配点会改变后者的目标并影响曲率近似，因此两阶段优化不是无需条件的默认保证。停止规则应结合独立验证误差、残差和预算，而非只看总训练损失。若验证误差持续恶化，应检查边界冲突、噪声模型或欠拟合，不应单纯提高物理权重。

\section{主要变体及其适用条件}
残差自适应采样先在候选点上评估误差，再增加高残差区域的配点。它适合尾流剪切层和局部变化较强的位置，但必须保留全域基础覆盖。高残差也可能来自错误方程或边界，反复加点不一定能够解决。若原目标是均匀体积积分，按非均匀分布采样后应进行相应的重要性修正，否则训练实际强调了新的空间权重。对独立均匀测试配点的评估可以区分局部拟合与全域改善。

弱形式或变分 PINN 将微分方程乘以测试函数并积分，适当分部积分可降低网络所需的导数阶数\cite{kharazmi2021hpvpinn}。例如对扩散项和在边界为零的测试函数$\varphi$，有
\begin{equation}
 \int_\Omega(-\nabla^2q)\varphi\,d\mathbf x
 =\int_\Omega\nabla q\cdot\nabla\varphi\,d\mathbf x.
\end{equation}
非零边界测试函数还会产生通量边界项。弱形式减少了直接计算二阶导数的需要，但测试空间和积分精度决定哪些残差能够被检测，不能把少量积分约束视为处处守恒的证明。

区域分解为不同子域训练局部网络，并在交界处约束状态与相应通量。Extended Physics-Informed Neural Network (XPINN) 将这种思路扩展到时空分解\cite{jagtap2020xpinn}。风电场可按风机邻域或尾流区域分区，以改善局部尺度表示并支持并行训练。接口仅匹配速度而忽略压力或应力一致性，仍可能产生动量跳变，分区数量增加也会增加接口优化成本。

时间分窗和因果训练强调先拟合早期状态，再逐步处理后续时间。一个常见的加权思路是令第$k$个时间段的残差权重取
\begin{equation}
 a_k=\exp\left(-\gamma\sum_{j<k}\mathcal L_{f,j}\right),\qquad \gamma>0.\label{eq:pinn-causal-weight}
\end{equation}
早期误差较大时，后续时间段暂时减小权重。使用式\eqref{eq:pinn-causal-weight}时应说明是否对权重停止梯度，以及如何逐渐开放后续时间段。权重长期接近零会造成后段完全未学到，必须单独验收每个时间窗。窗口交界还需状态传递，错误初值可能沿时间累积。

逆参数 PINN 将$\boldsymbol\eta$与$\theta$共同优化。正黏度可通过$\nu_{\mathrm{eff}}=\nu_{\min}+\log(1+\exp a)$参数化，但正值并不等于物理可辨识。应比较不同初值和先验下的参数稳定性，并分析参数变化是否会被速度或边界补偿。只有参数在独立工况上也能解释观测时，才适合将其解释为物理量，而不仅是有效拟合系数。

贝叶斯 PINN 则为参数或场表示设置先验，以后验表达部分不确定性\cite{yang2021bpinn}。形式上可写为
\begin{equation}
 p(\theta,\boldsymbol\eta\mid\mathbf y)
 \propto p(\mathbf y\mid\theta,\boldsymbol\eta)
 p(\text{物理约束}\mid\theta,\boldsymbol\eta)
 p(\theta,\boldsymbol\eta).
\end{equation}
将物理残差当作概率因子时，残差尺度也必须建模，不能任意解释成已知测量噪声。多个随机种子训练的集合能揭示部分优化敏感性，但不是严格的后验样本。无论采用哪种区间，都应在独立观测上检查覆盖率，并区分测量噪声、有限观测和方程失配造成的不确定性。

\section{风电场案例的验证流程}
\label{sec:pinn-validation}
一个可检验的案例应先写清场地、扫描几何和重建时间窗。每条观测至少记录时间戳、雷达位置、波束方向、距离门响应、径向速度和质量标记。风机位置与运行状态用于定义几何或体力，入口信息只能使用任务开始时确实能够得到的测风数据。若扫描覆盖多个高度，可以研究三维重建；若只有近水平的一层扫描，就应把任务限定为二维或准二维，并明确采用了怎样的垂直输运假设，不能据此宣称恢复完整三维湍流。

重建窗口内使用晚于目标时刻的观测属于离线平滑，不能称为实时预测。在线重建只能使用当时已到达的扫描，并计入扫描延迟、质量控制、优化和查询时间。滚动窗口可用上一窗口参数初始化，但要记录跨窗口传递的信息。跨日期或跨风况部署需要独立测试，单窗口拟合成功不代表已经学习到可直接迁移的通用求解器。

数据划分应按完整扫描、连续时间段或独立雷达组织，避免将同一扫描的相邻距离门随机分到训练与测试后夸大泛化。验证集用于选择权重、网络规模和停止时刻，测试观测在选择完成前保持隐藏。已知几何和公共物理方程可以用于测试区域配点，但测试风速及由其反推的边界不能进入训练。应区分窗口内未观测位置的插值重建与新窗口的泛化，两者难度不同。

仅有径向观测时，可以在留出集合$\mathcal T$上报告
\begin{equation}
 \operatorname{RMSE}_{\mathrm{LOS}}
 =\sqrt{\frac1{|\mathcal T|}\sum_{m\in\mathcal T}
 (\mathcal H_m[\mathbf u_\theta]-y_m)^2}.
\end{equation}
这一指标检验测量空间的一致性，不能证明每个速度分量都准确。若有独立三维风速参考，还应报告各分量误差、风向误差以及按空间区域分组的误差。低风速下风向不稳定，需要事先定义有效风速范围。对有限测点或探测体积，应将预测映射到参考仪器的采样尺度后比较。

合成试验可以从未参与训练的流场参考中模拟波束、距离门平均、扫描顺序和噪声，从而获得全场误差；若训练和参考都使用完全相同的简化方程，试验只检验理想匹配条件。为检查模型失配，还应改变入流、闭合或地形条件。真实试验可使用独立交叉雷达或测风仪器，但应说明其覆盖范围和自身误差。缺少全场真值时，不能用满足所假设方程来替代真实精度证据。

对照至少包括采用相同径向观测算子的无物理网络、适用的正则化反演，以及可获得时的数值模型同化。把已知完整速度分量交给某个基线，却只给 PINN 径向速度，会改变任务条件。所有方法应共享可用观测、边界来源和调参划分，同时报告训练或同化成本、单点查询成本和整个时间窗的交付延迟。

消融可以依次检查连续性约束、动量约束、真实距离门算子和边界信息的作用，并在匹配计算预算下比较。若去掉体积平均后径向误差变化很小而剪切误差增大，说明观测拟合不足以识别局部结构。若增加物理权重降低残差却提高独立观测误差，则应检查方程和边界失配。若更多扫描角度显著降低分量误差，收益可能主要来自观测几何改善，而不是网络规模。

下游应用还应验收轮毂高度速度、转子面平均入流、尾流亏损与位置等目标，并说明参考上游风速的定义。由稀疏平均场不能直接推断湍流强度或疲劳载荷；这些量需要相应时间分辨率和动力学信息。物理残差应在未参与训练的配点上报告，并辅以控制体净通量等积分检查。所有残差较小仍只是所选模型的一致性证据。

部署时可以联合监测径向创新、分区域物理残差、观测几何和不确定区间。观测大面积缺失、波束几何退化或边界状态超出训练范围时，应扩大不确定性说明并限制输出用途，必要时回退到已验证的测风或同化流程。自动微分能够精确计算网络导数，但不能保证网络代表了真实大气。

物理信息神经网络需要同时理解自动微分和方程残差，下面的材料提供相应讲解与代码示例。
\section*{辅助学习材料}
\begingroup
\emergencystretch=3em
\begin{itemize}
\item \textbf{Burgers equation}；作者/平台：DeepXDE；英文；官方代码教程。阅读 Problem setup 和 Implementation，以 Burgers 方程核对时空域、初边值条件、Jacobian/Hessian、PDE 残差与训练损失的代码对应。\href{https://deepxde.readthedocs.io/en/latest/demos/pinn\_forward/burgers.html}{在线阅读}。
\item \textbf{The Fundamentals of Autograd}；作者/平台：PyTorch 官方教程；英文；代码教程。重点看 requires\_grad、Jacobian、higher-order derivatives 与 autograd，帮助实现本章的 PDE 导数、残差构造和梯度检查。\href{https://docs.pytorch.org/tutorials/beginner/introyt/autogradyt\_tutorial.html}{在线阅读}。
\item \textbf{PINN——加入物理约束的神经网络}；知乎；中文解说。分别写出观测损失、方程残差与边界条件，对照本章约束设计；标题与主题据必应索引，原页受限。\href{https://zhuanlan.zhihu.com/p/544561165}{在线阅读}。
\item \textbf{物理信息神经网络（PINN）: 将物理知识融合到深度学习中}；CSDN；中文博客。核对自动微分与无量纲化；低方程残差不保证逆问题可辨识；标题与主题据必应索引，原页受限。\href{https://blog.csdn.net/qq\_28247201/article/details/136072595}{在线阅读}。
\item \textbf{物理信息神经网络 PINNs 融合物理定律与深度学习}；上传者：北游知；bilibili；中文技术讲解。按公开简介对照本章自动微分、方程残差、边界条件与参数反演；先核对损失各项，再检查尺度和可辨识性，不把教程中的前景描述当作精度保证。2026年10月4日公开元数据播放6,523次；在直接相关的中文解说候选中优先选用，未采用播放更高的机翻转载课程。已核对公开简介、标题和计数，未逐帧观看。\href{https://www.bilibili.com/video/BV1WfQsB6EKn/}{观看视频}。
\end{itemize}
\endgroup
\paragraph{本章小结。}
PINN 用方程帮助推断稀疏观测之间的风场。使用前要写清雷达的测量方式和重建尺度，再设置方程残差、初值和边界条件。多放一些配点，只是多检查一些位置是否符合方程；增加独立视角的观测，才可能带来新的测量信息。报告结果时应同时说明测量精度、方程与现场的偏差、数据能否确定所求风场，以及计算开销。一张平滑的图和较低的训练损失还不够。


\part{从学习方式看数据智能场景}
前面各章主要讨论如何用给定数据学习预测器。本部分进一步讨论数据与知识如何被组织：预训练构建可复用的表示，迁移学习处理源域与目标域，持续学习处理按时间到达的变化，主动学习决定哪些样本值得标注，联邦学习处理分散的数据所有权，流形学习利用样本几何，元学习利用任务经验。七章分别给出问题假设、关键算法、变体与工业评价边界，不能用其中一种方法的收益替代另一种设定的验证。
\chapter{预训练：构建可复用的数据表示}
\label{chap:pretraining}
\begin{quote}\small
\textit{本章摘要。}预训练通过先行学习任务获得可复用的参数或表示。本章从数据与目标的关系出发，解释监督预训练、掩码与生成式学习、对比学习、无负样本自蒸馏及领域继续预训练。工业振动信号案例用于说明如何选择增强与遮蔽方式，以及如何区分预训练损失下降、表示改善和下游收益。
\end{quote}

前面的模型章节回答序列、图像和文本应如何表示，上一章进一步让物理方程参与学习：当标签不足时，已知规律可以约束候选解。然而，不是所有设备都有足够准确的控制方程，许多可利用的知识只隐含在历史观测中。本章由此转向另一条路线：在目标故障标签不足时，能否先利用历史数据学习可复用结构，再用有限标签建立预测器？预训练的价值在于可复用性，其监督信号可以来自人工标签，也可以来自数据本身。它不等同于无监督学习，更不要求模型一定具有很大规模。

理解这条路线，需要区分三个并不等价的问题：辅助任务能否被完成，表示是否保留了任务需要的信息，有限标签能否把这些信息转成可靠判断。重构波形可以依靠平滑性，对比学习可以依靠设备身份，而故障诊断可能恰恰需要低能量的短时冲击。因此，本章不以罗列预训练模型为主线，而是逐步追问学习信号要求模型保留什么、丢弃什么：先由似然与重构解释预测目标，再由正负样本关系解释不变性，最后分析不使用负样本时如何避免所有表示趋同。贯穿这些方法的观点是，预训练目标是一种关于有用结构的假设，必须由独立的下游任务检验。

读完本章，可以按三个问题检查自己的方案：哪些历史记录允许参与训练，遮住的部分是否太容易猜中，模型是否把不同输入都编码成了相近的结果。例如波形补得很平滑，却丢掉短时冲击，就未必有利于故障诊断。下一章继续讨论拿到预训练参数以后，哪些可以保留，哪些需要用新设备的数据重新调整。

\section{预训练学到了什么}

% auxiliary-resource-guide
本章末的“辅助学习材料”列出与核心方法对应的讲解和实践入口，可在阅读相关推导后，按所列小节对照学习。

\begin{figure}[htbp]
\centering
\begin{tikzpicture}[node distance=0.75cm and 0.75cm]
 \node[idi input, minimum width=2.0cm] (data) {大规模数据\\$\mathcal D_{\rm pre}$};
 \node[idi process, right=of data, minimum width=2.0cm] (obj) {预训练目标\\重构/对比/生成};
 \node[idi state, right=of obj, minimum width=2.0cm] (enc) {可复用编码器\\$f_{\theta_0}$};
 \node[idi output, right=of enc, minimum width=2.0cm] (task) {目标任务\\线性探测或微调};
 \draw[idi arrow] (data)--(obj); \draw[idi arrow] (obj)--(enc); \draw[idi arrow] (enc)--(task);
 \draw[idi feedback] (task.south) .. controls +(0,-0.42) and +(0,-0.42) .. node[below,font=\scriptsize]{下游收益反查预训练目标} (obj.south);
\end{tikzpicture}
\caption{预训练的价值由下游任务检验：辅助目标下降并不自动意味着故障相关信息被保留。}
\label{fig:pretraining-downstream-validation}
\end{figure}
图\ref{fig:pretraining-downstream-validation}把无标签预训练与下游评估连成反馈路径，强调辅助任务学到的表示须通过独立故障任务检验，而不能只用预训练损失验收。反馈只允许使用训练与验证任务：若根据测试集表现反复改变预训练目标，该测试集就参与了模型选择，不能再充当最终独立证据。
阅读下面的目标函数，只需先掌握小批次训练、交叉熵和反向传播。与从零训练相比，本章多了一个先行任务：用什么训练信号产生可复用参数？先辨认每个目标的输入和答案来自哪里，再判断学到的特征能否帮助有限标签任务。

\paragraph{表示、编码器和头分别是什么。}
表示是从原始输入计算出的特征向量，而不是人工故障标签。例如一个窗口含$C=3$个通道、每通道$L=100$个采样值，可展平为$x\in\R^{300}$，编码器$f_\theta:\R^{300}\to\R^{32}$把它变成$32$维表示。重构头$g_\phi:\R^{32}\to\R^{300}$把表示映射回波形；下游三分类头则把同一表示映射到$3$个类别分数，经 softmax 变成概率。头是为特定目标服务的末端模块，骨干或编码器负责提取共享特征。不同任务的头输出尺寸不同，所以预训练结束时常需替换它，而非直接把重构波形当成故障概率。

在线性探测中，固定$h=f_{\theta_0}(x)\in\R^{32}$，只学习$W\in\R^{3\times32}$、$b\in\R^3$，计算类别分数$a=Wh+b$和概率$\softmax(a)$。这里“线性”指分数对表示$h$是仿射的，不是对原始波形一定线性，也不是把 softmax 删除。若表示已丢掉故障冲击，再训练这样的头不能恢复被删信息；全量微调则允许编码器重新学习，但需要更多目标证据控制过拟合。

“预训练”描述先学习可复用参数的阶段，“自监督”描述标签如何由数据本身构造，“迁移学习”描述把已有知识用于目标问题的做法。它们不互斥：用掩码恢复历史波形是自监督预训练，把得到的编码器用于新设备故障分类是迁移，训练新分类头或更新编码器则是具体适配操作。监督预训练也能用于迁移，迁移也不必使用单独设计的自监督阶段。

设预训练数据为$\mathcal D_{\mathrm{pre}}$，编码器为$f_\theta$，与预训练目标配套的预测头为$g_\phi$。一般训练形式为
\begin{equation}
 (\theta_0,\phi_0)=\argmin_{\theta,\phi}
 \E_{x\sim\mathcal D_{\mathrm{pre}}}
 \E_{\omega\sim P_\omega}
 \ell_{\mathrm{pre}}\bigl(g_\phi(f_\theta(v_\omega(x))),s_\omega(x)\bigr).
\end{equation}
式中$\argmin$表示寻找使平均损失尽量小的参数，外层期望平均不同记录，内层期望平均同一记录的随机处理方式，实际通常以一个小批次近似。下标$0$表示交付下游任务的预训练参数，并非所有参数取零。监督版本的外层抽样应读作$(x,y)\sim\mathcal D_{\mathrm{pre}}$，此时目标$s_\omega$取该记录的人工标签$y$；不能只凭输入$x$凭空得到人工标签。

其中$\omega$控制裁剪、遮蔽或目标构造，$v_\omega(x)$为模型输入，$s_\omega(x)$为学习目标。监督分类时目标来自人工标签，自监督任务则由样本结构构造目标，例如被遮蔽的片段或同一样本的另一视图。预训练完成后可保留编码器、丢弃辅助头，再接入目标任务头$h_\psi$。

若冻结$\theta_0$只训练线性$h_\psi$，检验的是该表示中的线性可分信息；全量微调则同时检验初始化能否支持适应。两者可能产生不同排序。模型也可能保留了设备编号和工况，却没有保留故障机理，因此预训练损失不能直接用作表示质量指标。下一章讨论微调、参数高效适配与域对齐，本章重点是初始化产生之前的学习过程。

预训练数据不一定与下游任务同分布。通用图像表示可能提供边缘和纹理结构，但工业细裂纹要求保留更细的空间信息；自然语言模型能够学习句法，却未必知道某台设备的有效规程。数据规模、任务相关性和目标设计共同决定迁移潜力，增加其中一项不能保证补偿另一项的缺失。

\section{监督预训练与重建式目标}
监督预训练利用较大源数据集的标签训练骨干，例如图像分类网络中的交叉熵目标\cite{he2016resnet}：
\begin{equation}
 \mathcal L_{\mathrm{sup}}=-\frac1N\sum_{i=1}^{N}
 \log p_{\theta,\phi}(y_i\mid x_i).
\end{equation}
随后保留骨干并替换分类头。源标签可以来自其他设备或相关任务，但其类别定义决定了哪些差异受到强调。只以设备类别预训练，可能更容易区分传感器安装位置而不是早期故障。弱标签和伪标签能扩大覆盖，却需要单独处理标签噪声，不能把它们等同于准确人工标注。

自编码器从压缩表示重构输入，去噪自编码器则从受扰动输入$\widetilde x$恢复原始$x$\cite{vincent2008denoising}：
\begin{equation}
 \mathcal L_{\mathrm{denoise}}
 =\E_{x\sim\mathcal D_{\mathrm{pre}}}\E_{\widetilde x\sim q(\cdot\mid x)}
 \|g_\phi(f_\theta(\widetilde x))-x\|_2^2.
\end{equation}
扰动机制决定模型需要恢复什么。噪声过小可能只要求近似恒等映射，噪声过大可能使恢复目标含糊。重构均方误差倾向于优先解释高能量成分，低幅值但诊断价值较高的冲击不一定受到足够重视。频谱或多尺度损失可以改变关注点，但要明确引入的频段先验，避免利用测试故障选择损失。

\subsection{从条件似然理解重构为何产生平滑结果}
把一个窗口展平为$x\in\R^D$，令$u=\widetilde x$表示受扰动观测，编码器输出$f_\theta(u)\in\R^d$，解码器输出$m_{\theta,\phi}(u)\in\R^D$。若建模假设为
\begin{equation}
 p_{\theta,\phi}(x\mid u)=\mathcal N(x;m_{\theta,\phi}(u),\sigma^2 I_D),
\end{equation}
则单个样本的负对数似然是
\begin{equation}
 -\log p(x\mid u)=\frac D2\log(2\pi\sigma^2)
 +\frac1{2\sigma^2}\|x-m_{\theta,\phi}(u)\|_2^2.
\end{equation}
只有在$\sigma^2$固定时，最大似然才等价于普通均方误差最小化；若同时学习方差，就不能丢弃对数方差项。不同通道单位差异很大时，采用对角协方差$\Sigma$会得到$(x-m)^\top\Sigma^{-1}(x-m)$，即按噪声尺度加权的误差。这样的权重表达的是可靠性或尺度假设，不能任意解释为故障重要性。

进一步令$m^*(u)=\E[X\mid U=u]$，对任意候选重构$a$展开平方：
\begin{align}
 \E[\|X-a\|^2\mid u]
 &=\E[\|X-m^*+m^*-a\|^2\mid u]\\
 &=\E[\|X-m^*\|^2\mid u]+\|m^*-a\|^2.\label{eq:pretraining-conditional-mse}
\end{align}
式\eqref{eq:pretraining-conditional-mse}中交叉项消失，是因为$\E[X-m^*\mid u]=0$。因此无限容量、总体风险最优的平方误差预测器输出条件均值。假如遮蔽区的冲击可能出现在两个不同位置，条件均值可能在两处各产生半个冲击，虽然平均误差最小，却不是任何一次真实波形。这不是单纯增加网络容量就能解决的问题；需要缩短遮蔽区、加入可靠上下文，或用能表达多峰条件分布的预测头。

用两个等可能的真实片段$(2,0)$与$(0,2)$可检查条件均值结论。预测$(1,1)$时，两种情况的平方误差都为$2$；固定预测$(2,0)$时，误差分别为$0$和$8$，平均为$4$。因此均值确实赢得平方损失，却把一个完整冲击拆成两个半冲击。重构指标和故障语义之间的差异由此可见。

\subsection{遮蔽任务与因果预测的边界}
掩码建模只显示部分输入，并预测被遮蔽区域。令$M$为目标位置集合、$V$为可见位置集合，可写为
\begin{equation}
 \mathcal L_{\mathrm{mask}}
 =\E_{x,M}\frac1{|M|}\sum_{j\in M}
 \ell\bigl(\widehat x_j(x_V,M),x_j\bigr).
\end{equation}
离散词元通常采用交叉熵，连续图像块或信号片段可采用重构误差。BERT 使用掩码语言建模，MAE 则以可见图像块训练编码器，并由轻量解码器重构被遮蔽块\cite{devlin2019bert,he2022mae}。MAE 的编码器只处理可见块，减少了高遮蔽率下的编码开销；解码器仍需接收位置信息与遮蔽标记。不能仅按遮蔽比例推断所有架构都具有相同节省。

这里假定输入由$L$个元素或片段组成，每个连续片段$x_j\in\R^p$，$M\subset\{1,\ldots,L\}$非空，$V$为其补集；解码器需为每个$j\in M$输出同维预测。若采用逐位置条件分布$q_{\theta,\phi}(x_j\mid x_V,M)$，掩码损失就是条件负对数似然之和。逐位置相加相当于用条件独立输出头拟合各位置的边缘分布，不表示真实被遮蔽片段给定可见输入后必然独立；若要生成相互协调的一整段，还需联合或自回归解码器。$1/|M|$避免遮蔽更多位置时仅因项数增加而提高损失，但它不会消除不同遮蔽难度带来的梯度差异。

离散监督分类和掩码词元分类共享交叉熵结构。令真实条件分布为$r(k\mid u)$、预测为$p(k\mid u)$，则$-\sum_kr_k\log p_k=H(r)+\KL(r\|p)$。固定输入条件下，总体最优预测应匹配$r$；有限标签训练以单个类别指示向量代替未知$r$。因此监督预训练强调源标签可区分的信息，掩码训练强调由上下文可预测的信息，二者都没有数学理由自动保留与自身目标无关的下游细节。

例如$x=(1,2,3,4)$，遮蔽集合$M=\{2,3\}$，输入只显示位置$1,4$及其数值。若模型对遮蔽位置输出$(2.5,2.5)$，则按本节平方误差约定，损失为$[(2.5-2)^2+(2.5-3)^2]/2=0.25$；可见位置不参与这次目标平均。监督值来自原始完整记录，但不能先把完整记录编码后才遮蔽目标，否则模型已经看过答案。连续片段含$p$个数时，单位置损失可用平方范数，是否再除以$p$也应统一。

独立随机点遮蔽可能被局部平滑插值轻易解决。时间序列可使用连续片段遮蔽、跨通道遮蔽或多尺度遮蔽，以迫使模型利用更长上下文。遮蔽过长又可能使目标不可预测，网络只能输出条件平均。应根据信号周期、传感器耦合与任务时间尺度选择策略，并检查模型是否只是复制相邻样本。

自回归预训练预测给定过去的下一个元素，目标为
\begin{equation}
 \mathcal L_{\mathrm{AR}}=-\E_x\sum_{j=1}^{L}
 \log p_\theta(x_j\mid x_{<j}).
\end{equation}
该目标来自联合分布的链式分解$p_\theta(x_1,\ldots,x_L)=\prod_{j=1}^Lp_\theta(x_j\mid x_{<j})$；取负对数将连乘转为求和，不要求序列元素相互独立。训练时使用真实前缀计算条件似然，而生成时使用模型已生成的前缀，两者输入分布可能不同。对总体数据分布$P$，期望负对数似然等于$H(P)+\KL(P\|P_\theta)$；模型族足够表达且优化达到总体最优时，才有匹配数据分布的解释。有限数据、有限上下文或模型失配下，即便整体似然改善，稀有故障的条件概率也可能仍很差。对文本这是下一词元预测，对连续信号需要明确预测分布或回归损失。双向掩码表示可以利用目标位置两侧上下文，自回归表示受因果边界约束。在线故障预测若采用窗口表示，应保证窗口中没有预测时刻之后的数据，不能因任务名称为自监督就忽略时间泄漏。

\section{对比学习：哪些样本应当接近}
对比学习用正样本对定义应保留的一致性，用负样本提供区分信号。A Simple Framework for Contrastive Learning of Visual Representations (SimCLR) 对一个批次中的每个样本独立增强两次，经编码器与投影头得到归一化表示$z_i$\cite{chen2020simclr}。对锚点$i$及其正样本$j$，损失为
\begin{equation}
 \ell_{i,j}=-\log
 \frac{\exp(z_i^\top z_j/\tau)}
 {\sum_{k\ne i}\exp(z_i^\top z_k/\tau)},\qquad \tau>0.
\end{equation}
含$B>1$个原始样本的批次产生$2B$个视图；若$B=1$，只有正候选，损失恒为零，无法提供这里的对比信号。分母含一个正样本和$2B-2$个负样本，训练对两个方向求平均。投影头使对比目标作用于专门的空间，下游任务常使用投影前的编码器表示。温度$\tau$改变相似度差异对梯度的影响，较小温度强调难区分样本，也可能放大假负样本的影响。

“视图”就是同一条记录经不同允许处理得到的输入，正样本不是“正常设备”，负样本也不是“故障设备”，而是当前配对任务中匹配与不匹配的候选。例如$B=2$，某锚点有一个正候选和两个负候选，若与三者的相似度分别为$1,0,0$且$\tau=1$，则正候选概率为$e/(e+2)\approx0.576$，损失约$0.551$。若三者都同样相似，则概率为$1/3$、损失为$\log3\approx1.099$。这说明对比训练不是单独提高正相似度，而是提高它相对于其他候选的辨识度。

\paragraph{把对比损失展开为一个分类问题。}
设投影前表示为$h_i\in\R^d$，投影后归一化向量为$z_i\in\R^m$，$\|z_i\|_2=1$。固定锚点$i$时，其余$2B-1$个视图构成候选集合。令$s_{ik}=z_i^\top z_k$及
\begin{equation}
 q_{ik}=\frac{e^{s_{ik}/\tau}}{\sum_{l\ne i}e^{s_{il}/\tau}}.
\end{equation}
任务就是从候选中识别真正配对的视图$j$。把$\ell_{i,j}=-s_{ij}/\tau+\log\sum_{k\ne i}e^{s_{ik}/\tau}$逐项求导，有
\begin{equation}
 \frac{\partial\ell_{i,j}}{\partial s_{ik}}
 =\frac1\tau(q_{ik}-\mathbf1[k=j]),\qquad
 \nabla_{z_i}\ell_{i,j}
 =\frac1\tau\left(\sum_{k\ne i}q_{ik}z_k-z_j\right).
\end{equation}
这里第二式暂把其他向量视为常量；完整批次梯度还包含$z_i$作为别的锚点的候选时收到的梯度。下降方向拉近正样本并远离概率加权的候选均值，而非均匀排斥每个负样本。若未归一化投影为$u_i\ne0$，还需乘归一化雅可比$(I-z_iz_i^\top)/\|u_i\|_2$，不能把单位球面上的变量当成任意自由向量。零向量不能直接归一化；实现中若用小常数保护分母，保护区间内的导数也应按实际公式计算。

这个推导解释了温度的两个作用：显式的$1/\tau$改变梯度尺度，softmax 又把权重集中到最相似候选。因此减小温度不只是“加大学习率”。当一个高相似负样本实际属于同类故障时，它会得到很大的排斥权重。另一方面，若所有$z$完全相等，则$q_{ik}=1/(2B-1)$，损失为$\log(2B-1)$，没有完成配对识别；但对称构型仍可能有退化驻点，不能仅从负样本存在推断优化必然成功。

大批次可以提供更多负样本，但增加内存成本。Momentum Contrast (MoCo) 使用键队列复用近期表示，并用动量编码器降低队列中新旧表示的不一致\cite{he2020moco}。若查询编码器参数为$\theta_q$、键编码器参数为$\theta_k$，则
\begin{equation}
 \theta_k\leftarrow\mu\theta_k+(1-\mu)\theta_q,\qquad 0\leq\mu<1.
\end{equation}
查询分支通过梯度更新，键分支按指数移动平均更新，当前键入队并移除最旧键。队列增大并不总有益：分布变化较快时，陈旧键可能形成不合适的比较对象。

正样本定义比损失名称更重要。图像颜色扰动可以弱化照明差异，但在变色故障中颜色本身就是标签信息。振动信号的强幅值归一化可能删除故障严重度，时间伸缩可能改变与转速相关的频率结构。只有在目标任务对某种变化确实不敏感时，才应训练这种不变性。状态变化对应的敏感方向应保留，而不是一律压缩。

同一故障的两个不同记录可能被当作负样本，造成假负样本排斥。监督对比学习利用已有类别标签为一个锚点构造多个正样本\cite{khosla2020supcon}，但其适用范围由标签覆盖决定。工业采样还应控制设备和批次：若一个批次中每台设备只出现一种工况，模型可能靠设备身份完成对比任务，而非学习可迁移的状态表示。

跨模态对比预训练将同步波形、图像或文本对齐。CLIP 使用图像与文本的双向批次对比目标\cite{radford2021clip}，其思想可以扩展到工业配对数据，但工单时间不一定等于故障发生时间，粗糙配对会引入错误正样本。跨模态目标学到的是给定配对规则下的关联，不会自动建立因果对应。

\section{无负样本学习与表示坍缩}
如果只令两个增强视图的表示接近，所有输入都映射到同一个常量也能使距离为零，这称为表示坍缩。Bootstrap Your Own Latent (BYOL) 通过在线编码器、预测头和移动平均目标编码器构成不对称训练\cite{grill2020byol}。用$\operatorname{sg}$表示停止梯度，归一化后的预测与目标分别为$\overline q_\theta$和$\overline z_\xi$，则一个方向的损失为
\begin{equation}
 \ell_{\mathrm{BYOL}}=
 \left\|\overline q_\theta(v_1)
 -\operatorname{sg}\bigl(\overline z_\xi(v_2)\bigr)\right\|_2^2,
 \qquad \xi\leftarrow\mu\xi+(1-\mu)\theta.
\end{equation}
实际训练交换视图后再平均。这里$\theta$和$\xi$表示两分支中对应编码与投影模块的参数，在线预测头单独通过梯度更新。停止梯度和移动平均控制目标变化，并不使任意网络与任意数据增强都自动免于坍缩。

Simple Siamese (SimSiam) 去掉移动平均目标更新，使用共享编码器、预测头和停止梯度构造更简洁的孪生学习\cite{chen2021simsiam}。Self-Distillation with No Labels (DINO) 则让学生匹配教师在不同裁剪视图上的概率输出，并以移动平均教师、输出中心化和温度锐化调节目标分布\cite{caron2021dino}。这些方法的教师通常来自学生训练轨迹，而不是额外人工标签，因此属于自蒸馏。训练中要检查特征方差、有效秩和样本间距离，不能只看两个分支越来越一致。

另一类方法直接约束表示统计量。VICReg 组合视图一致性、各维方差下限和非对角协方差惩罚\cite{bardes2022vicreg}：
\begin{equation}
 \mathcal L=\lambda_s\mathcal L_{\mathrm{inv}}
 +\lambda_v\mathcal L_{\mathrm{var}}
 +\lambda_c\mathcal L_{\mathrm{cov}}.
\end{equation}
对中心化批次表示$Z\in\R^{B\times d}$，协方差为$C=Z^\top Z/(B-1)$，协方差项惩罚$C$的非对角元素平方，方差项则惩罚标准差低于给定阈值的维度。两项通常分别作用于两个视图分支。它以显式统计约束抑制常量表示和维度冗余，但小批次统计不稳定，低相关性也不保证所有维度都具有诊断意义。

\paragraph{停止梯度和统计约束究竟改变了什么。}
对单位向量$q,z$，$\|q-z\|^2=2-2q^\top z$，所以 BYOL 的距离最小化也在最大化方向一致性。若目标$z$停止梯度，单次更新只改变在线分支，而目标编码器通过移动平均间接改变。把动量递推展开，得到
\begin{equation}
 \xi_t=\mu^t\xi_0+(1-\mu)\sum_{s=1}^t\mu^{t-s}\theta_s.
\end{equation}
教师因此是历史学生的指数加权组合：较大的$\mu$平滑噪声，也增加对新分布的响应延迟。停止梯度不是修改距离的数值，而是修改计算图的导数路径；它与预测头、归一化和优化共同作用，单独的常量映射依然能让一致性项为零。

为了明确 VICReg 如何额外排斥这个解，令未中心化批次表示为$Z^{(a)}\in\R^{B\times d}$、$a\in\{1,2\}$，$B>1$，中心化后协方差为$C^{(a)}$。一组明确的损失约定是
\begin{align}
 \mathcal L_{\mathrm{inv}}&=\frac1B\sum_{i=1}^B\|z_i^{(1)}-z_i^{(2)}\|_2^2,\\
 \mathcal L_{\mathrm{var}}&=\frac1{2d}\sum_{a=1}^2\sum_{j=1}^d
 \max\{0,\gamma-\sqrt{C_{jj}^{(a)}+\epsilon}\},\\
 \mathcal L_{\mathrm{cov}}&=\frac1{2d}\sum_{a=1}^2\sum_{j\ne k}(C_{jk}^{(a)})^2.
\end{align}
这里$\gamma>\sqrt\epsilon>0$规定标准差下限。常量表示使一致性和协方差项都为零，却令每一维方差惩罚严格为正；所有维度复制同一个变化信号能够通过方差检查，却会产生很大的非对角协方差。这说明两种统计惩罚分别针对无变化和重复编码，缺一不可。不过$\operatorname{rank}(C)\leq\min(d,B-1)$，当$d>B-1$时不可能在一个批次内同时达到全维正方差与严格零非对角协方差。损失权重是在这些目标间折中，而不是保证有限批次出现理想白化表示。

可把停止梯度先看成标量操作：$L=(a-\operatorname{sg}(b))^2$在$a=1,b=2$时值为$1$，对$a$的导数为$-2$，而沿被停止的目标路径对$b$的导数为零。若移去$\operatorname{sg}$，数值仍为$1$，但对$b$的导数变为$2$。教师不是“知道真相”的专家，只是提供当前更新暂不追赶的目标。比如教师参数$\xi=0$、学生对应参数$\theta=1$、$\mu=0.9$时，移动平均把教师更新到$0.1$，不是立刻复制成$1$。

无负样本方法避免了显式的负样本队列或大规模两两比较，却仍依赖增强、架构和优化。对比、蒸馏与重建可以组合，但新增目标会引入权重与梯度冲突。应先验证单一目标，再判断组合是否保留了不同的有效信息，而不是因目标数量增加就推断表示更通用。

\section{数据组织与领域继续预训练}
预训练数据应登记来源、设备、时间、工况和使用许可。没有标签不等于可以随意使用：历史记录可能含有机密信息，也可能已经包含后续测试期的数据。对连续波形，应先按设备或时间段划分，再切窗口；否则同一段波形的重叠片段会落入不同集合，让表示看起来比实际更好。归一化统计和滤波参数也只能用规则允许的数据估计。

领域继续预训练在已有参数上使用领域数据延续学习。领域自适应预训练与任务相关语料预训练在语言任务中已有系统研究\cite{gururangan2020dontstop}。对工业波形，相应思路是在通用信号表示上使用目标设备族的历史无标签数据，但应区分部署前可用历史与未来测试期数据。使用无标签测试输入属于不同的信息使用规则，不能悄悄并入普通归纳评估。

只使用狭窄领域数据继续训练可能遗忘通用能力，也可能过拟合设备身份。可比较较小学习率、较短训练、混合原始数据或限制可训练参数等策略。若采用数据混合，可令来源$s$以概率$q_s$被采样，$\sum_sq_s=1$。按样本量取$q_s$使大来源占主导，均衡来源则会重复小数据集；两者改变了优化分布，需要记录而不能仅报告总样本数。

例如两个来源分别有$900$和$100$个窗口，按样本量抽取对应$q_1=0.9,q_2=0.1$；均衡来源则取$q_1=q_2=0.5$。若总共抽取$1000$次，后一种策略平均各抽$500$次，小来源的单个窗口会被重复使用得更多，但并未产生$500$条独立信息。采样比例需要围绕目标设备群确定，不能只以损失更低为由过度重复容易的来源。

固定算力预算下，模型规模、有效样本量、输入长度和训练步数存在取舍。遮蔽率改变任务难度和部分架构的计算量，队列长度改变负样本覆盖，继续训练改变历史知识的保留。比较时应记录实际处理样本或词元数、去重比例、训练时间与峰值显存。跨领域缩放规律不能直接当作当前工业数据的性能承诺。

预训练产物应同时保存模型、输入规范、数据来源和训练截止时间。投影头是否保留、归一化层统计是否更新、通道顺序如何定义，都会影响复用结果。所谓冻结编码器还应说明是否固定归一化运行统计，否则无参数梯度更新也可能改变模型行为。预训练模型版本应与适用设备和数据使用规则绑定。

\section{工业振动信号的完整学习任务设置}
考虑利用多个设备的历史振动记录预训练编码器，再以有限故障标签进行诊断。首先限定预测对象：窗口状态分类、未来故障预测和故障严重度回归需要保留的信息不同。本例以窗口状态分类为教学任务，输入为窗口内多通道波形及允许获得的运行条件，不假定已有实验结果。

按设备划分预训练与下游评估范围，再按连续记录或运行批次组织窗口。跨设备泛化试验中，测试设备的无标签记录也不进入预训练；若研究目标设备无标签适应，则另设一组明确允许该信息的规则。故障状态标签只能来自可核查记录，不能将未标注窗口全部视为正常。

掩码方案可将波形分成连续片段，只向编码器显示可见片段，再预测被遮蔽片段。对比方案可使用不改变窗口状态的裁剪和经过幅值校准的小噪声，增强幅度在训练与验证范围内选择。若裁剪可能移除唯一故障冲击，两视图便不再可靠共享状态，应限制裁剪或改用其他目标。转速变化明显时，可以在明确的工况条件下比较，避免把转速差异一律当作应消除的干扰。

\begin{algorithm}[htbp]
\caption{工业信号预训练与下游表示评价}
\label{alg:industrial-pretraining}
\begin{algorithmic}[1]
\Require 按设备和时间划分的数据、预训练目标、有限下游标签
\State 固定允许使用的数据、窗口边界、归一化与增强规则
\State 初始化编码器及预训练辅助头
\For{每个预训练批次}
 \State 按设备和工况采样窗口，生成掩码或增强视图
 \State 计算选定的重建、对比或自蒸馏损失
 \State 更新在线参数，并按算法更新教师或队列
 \State 监测目标损失、表示方差和不同来源的训练贡献
\EndFor
\State 在固定标签子集上训练线性探测头，并单独进行微调对照
\State 用验证集选择配置，冻结选择后在独立测试集上评价
\end{algorithmic}
\end{algorithm}

对照应包括随机初始化、适用的监督预训练，以及在相同数据范围下的代表性自监督目标。线性探测与全量微调分别报告，至少在多个标签预算下绘制学习曲线，比较每类故障召回、宏平均指标、误报及校准。若使用同一个预训练模型服务多个任务，可另外报告摊销成本，但不能把预训练计算完全隐藏。

消融应保持数据和计算条件尽量一致。例如固定编码器比较随机点遮蔽与连续片段遮蔽，检验收益是否来自更长上下文；固定增强比较是否使用投影头，检验下游表示与训练目标空间的区别；固定预训练样本量比较来源混合，检验模型是否依赖某个设备族。注意单纯匹配训练轮数不一定匹配算力，尤其是不同遮蔽编码和对比批次设置之间。

预训练损失下降但线性探测不改善，可能说明任务主要学到了插值或身份线索。线性探测改善而微调收益消失，则需检查标签量、适应策略和优化预算。某些设备收益明显、另一些退化时，应报告子群差异，而非只用平均分掩盖负迁移。表示可视化只提供诊断线索，不能代替独立分类或预测评价。

监督预训练、自监督目标和表示坍缩可结合下面的教程复习，分别关注训练目标与对比/非对比表示学习。
\section*{辅助学习材料}
\begingroup
\emergencystretch=3em
\begin{itemize}
\item \textbf{Self-Supervised Representation Learning}；作者/平台：Lilian Weng；英文；技术教程。重点阅读 pretext tasks、上下文预测和 Contrastive Predictive Coding，帮助理解本章的重建、对比目标和表示质量。\href{https://lilianweng.github.io/posts/2019-11-10-self-supervised/}{在线阅读}。
\item \textbf{Transfer Learning for Computer Vision}；作者/平台：PyTorch 官方教程；英文；代码教程。重点看 pretrained weights、fine-tuning、feature extraction 和 augmentation，帮助理解本章从通用预训练到领域继续训练的实验边界。\href{https://docs.pytorch.org/tutorials/beginner/transfer\_learning\_tutorial.html}{在线阅读}。
\item \textbf{通俗易懂讲AI--自监督学习}；知乎；中文解说。区分预训练辅助目标与下游目标，比较遮蔽、重建和对比学习；标题与主题据必应索引，原页受限。\href{https://zhuanlan.zhihu.com/p/676130861}{在线阅读}。
\item \textbf{【深度学习】自监督学习详解（self-supervised learning）}；CSDN；中文博客。阅读对比学习和生成学习的入门说明，表示质量仍须由独立下游任务评价；已读取公开页面引言与方法概述。\href{https://blog.csdn.net/qq\_51392112/article/details/128806998}{在线阅读}。
\item \textbf{2021 - 自监督式学习 (二) – BERT简介}；李宏毅；bilibili；中文课程。选看第72集，以语言遮蔽预训练为例对照本章自监督目标。2026年10月4日官方公开元数据显示整套课程累计播放约299.5万，并非本分集播放量；已核对目录与计数，未逐帧观看。\href{https://www.bilibili.com/video/BV1Wv411h7kN/?p=72}{观看分集}。
\end{itemize}
\endgroup
\paragraph{本章小结。}
预训练的设计由三件事共同决定：哪些数据可以使用、什么变化应当保持一致、什么信息必须保留。监督目标利用已有标签，掩码与生成目标恢复缺失或后续内容，对比目标定义样本关系，自蒸馏和统计约束抑制表示退化。这些目标的数学结构对应不同偏好：平方误差输出条件均值，对比交叉熵放大难区分候选的作用，停止梯度改变优化路径，方差与协方差项分别约束常量坍缩和维度重复。理解这些偏好，比记住某个模型名称更有助于判断故障信息会不会被保留。

判断预训练是否有用，要在独立任务上给出相同标签预算，再比较线性探测或微调后的结果。报告中应写清用了哪些历史数据，怎样增强或遮蔽，表示有没有退化，以及故障识别究竟改善多少。辅助损失下降只能说明辅助任务学得更好。下一章接着检查这些表示换到新设备后是否仍适用，并讨论重加权、表示对齐、参数适配和负迁移。

\chapter{迁移学习}
\label{chap:transfer}
\begin{quote}\small
\textit{本章摘要。}迁移学习利用已有数据、表示或参数降低目标任务的学习成本。本章先区分任务变化与分布偏移，再从风险重加权、表示对齐和参数适配推导主要算法，最后讨论测试时适配与负迁移。选择依据是哪些假设仍成立，而不是预训练模型有多大。
\end{quote}

上一章把大量历史观测压缩成一个可复用的编码器，但“可复用”只是潜力，不是新设备上的性能承诺。旧相机学到的边缘结构可能仍然有用，照明统计可能需要更新，而不同产品的缺陷定义甚至可能不再相同。本章因此从参数初始化走向目标风险：不是询问源模型有多强，而是询问源数据中的哪些关系在目标环境中仍成立。

这里存在三种不同的干预。若预测关系不变而采样频率改变，可重加权样本；若不同观测方式背后存在共享语义，可尝试对齐表示；若目标任务需要新的决策规则，则必须更新分类头或部分骨干。三种干预对应不同假设，不能把它们看成可互换的技巧。本章先推导风险如何从源分布转换到目标分布，再解释统计对齐与域对抗到底消除了什么差异，随后讨论有限标签与计算预算下的参数适配，并用负迁移说明这些方法的边界。

假设两台相机的图像颜色变得更接近了，还要检查裂纹和正常纹理是否仍能区分。对齐指标下降，只说明某种分布差异缩小，不能直接说明目标设备上的故障更容易识别。章末再考虑连续更新：用新设备数据改好模型后，旧设备上的故障还能不能识别？下一章的持续学习就从这里开始。

\section{源域知识与目标域风险}

% auxiliary-resource-guide
本章末的“辅助学习材料”列出与核心方法对应的讲解和实践入口，可在阅读相关推导后，按所列小节对照学习。

\begin{figure}[htbp]
\centering
\begin{tikzpicture}[node distance=0.75cm and 0.7cm]
 \node[idi input, minimum width=2.0cm] (source) {源域\\$P_s(X,Y)$};
 \node[idi process, right=of source, minimum width=2.0cm] (repr) {共享表示\\$F_\theta(X)$};
 \node[idi output, right=of repr, minimum width=2.0cm] (target) {目标域\\$P_t(X,Y)$};
 \node[idi state, below=0.8cm of repr, minimum width=2.3cm] (adapt) {重加权、对齐\\或参数适配};
 \draw[idi arrow] (source)--(repr); \draw[idi arrow] (repr)--(target);
 \draw[idi arrow] (source)--(adapt); \draw[idi arrow] (target)--(adapt);
 \draw[idi feedback] (adapt)--(repr);
\end{tikzpicture}
\caption{迁移学习的核心是把源域知识送入目标域，同时检验哪些分布或任务假设仍然成立。}
\label{fig:transfer-source-target}
\end{figure}
图\ref{fig:transfer-source-target}将源域知识、适配过程和目标域评价分开；共享参数并不意味着分布相同，目标域验证必须检验迁移后是否降低了实际风险。
设源域与目标域联合分布分别为$P_s(X,Y)$与$P_t(X,Y)$，目标是降低$R_t(f)=\E_{P_t}\ell(f(X),Y)$。迁移学习允许域或任务发生变化，领域自适应通常保持预测任务并处理域差异，领域泛化则要求在训练时未见的域上工作\cite{pan2010transfer}。新相机质检可保持缺陷类别不变，跨产品迁移却可能改变标签空间；二者不能使用同一适配规则。

读到风险与密度比时，可先把期望看成按出现概率加权的平均损失；后文的两工况例子先用有限个数说明这个平均如何改变，再推广到积分。使用迁移方法前还要列明手头有什么：源标签、预训练参数、目标无标签输入和目标标签是四种不同资源，不能默认全部可用。

这里的“域”是数据产生环境及其分布，不只是文件夹名称；“任务”包括要预测什么和怎样计算错误。$X$是输入，$Y$是目标标签，$f(X)$是预测，$\ell$把单次错误变成一个数，风险$R_t$是目标环境中这些错误的期望。例如零一损失在判对时为$0$、判错时为$1$，风险就是错误率；若使用交叉熵，风险衡量的还包括预测概率的好坏，不能直接当成错误率。联合分布$P(X,Y)$同时描述输入和标签；边缘分布$P(X)$忽略标签；条件分布$P(Y\mid X)$描述给定输入后的标签规则。

预训练与迁移的分工也由此明确：前者产生一组参数，后者要证明使用这些参数或源数据确实降低了目标风险。即使无需更新参数、直接复用源模型，也是一种需要目标评价的复用方式；微调只是迁移的一种手段。

协变量偏移假定$P_s(Y\mid X)=P_t(Y\mid X)$而输入边缘分布变化；标签偏移假定$P_s(X\mid Y)=P_t(X\mid Y)$而类别比例变化；概念变化则涉及预测关系本身。实际数据可能同时存在多种变化。仅观察无标签目标输入通常无法确定标签机制是否改变，因此无监督适配不能替代目标标签审核。

\section{实例重加权与分布对齐}
\paragraph{重要性加权。}
在协变量偏移且目标分布支撑被源分布覆盖时，目标风险可写为
\begin{equation}
 R_t(f)=\E_{P_s}\left[w(X)\ell(f(X),Y)\right],\qquad
 w(x)=\frac{p_t(x)}{p_s(x)}.
\end{equation}
先用源、目标输入估计密度比，再最小化加权源域损失。核均值匹配直接使加权源域特征均值接近目标均值，避免分别估计高维密度\cite{huang2007kmm}。若用域分类器估计比值，必须校正训练时源、目标样本的先验比例。大权重会使少数样本主导梯度，截断权重可以降低方差但引入偏差。有效样本量$(\sum_iw_i)^2/\sum_iw_i^2$可用于判断加权是否退化。

\paragraph{先用两个工况手算权重。}
设输入只有低速、高速两种工况，源比例为$(0.8,0.2)$，目标比例为$(0.4,0.6)$，并假定同一工况中的预测关系不变。两组权重分别为$0.4/0.8=0.5$和$0.6/0.2=3$。若某模型两组平均损失为$0.1,0.5$，源风险为$0.8\times0.1+0.2\times0.5=0.18$，目标风险为$0.4\times0.1+0.6\times0.5=0.34$。加权源风险$0.8\times0.5\times0.1+0.2\times3\times0.5=0.34$恰好一致。权重不是把高速样本改成新样本，而是增加它们对平均损失和梯度的贡献。若源数据完全没有高速工况，分母为零，重加权就无法替代采集数据。

\paragraph{风险恒等式、密度比估计与方差。}
以离散标签和连续输入为例，协变量偏移要求对目标几乎处处的$x$，两域条件分布相同，且$p_t(x)>0$时$p_s(x)>0$。于是
\begin{align}
 R_t(f)&=\int\sum_y\ell(f(x),y)p_t(y\mid x)p_t(x)\,\mathrm dx\\
 &=\int\sum_y\ell(f(x),y)p_s(y\mid x)p_s(x)
 \frac{p_t(x)}{p_s(x)}\,\mathrm dx.\label{eq:transfer-source-risk-integral}
\end{align}
式\eqref{eq:transfer-source-risk-integral}才是可从源标签估计的量。支撑覆盖不是技术性装饰：源域从未出现的输入区域没有标签项可重用，任何有限权重都不能补上这部分风险。

令域指示$D=1$表示目标，域分类器训练混合中的先验为$\pi_t$与$\pi_s=1-\pi_t$，其校准输出为$q(x)=P(D=1\mid x)$。由贝叶斯公式
\begin{equation}
 \frac{q(x)}{1-q(x)}=\frac{\pi_t p_t(x)}{\pi_s p_s(x)},\qquad
 w(x)=\frac{\pi_s}{\pi_t}\frac{q(x)}{1-q(x)}.
\end{equation}
若训练时人为平衡两域，先验就是该平衡采样的比例，而不是原始数据库大小的比例。$q$接近$1$使估计非常不稳定；高域分类准确率此时可能表明重叠不足，而不是重加权更可靠。

对独立源样本，使用真实权重的估计$\widehat R=n_s^{-1}\sum_iw_i\ell_i$在可积时无偏；有限方差时其方差为$\Var_s(w\ell)/n_s$。若截断为$\bar w=\min(w,c)$，且$0\leq\ell\leq L_{\max}$，则
\begin{equation}
 0\leq R_t-\E_s[\bar w\ell]
 =\E_s[(w-c)_+\ell]\leq L_{\max}\E_s[(w-c)_+].
\end{equation}
这显式展示了截断的偏差来源。自归一化估计$\sum_iw_i\ell_i/\sum_iw_i$也通常具有有限样本偏差，但可改善数值稳定性。有效样本量来自归一化权重$\alpha_i=w_i/\sum_jw_j$的$1/\sum_i\alpha_i^2$：全部相等时为$n_s$，一个样本几乎占全重时接近$1$；它是权重集中的诊断量，不是所有损失下精确的独立样本数。

标签偏移下，权重改为类别先验比$p_t(y)/p_s(y)$。这个结论来自$p_t(x,y)=p_s(x\mid y)p_t(y)=p_s(x,y)p_t(y)/p_s(y)$，要求目标出现的类别在源域具有正概率。若希望从无标签目标数据估计先验，可先固定一个分类器，定义源验证集上的混淆矩阵$C_{ky}=P_s(\widehat Y=k\mid Y=y)$。标签偏移意味着该矩阵在目标域不变，目标预测类别频率$\mu_k=P_t(\widehat Y=k)$满足$\mu=C\pi_t$，其中$\pi_t\in\R^K$是目标类别先验。可在$\pi_t\geq0$、$\mathbf1^\top\pi_t=1$约束下拟合这个线性关系。矩阵列若近似相同，分类器不能区分类别，先验反演便病态；若$P(X\mid Y)$也改变，则关系本身失效。它不能直接与协变量偏移权重互换。若目标域出现源域从未包含的故障，仅靠重加权不能产生缺失的类别证据。

标签偏移可手算为另一种情形：两个真实类别上的混淆矩阵为$C=\begin{bmatrix}0.9&0.2\\0.1&0.8\end{bmatrix}$，每列对应一个真实类别，每列和为$1$。若目标预测频率为$\mu=(0.55,0.45)^\top$，写真实类别比例为$(a,1-a)^\top$，第一行给出$0.55=0.9a+0.2(1-a)$，故$a=0.5$。这个反演只有在目标域仍保留这些类条件识别率时成立，不能把预测频率直接当成真实类别频率。

\paragraph{Maximum Mean Discrepancy (MMD) 与 Correlation Alignment (CORAL)。}
深度适配网络在特征$z=F_\theta(x)$上同时优化源域监督损失和域差异\cite{long2015dan}：
\begin{equation}
 \mathcal L=\mathcal L_s+\lambda\left\|\frac1{n_s}\sum_i\varphi(z_i^s)
 -\frac1{n_t}\sum_j\varphi(z_j^t)\right\|_{\mathcal H}^2.\label{eq:transfer-mmd-objective}
\end{equation}
核技巧将平方范数展开为源源、目标目标和源目标三组核值，无需显式构造$\varphi$。训练时从两域抽取批次，计算分类与对齐损失，再共同更新编码器。核带宽决定比较的尺度，多核 MMD 组合不同带宽以减少对单一尺度的依赖。

$\varphi$把表示映射到用于比较分布的特征空间$\mathcal H$，核函数$k(u,v)$直接计算该空间中的内积，因此不必真的存储可能无限维的$\varphi(u)$。例如线性核$k(u,v)=u^\top v$就是在原特征空间比较均值；高斯核$k(u,v)=\exp(-\|u-v\|^2/(2b^2))$用带宽$b>0$控制多远的两个样本仍被视为相似。MMD 是最大均值差异，CORAL 是相关性对齐；本节给出的 CORAL 实际比较协方差矩阵，而不是逐个样本配对。

\paragraph{核均值差的实际计算。}
设$k(u,v)=\langle\varphi(u),\varphi(v)\rangle_{\mathcal H}$为正定核，$z\in\R^d$。将式\eqref{eq:transfer-mmd-objective}中的平方范数按$\|a-b\|^2=\langle a,a\rangle+\langle b,b\rangle-2\langle a,b\rangle$展开，得到
\begin{align}
 \widehat{\mathrm{MMD}}_b^2
 &=\frac1{n_s^2}\sum_{i,i'}k(z_i^s,z_{i'}^s)
 +\frac1{n_t^2}\sum_{j,j'}k(z_j^t,z_{j'}^t)\\
 &\quad-\frac2{n_sn_t}\sum_{i,j}k(z_i^s,z_j^t).
\end{align}
这是含对角项的非负、有偏估计。若两域样本分别独立同分布且$n_s,n_t>1$，去掉前两项中相同下标的配对，并将分母换为$n_s(n_s-1)$和$n_t(n_t-1)$，可得到总体平方 MMD 的无偏估计；后者有限样本下可能为负，不能误认为程序必定出错。采用特征核（characteristic kernel）时，核均值能够唯一确定其适用分布类中的分布，因此总体 MMD 为零可以识别分布相同；这里不是泛指作用于特征的任意核，线性核就只检查均值相同。无论采用哪种核，比较对象都是$P_s(Z)$与$P_t(Z)$，并没有自动检查$P_s(Y\mid Z)$与$P_t(Y\mid Z)$。

CORAL 对齐二阶统计量\cite{sun2016coral}：
\begin{equation}
 \mathcal L_{\mathrm{CORAL}}=\frac{1}{4d^2}\|C_s-C_t\|_F^2,
\end{equation}
其中$C_s,C_t$为$d$维表示的批次协方差。具体地，把样本按行写成$Z_s\in\R^{n_s\times d}$，令$H_s=I_{n_s}-\mathbf1\mathbf1^\top/n_s$，则$C_s=Z_s^\top H_sZ_s/(n_s-1)$，目标域同理。由于$C_s-C_t$对称，矩阵微分给出
\begin{equation}
 \nabla_{Z_s}\mathcal L_{\mathrm{CORAL}}
 =\frac{H_sZ_s(C_s-C_t)}{d^2(n_s-1)}.
\end{equation}
可见梯度来自中心化表示与协方差差异的乘积，源、目标均值本身不在该损失中；如果还需均值对齐，应另行添加对应项。$n_s=1$时样本协方差未定义，小批次下秩至多为$n_s-1$，也解释了协方差估计为何依赖批次组织。它计算直接，但小批次协方差不稳定，而且二阶一致不代表类别条件分布一致。MMD 与 CORAL 都可能把不同类别拉近；存在目标标签时应检查类条件结构，使用伪标签时则需控制错误对齐的累积。

一维源表示$(-1,1)$与目标表示$(-2,2)$的均值都为零，线性核 MMD 因而为零；样本方差却分别为$2$与$8$，CORAL 损失为$(2-8)^2/4=9$。相反，把$(-1,1)$整体平移为$(4,6)$后，方差不变，CORAL 为零，均值差却为$5$。这两个手算例子解释了为何均值、协方差检查互不替代，更不能替代标签语义检查。$\|A\|_F^2$就是矩阵所有元素的平方和。

\paragraph{域对抗适配。}
Domain-Adversarial Neural Network (DANN) 用域判别器$D$识别表示来自哪一个域\cite{ganin2016dann}。若域损失为交叉熵，目标写为
\begin{equation}
 \min_{F,C}\max_D\;\mathcal L_s(C(F(x_s)),y_s)
 -\lambda\mathcal L_d(D(F(x)),d).
\end{equation}
域判别器最小化自己的交叉熵，编码器通过梯度反转最大化该损失，同时维持源域分类能力。实现时不能将两个模块都沿同一符号更新。域准确率接近随机只是对齐的诊断信号，并非目标分类正确的证据；特征全部坍缩也可使域不可区分。

\paragraph{域不可区分为何不是语义正确。}
为了看清目标的符号，采用等域权重的域损失
\begin{equation}
 \mathcal L_d=-\E_s\log(1-D(Z))-\E_t\log D(Z).
\end{equation}
固定编码器，在具有特征密度且$p_s(z)+p_t(z)>0$的点逐点最小化，得到
\[
 D^*(z)=\frac{p_t(z)}{p_s(z)+p_t(z)}.
\]
两域密度都为零的点不贡献损失；仅一方密度为零时，最优值在边界，严格取值于$(0,1)$的判别器只能趋近它。下式的最小值按允许边界输出的理想判别器理解。代回可得
\begin{equation}
 \min_D\mathcal L_d=\log4-2\operatorname{JS}(P_s^Z\|P_t^Z),
\end{equation}
其中$\operatorname{JS}(P\|Q)=\tfrac12\KL(P\|M)+\tfrac12\KL(Q\|M)$，$M=(P+Q)/2$。因此在判别器充分优化这一理想条件下，编码器最大化域损失相当于减小特征边缘分布差异。若实现中再把两域损失除以$2$，常数和系数也相应减半，方向不变。

但令两域特征都均匀取$-1,+1$，源域标签满足$Y=\mathbf1[Z=1]$，目标域却满足$Y=\mathbf1[Z=-1]$。两域特征分布完全相同，任何域判别器都只能猜测，而源分类器在目标域全部判错。这个反例说明真正需要的是可共享的预测关系，而不只是看不出设备来源。二分类的零一风险可用理论形式
\begin{equation}
 R_t(h)\leq R_s(h)+\tfrac12d_{\mathcal H\Delta\mathcal H}(P_s^X,P_t^X)+\lambda^*,
 \qquad \lambda^*=\min_{h\in\mathcal H}(R_s(h)+R_t(h))
\end{equation}
理解这一点：差异量定义为$2\sup_{h,h'\in\mathcal H}|P_s(h\ne h')-P_t(h\ne h')|$，检查假设对在两域的分歧概率；$\lambda^*$则表达同一分类器同时做好两个域的可能性。这是总体风险界，有限样本估计还需要复杂度误差项；域对抗损失也不等于这个差异量的精确值。对齐可以降低某种域差异，却不能从无标签输入中证明$\lambda^*$很小。

条件对齐将类别预测纳入域判别输入，以区分同一域中的不同语义。多源适配对不同源域分别估计可信度，避免样本最多的源主导目标；部分域适配需降低源域独有类别的权重，开放集适配则必须识别目标域未知类别。这些变体改变的是标签空间假设，不能只通过调大域对齐系数实现。

\section{预训练、微调与参数高效适配}
第\ref{chap:pretraining}章讨论如何从预训练数据构造可复用表示，本节转向已有表示如何适应目标任务。
预训练提供初始化，但目标任务仍决定哪些层需要改变。线性探测冻结编码器只训练头，可检验原表示是否已具有判别性；逐层解冻先调整高层，再判断低层滤波是否需要变化；全量微调自由度最大，也更容易在少量目标样本上过拟合。比较时应固定目标标签量、预训练来源与超参数搜索预算。

Adapter 在冻结骨干中插入低维瓶颈。对输入表示$h$，一个残差适配器可写为$h'=h+W_{up}\sigma(W_{down}h)$，只训练新增参数\cite{houlsby2019adapter}。LoRA 则令权重更新$\Delta W=BA$，以秩$r$限制更新自由度\cite{hu2021lora}。二者都降低每个目标域的存储开销，但低秩不意味着任何域变化都能被表达。提示学习优化输入或中间层的提示向量，适合保持基础模型不变的条件化使用；其可迁移性仍取决于基础模型是否理解目标语义。

\paragraph{低秩更新的维度、梯度与容量。}
令冻结的线性层为$W_0\in\R^{d_{out}\times d_{in}}$，输入$x\in\R^{d_{in}}$。LoRA 使用$A\in\R^{r\times d_{in}}$、$B\in\R^{d_{out}\times r}$及尺度$s=\alpha/r$，输出为
\begin{equation}
 y=(W_0+sBA)x.
\end{equation}
可训练参数数目为$r(d_{out}+d_{in})$，更新秩不超过$r$。记$W_{eff}=W_0+sBA$，则$G=\partial\mathcal L/\partial W_{eff}\in\R^{d_{out}\times d_{in}}$，矩阵链式法则给出
\begin{equation}
 \nabla_B\mathcal L=sGA^\top,\qquad
 \nabla_A\mathcal L=sB^\top G.
\end{equation}
这解释了常见的随机初始化$A$、零初始化$B$：初始$BA=0$保证原输出不变，同时$B$可以收到非零梯度；若两者都初始化为零，则两条梯度都为零，低秩分支无法启动。假如理想全量更新$\Delta W^*$的奇异值为$\sigma_1\geq\sigma_2\geq\cdots$，最佳秩$r$矩阵近似的平方 Frobenius 误差为$\sum_{j>r}\sigma_j^2$。因此低秩适配隐含的是有用更新可被少数方向近似，而非原权重本身低秩；该矩阵逼近结论也不等于网络任务风险的误差界。

对 Adapter，若$h\in\R^d$、瓶颈宽度为$r$，则$W_{down}\in\R^{r\times d}$、$W_{up}\in\R^{d\times r}$，忽略偏置时有$2dr$个参数。它的非线性瓶颈通常不能像单个线性 LoRA 更新那样直接合并进原权重。选择应同时看可表达变化、路由和推理限制，而不只看参数比率。

可用一个小矩阵检查低秩的含义。令$d_{in}=d_{out}=2,r=1,s=1$，$A=[1\ \ 2]$、$B=[3\ \ 4]^\top$，则$BA=\begin{bmatrix}3&6\\4&8\end{bmatrix}$；对$x=(1,0)^\top$，新增输出为$(3,4)^\top$。两列彼此成比例，所以只有一个独立更新方向。这个极小例中参数数目仍为$4$，并不省参数；若原层为$100\times100$且$r=2$，才从$10000$个参数降为$400$个。实际省参数需要$r(d_{out}+d_{in})<d_{out}d_{in}$，低秩的名字本身不是节省保证。

参数量、训练显存与推理开销应分别报告。冻结权重仍可能需要存储用于反向传播的激活；多个 Adapter 需要路由规则，LoRA 合并则需要确认目标域选择是否明确。不能把参数高效方法自动视为持续学习算法：若所有域共享并反复覆盖同一个低秩更新，仍会遗忘旧域。

\section{测试时适配与负迁移}
测试时适配只使用部署输入调整有限参数。Tent 最小化预测熵，并在其设定下更新归一化层仿射参数\cite{wang2021tent}：
\begin{equation}
 \min_{\eta}\frac1{|B|}\sum_{x\in B}
 \left[-\sum_kp_\eta(k\mid x)\log p_\eta(k\mid x)\right].
\end{equation}
先保留训练好的模型，按部署批次计算熵，再对允许更新的参数执行梯度步。低熵可能是正确判断，也可能是错误自信。类别极不平衡、批次很小或流中含未知故障时，熵最小化可能将预测推向单一类别。逐批重置与持续累积更新是不同规则，必须明确是否保留适配状态。

用二分类数值理解熵：$p=(0.5,0.5)$时$H=\log2\approx0.693$，$p=(0.9,0.1)$时$H\approx0.325$，确定预测时熵趋于零。若真实类别恰是第二类，后一个预测虽然熵更低，却错得更自信。因此无标签目标衡量的是预测集中程度，不是准确率。

部署评估还要规定信息到达顺序。严格在线的“先预测、后更新”先用当前模型记录本批结果，再利用本批无标签输入更新供后续批次使用；“先适配、后预测”允许本批输入参与自身预测前的更新，是另一种协议，不能混报。两种协议都不得访问未来批次或测试标签；重置频率、步数和学习率应由先前验证数据确定，而不是查看测试准确率后调整。

\paragraph{熵下降为什么会强化错误。}
令类别 logits 为$a\in\R^K$、$p=\softmax(a)$。由$\partial p_k/\partial a_j=p_k(\mathbf1[k=j]-p_j)$，可逐项得到
\begin{equation}
 \frac{\partial H(p)}{\partial a_j}=-p_j(\log p_j+H(p)).
\end{equation}
二分类时令$p=\sigma(a)$，则$\mathrm dH/\mathrm da=p(1-p)\log((1-p)/p)$。若$p>1/2$，梯度下降会继续增大$a$和$p$；若$p<1/2$则进一步压低$p$。没有标签参与时，这个更新没有能力识别当前偏向是否正确，只会强化已有偏向。熵适配因而依赖初始预测大体可信、低密度区域适合作决策边界等额外条件。应限制步数和漂移幅度，保留重置机制，并使用独立监测数据检验校准与少数类召回；这些操作降低风险，但仍不能构成无标签正确性保证。

负迁移指相对于合适的无迁移基线，使用源知识反而提高目标风险。应比较从零训练、线性探测、全量微调和轻量适配，并绘制随目标标签量变化的学习曲线。领域自适应理论中的分布差异项之外，还存在源、目标任务能否被同一假设同时解释的误差项\cite{bendavid2010theory}；只减小域差异并不足够。

换了新相机后，先查清楚是亮度、镜头和噪声变了，还是“什么算缺陷”的定义变了。前一类问题可以先试归一化适配和线性探测；后一类问题需要重新核对标签。若规则允许，可以用新相机的无标签图像做适配，但测试标签只能留到最终评价。报告各类缺陷召回、未知输入拒识、校准、训练成本和发生负迁移的设备组。若还要长期维护旧相机能力，就需要下一章的持续学习机制。

迁移学习中的分布偏移、对齐与参数适配可用以下材料复习，分别覆盖数据变换和模型复用。
\section*{辅助学习材料}
\begingroup
\emergencystretch=3em
\begin{itemize}
\item \textbf{Transfer Learning}；作者/平台：Stanford CS231n；英文；课程笔记。阅读固定特征提取、微调以及数据规模和领域相似性对应的选择建议，理解迁移范围与过拟合风险。\href{https://cs231n.github.io/transfer-learning/}{在线阅读}。
\item \textbf{Transfer Learning for Computer Vision Tutorial}；作者/平台：Sasank Chilamkurthy / PyTorch；英文；代码教程。比较微调整个网络与冻结特征提取器两种训练流程，核对训练参数范围和验证步骤；该教程不是 LoRA 实现。\href{https://docs.pytorch.org/tutorials/beginner/transfer\_learning\_tutorial.html}{在线阅读}。
\item \textbf{什么是迁移学习（Transfer Learning）？}；知乎；中文解说。比较冻结特征与微调，记录可训练参数；标题与主题据必应索引，原页受限。\href{https://zhuanlan.zhihu.com/p/1899236885408256815}{在线阅读}。
\item \textbf{迁移学习(Transfer Learning)概述及代码实现}；CSDN；中文博客。对照源域与目标域划分，用无迁移基线识别负迁移；标题与主题据必应索引，原页受限。\href{https://blog.csdn.net/weixin\_49272172/article/details/115791703}{在线阅读}。
\item \textbf{第十一节 2021 - 概述领域自适应 (Domain Adaptation)}；李宏毅；bilibili；中文课程。选看第110集，对照本章分布差异与适应条件。2026年10月4日官方公开元数据显示整套课程累计播放约299.5万，并非本分集播放量；已核对目录与计数，未逐帧观看。\href{https://www.bilibili.com/video/BV1Wv411h7kN/?p=110}{观看分集}。
\end{itemize}
\endgroup
\paragraph{本章小结。}
迁移学习依次回答源知识是否相关、分布差异能否被估计、哪些参数应更新。重加权改变样本贡献，成立前提是正确的偏移假设和支撑覆盖；分布对齐改变表示，却必须额外保护类别语义；微调和低秩适配改变预测器，其容量与标签预算应匹配。测试时熵下降进一步说明，没有标签的自洽改进可能只是错误自信。所有路线都需要目标域证据与无迁移基线，不能以对齐损失下降代替目标风险下降。

实际选方法时，先核对新旧数据的标签含义是否相同，再看目标设备有多少标签、能否使用它的无标签记录，最后选择足够解决问题的适配办法。如果系统仍需服务旧相机，更新后就要一起复测；新相机得分提高、旧相机却明显退步，不能算整个更新任务已经完成。下一章据此讨论持续学习：数据不断到来时，怎样学习新情况、保留旧能力，并控制存储和计算开销。

\chapter{持续学习}
\label{chap:continual}
\begin{quote}\small
\textit{本章摘要。}持续学习把更新过程放到时间轴上，同时约束新知识获取、旧知识保持和资源增长。本章从增量设置出发，推导回放、梯度约束、参数正则和知识蒸馏，再比较隔离与扩展方法。评价对象是整个学习轨迹，而非最后一个任务。
\end{quote}

上一章研究如何借助源知识降低一个目标域的风险。实际系统却往往没有永久不变的“最终目标域”：新相机、新工况和新故障按时间到达，上一轮的目标域很快又成为需要保护的历史域。如果每次适应都只考核最新数据，系统可能一边发布更好的新版本，一边失去识别旧故障的能力。本章将学习对象从一个最终模型改为一串带资源约束的更新过程；因此，读者要问的不是“最后一次训练是否收敛”，而是“沿时间轴是否持续服务所有仍需服务的分布”。

持续学习并非要求所有参数永远不变，而是要求区分有益的知识修正与不必要的历史破坏。旧数据能否保留、推理时是否知道任务身份、类别是否持续增加，会直接改变可用算法和评价标准。本章先说明回放如何近似历史风险，再从损失的局部变化推出梯度约束，从后验近似推出参数重要性惩罚，并解释蒸馏和参数隔离分别保留了什么。它们共同处理稳定性与可塑性的张力，但保护对象并不相同：样本风险、局部方向、参数邻域和可调用模块不能互相替代。

每次更新以后，都要重新测一遍仍需服务的旧设备，记录哪些故障开始漏检。只看最后一个阶段的总分，会看不到中途丢失的能力。另一个实际困难是，新到的记录往往还没有标签，专家不可能逐条检查。下一章的主动学习讨论怎样挑选优先标注的记录，再把得到的标签交给更新算法。

\section{数据流、任务边界与遗忘}

% auxiliary-resource-guide
本章末的“辅助学习材料”列出与核心方法对应的讲解和实践入口，可在阅读相关推导后，按所列小节对照学习。

\begin{figure}[htbp]
\centering
\begin{tikzpicture}[node distance=0.7cm and 0.65cm]
 \node[idi input, minimum width=1.7cm] (d1) {$\mathcal D_1$};
 \node[idi input, right=of d1, minimum width=1.7cm] (d2) {$\mathcal D_2$};
 \node[idi input, right=of d2, minimum width=1.7cm] (d3) {$\mathcal D_3$};
 \node[idi state, below=0.8cm of d2, minimum width=2.1cm] (model) {共享模型\\更新参数};
 \node[idi output, below=0.8cm of d3, minimum width=2.1cm] (old) {旧任务性能\\可能下降};
 \draw[idi arrow] (d1)--(model); \draw[idi arrow] (d2)--(model); \draw[idi arrow] (d3)--(model);
 \draw[idi arrow] (model)--(old);
 \draw[idi feedback] (old.south) .. controls +(0,-0.6) and +(0,-0.6) .. node[below,font=\scriptsize]{回放、蒸馏或参数隔离} (model.south);
\end{tikzpicture}
\caption{持续学习沿数据流更新模型；新任务的可塑性与旧任务的稳定性形成直接张力。}
\label{fig:continual-stability-plasticity}
\end{figure}
图\ref{fig:continual-stability-plasticity}沿任务流展示更新过程，并将旧知识保持与新任务适应并列；评价更新效果时须同时检查历史任务和当前任务，不能只看新数据损失。
设数据依次到达$\mathcal D_1,\ldots,\mathcal D_T$，学习第$t$阶段时通常无法完整访问历史数据。直接在新数据上继续优化会改变共享参数，旧分布风险可能增加，这就是灾难性遗忘的典型表现。持续学习的核心矛盾是可塑性与稳定性，而非单纯限制更新幅度\cite{delange2022continual}。

任务增量通常在推理时提供任务身份，可选择相应输出头；类增量要求在已出现的所有类别中统一预测，测试时不提供任务身份；域增量通常保持标签空间而改变输入分布。类增量训练可以有已知阶段边界，不能把它与无边界数据流混为一谈。任务无关的在线持续学习进一步要求算法自行处理连续到达的数据，标签还可能延迟。

\paragraph{先分清一条样本、一个任务和一个阶段。}
例如一台电机连续记录振动，把一段互不重叠的时间窗提取成$x\in\R^d$，复检给出标签$y\in\{1,\ldots,K\}$，这一对$(x,y)$才是一条有标签样本。设备是产生记录的物理对象；任务是规定输入、标签含义和评价分布的预测问题；阶段$t$是允许更新的时间区间。一个阶段可以有多台设备，一台设备也可以经历多个阶段，三种编号不能混用。本章参数写作$\theta\in\R^p$，模型输出$K$类概率，损失$\ell_\theta(x,y)$是一个标量，例如真实类别概率为$0.8$时交叉熵为$-\log0.8$。梯度是损失对$p$个参数的偏导数组成的向量，更新$\theta\leftarrow\theta-\eta\nabla_\theta\ell$才是实际改变模型的步骤。

\paragraph{为什么学新任务会影响旧任务。}
用只有一个参数的例子，旧损失为$(\theta-1)^2/2$、新损失为$(\theta+1)^2/2$。旧任务训练到$\theta=1$后，新损失梯度为$2$；步长$\eta=1/2$使参数变为$0$，新损失从$2$降为$1/2$，旧损失却从$0$升为$1/2$。两个任务争用同一个参数，因而只优化当前损失不足以保留历史能力。真实网络不一定如此直接冲突，但也没有自动保存旧功能的机制。若标签定义本身被纠正，则旧测试标准也应按预定规则重审，不能把坚持错误旧标签当作成功保持。

\section{经验回放与梯度约束}
\paragraph{有限缓存如何近似历史风险。}
经验回放保存大小为$M$的缓存$\mathcal B$，每次从当前数据与缓存分别采样：
\begin{equation}
 \mathcal L_t=\frac1{|B_t|}\sum_{(x,y)\in B_t}\ell_\theta(x,y)
 +\lambda\frac1{|B_m|}\sum_{(x,y)\in B_m}\ell_\theta(x,y).\label{eq:continual-replay-objective}
\end{equation}
先混合计算梯度并更新模型，再按既定策略更新缓存。蓄水池采样在观察到第$n>M$个样本时，以$M/n$概率将其放入缓存，并均匀替换一个槽位，使历史每个样本有相同保留概率。类别均衡缓存则改变采样目标，保障少数类但不再保持原始频率。iCaRL 使用代表样本与蒸馏维护类增量表示，并以类均值进行预测\cite{rebuffi2017icarl}。

\paragraph{把回放公式执行一次。}
$\mathcal B$是跨轮保存的缓存，$B_m\subseteq\mathcal B$是本次抽取的旧样本小批，$B_t$是当前小批；这里的竖线表示集合大小（样本数），不是向量范数。若当前两条损失为$0.2,0.6$，回放两条为$0.1,0.3$，且$\lambda=2$，目标值为$0.4+2\times0.2=0.8$。更新时对这个加权和求导，不是把损失值直接从参数上减去。第一次还没有旧缓存时只计算当前项，不能除以$|B_m|=0$。容量$M=2$、已见两条记录时，第三条以$2/3$概率入选，入选后在两个槽位中各以$1/2$概率替换；所以每条旧记录被替换的概率为$1/3$。这说明随机替换与每次总删最旧记录是不同算法。

\paragraph{缓存无偏性与风险权重不是同一件事。}
设此前共观察到$n$个样本且缓存容量为$1\leq M\leq n$；$M=0$表示停用回放，不计算缓存均值。蓄水池采样的等概率性质可由归纳法证明：第$n$个新样本的入选概率是$M/n$；任一旧样本在$n-1$时刻位于缓存的概率为$M/(n-1)$，本轮被替换的条件概率为$(M/n)(1/M)=1/n$。所以它本轮后仍在缓存的概率为
\begin{equation}
 \frac{M}{n-1}\left(1-\frac1n\right)=\frac Mn.
\end{equation}
初始$n=M$时所有样本都被保留，归纳因此成立。对固定参数$\theta$，缓存平均损失对采样随机性取期望等于历史经验平均损失：
\begin{equation}
 \E_{\mathcal B}\left[\frac1M\sum_{i\in\mathcal B}\ell_\theta(x_i,y_i)\right]
 =\frac1M\sum_{i=1}^n\frac Mn\ell_\theta(x_i,y_i)
 =\frac1n\sum_{i=1}^n\ell_\theta(x_i,y_i).
\end{equation}
这是固定函数的抽样性质，不表示在同一缓存上反复优化后，所选模型的缓存分数仍是无偏泛化估计；模型会适应缓存，所以还需要固定的历史测试集。

此外，式\eqref{eq:continual-replay-objective}的训练目标除以$1+\lambda$后对应新数据权重$1/(1+\lambda)$与历史权重$\lambda/(1+\lambda)$。若真要拟合所有已见样本的等权经验风险，当前阶段有$n_t$个样本、历史有$n_{old}$个样本时，应取$\lambda=n_{old}/n_t$，并确保两侧批次各自代表相应经验分布。整体除以$1+\lambda$不改变极小点，却会把梯度缩小同样倍数；对于无额外梯度变换的普通梯度下降，归一化后把学习率相应放大才能保持同一步更新；带自适应缩放、梯度裁剪或独立权重衰减的优化器不能直接套用这一等价关系。固定$\lambda$实际是在规定新旧风险的政策权重，而不是自动复现全历史训练。对逐渐变化的设备，近期分布优先也许更合理，但必须说明这是主动选择的目标。

缓存按样本数公平并不等于按字节公平：原图、压缩特征与 logits 占用不同。特征回放节省空间，却依赖编码器版本兼容；生成回放用生成器代替原样本，但生成器本身也会遗忘，偏差可能逐阶段累积\cite{shin2017replay}。保存合成样本也不能自动免除隐私审计。

\paragraph{梯度情景记忆（Gradient Episodic Memory，GEM）与 A-GEM。}
GEM 将当前梯度$g$投影为与旧任务梯度$g_k$不冲突的方向\cite{lopezpaz2017gem}：
\begin{equation}
 \min_{\tilde g}\frac12\|\tilde g-g\|_2^2,
 \qquad \tilde g^{\mathsf T}g_k\geq0\quad\text{对所有旧任务 }k.
\end{equation}
因为梯度下降的旧损失一阶变化约为$-\eta g_k^{\mathsf T}\tilde g$，该约束限制局部增损，而非保证有限步后的旧风险不升。A-GEM 用回放批次的平均参考梯度$g_r$替代全部约束\cite{chaudhry2019agem}。若$g^{\mathsf T}g_r<0$且$\|g_r\|>0$，投影为
\begin{equation}
 \tilde g=g-\frac{g^{\mathsf T}g_r}{\|g_r\|_2^2}g_r;
\end{equation}
否则保留原梯度。它降低计算开销，但只约束平均回放方向，少数旧任务仍可能受损。

\paragraph{从约束优化推出投影及其步长限制。}
令参数维度为$p$，$g,g_r,\tilde g\in\R^p$。单个约束的拉格朗日函数为
\begin{equation}
 \mathcal J(\tilde g,\alpha)=\tfrac12\|\tilde g-g\|^2-\alpha g_r^\top\tilde g,
 \qquad\alpha\geq0.
\end{equation}
驻点条件给出$\tilde g=g+\alpha g_r$。当$g^\top g_r\geq0$时约束已满足，$\alpha=0$；否则最优点位于边界$g_r^\top\tilde g=0$，故$\alpha=-g_r^\top g/\|g_r\|^2$，代回即得 A-GEM 公式。若$g_r=0$，该一阶约束没有信息，不应执行除法。对于$K_{\mathrm{old}}$个旧任务，把梯度写成矩阵$G=[g_1,\ldots,g_{K_{\mathrm{old}}}]\in\R^{p\times K_{\mathrm{old}}}$；这里旧任务数不同于类别数$K$。此时$\alpha\in\R^{K_{\mathrm{old}}}$为逐任务的非负乘子向量，同样有$\tilde g=g+G\alpha$，对偶问题成为
\begin{equation}
 \min_{\alpha\geq0}\ \tfrac12\alpha^\top G^\top G\alpha+\alpha^\top G^\top g.
\end{equation}
这把高维参数空间中的投影转换为以旧任务数为变量数的二次规划，但需要计算和保存参考梯度。

即使投影精确，有限步长仍可能增加旧损失。若旧损失梯度为$L_k$-Lipschitz，则
\begin{equation}
 \mathcal L_k(\theta-\eta\tilde g)
 \leq\mathcal L_k(\theta)-\eta g_k^\top\tilde g
 +\tfrac12L_k\eta^2\|\tilde g\|^2.
\end{equation}
约束仅保证一次项不为正；当内积恰好为零时，二次项仍允许增损。当内积严格为正且梯度确实来自该损失时，$\eta\leq2g_k^\top\tilde g/(L_k\|\tilde g\|^2)$才给出这个上界下的非增条件。实际使用的是小缓存梯度，既有抽样误差又有曲率误差，因此仍须检查历史风险。若优化器再对投影梯度施加动量或自适应预条件，实际更新方向也可能不再满足约束；约束应针对实际步向量解释。

\paragraph{手算一次梯度投影。}
取$g=(-1,2)^{\mathsf T}$、$g_r=(1,0)^{\mathsf T}$，内积为$-1$。原始下降步为$-\eta g=(\eta,-2\eta)^{\mathsf T}$，旧损失的一阶变化为$+\eta$，所以发生冲突。投影后$\tilde g=(-1,2)^{\mathsf T}-(-1)(1,0)^{\mathsf T}=(0,2)^{\mathsf T}$，内积变为零，只去掉冲突分量。注意算法投影的是梯度，执行时仍要取负号；把约束误写成梯度内积小于零，会恰好鼓励旧损失上升。

\section{参数重要性与知识蒸馏}
\paragraph{Elastic Weight Consolidation (EWC) 及其在线变体。}
EWC 将旧任务后验在最优参数附近近似为对角二次项\cite{kirkpatrick2017ewc}：
\begin{equation}
 \mathcal L_t(\theta)=\mathcal L_{new}(\theta)
 +\frac\lambda2\sum_jF_j(\theta_j-\theta_j^*)^2,
 \quad
 F_j\approx\frac1N\sum_n[\partial_{\theta_j}\log p_{\theta^*}(y_n\mid x_n)]^2.
\end{equation}
先训练旧阶段、保存$\theta^*$并估计$F$，再在新阶段加入惩罚。经验 Fisher 使用观测标签，和对模型预测分布取期望的 Fisher 并不完全相同。对角近似忽略参数耦合，旧任务局部曲率也未必能描述大范围更新风险。在线变体将历史重要性按衰减系数累积为一个统计量，控制内存，但衰减速度决定更早知识被保留多少。

\paragraph{从贝叶斯递推到对角二次保护。}
若给定参数后各阶段数据条件独立，后验满足
\begin{equation}
 p(\theta\mid\mathcal D_{1:t})\propto p(\mathcal D_t\mid\theta)
 p(\theta\mid\mathcal D_{1:t-1}).
\end{equation}
对旧后验的负对数$V(\theta)$在局部极小点$\theta^*$二阶展开，若梯度为零，有
\begin{equation}
 V(\theta)\approx V(\theta^*)+\tfrac12(\theta-\theta^*)^\top H(\theta-\theta^*),
 \qquad H=\nabla^2V(\theta^*).
\end{equation}
这里$H\in\R^{p\times p}$，完整存储需要$O(p^2)$。EWC 用似然 Fisher 的对角近似作为局部精度信息，忽略参数间相关性后只需$O(p)$存储。模型 Fisher 的定义为
\begin{equation}
 F=\E_{X}\E_{Y\sim p_{\theta^*}(\cdot\mid X)}[s(X,Y)s(X,Y)^\top],
 \qquad s=\nabla_\theta\log p_{\theta^*}(Y\mid X).
\end{equation}
在可交换微分和积分等正则条件下，模型分布内的 Fisher 等于负对数似然 Hessian 的期望；把模型标签换为观测标签得到的是经验 Fisher，不必等于实际损失 Hessian。若新损失采用样本平均而旧后验来自样本似然之和，还要计入样本量尺度；实践中的$\lambda$同时承担尺度转换与保持强度调节，不能把任意$\lambda$都当作精确贝叶斯更新。

二次惩罚对第$j$个参数的梯度为$\lambda F_j(\theta_j-\theta_j^*)$。在单个坐标上，若新目标局部为$\tfrac12a(\theta_j-b)^2$，$a>0$，则合并后的极小点为
\begin{equation}
 \theta_j^{new}=\frac{ab+\lambda F_j\theta_j^*}{a+\lambda F_j}.
\end{equation}
这明确展示了折中：新证据曲率$a$越强，参数越靠近$b$；旧重要性越强，越靠近$\theta_j^*$。但两个参数可能需要协同移动才能不改变旧函数，对角惩罚看不到这种安全方向。在线累积$\bar F_t=\rho\bar F_{t-1}+F_t$还应同时说明锚点如何重置；把多阶段局部曲率压缩到最新锚点，是进一步近似，不等于保留所有历史后验。

Synaptic Intelligence 沿训练轨迹累计参数更新对损失下降的贡献，再按总位移归一化估计重要性\cite{zenke2017si}。例如令一次参数增量为$\Delta\theta_j^{(s)}$、更新前梯度为$g_j^{(s)}$，可累计$\omega_j=\sum_s-g_j^{(s)}\Delta\theta_j^{(s)}$，再以$\Omega_j=\max(0,\omega_j)/((\theta_j^{end}-\theta_j^{start})^2+\epsilon)$作为非负重要性的一种实现约定。分子来自一阶损失下降的逐坐标贡献，分母避免仅因坐标位移较大就给出高重要性；$\epsilon>0$防止小位移导致数值爆炸。噪声、动量和坐标间耦合会使逐坐标贡献不稳定，因此需记录截断与累积约定。它与 EWC 的区别是使用优化路径而非阶段末的 Fisher。两者都可能在任务冲突较强时过度限制可塑性，因此要同时报告新任务学习不足与旧任务遗忘。

\paragraph{参数重要性可以怎样直观核算。}
在上面的一维折中式中，取$a=1,b=-1,\theta_j^*=1,\lambda F_j=3$，新最优参数为$(-1+3)/4=1/2$，处于新旧偏好之间。$F_j$大不是说参数数值大，而是旧预测对该坐标变化敏感；即使$\theta_j^*=0$也可能很重要。$\lambda F_j=0$时完全服从新目标，趋于无穷时近乎冻结，这两端分别可能遗忘旧任务和学不会新任务。Synaptic Intelligence 中若梯度为$2$、本次位移为$-0.1$，一阶下降贡献为$-2(-0.1)=0.2$；它记录的是沿路径的贡献，与只在末点计算的 Fisher 不是同一统计量。

\paragraph{输出蒸馏与特征保持。}
Learning without Forgetting 保存旧模型输出，在新数据上约束新模型\cite{li2017lwf}：
\begin{equation}
 \mathcal L=\mathcal L_{new}+\lambda\tau^2
 \E_x\KL(p_{old}^{(\tau)}(\cdot\mid x)\|p_\theta^{(\tau)}(\cdot\mid x)).
\end{equation}
温度$\tau$软化类别概率，旧教师应停止梯度更新。蒸馏只能约束所观察输入上的行为；若新输入不覆盖旧分布，保持这些输出不保证保持旧任务。类增量场景还需统一输出空间并处理新旧类别的偏置。特征蒸馏可约束中间表示，但应明确对齐层、距离和容量，否则会阻碍必要的表示变化。

\paragraph{蒸馏温度、梯度缩放与类增量缺口。}
令旧教师与新学生在相同$K$个旧类别上的 logits 分别为$u,v\in\R^K$，$p_k=\softmax(u/\tau)_k$、$q_k=\softmax(v/\tau)_k$。停止教师梯度后，$\KL(p\|q)=\sum_kp_k\log p_k-\sum_kp_k\log q_k$中的第一项为常数，所以
\begin{equation}
 \frac{\partial\KL(p\|q)}{\partial v_j}=\frac{q_j-p_j}{\tau},\qquad
 \frac{\partial[\tau^2\KL(p\|q)]}{\partial v_j}=\tau(q_j-p_j).
\end{equation}
较高温度下，令$\bar u=K^{-1}\sum_ku_k$、$\bar v$同理，有$p_k\approx K^{-1}+(u_k-\bar u)/(K\tau)$，故$q_k-p_k$本身约为$1/\tau$量级。乘$\tau^2$使高温近似下的梯度尺度不至于随温度平方衰减，主要约束中心化 logits 的相对差异；它不是对所有温度都精确保持梯度范数。

类增量时必须说明 softmax 的归一化范围。如果只在旧类别上归一化，蒸馏保存旧类之间的相对关系，却不控制它们与新类之间的整体偏移；如果在全部类别归一化并把教师对新类概率设为零，就会额外压低新类概率，可能与新任务监督冲突。回放样本、分类校准和权重调节可以缓解，但不能省略这个输出空间选择。

特征保持也有明确几何假设：若教师特征$h_o\in\R^{d_o}$、学生特征$h_n\in\R^{d_n}$维度不同，应通过映射$P\in\R^{d_o\times d_n}$再比较$\|Ph_n-h_o\|^2$，或比较关系结构；直接相减没有定义。可训练的$P$可能吸收一部分变化，过强的逐点保持则会限制新任务需要的表示重组。

\paragraph{蒸馏保留的是答案分布，不是参数副本。}
logit 指 softmax 之前的实数得分，并非概率；对$u=(2,0)$，温度$1$下概率约为$(0.881,0.119)$，温度$2$下约为$(0.731,0.269)$，次要类别更容易被看见。实际一步先用冻结教师对同一批输入算概率，再用学生算概率，组合监督损失与蒸馏损失后仅更新学生。这里$\KL(p\|q)=\sum_kp_k\log(p_k/q_k)$衡量两分布差异，不是两个类别标签的距离。教师若把新故障误判为正常，照搬其输出也会保留错误，因此蒸馏必须与新标签监督共同权衡。

\section{参数隔离、扩展与混合策略}
Progressive Networks 为新任务增加网络列，并通过横向连接复用旧表示，旧参数冻结以避免直接覆盖\cite{rusu2016progressive}。代价是参数随任务增长，并依赖任务路由。掩码方法为不同任务选择不同参数子集，容量耗尽后仍需要压缩或扩展。Adapter 和 LoRA 可以成为隔离模块，但只有旧模块被冻结且路由可靠时，才提供对应的保持机制。

冻结模块只有在它实际被调用时才保护输出。设路由器为$r(x)$、旧任务$i$的固定专家为$f_i$，最终预测为$f_{r(x)}(x)$。对零一损失，记错误路由概率为$\epsilon_r$，旧专家单独使用时风险为$R_i$，则由事件的并集界有$R_{system,i}\leq\epsilon_r+R_i$。即使$R_i$完全不变，路由错误也足以损坏系统性能；若另有共享归一化统计或公共前处理被更新，冻结专家权重甚至不保证$f_i$这一函数不变。增加$T$个每个含$P_0$参数的专家通常需$O(TP_0)$存储，隔离是在保持与资源之间交换，不是免费消除遗忘。

回放、蒸馏和隔离并不互斥。可用小缓存维持真实旧分布，用蒸馏保留相对类别信息，再通过新模块吸收变化。比较混合算法时必须同时限制缓存、模型参数、教师副本与训练计算，避免用更大资源获得的收益被误认为算法优势。

\paragraph{如何选择并调用被隔离的模块。}
输出头是把共享特征变成类别得分的最后一段网络；掩码是与参数同形状的零一开关，决定某任务使用哪些参数；路由器是决定调用哪个头或专家的规则。若生产系统已提供相机型号，可以按该元数据路由；若测试不提供任务身份，就必须从输入估计路由，不能用真实标签代选正确专家。Adapter 是插入主干的小型可训练模块，LoRA 以低秩矩阵乘积表示权重增量；二者只是节约任务专用参数的方式，仍须冻结对应旧路径以及其前处理。选择混合策略时应先问能否存旧数据、能否知道任务身份、最多能增加多少参数，再确定保护对象，而不是默认把所有方法相加就最好。

\section{评价矩阵与工业更新规则}
定义$A_{t,i}$为学习完阶段$t$后在任务$i$固定测试集上的性能，行对应学习时间，列对应测试任务。最终平均性能和平均遗忘为
\begin{equation}
 \mathrm{ACC}=\frac1T\sum_{i=1}^TA_{T,i},\qquad
 F=\frac1{T-1}\sum_{i=1}^{T-1}
 \left(\max_{i\leq t<T}A_{t,i}-A_{T,i}\right).
\end{equation}
遗忘量可以为负，表示最终性能超过此前最佳值；若采用截断为零的定义，应显式声明。后向迁移为
\begin{equation}
 \mathrm{BWT}=\frac1{T-1}\sum_{i=1}^{T-1}(A_{T,i}-A_{i,i}).
\end{equation}
前向迁移可比较学习任务$i$之前的$A_{i-1,i}$与未迁移基线，但仅在输出空间与任务可访问性允许该测试时有意义\cite{lopezpaz2017gem}。类增量应同时测试所有已见类别，不能利用真实任务身份缩小候选集合。

这些指标中的$A$默认越大越好，例如准确率或召回；若填入损失，需要相应反转差分方向。遗忘与后向迁移并不相反等价：若某旧任务初学得分为$0.70$，中期最高为$0.90$，最终为$0.80$，它的后向迁移为$+0.10$，但从历史最佳计算的遗忘仍为$0.10$。数值仅为解释定义的教学例。只报告最终平均性能可能同时掩盖某旧域崩溃与某新域改善，因此应保留完整矩阵、最差子群指标和新阶段初学性能；没有学会旧任务的模型也可能表现为“几乎不遗忘”。当$T=1$时，含$T-1$分母的指标不定义。

\paragraph{从完整矩阵读出最终指标。}
例如三阶段的下三角成绩可逐行写为
\begin{align*}
 (A_{1,1})&=(0.8),\\
 (A_{2,1},A_{2,2})&=(0.7,0.9),\\
 (A_{3,1},A_{3,2},A_{3,3})&=(0.6,0.85,0.9).
\end{align*}
未学任务的格子留空而非填零。最终平均性能为
\begin{equation*}
 \mathrm{ACC}=\frac{0.6+0.85+0.9}{3}\approx0.783.
\end{equation*}
前两任务遗忘分别为$0.8-0.6=0.2$和$0.9-0.85=0.05$，平均为$0.125$，后向迁移为$-0.125$。这里两者恰好互为相反数，是因为旧任务历史最佳都在初学时取得，不是一般规律。$0.2$表示二十个百分点的绝对下降，不是相对下降百分之二十。

新相机按周接入时，应提前为每台相机固定测试集，每次更新后都在这些集合上填写评价矩阵。缓存不能混入测试样本，超参数也不能看到未来阶段的结果后反复修改。线上流程应先用当前模型预测，等工单或复检标签到达后再学习，并记录标签延迟、漂移响应时间、每阶段训练时长和累计存储。旧相机或安全关键缺陷一旦超过风险阈值，应回滚到已验收版本，而不是自动覆盖生产模型。

\paragraph{自测与答案。}
若某旧任务初学、历史最佳和最终准确率依次为$0.70,0.90,0.85$，其遗忘和后向迁移各是多少？答案：遗忘为$0.90-0.85=0.05$，后向迁移为$0.85-0.70=0.15$。二者可以同时为正，因为参照时刻不同。若只发生一个阶段，则本章的平均遗忘和后向迁移不定义，不能把分母为零的结果记成零。

以下教程与讲解可用于复习持续学习设置，并组织回放、模块保存和遗忘评价实验。
\section*{辅助学习材料}
\begingroup
\emergencystretch=3em
\begin{itemize}
\item \textbf{Avalanche: an End-to-End Library for Continual Learning}；作者/平台：ContinualAI；英文；开源框架教程。重点阅读 scenarios、benchmarks、replay plugins 和 continual evaluation，帮助把本章的任务流、缓存、类增量与遗忘矩阵转成可重复实验。\href{https://avalanche.continualai.org/}{在线文档}。
\item \textbf{Saving and Loading Models}；作者/平台：PyTorch 官方教程；英文；代码教程。重点看 checkpoint、optimizer state 与模型版本保存，帮助组织本章的回放、蒸馏和参数隔离对照。\href{https://docs.pytorch.org/tutorials/beginner/saving\_loading\_models.html}{在线阅读}。
\item \textbf{持续学习/终身学习/增量学习基础知识}；知乎；中文解说。区分任务增量、类增量与数据流设定；标题与主题据必应索引，原页受限。\href{https://zhuanlan.zhihu.com/p/696297033}{在线阅读}。
\item \textbf{持续学习（Continual Learning）：让AI像人类一样终身成长}；CSDN；中文博客。比较回放与参数约束，同时检查旧任务遗忘和新任务学习；标题与主题据必应索引，原页受限。\href{https://blog.csdn.net/daqianai/article/details/152846877}{在线阅读}。
\item \textbf{2021 - 机器終身学习 (二) - 灾难性遗忘(Catastrophic Forgetting)}；李宏毅；bilibili；中文课程。选看第132集，对照本章序贯训练中的遗忘问题。2026年10月4日官方公开元数据显示整套课程累计播放约299.5万，并非本分集播放量；已核对目录与计数，未逐帧观看。\href{https://www.bilibili.com/video/BV1Wv411h7kN/?p=132}{观看分集}。
\end{itemize}
\endgroup
\paragraph{本章小结。}
回放重建历史风险，梯度约束限制局部干扰，参数正则保存重要方向，蒸馏保存所观察输入上的函数行为，隔离减少共享参数冲突。这些保护机制各有缺口：缓存未必覆盖稀有旧故障，一阶约束不控制有限步长曲率，对角重要性忽略参数耦合，新输入上的蒸馏不保证旧输入行为，冻结模块还依赖正确路由。没有一种方法能够脱离任务身份、历史访问和资源预算单独比较。

持续学习要逐阶段记录新旧任务的成绩。先列出仍需识别的设备和故障，再选择保留历史样本、约束重要参数，或保存独立模块。同时检查新任务是否学会：完全不更新当然不会因更新而遗忘，但也可能一直学不会新故障。现场还会受标注速度限制。回放与正则化解决拿到标签以后怎样更新，下一章的主动学习则决定有限的专家时间先花在哪些记录上。

\chapter{主动学习}
\label{chap:active}
\begin{quote}\small
\textit{本章摘要。}主动学习在有限标注预算下选择最有价值的查询。本章从候选池与标注者的交互规则出发，比较不确定性、委员会分歧、预期模型变化和批量多样性，再讨论成本、噪声与流式查询。核心不是找到最难的样本，而是使下一次标注改善部署风险。
\end{quote}

上一章讨论在新数据到来后如何更新而不遗忘，但这通常默认训练所需标签能够取得。工业现场中，更稀缺的资源可能是能够判断裂纹性质、核对工单因果或识别异常波形的专家时间。面对成千上万条未标注记录，更新算法之前还有一个决策：先看哪一条，为什么值得看，得到答案后能够改变什么？本章把这种决策本身纳入学习过程：标注预算是有限资源，查询规则的目标是降低部署风险，而不是让训练集看起来更大。

主动学习不是把最难样本交给专家，而是在当前信息和预算下选择预期有用的观测。模型拿不准可能源于知识不足，也可能源于输入本身无法判读；两个都很有价值的单点放在同一批次中，也可能提供重复信息。因此，本章依次区分预测不确定性与可减少的不确定性，推导标签带来的信息增益和预期风险变化，再用几何覆盖与梯度表示处理批量冗余，最后把费用、标注错误和抽样偏差纳入查询与标注流程。

获取函数给记录排出的次序，只是在估计先标哪一条可能更有帮助。它的分数再高，也要通过实际标注和重新训练来检验。比较方法时应给它们相同的总预算，并计入专家核对、返工和等待时间，再看哪一种更快减少漏报与误报。如果记录分散在不能交换原始数据的工厂中，还需要安排各处怎样共同训练，这就是下一章联邦学习的问题。

\section{查询规则与标注预算}

% auxiliary-resource-guide
本章末的“辅助学习材料”列出与核心方法对应的讲解和实践入口，可在阅读相关推导后，按所列小节对照学习。

\begin{figure}[htbp]
\centering
\begin{tikzpicture}[node distance=0.7cm and 0.75cm]
 \node[idi input, minimum width=1.8cm] (pool) {未标注池\\$U$};
 \node[idi process, right=of pool, minimum width=2.0cm] (score) {获取函数\\$a(x;L)$};
 \node[idi state, right=of score, minimum width=1.8cm] (query) {查询集\\$Q$};
 \node[idi output, right=of query, minimum width=1.8cm] (label) {专家标签};
 \node[idi process, below=0.8cm of score, minimum width=2.2cm] (train) {更新模型\\加入有效标注对};
 \draw[idi arrow] (pool)--(score); \draw[idi arrow] (score)--(query); \draw[idi arrow] (query)--(label); \draw[idi arrow] (label)--(train);
 \draw[idi feedback] (train) to[bend right=35] (score);
\end{tikzpicture}
\caption{池式主动学习在每轮用获取函数分配有限标注预算，再用新标签更新模型和下一轮查询。}
\label{fig:active-query-loop}
\end{figure}
图\ref{fig:active-query-loop}的反馈路径说明查询与训练相互影响：获取函数选择待标注样本，新标签改变模型，也改变下一轮的查询分布。
设已标注集为$L$，未标注池为$U$，标注成本为$c(x)$，总预算为$B$。每轮训练模型，按获取函数$a(x;L)$选择查询集$Q\subset U$，获取标签后更新$L\leftarrow L\cup\{(x,y):x\in Q\}$、$U\leftarrow U\setminus Q$。查询满足$\sum_{x\in Q}c(x)\leq B_{remaining}$。获取函数（acquisition function）是对未来收益的近似，不是测试集上的真实增益\cite{settles2009active}。

池式主动学习可比较整个$U$；流式方法只能对到达样本决定查询或跳过；查询合成允许模型构造新输入，但工业中未必存在可由专家可靠判读的合成波形。无论采用哪种设定，测试集都不属于候选池。初始化标签过少时，模型置信度尚不可靠，因此随机或分层覆盖的种子集十分重要。

\paragraph{一次查询究竟买到了什么。}
一张巡检图像或一段振动窗口是输入$x\in\R^d$，故障类别$y\in\{1,\ldots,K\}$是标签；未标注不是没有输入，而是没有可供训练的可靠$y$。设备是输入的来源，不是一个类别，也不是一次查询。任务是要解决的诊断问题；本章可以始终只有一个任务，却进行很多轮查询。$L$存样本对，$U$只存输入，$Q$在送标注时也只含输入；因此 $|L|$ 计已返回的有效标签，而不计待处理记录。若初始$|L|=4,|U|=10$，本轮查询两条并都得到有效标签，则下一轮为$|L|=6,|U|=8$；标签尚未返回的记录应标为待处理，避免重复派单。拒标仍可能消耗费用，但不能伪造标签加入$L$。

\paragraph{区分训练、选择和验收。}
训练把$L$变成参数$\theta$；获取函数把每个候选变成标量分数；专家查询才产生真实标签；重新训练后独立测试才衡量收益。获取分数不是标签，也不能直接当作训练目标。预算单位应先固定，例如专家分钟数：剩余$10$分钟、每条$3$分钟时最多查询三条，而不是十条。停止条件可以是预算用完或预先规定的验证收益不足，不能反复看最终测试集来决定何时停止。

\section{不确定性、委员会与贝叶斯查询}
\paragraph{最小置信度、间隔与熵。}
设模型类别概率按大小排列为$p_{(1)}(x)\geq p_{(2)}(x)$。常见获取函数为
\begin{align}
 a_{LC}(x)&=1-p_{(1)}(x),\qquad
 a_{margin}(x)=-(p_{(1)}(x)-p_{(2)}(x)),\\
 a_H(x)&=-\sum_kp_k(x)\log p_k(x).
\end{align}
三者均选择得分较大的样本。最小置信度只看最大类别，间隔看前两类竞争，熵考虑全部概率。多分类下它们并不等价。Lewis 与 Gale 的不确定性采样工作展示了基于当前分类器选择标注对象的基本思想\cite{lewis1994uncertainty}。

\paragraph{何时这些分数等价，何时不等价。}
对二分类，记$p=P(Y=1\mid x)$，则最小置信度为$\min(p,1-p)$，负间隔为$-|2p-1|$，熵为$h(p)=-p\log p-(1-p)\log(1-p)$。在$0<p<1/2$时，$h'(p)=\log((1-p)/p)>0$，三者都随$p$靠近$1/2$而增加，故产生相同排序。多分类没有这个结论。例如两个候选概率为$(0.50,0.49,0.01)$与$(0.50,0.25,0.25)$，最小置信度相同，间隔优先选前者，熵却优先选后者；前者突出两类决策竞争，后者突出整体分散程度。这里及下文取自然对数，约定$0\log0=0$。

边界附近可能有可学习样本，也可能是永久模糊的记录、传感器损坏或标签定义冲突。反复查询后一类样本会消耗预算却不减少模型不确定性。应按设备与工况审计获取分数，并检查模型校准，而不是默认低置信度必有高信息价值。

\paragraph{分数的数值和含义。}
二分类候选甲预测为$(0.9,0.1)$，乙为$(0.5,0.5)$，则最小置信度分别为$0.1,0.5$，负间隔为$-0.8,0$，熵约为$0.325,0.693$，三种规则都优先乙。负间隔越接近零越不确定，不能因为它有负号而改成选最小值。熵$H(p)$就是按概率加权的意外程度$-\log p_k$；这里的自然对数使单位为 nat，换成以$2$为底只改变单位和统一缩放，不改变单点排序。预测概率来自当前模型而非真实频率；自信但错误的模型可能永远不主动查询甲。

\paragraph{委员会查询（Query-by-Committee）。}
委员会方法训练多个与当前数据相容的预测器，选择其分歧最大的样本\cite{seung1992qbc}。若有$M$个成员，类别投票比例为$v_k(x)$，投票熵为$-\sum_kv_k\log v_k$。也可用成员概率分布到平均预测的 KL 散度。Bootstrap、不同初始化或后验样本可形成委员会，但高度相关的成员会低估分歧，毫无约束的弱模型则可能制造无效分歧。

\paragraph{Bayesian Active Learning by Disagreement (BALD) 与联合批量查询。}
BALD 选择标签与模型参数间条件互信息最大的输入\cite{houlsby2011bald}：
\begin{equation}
 a_{BALD}(x)=H\!\left[\E_{p(\theta\mid L)}p(y\mid x,\theta)\right]
 -\E_{p(\theta\mid L)}H[p(y\mid x,\theta)].
\end{equation}
第一项是总体预测熵，第二项扣除各模型自身仍有的不确定性。用$M$次后验近似采样计算类别概率、先平均再取熵，并减去逐次熵的平均。它试图区分认知不确定性与数据歧义，但效果取决于后验近似。Monte Carlo (MC) dropout 是一种近似手段，不是精确贝叶斯采样\cite{gal2017deepactive}。

\paragraph{先把两种平均分开计算。}
两个模型分别输出$(0.9,0.1)$和$(0.1,0.9)$。先平均概率得到$(0.5,0.5)$，熵为$0.693$；各自先算熵再平均得到$0.325$，相减为$0.368$。若两个模型都输出$(0.5,0.5)$，相减为零。前者是模型意见相左，标签可能帮助选择模型；后者是每个模型自身都模糊。贝叶斯后验$p(\theta\mid L)$表示看到$L$以后对参数的分布，$\E_{p(\theta\mid L)}$指对不同可能参数求平均，而非对候选样本求平均。认知不确定性指可望随证据减少的参数分歧，数据歧义指给定解释后仍可能出现多个标签；这种分解始终依赖所选模型，不能当作对现实噪声来源的直接测量。

\paragraph{从后验信息增益推出 BALD。}
固定已标注集$L$和候选$x$，把参数$\Theta$看成后验随机变量、标签$Y$看成尚未观测的离散变量。理想查询的信息价值是观察标签后后验改变多少：
\begin{align}
 I(Y;\Theta\mid x,L)
 &=\E_{Y\mid x,L}\KL(p(\theta\mid L,x,Y)\|p(\theta\mid L))\label{eq:active-posterior-information}\\
 &=\E_{\Theta\mid L}\sum_yp(y\mid x,\Theta)
 \log\frac{p(y\mid x,\Theta)}{p(y\mid x,L)}\label{eq:active-predictive-information}\\
 &=H[p(y\mid x,L)]-\E_{\Theta\mid L}H[p(y\mid x,\Theta)].
\end{align}
从式\eqref{eq:active-posterior-information}到式\eqref{eq:active-predictive-information}使用贝叶斯公式$p(\theta\mid L,x,y)/p(\theta\mid L)=p(y\mid x,\theta)/p(y\mid x,L)$，这里采用条件预测模型并将候选$x$作为固定设计。互信息非负且至多为$\log K$，其中$K$为类别数，因为条件熵非负且$K$类熵至多为$\log K$。

取$M$个后验近似样本，令$p^{(m)}\in\R^K$为其预测、$\bar p=M^{-1}\sum_mp^{(m)}$，则估计为
\begin{equation}
 \widehat I=H(\bar p)-\frac1M\sum_mH(p^{(m)})
 =\frac1M\sum_m\KL(p^{(m)}\|\bar p).
\end{equation}
最后一个等式也说明概率委员会分歧与 BALD 的联系，但普通随机初始化委员会未必代表后验。二分类中，若每个模型都输出$(1/2,1/2)$，预测熵为$\log2$而互信息为零：再看标签未必能解决输入本身歧义。若一半模型输出$(1,0)$、另一半输出$(0,1)$，平均预测同样为$(1/2,1/2)$，但每个模型熵为零，互信息达到$\log2$：标签可以区分两组假设。这是两种不确定性的理想示例，不保证近似后验能在真实数据中正确区分它们。

逐点取 BALD 前若干名可能选出几乎相同的相邻窗口。BatchBALD 对一批标签与参数的联合互信息进行近似优化，考虑查询间冗余\cite{kirsch2019batchbald}。批量规模增加使联合计算更困难，因此应报告候选池大小、后验采样次数和获取时间。

对查询批次$Q=\{x_1,\ldots,x_b\}$，若给定参数后标签条件独立，联合信息增益可写为
\begin{equation}
 I(Y_Q;\Theta\mid L)=H(Y_Q\mid L)-\sum_{j=1}^b\E_{\Theta\mid L}H(Y_j\mid x_j,\Theta).
\end{equation}
不能把第一项换成单点熵之和，因为对参数积分后各标签往往相关。两点时，条件独立假设给出$I(Y_1,Y_2;\Theta)=I(Y_1;\Theta)+I(Y_2;\Theta)-I(Y_1;Y_2)$，省略的条件均为固定输入与$L$。这正是冗余扣除：两条记录若几乎回答同一个参数问题，单点相加会重复计算收益。联合标签空间有$K^b$种配置，直接枚举随批量指数增长，所以近似采样和贪心批量构造是计算折中，而不是改变查询目标。

\section{预期风险、梯度变化与多样性}
\paragraph{预期误差减少。}
理想查询希望标注后风险下降。由于未知$y$，可按当前预测枚举候选标签并模拟更新\cite{roy2001error}：
\begin{equation}
 a_{EER}(x)=\widehat R(\theta)
 -\sum_yp_\theta(y\mid x)\widehat R(\theta_{L\cup\{(x,y)\}}).
\end{equation}
风险代理只能在训练侧保留集或规定的无标签代理上计算，不能使用最终测试标签。每个候选、每个标签都重训代价很高，可用少量梯度步近似，但近似可能误判长期收益。若当前模型给真实类别极低概率，预期收益也可能被低估。

\paragraph{怎样模拟还没拿到的标签。}
假设某候选当前预测两类概率为$0.75,0.25$，当前代理风险为$0.20$。分别暂设标签为第一类和第二类，复制模型后独立模拟更新，得到代理风险$0.16,0.24$，则期望更新后风险为$0.75\times0.16+0.25\times0.24=0.18$，获取分数为$0.02$。这两个假想标签都不能写入正式训练集；真正查询后只用专家返回的标签更新。风险代理是训练侧可计算的近似风险，不是预知未来测试误差的工具。若没有代表性验证标签，必须说明无标签代理如何定义，不能把高置信度直接等同于低真实风险。

\paragraph{一次梯度更新能否预测风险收益。}
令$\theta\in\R^p$，候选标签为$y$时的梯度为$g_y(x)=\nabla_\theta\ell(x,y)$，模拟一步$\theta'_y=\theta-\eta g_y$。若风险代理$\widehat R$二阶可微，在当前参数处梯度为$r$、Hessian 为$H$，则
\begin{equation}
 \widehat R(\theta)-\widehat R(\theta'_y)
 \approx\eta r^\top g_y-\frac{\eta^2}{2}g_y^\top Hg_y.
\end{equation}
所以有效更新取决于方向对齐和曲率，并非只取决于$\|g_y\|$。若代理风险在当前模型处已近驻点，一阶项很小，大梯度甚至可能主要增加二阶风险。

还有一个容易忽略的抵消现象。对交叉熵$\ell=-\log p_\theta(y\mid x)$，若用同一个当前模型对未知标签求平均，则
\begin{equation}
 \sum_yp_\theta(y\mid x)g_y(x)
 =-\sum_y\nabla_\theta p_\theta(y\mid x)
 =-\nabla_\theta1=0.
\end{equation}
这里先对每个假想标签的损失求梯度，再用当前概率加权；不是对$\sum_yp_\theta(y\mid x)\ell(x,y)$整体求导，后者还会对概率权重求导。因此简单的一步、固定代理风险、一阶近似组合会给所有候选零平均收益；它不是完整 Expected Error Reduction (EER) 的替代。完整重训会改变最优解、后验或风险估计，高阶项也不会一般消失。这个结果解释了为什么预期梯度长度使用范数的平均，而不是平均梯度的范数，二者由三角不等式并不相同。

预期梯度长度以$\sum_yp_\theta(y\mid x)\|\nabla_\theta\ell(x,y)\|$估计模型变化。较大梯度说明会产生较大更新，却不保证更新方向有益，异常值可能占据排名前列。

\paragraph{Core-set 与 BADGE。}
BADGE 的英文全称为 Batch Active Learning by Diverse Gradient Embeddings，即利用多样化梯度嵌入进行批量主动学习。
Core-set 在表示空间选择覆盖候选池的样本\cite{sener2018coreset}：
\begin{equation}
 \min_{Q\subseteq U:\,|Q|=b}\max_{x\in U}\min_{z\in L_X\cup Q}
 \|f(x)-f(z)\|_2.
\end{equation}
这里$L_X=\{x:(x,y)\in L\}$为已标注集的输入集合，$b$为本轮查询数，$f$为当前特征编码器。这样距离比较的两端都是输入，而不是将带标签样本对与输入混合。
常用贪心近似每次选离已选中心最远的点，并更新所有候选的最近距离。它鼓励覆盖，但若表示由背景或设备身份主导，就会覆盖无关差异；极端离群点也可能被优先选择。

\paragraph{覆盖半径为什么有意义，又为什么不够。}
记当前已选中心集合为$S$，所有候选到中心的最近距离为$d_S(x)$。贪心最远点步骤选择$x^*=\argmax_xd_S(x)$，随后仅需更新$d_{S\cup\{x^*\}}(x)=\min\{d_S(x),\|f(x)-f(x^*)\|\}$，而不必每轮重新比较所有中心。在固定度量、等成本、中心从候选中选取的离散覆盖问题中，最远点策略有经典的两倍半径近似解释：若最终半径超过最优半径的两倍，依次选出的点和最终最远点必须彼此相距超过该值；但最优的$b$个新增半径球无法容纳$b+1$个这样的点，因为同一球内任意两点的距离至多为两倍半径。已有固定中心时，这些点也不可能由已有中心在最优半径内覆盖，故相同计数逻辑仍成立。

这个结论只保证几何覆盖。若条件期望损失$\bar\ell(z)=\E[\ell\mid Z=z]$在表示空间是$L_\ell$-Lipschitz，其中$L_\ell\geq0$是常数而非已标注集$L$，则$|\bar\ell(z)-\bar\ell(z_c)|\leq L_\ell\|z-z_c\|$，覆盖半径$r$才能被解释为至多$L_\ell r$的局部损失差异控制；但故障边界可能不平滑，表示可能遗漏关键变量，而且一个中心标签不等于精确的条件期望损失。因此几何近似保证不能直接改写为标注后分类风险保证。

BADGE 将不确定性与多样性放到最后一层的梯度表示中\cite{ash2020badge}。先取预测概率最大的类别为伪标签，再计算各类别的梯度块：
\begin{equation}
 \hat y=\arg\max_kp_k(x),\qquad
 g_k(x)=(p_k(x)-\mathbf1[k=\hat y])h(x),\label{eq:active-badge-embedding}
\end{equation}
其中$h(x)$是分类头输入。拼接各类梯度后，用类似$k$-means++的距离采样选择一批点。梯度长度反映当前预测的不确定性，梯度方向帮助减少冗余。该方法不需要候选真实标签，但依赖伪标签与当前表示质量。

\paragraph{从覆盖与梯度各算一个小例子。}
在一维表示上，已标注中心为$0$，候选为$1,2,8,9$，预算为两条。最远点法先选$9$，最近中心距离更新为$1,2,1,0$，再选$2$，而不是继续选靠近$9$的$8$；这就是批量多样性的作用。若$9$只是坏传感器的读数，则覆盖仍可能浪费一次查询，故离群审核不能省略。BADGE 中若$h=(1,2)^{\mathsf T}$、$p=(0.6,0.4)$，伪标签为第一类，两个梯度块为$(-0.4,-0.8)^{\mathsf T}$和$(0.4,0.8)^{\mathsf T}$，拼成四维向量。伪标签仅用于算选择特征，不表示已经获得专家标签。

\paragraph{梯度嵌入为何同时包含置信度与方向。}
令线性分类头$W\in\R^{K\times d_h}$、特征$h(x)\in\R^{d_h}$，其中$d_h$是编码后的特征数，不必等于输入维数$d$；logits 为$Wh\in\R^K$，概率为$p\in\R^K$。对标签$y$，交叉熵梯度为
\begin{equation}
 \nabla_W\ell=(p-e_y)h^\top\in\R^{K\times d_h},
\end{equation}
其中$e_y\in\R^K$为第$y$个标准基向量。逐行展平这张矩阵得到$Kd_h$维嵌入，便得到式\eqref{eq:active-badge-embedding}中的$g_k$。其范数满足
\begin{equation}
 \|g(x)\|_2^2=\|p-e_{\hat y}\|_2^2\|h(x)\|_2^2.
\end{equation}
若预测接近某个确定类别，第一个因子接近零；若多个类别竞争，则残差通常更大。但第二个因子也会影响长度，所以不能把它当作纯粹的概率不确定性。两个样本的梯度内积为$(p-e_{\hat y})^\top(p'-e_{\hat y'})\,h^\top h'$：类别竞争结构与特征相似性共同决定方向。

一种可复现的初始化约定是先选梯度范数最大的候选作为首中心，并列时按固定种子的随机规则选择。随后类似$k$-means++，在已有中心集合$S$下，以$D(x)^2/\sum_{u\in U\setminus S}D(u)^2$从未选候选中抽取新中心，其中$D(x)=\min_{z\in S}\|g(x)-g(z)\|$。每选一条即更新最近距离，直到达到本轮预算。这既避免单纯选取最大梯度的一簇近重复样本，又保留随机探索。若所有剩余距离为零，可在未选候选中均匀无放回抽样，不能除以零或重复选择同一条。它是在梯度空间中寻找代表性，并不保证伪标签正确或最终风险最小。

\section{成本、噪声与工业流式变体}
等成本预算下按获取分数排序即可形成简单策略；成本差异明显时，对非负收益分数且$c(x)>0$的情形，可用$a(x)/c(x)$作启发式，或直接求带预算约束的批量选择。前面的负间隔不能未经处理就按此解释为收益：例如固定加一个常数虽不改变单点排序，却会改变除以不同成本后的排序，所以需先明确收益代理的零点与尺度。比值排序不是一般情况下的最优解。长工单、复杂图像与短波形的标注时间不同，应使用专家分钟数或费用，而不只用样本数比较。

若先把收益近似为可加的$a_i\geq0$，批量选择就是$\max_{q_i\in\{0,1\}}\sum_iq_ia_i$，约束$\sum_iq_ic_i\leq B$。比值排序解的是可分割背包的思想，而样本标签不可分割。例如预算为$10$，一个候选成本$6$、收益$12$，另外两个各成本$5$、收益$9$；按比值会先选第一个而得到$12$，选后两个却得到$18$。数值只是算法反例，不是实验结果。实际批量收益还因冗余不满足可加性，故应把多样性和成本联合考虑。

多标注者场景还需选择由谁标注，并允许“不确定”或“无法判读”。对高风险少数故障可配置双人复核；对重复分歧应先修改标签定义。密度加权可以降低孤立噪声的优先级，但也可能排斥真正罕见的新故障，因此应保留随机探索和新颖性审核预算。

流式方法可在$a(x_t)>\tau_t$时查询，并根据剩余预算调整阈值。漂移后旧阈值可能失效，不能仅依赖过去的分数分布。持续学习决定如何更新模型，主动学习决定获取哪些标签，二者可联合但应分别消融。查询高不确定样本产生选择偏差，性能监测仍需要独立的代表性抽样，不能直接用已查询数据的准确率估计上线风险。

\paragraph{查询偏差与代表性监测。}
设一个固定候选总体有$N$个样本，$S_i\in\{0,1\}$表示是否被查询，已知边际入选概率为$\pi_i>0$。对预先固定的评估器，若损失为$\ell_i$，则逆概率估计
\begin{equation}
 \widehat R_{IPW}=\frac1N\sum_{i=1}^N\frac{S_i\ell_i}{\pi_i},\qquad
 \E[\widehat R_{IPW}]=\frac1N\sum_{i=1}^N\ell_i
\end{equation}
利用的是$\E S_i=\pi_i$。查询越少的区域权重越大，故方差可能很高；确定性只选最高分时，未入选样本$\pi_i=0$，无法作此修正。自适应多轮查询的条件概率还依赖历史，最终入选概率不能随意替换为某轮的查询概率；同一批查询标签又参与训练时，固定评估器的无偏论证也不再直接适用。

一个可操作的补充是将部分预算留给概率探索：每轮候选数为$|U|$时，按$q_t(x)=(1-\epsilon)q_{a,t}(x)+\epsilon/|U|$查询，$0<\epsilon\leq1$，并记录当轮概率与历史状态。这保证当轮每个候选有非零机会，却不自动保证精确最终风险估计。最清晰的验收仍是保留不参与获取排序和训练的代表性监测样本，其成本应计入整体标注预算，而不是被当作免费信息。

标注噪声还会改变信息价值。若标注者$a$的混淆矩阵为$T^{(a)}_{ky}=P(\widetilde Y=k\mid Y=y,a)$，并假定给定真实标签后错误与$x,\theta$无关，则观测标签预测为$p(\widetilde y=k\mid x,\theta,a)=\sum_yT^{(a)}_{ky}p(y\mid x,\theta)$。查询信息增益应针对实际收到的$\widetilde Y$计算，而非假设完美的$Y$。矩阵本身未知、错误随图像质量变化或多人错误相关时，需要审核样本和更细模型；简单多数投票不能消除共同误判。

\paragraph{查询偏差与标注噪声的小例子。}
某记录以$\pi_i=1/4$概率被查询，则它被查询时在逆概率求和中贡献$4\ell_i$，未查询时贡献零，期望仍为$\ell_i$；除以总体数$N$后才是平均风险贡献。这不允许在$\pi_i=0$时补权重，也不意味着一次抽样结果一定接近真实平均。对于二分类标注者，若每类都有$0.1$概率被误标为另一类，则$T=\left(\begin{smallmatrix}0.9&0.1\\0.1&0.9\end{smallmatrix}\right)$，列对应真标签且每列和为一。模型真标签概率为$(0.8,0.2)^{\mathsf T}$时，收到标签的预测为$T(0.8,0.2)^{\mathsf T}=(0.74,0.26)^{\mathsf T}$，并非原概率。应把标注质量也纳入查询价值，不能只增加重复标注数量。

\section{算法流程与验收}
在巡检图像标注中，可以先按设备划分候选池和固定测试集，再从训练池取一批共同的初始样本。每轮用相同训练预算更新模型，给候选图像打获取分数，去掉连续拍摄的近重复图像后送专家标注，同时记录标签、耗时、费用和拒标原因。新标签加入训练集后再开始下一轮，不能把本轮查询结果提前用于本轮的分数计算。

比较随机、分层随机、不确定性和一种多样性方法，绘制真实测试召回随累计标注成本的曲线。重复不同种子集与查询随机种子，同时报告获取耗时、训练耗时、少数故障覆盖和失败查询比例。主动学习只有在同等总成本下提高真实风险表现时才有价值；节省标签却需要无法承受的反复重训，并非完整收益。

\paragraph{自测与答案。}
两个委员会成员都预测$(0.5,0.5)$，其预测熵和经验 BALD 分数各是多少？答案：预测熵为$\log2$，BALD 为$\log2-(\log2+\log2)/2=0$；高熵未必表示成员分歧。若本轮送标三条、两条有效返回、一条拒标，已标注集只增加两条，但三条都可能产生费用，不能把拒标记录当作已获得标签。

主动学习的获取函数、批量多样性和预算评价可用以下材料复习，分别对应算法直觉和工程实验。
\section*{辅助学习材料}
\begingroup
\emergencystretch=3em
\begin{itemize}
\item \textbf{Active Learning}；作者/平台：Small-Text；英文；官方教程。阅读 Components 和池式主动学习流程，理解初始化、查询策略、人工标注与停止准则如何组成迭代过程。\href{https://small-text.readthedocs.io/en/latest/active\_learning.html}{在线阅读}。
\item \textbf{Validation curves: plotting scores to evaluate models}；作者/平台：scikit-learn Developers；英文；官方示例文档。重点看训练集大小、验证曲线和重复评估，帮助把本章的累计标注预算、风险曲线与独立测试集评价落到实验设计。\href{https://scikit-learn.org/stable/modules/learning\_curve.html}{在线阅读}。
\item \textbf{通俗易懂讲AI--主动学习}；知乎；中文解说。沿查询、标注与更新循环理解预算分配；标题与主题据必应索引，原页受限。\href{https://zhuanlan.zhihu.com/p/676122996}{在线阅读}。
\item \textbf{用大白话讲主动学习（Active Learning）}；CSDN；中文博客。比较不确定性和样本多样性，评价必须使用独立测试集；标题与主题据必应索引，原页受限。\href{https://blog.csdn.net/qq\_51323826/article/details/144357831}{在线阅读}。
\item \textbf{主动学习(Active Learning，AL)的步骤、代码流程讲解及实战}；上传者：来包番茄沙司；bilibili；中文技术讲解。对照本章查询策略、不确定性度量与标注循环，检查每轮预算；评价仍使用独立测试集，不以查询样本上的准确率作为停止依据。2026年10月4日公开元数据播放14,386次，高于本次另一直接相关入门候选的1,501次；已核对公开流程简介、标题和计数，未逐帧观看。\href{https://www.bilibili.com/video/BV1y34y147xN/}{观看视频}。
\end{itemize}
\endgroup
\paragraph{本章小结。}
主动学习把标签视为需要分配的资源。不确定性寻找尚未确定的判断，BALD 进一步区分模型间分歧与单模型仍有的歧义，联合信息增益扣除批量冗余，预期风险方法直接考虑更新后的收益，而几何覆盖与梯度嵌入提供较便宜的批量代理。它们各自近似不同对象，因此获取得分不能跨方法直接解释为真实收益；大梯度不一定有益，高熵不一定可学，远距离也不一定对应新故障。

模型会判断错，专家也可能标错，同一批查询还可能重复，所以选样策略必须连同费用一起检验。上线效果应在独立且有代表性的数据上估计，不能只用主动挑出的难例测试。与持续学习结合时，可以分别替换选样方法和更新方法，查清收益来自哪里。如果多个工厂无法汇集原始记录，还要安排本地计算、更新通信与聚合。下一章讨论这种联邦协作；标签缺失、选择偏差和历史遗忘仍需各自处理。

\chapter{联邦学习}
\begin{quote}\small
\textit{本章摘要。}本章把“数据不能集中”作为约束，重新审视训练、通信、隐私和公平性。先定义联邦优化与非独立同分布，再分析 Federated Averaging (FedAvg)、本地更新、差分隐私、安全聚合和个性化方法，最后回到系统延迟、参与方差异和攻击面。联邦学习解决的是数据协作边界，不会自动解决统计偏差。
\end{quote}
上一章主动学习讨论把有限标注预算花在哪些样本上，隐含的问题是所选样本能否进入同一个学习流程。本章进一步撤去数据可集中这一便利条件：不同工厂或设备群既可能具有不同工况，也可能不能把原始记录汇集到同一处。故障预测目标没有改变，改变的是谁持有数据、谁计算梯度，以及谁能够观察更新。

假设几家工厂不能交换原始记录，但希望一起改进故障模型。首先要决定大工厂是否因样本多而占更大权重，再看各厂独立训练多步以后，更新为什么可能朝不同方向走。通信中断、设备掉线和更新中泄露的信息，也都要考虑。本章从加权经验风险推导 FedAvg，再解释近端项、控制变量和个性化方法。评价时既看平均成绩，也看表现最差的工厂，并记录通信成本和隐私预算。原始数据留在本地，并不自动保证上传的更新安全；本章因此同时追踪统计目标、通信过程和可观察信息，而不把“分散存储”误当成隐私结论。
\section{问题定义}

% auxiliary-resource-guide
本章末的“辅助学习材料”列出与核心方法对应的讲解和实践入口，可在阅读相关推导后，按所列小节对照学习。

\begin{figure}[htbp]
\centering
\begin{tikzpicture}[node distance=0.8cm and 1.0cm]
 \node[idi state, minimum width=2.4cm] (server) {服务器\\全局参数 $\theta^t$};
 \node[idi input, below left=1.0cm and 1.0cm of server, minimum width=2.0cm] (c1) {工厂 1\\本地数据};
 \node[idi input, below=1.0cm of server, minimum width=2.0cm] (c2) {工厂 2\\本地数据};
 \node[idi input, below right=1.0cm and 1.0cm of server, minimum width=2.0cm] (c3) {工厂 $K$\\本地数据};
 \draw[idi arrow,<->] (server)--(c1); \draw[idi arrow,<->] (server)--(c2); \draw[idi arrow,<->] (server)--(c3);
 \node[below=0.3cm of c2,font=\scriptsize] {双向通信：下发模型、回传更新；服务器按权重聚合};
\end{tikzpicture}
\caption{联邦学习的基本循环。原始数据留在客户端，服务器只接收模型参数或更新。}
\label{fig:federated-basic-loop}
\end{figure}
图\ref{fig:federated-basic-loop}区分本地数据路径与服务器聚合路径；客户端在原始记录上训练，服务器通过汇总更新协调模型，而不是集中收集训练样本。
FedAvg 通过客户端本地更新与服务器加权聚合实现分散训练\cite{mcmahan2017fedavg}。Federated Proximal Optimization (FedProx) 加入近端约束，Stochastic Controlled Averaging for Federated Learning (SCAFFOLD) 使用控制变量校正客户端漂移\cite{li2020fedprox,karimireddy2020scaffold}；它们改变优化机制，并不单独提供隐私保证。
联邦学习在多个客户端保留本地数据，仅交换模型更新。客户端（client）$k$有目标$F_k(\theta)$和权重$p_k$，服务器（server）协调全局目标为
\begin{equation}
 \min_\theta F(\theta)=\sum_{k=1}^{K}p_kF_k(\theta).
\end{equation}
FedAvg在每轮下发全局参数，客户端执行若干步本地 SGD，再按样本量聚合：
\begin{equation}
 \theta^{t+1}=\sum_{k=1}^{K}p_k\theta_k^{t+1}.
\end{equation}

\paragraph{客户端不是一条样本。}
一条样本$(x_{ki},y_{ki})$可以是一段振动窗口及其故障标签，$i$编号记录，$k$编号持有这些记录并能执行本地训练的客户端。一家工厂可以作为一个客户端并管理多台物理设备，也可以按设备群拆成多个客户端；划分必须与权限、通信和评价一致。任务是预测故障的规则，各客户端可以共同执行同一任务，不能把客户端数直接当成类别数或任务数。本章讨论各方具有相容输入与模型结构的横向联邦设定；同一对象的不同特征分散在不同机构的纵向联邦需要另行设计样本对齐和计算流程，不能直接平均本章的参数。

模型参数$\theta\in\R^d$是所有可训练权重按固定顺序排成的向量，此处$d$是参数数而非传感器通道数。$F_k$是客户端平均损失，$p_k$是其全局权重；上标$t$表示通信轮次而非幂次。所有客户端必须从同一个版本开始，且同一参数坐标和类别编号含义一致，参数逐坐标平均才有定义。SGD 即随机梯度下降：从本地数据抽小批算梯度，再减去步长乘梯度，不必每步遍历全部记录。

\subsection{权重、抽样与集中式基线}
设参数$\theta\in\R^d$，客户端$k$持有$n_k$个样本，且
\begin{equation}
 F_k(\theta)=\frac1{n_k}\sum_{i=1}^{n_k}\ell(\theta;x_{ki},y_{ki}),
 \qquad p_k\geq0,\quad\sum_kp_k=1.
\end{equation}
取$p_k=n_k/\sum_jn_j$得到所有样本等权的经验风险；取$p_k=1/K$得到所有工厂等权的风险。二者不是实现上的可互换选项：一家拥有大量正常记录的大工厂，在第一种目标中可能压过数据稀少但故障代价高的小工厂。权重应先由业务目标定义，而不是由方便上传的数据量事后决定。

若所有客户端参与、从同一参数出发且以共同步长$\eta$只做一步全梯度下降，则
\begin{align}
 \sum_kp_k(\theta^t-\eta\nabla F_k(\theta^t))
 &=\theta^t-\eta\sum_kp_k\nabla F_k(\theta^t)\\
 &=\theta^t-\eta\nabla F(\theta^t).\label{eq:federated-centralized-equivalence}
\end{align}
式\eqref{eq:federated-centralized-equivalence}表明，这个特例与集中式全梯度更新严格相同。若使用随机梯度$g_k$，需要条件无偏性$\E[g_k\mid\theta^t]=\nabla F_k(\theta^t)$才有期望意义上的相同；模型中的随机状态、归一化统计量等还应采用相容的处理方式。

部分参与使权重问题更复杂。令$I_k^t\in\{0,1\}$表示是否参与，已知参与概率$\pi_k>0$，则在更新与抽样条件相容时，
\begin{equation}
 \widehat g^t=\sum_k I_k^t\frac{p_k}{\pi_k}g_k^t,
 \qquad \E[\widehat g^t\mid\theta^t]=\nabla F(\theta^t).\label{eq:federated-inclusion-gradient}
\end{equation}
式\eqref{eq:federated-inclusion-gradient}中的期望等式来自$\E[I_k^t]=\pi_k$。这种逆概率校正不要求各客户端独立入选，但小$\pi_k$会放大方差。常用的参与集合内重新归一化权重是随机比值，通常不能声称对原目标严格无偏。如果网络故障与设备状态相关，又没有可靠参与概率，仅靠增加训练轮数不能恢复缺失工厂的代表性。

\paragraph{手算一轮聚合和一种掉线后果。}
两客户端分别有$1$条和$3$条样本，按样本量取权重$1/4,3/4$。从标量参数$\theta^t=0$出发，本地梯度分别为$2,-2$，步长$0.1$，一步后参数为$-0.2,0.2$，服务器得到标量参数$\theta^{t+1}=0.1$；这里的数字是参数值，不是客户端编号。直接算全局梯度$(1/4)2+(3/4)(-2)=-1$也得到$0-0.1(-1)=0.1$。若改为客户端等权，结果为$0$，这不是算错，而是目标变了。若本轮只有第一客户端上线并把其权重归一成$1$，结果为$-0.2$；掉线改变了被代表的人群，不能把此次更新称为原全局梯度。

上传差分$u_k=\theta_{k,E}^t-\theta^t$与上传完整参数在全参与且权重和为一时等价：服务器执行$\theta^{t+1}=\theta^t+\sum_kp_ku_k$。一个轮次包含选择、下发、本地计算、上传、聚合；$E$表示本地梯度步数，不自动等于遍历数据的轮数。一遍数据含多个小批时，一个本地 epoch 就有多步，比较通信与计算时要分别记录。

\section{联邦优化、隐私与系统约束}
\paragraph{非独立同分布与通信。}
当$P_k(X,Y)$不同，本地更新方向可能冲突，客户端数量、参与率和本地步数会影响收敛。通信压缩、个性化联邦学习和聚类联邦学习分别处理带宽、客户端差异和多群体结构。服务器不能直接看到原始数据，但更新仍可能泄露信息，因此需要安全聚合、差分隐私或加密协议。这里的协议指参与方之间明确规定的消息交换和计算步骤。

训练完成后，共享完整模型的客户端可直接对新输入做前向预测，不需要每预测一次就重新聚合。个性化模型还须加载该客户端自己的预测头及预处理状态，并核对它们与共享编码器的版本；新客户端没有本地头时，应先执行约定的适应流程。

\paragraph{评价与边界。}
联邦系统应报告全局性能、最差客户端性能、通信轮数、隐私预算与掉线鲁棒性。联邦学习解决的是数据不能集中这一约束，并不会自动解决标签偏差、恶意客户端或模型公平性问题。

\begin{table}[htbp]
\centering
\caption{联邦学习方法的约束匹配}
\label{tab:federated-method-constraints}
\scriptsize
\setlength{\tabcolsep}{3pt}
\resizebox{\linewidth}{!}{%
\begin{tabular}{p{2.5cm}p{3.8cm}p{3.6cm}p{4.0cm}}
\hline
方法 & 主要修正对象 & 额外代价 & 适用边界 \\
\hline
FedAvg & 基本全局聚合 & 本地漂移未被显式校正 & 分布接近、通信受限的基线 \\
FedProx & 客户端更新偏离 & 近端系数和局部适应之间需调节 & 非 Independent and Identically Distributed (IID) 与算力异构 \\
SCAFFOLD & 系统性梯度漂移 & 控制变量存储和同步 & 客户端长期参与、通信可承受 \\
个性化联邦 & 客户端条件差异 & 共享层与本地层的划分 & 工厂、设备或人群差异明显 \\
安全聚合与 Differential Privacy (DP) & 更新可见性和样本影响 & 密码学开销或噪声导致效用下降 & 需要隐私保证的跨组织协作 \\
\hline
\end{tabular}}
\end{table}

表\ref{tab:federated-method-constraints}将联邦方法与数据异质性、通信及个性化约束对应；比较方法时须先确定主要瓶颈，不能仅按聚合后平均精度选型。
\paragraph{FedAvg的本地更新。}
客户端$k$从全局参数$\theta^t$开始执行$E$步本地更新：
\begin{equation}
 \theta_{k,e+1}^{t}=\theta_{k,e}^{t}-\eta\nabla F_k(\theta_{k,e}^{t}),
 \qquad \theta_{k,0}^{t}=\theta^t.
\end{equation}
服务器聚合$\theta^{t+1}=\sum_kp_k\theta_{k,E}^t$。当各客户端数据同分布且目标函数平滑时，增加本地步数可减少通信；非 IID 情况下，本地漂移会使聚合方向偏离全局梯度。

\paragraph{从本地多步展开看漂移来源。}
先考虑全参与、确定性梯度与共同步长，以免将抽样噪声和漂移混在一起。把$E$步递推求和可得
\begin{align}
 \theta^{t+1}-\theta^t
 &=-\eta\sum_kp_k\sum_{e=0}^{E-1}\nabla F_k(\theta_{k,e}^t)\\
 &=-\eta E\nabla F(\theta^t)-\eta R_t,\\
 R_t&=\sum_kp_k\sum_{e=0}^{E-1}
 [\nabla F_k(\theta_{k,e}^t)-\nabla F_k(\theta^t)].
\end{align}
这里$R_t\in\R^d$就是相对于固定起点梯度的累计偏离。假设各$F_k$梯度为$L$-Lipschitz，且本轮轨迹上$\|\nabla F_k\|\leq G$，由三角不等式有
\begin{equation}
 \|\theta_{k,e}^t-\theta^t\|\leq\eta eG,
 \qquad \|R_t\|\leq L\eta G\frac{E(E-1)}2.
\end{equation}
于是参数误差上界为$L\eta^2GE(E-1)/2$。这不是完整收敛定理，也未区分同分布曲率与异质性；它明确说明：少通信带来的多步本地训练会累积偏离，而不是无代价地模拟集中式梯度。

异质性可用$\zeta^2(\theta)=\sum_kp_k\|\nabla F_k(\theta)-\nabla F(\theta)\|^2$描述。全局最优点上$\nabla F=0$并不要求各$\nabla F_k=0$，因此本地继续优化仍可能把参数拉向不同方向。一个可直接核算的例子是$F_k(\theta)=a_k(\theta-b_k)^2/2$，其中$\theta\in\R$、$a_k>0$。递推给出
\begin{equation}
 \theta_{k,E}=b_k+(1-\eta a_k)^E(\theta^t-b_k).
\end{equation}
当$|1-\eta a_k|<1$且$E\to\infty$，聚合趋于$\sum_kp_kb_k$，但全局最优点是$\sum_kp_ka_kb_k/\sum_kp_ka_k$。除特殊情形外二者不同；本地模型各自训练得更充分，不代表全局模型更接近目标最优点。

\paragraph{把漂移例子代入数字。}
取两客户端等权，$a_1=1,b_1=0,a_2=3,b_2=2$，即第二客户端的损失曲线更陡。各自充分训练后参数分别为$0,2$，平均为$1$；全局最优却是$(0+6)/(1+3)=1.5$。在$\theta=1$处，全局梯度为$[1+3(1-2)]/2=-1$，并未达到驻点。非 IID 意为不同客户端的记录不服从同一个联合分布，例如故障比例、输入条件或标签含义有差异；它与本地训练时间不同的系统异构不是同一个概念。即使没有异质性，多步梯度也在移动后的点求值，通常不等于把起点梯度重复$E$遍。

\paragraph{隐私的三个层次。}
数据不上传不等于隐私完好。更新可能通过梯度反演暴露样本特征，恶意客户端可能投毒，服务器也可能推断客户端属性。安全聚合保护服务器免于看到单个更新；差分隐私在更新上加入噪声，提供形式化的$(\epsilon,\delta)$保证；鲁棒聚合则针对异常更新，但不能自动提供隐私。

\paragraph{个性化联邦学习。}
全局模型适合共享规律，个性化模型允许客户端保留本地参数：
\begin{equation}
 \theta_k=\theta_{shared}+\theta_{k,local}.
\end{equation}
共享层、个性化头、元学习初始化和客户端聚类是常见路线。工业设备差异很大时，追求一个平均最优模型可能不如学习一组有结构的个性化模型。

\paragraph{个性化目标为什么需要耦合。}
仅写$\theta_k=\theta_{shared}+\theta_{k,local}$还没有定义算法，因为给共享参数加一个向量、给每个本地参数减去同一向量不会改变预测。一个明确的软共享目标是
\begin{equation}
 \min_{\theta,\{v_k\}}\sum_kp_k
 \left[F_k(v_k)+\frac\lambda2\|v_k-\theta\|^2\right],
 \quad \theta,v_k\in\R^d,\quad\lambda>0.
\end{equation}
固定$\theta$，各客户端独立优化其正则化风险；固定$v_k$，对$\theta$求导得到$\lambda\sum_kp_k(\theta-v_k)=0$，故$\theta=\sum_kp_kv_k$。这一交替过程把共享信息理解为各个性化参数的参照，而不是强迫它们完全相同。$\lambda$很大时接近共享模型，很小时接近独立训练，但数据少的客户端可能因此过拟合。这里的$v_k$是长期保留的个性化模型，与 FedProx 中每轮重置的局部变量不同。

若不同工厂的输出标签或传感器结构并不相同，可令共享编码器为$f_\theta$、本地头为$h_{\psi_k}$，优化$\sum_kp_kF_k(\theta,\psi_k)$。此时只能聚合语义和维度一致的共享参数，不能把不同类别顺序的分类头逐坐标平均。对从未参与的新工厂，还必须规定支持样本和本地适应步骤；这个问题将在元学习章进一步展开。

\paragraph{通信和系统约束。}
量化、稀疏化、低秩更新和客户端选择降低通信量。实际系统还面对带宽波动、设备掉线、异构算力和更新延迟。实验应记录每轮参与客户端数、有效样本量、字节数、端侧训练时间和掉线率，而不能只报告最终准确率。

\paragraph{压缩误差与误差反馈。}
设待上传差分$u_k\in\R^d$，压缩器$Q$可返回量化值或稀疏向量。若$\E[Q(u)\mid u]=u$且$\E\|Q(u)-u\|^2\leq\omega\|u\|^2$，压缩不改变条件均值，却增加更新方差；确定性的最大幅度稀疏化通常不满足无偏性。误差反馈保存残差$r_k^t$，令
\begin{equation}
 s_k^t=Q(u_k^t+r_k^t),\qquad
 r_k^{t+1}=u_k^t+r_k^t-s_k^t.
\end{equation}
由此$s_k^t=u_k^t+r_k^t-r_k^{t+1}$，对轮次求和后中间残差相消：$\sum_{t=0}^{T-1}s_k^t=\sum_{t=0}^{T-1}u_k^t+r_k^0-r_k^T$。它解释了为何未发送信息可以延后补偿，但残差有界及收敛仍需压缩器和目标的额外条件。长期掉线、状态丢失或模型版本变化都可能破坏这种补偿；稀疏上传还要计入索引字节，不能只数非零数值。

\paragraph{压缩发送的是近似更新，残差不能丢。}
取二维更新$u=(0.6,0.4)^{\mathsf T}$，只保留绝对值最大的一个坐标，首次发送$(0.6,0)^{\mathsf T}$，残差为$(0,0.4)^{\mathsf T}$。下一次相同更新先加残差，得到$(0.6,0.8)^{\mathsf T}$，发送$(0,0.8)^{\mathsf T}$，新残差为$(0.6,0)^{\mathsf T}$。累计发送加最终残差恰等于累计原更新$(1.2,0.8)^{\mathsf T}$。这说明信息可以延后发送，但不保证每一轮的模型与未压缩模型相同。个性化模型还需保存本地头，控制变量方法需保存$c_k$，所以通信减少与状态内存增加应一起核算。

\section{联邦诊断案例与验收}
将不同工厂作为客户端，训练共享的故障表征和个性化的告警头。一轮 FedAvg 与集中式梯度下降的差异来自本地多步更新和客户端数据权重；安全聚合只能隐藏单个更新，不能防止模型投毒，因此验收还应约束最差客户端性能、异常更新比例和通信失败后的恢复行为。

\paragraph{一轮联邦训练的完整过程。}
服务器首先选择客户端子集并发送当前模型。客户端在本地数据上训练若干步，计算更新后上传参数或差分；服务器按样本量或其他权重聚合，再将新模型发送给下一轮。这个过程的关键不是“把 SGD 分散执行”这么简单，而是每轮参与者、数据量、本地步数和网络状态都会改变优化轨迹。

\begin{figure}[htbp]
\centering
\begin{tikzpicture}[>=Latex, node distance=0.9cm, every node/.style={font=\small}]
 \node[draw, rounded corners, fill=yellow!20, minimum width=2.5cm] (server) {\shortstack{服务器\\全局模型}};
 \node[draw, rounded corners, fill=blue!10, below left=1.2cm and 1.8cm of server, minimum width=2.2cm] (c1) {\shortstack{客户端 1\\本地数据}};
 \node[draw, rounded corners, fill=green!12, below=1.2cm of server, minimum width=2.2cm] (c2) {\shortstack{客户端 2\\本地数据}};
 \node[draw, rounded corners, fill=orange!12, below right=1.2cm and 1.8cm of server, minimum width=2.2cm] (c3) {\shortstack{客户端 3\\本地数据}};
 \draw[<->, thick] (server)--(c1); \draw[<->, thick] (server)--(c2); \draw[<->, thick] (server)--(c3);
 \node[below=0.3cm of c2,font=\small] {双向通信：下发模型 / 上传更新 / 聚合};
\end{tikzpicture}
\caption{联邦学习的服务器--客户端循环；原始数据留在客户端，但更新仍需保护。}
\label{fig:federated-training-round}
\end{figure}
图\ref{fig:federated-training-round}展开一轮下发、本地训练和回传聚合的关系；客户端计算与通信构成不同成本，且回传更新本身仍可能包含敏感信息。

\paragraph{客户端漂移与 FedProx。}
非 IID 数据下，本地目标的最优方向可能不同。客户端执行过多本地步数后，更新会远离全局模型，聚合后产生客户端漂移。FedProx 在本地目标中加入近端项：
\begin{equation}
 \min_\theta F_k(\theta)+\frac{\mu}{2}\|\theta-\theta^t\|_2^2.
\end{equation}
近端系数限制本地模型偏离服务器模型的幅度，适合数据和算力差异较大的系统，但过大的$\mu$会降低本地适应能力。

对该目标求梯度，本地第$e$步为
\begin{equation}
 \theta_{k,e+1}=\theta_{k,e}-\eta
 [\nabla F_k(\theta_{k,e})+\mu(\theta_{k,e}-\theta^t)].
\end{equation}
起点$\theta_{k,0}=\theta^t$处近端梯度为零，所以第一步与 FedAvg 相同；它主要约束之后的偏离。对二次目标$F_k(v)=\tfrac12(v-b_k)^{\mathsf T}A_k(v-b_k)$，其中$A_k\in\R^{d\times d}$对称半正定、$\mu>0$，驻点方程为$(A_k+\mu I)v=A_kb_k+\mu\theta^t$，从而
\begin{equation}
 v-\theta^t=-(A_k+\mu I)^{-1}\nabla F_k(\theta^t).
\end{equation}
沿$A_k$的特征方向，位移被$1/(a+\mu)$缩放，展示了近端项如何抑制偏离。对非凸神经网络，近端项不自动使目标凸；若 Hessian 特征值下界为$-\rho$，只有$\mu>\rho$才足以使该区域的正则化 Hessian 正定。实际内层通常仅近似求解，应同时报告本地停止准则，而非默认获得精确极小点。

\paragraph{近端项不是直接把更新裁短。}
取$F_k(v)=(v-2)^2/2$、服务器起点$0$、$\mu=3$，正则化驻点满足$(v-2)+3v=0$，所以$v=0.5$而非本地独立最优$2$。它改变了整个本地目标；每步都受到返回起点的力，但并非强制要求位移小于某固定半径。若$\eta=0.1$，第一步仍为$0.2$，第二步梯度为$(0.2-2)+3(0.2)=-1.2$，更新至$0.32$。这个例子也说明有限步结果不等于精确近端解。

\paragraph{控制变量与 SCAFFOLD。}
SCAFFOLD 为服务器和客户端维护控制变量，用于估计并校正本地更新中的系统性偏差。其思想是把“客户端自身的梯度偏移”和“当前模型真正的全局下降方向”区分开。该类方法通常以额外状态和通信为代价改善非 IID 收敛，实验时必须报告控制变量的存储和同步开销。

具体地，设$c,c_k\in\R^d$分别估计全局和本地梯度方向，本地更新改为
\begin{equation}
 v_{k,e+1}=v_{k,e}-\eta(g_k(v_{k,e})-c_k+c),
 \qquad v_{k,0}=\theta^t.
\end{equation}
若在同一参考点$c_k=\nabla F_k(\theta^t)$、$c=\sum_kp_kc_k$，则第一步修正方向的条件期望恰为$\nabla F(\theta^t)$，不再随客户端改变。后续各本地点不同，校正并不意味着始终拥有精确全局梯度。一种固定$E,\eta$的控制变量更新为
\begin{equation}
 c_k^{new}=c_k-c+\frac{\theta^t-v_{k,E}}{E\eta}
 =\frac1E\sum_{e=0}^{E-1}g_k(v_{k,e}).
\end{equation}
第二个等号由本地递推求和得到，说明新控制变量代表本轮沿轨迹的平均梯度。若维护加权关系$c=\sum_kp_kc_k$，服务器应更新$c^{new}=c+\sum_{k\in S_t}p_k(c_k^{new}-c_k)$，未参与者保留旧状态。经典均匀客户端设定取$p_k=1/K$；不同权重、服务器步长和参与方案需要配套实现，不能混用聚合规则。陈旧控制变量以及额外的$d$维上传，是该方法的实际代价。

\paragraph{控制变量为什么有减有加。}
若客户端当前梯度为$5$，保存的本地方向估计为$c_k=4$，全局方向估计为$c=1$，修正后是$5-4+1=2$。减去$4$是在扣除本地系统偏移，加上$1$是在补回共享参考方向，不是把梯度都强制为零。若当前梯度仍恰为$4$，则得到$1$；若本地分布已经改变，陈旧的$4$也可能校正过头，所以状态要随参与更新并与模型版本对应。

\paragraph{差分隐私与隐私预算。}
必须先指定保护单位。样本级邻接表示训练集仅相差一条记录；客户端级邻接表示相差一个客户端的全部贡献。单样本梯度裁剪通常用于样本级保护，不能直接称为客户端级保护。随机机制$M$满足$(\epsilon,\delta)$差分隐私，是指任意相邻数据集$D,D'$和可测输出事件$A$都有
\begin{equation}
 \Pr[M(D)\in A]\leq e^\epsilon\Pr[M(D')\in A]+\delta.
\end{equation}
它约束观察结果因一单位数据变化而改变的程度，不是“攻击一定失败”的概率声明。

以客户端级发布为例，先把每个上传向量$u_k\in\R^d$裁剪为
\begin{equation}
 \bar u_k=u_k\min\{1,C/\|u_k\|_2\},\qquad
 M=\frac1m\left(\sum_k\bar u_k+\xi\right),\quad
 \xi\sim\mathcal N(0,\sigma^2C^2I_d),
\end{equation}
其中零向量保持为零，$m$是固定公开归一化常数，缺席贡献以零计。加入或删除一个客户端时，和的$\ell_2$敏感度至多为$C$；替换一个客户端时由三角不等式得到$2C$。因此两种邻接定义不可使用相同敏感度而不作说明。对一次发布、加入删除邻接、$0<\epsilon<1$，经典高斯机制的一个充分条件是$\sigma\geq\sqrt{2\log(1.25/\delta)}/\epsilon$；该式是保守校准条件，不是多轮训练的最终预算。

除以$m$后，敏感度与噪声标准差同时缩小，不能据此声称隐私自动增强。裁剪把过大更新的影响截断，也引入偏差；噪声期望为零，但其平均平方范数为$d\sigma^2C^2/m^2$，高维、少参与者时可能明显损害效用。若服务器先看到未扰动单个更新，再在总和上加噪声，则保证针对发布结果的观察者，而不保护客户端免受该服务器观察；必须明确可信聚合者，或使用安全聚合与按约定步骤进行的分布式噪声生成。

对于自适应的$T$轮发布，基本组合给出$\epsilon_{tot}\leq\sum_t\epsilon_t$、$\delta_{tot}\leq\sum_t\delta_t$。实际可用更紧的隐私会计，但需声明抽样方式、采样率、噪声乘子、邻接定义和会计方法。抽样放大结论依赖抽样独立性及可见信息，不能对按故障分数挑选客户端的过程直接套用均匀抽样公式。

\paragraph{用二维向量理解裁剪、噪声和保护单位。}
若上传$u=(3,4)^{\mathsf T}$，长度为$5$，阈值$C=2$，裁剪后为$(1.2,1.6)^{\mathsf T}$，方向不变而长度变为$2$；小于阈值的向量不变。噪声不是每次固定加同一个数，而是按声明分布独立生成新的随机向量，再按发布机制聚合。$\epsilon$控制相邻数据发布概率比的容许变化，$\delta$是概率不等式的松弛量，二者都不是准确率。保护一条窗口与保护整家工厂的所有窗口差别很大，应先回答要隐藏哪一种贡献，再选择裁剪与会计方式。

\paragraph{安全聚合与攻击面。}
安全聚合使服务器只能看到聚合结果，不能直接读取单个客户端更新，但它不能防止恶意客户端上传投毒或后门更新。鲁棒聚合可以按坐标中位数、截断均值或相似性筛除异常，却可能误伤来自真实少数群体的正常更新。安全性需要从服务器、客户端、通信通道和模型输出四个方面分别分析。

\paragraph{安全聚合为何能只显露总和。}
下面只用代数说明掩码抵消，不把它当作可直接部署的完整密码协议。设更新经约定量化后为有限域$\mathbb F_q^d$上的向量$u_k$，每对客户端共享随机掩码$r_{kj}$，并约定$r_{jk}=-r_{kj}$。客户端发送
\begin{equation}
 \widetilde u_k=u_k+\sum_{j\ne k}r_{kj},\qquad
 \sum_k\widetilde u_k=\sum_ku_k+
 \sum_{k<j}(r_{kj}+r_{jk})=\sum_ku_k.
\end{equation}
掩码在总和中抵消，单个带掩码消息却不直接等于原更新。实际协议还需要认证、密钥协商、门限恢复及串谋假设；若某客户端掉线，它对应的掩码未必成对抵消，需要恢复机制，且不能同时恢复足以解开其原始贡献的全部秘密。量化范围和模数还要防止解码溢出。

服务器得到总和仍可能通过参与集合差分或很小的聚合组推断贡献，因此还需最小聚合规模、查询限制及必要的差分隐私。另一个根本冲突是：坐标中位数不是总和的函数，服务器若只能获得总和，就不能直接执行基于单个更新的筛选。安全聚合与鲁棒聚合必须联合设计，而不是把两个算法名称并列便认为两项保证同时成立。

\paragraph{掉线、异构与公平性。}
慢客户端会拖延同步，掉线客户端可能造成样本代表性变化；按样本量加权又可能让大客户端压过小群体。评价时应报告平均性能、最差客户端、不同客户端规模区间和掉线条件下的性能。对于工业多工厂系统，公平性不仅是统计指标，也关系到不同工厂是否都能获得足够的故障检测能力。

在多工厂轴承诊断中，可把工厂、产线或设备群定义为客户端。首先固定一个不参与训练的跨工厂测试集，并在每轮记录参与客户端、样本量、本地步数、上传字节数和训练时长。随后构造三种非 IID 情形：不同工厂故障比例不同、不同工厂只观测部分类别、不同工厂的相机或传感器分布不同。FedAvg、FedProx、SCAFFOLD 和个性化头应使用相同总通信预算比较，同时报告全局平均、宏平均、最差工厂和各类故障召回。

隐私验收不能只检查原始文件是否离开工厂。应分别测试更新反演、成员推断、恶意客户端投毒和服务器串谋；对差分隐私记录每轮裁剪、噪声和累计$(\epsilon,\delta)$，对安全聚合记录阈值、掉线恢复和密钥失效行为。若加入噪声后小工厂的少数故障召回显著下降，平均准确率可能仍然稳定，但系统已经产生了不公平的风险分配，因此最差客户端和少数类指标必须进入停止条件。

\paragraph{自测与答案。}
两客户端等权，从$\theta^t=0$出发，各做一步，梯度为$2,-2$、步长为$0.1$，聚合参数是多少？答案：$(-0.2+0.2)/2=0$。若只第二家参与并将其权重归一为$1$，则得到$0.2$，不再是上述全参与更新。若采用 FedProx，第一步是否变化？不变，因为起点的近端梯度为零；差别从后续本地步开始出现。

以下框架教程、隐私指南与入门讲解可用于复习联邦训练流程及其保证边界。
\section*{辅助学习材料}
\begingroup
\emergencystretch=3em
\begin{itemize}
\item \textbf{Quickstart PyTorch}；作者/平台：Flower；英文；官方框架教程。阅读 ClientApp、ServerApp、FedAvg 与本地模拟流程，观察数据分区、轮次聚合和客户端评估。\href{https://flower.ai/docs/framework/tutorial-quickstart-pytorch.html}{在线阅读}。
\item \textbf{TensorFlow Privacy：隐私保护指南}；作者/平台：TensorFlow；中文；官方指南。阅读 DP-SGD 的介绍，理解限制单个样本影响的动机及其与联邦学习用户级隐私的区别；隐私预算还需结合本章的会计分析。\href{https://tensorflow.google.cn/responsible\_ai/privacy}{在线阅读}。
\item \textbf{详解联邦学习Federated Learning}；知乎；中文解说。画出本地训练与服务器聚合流程，检查客户端权重；标题与主题据必应索引，原页受限。\href{https://zhuanlan.zhihu.com/p/79284686}{在线阅读}。
\item \textbf{终于有人把联邦学习讲明白了}；CSDN；中文博客。对照非独立同分布和通信约束；数据不出本地不等于已经满足隐私保证；标题与主题据必应索引，原页受限。\href{https://blog.csdn.net/hzbooks/article/details/118097998}{在线阅读}。
\item \textbf{《联邦学习技术介绍、应用和FATE开源框架》第1课！了解一下AI的新风口？}；上传者：FedAI联邦学习；bilibili；中文课程。对照本章联邦协作流程与应用边界，区分框架实现和隐私保证；旧课程的软件接口以当前官方文档为准。2026年10月4日公开元数据播放12,741次，高于本次异步联邦教程候选的6,701次；已核对公开课程简介、标题和计数，未逐帧观看。\href{https://www.bilibili.com/video/BV1JE411p72s/}{观看视频}。
\end{itemize}
\endgroup
\paragraph{本章小结。}
联邦学习要求数据留在各工厂，但这不会替我们决定应该最小化哪一种风险。所有工厂各做一步、再按样本聚合时，FedAvg 可还原集中式梯度；每个工厂本地走多步后，各自停在不同位置，聚合方向就会产生偏差。FedProx 用近端项限制本地模型离开全局模型太远，SCAFFOLD 用控制变量修正客户端漂移，个性化方法则允许不同工厂保留不同的预测头。它们解决的不是同一个优化问题，不能只按通信轮数或平均准确率排顺序。

隐私与可靠性同样必须拆开：安全聚合限制单个更新的可见性，差分隐私限制单个保护单位对发布分布的影响，鲁棒聚合处理异常贡献。有效系统还要在掉线、少数故障和最差工厂上通过验收。本章的基本立场不是尽可能多地分散计算，而是在可说明的统计目标、信任边界与资源预算之间选择可执行的协作方式。

能够共同训练以后，还要检查怎样衡量样本相似性。高维读数距离近，可能只是温度量纲较大，不一定代表故障相似。下一章的流形学习研究局部邻近关系能否反映转速、负载等少量变化因素。理解这些关系，也有助于判断联邦学习中所谓相似工厂究竟是工况相似，还是仅仅表示中的数值接近。

\chapter{流形学习}
\label{chap:manifold}
\begin{quote}\small
\textit{本章摘要。}流形学习假定高维观测由少量连续因素控制。本章从线性子空间基线出发，分别推导保持测地距离、局部重构和图平滑的算法，再讨论 t-distributed Stochastic Neighbor Embedding (t-SNE)、Uniform Manifold Approximation and Projection (UMAP) 与流形正则化。低维图形只是一种表示，不能直接证明故障类别可分。
\end{quote}

上一章解决了数据分散时如何协同训练，却没有保证传递给模型的距离和表示本身合理。相同故障在不同转速下可能相距很远，不同故障也可能因传感器尺度相似而靠得很近。本章把问题从“数据在哪里”转向“数据之间什么关系应被保留”：如果观测的高维变化主要由少量连续因素驱动，能否在低维坐标中保留对判断有用的结构？

把许多传感器读数画到平面上，可以有不同的画法。PCA 尽量保留线性重构能力，Isometric Mapping (Isomap，等距映射) 保留沿邻域图行走的测地距离，Locally Linear Embedding (LLE) 保留局部重构关系，谱方法关注图上的平滑变化；概率可视化更重视谁与谁相邻，并不保证远处两点的距离准确。本章分别推导这些目标。看图时还要问：聚在一起的是同一种故障，还是同一种转速？换一台设备后，这种表示是否仍有助于识别故障？本章的几何目标只有在服务这个预测问题时才有意义，漂亮的二维图不是终点。

\section{几何假设与邻域构造}

% auxiliary-resource-guide
本章末的“辅助学习材料”列出与核心方法对应的讲解和实践入口，可在阅读相关推导后，按所列小节对照学习。

\begin{figure}[htbp]
\centering
\begin{tikzpicture}[x=0.9cm,y=0.75cm]
 \draw[gray!35] (0,0) grid (6,4);
 \draw[blue!75!black, very thick] plot[smooth] coordinates {(0.4,0.7)(1.1,1.5)(1.8,2.8)(2.7,3.1)(3.5,2.0)(4.3,0.9)(5.4,1.2)};
 \foreach \x/\y in {0.4/0.7,1.1/1.5,1.8/2.8,2.7/3.1,3.5/2.0,4.3/0.9,5.4/1.2} {\fill[red!70!black] (\x,\y) circle (2pt);}
 \draw[orange!80!black, dashed] (1.1,1.5) circle (0.55cm);
 \draw[orange!80!black, dashed] (3.5,2.0) circle (0.55cm);
 \node[font=\scriptsize, anchor=west] at (6.15,3.2) {局部邻域};
 \node[font=\scriptsize, anchor=west] at (6.15,2.65) {高维观测};
 \node[font=\scriptsize, anchor=west] at (6.15,2.1) {低维流形};
\end{tikzpicture}
\caption{流形学习把高维观测看作嵌入空间中的低维结构；邻域大小决定局部几何是否可靠。}
\label{fig:manifold-local-geometry}
\end{figure}
图\ref{fig:manifold-local-geometry}用曲线上的邻域说明局部距离与整体嵌入形状的区别；邻域过大会跨越不同局部区域，过小则可能使几何估计缺乏样本支持。
设$x=g(z)+\epsilon\in\R^d$，$z\in\R^r$且$r\ll d$，$g$在局部光滑。目标可以是恢复坐标、保持邻域或改善预测，这三者不是同一个优化问题。不同载荷、转速和传感器尺度会改变距离；先按训练数据确定标准化和度量，再构造$k$近邻或半径邻域图。

邻域太小会使图断裂，太大会跨越流形折叠形成捷径。互为近邻与并集近邻的连通性不同，应记录图的连通分量和边长分布。故障突变可能跨越不连续机制，不能因为原始欧氏距离近就强迫标签平滑。

\paragraph{低维并不等于只留下少数传感器。}
每个$x_i\in\R^d$是一条记录的$d$个特征，$n$条记录堆成矩阵$X$，行是样本、列是特征。设备是记录来源；任务是用这些记录作何种预测，不是矩阵的一行。潜在变量$z\in\R^r$是少量连续因素，$\epsilon\in\R^d$是观测扰动，$g:\R^r\to\R^d$把因素变成观测。比如$x=(\cos z,\sin z)$有两个读数，却只由一个角度控制；这是局部一维圆周，不能靠删掉其中一个读数无歧义地表示整圈。流形是局部像低维平面一样可用坐标描述的集合，通常不要求整体是平面。实际算法只能依据有限且有噪声的点猜测这种结构。

构图时每个样本是一节点，近邻之间连边；边长通常是距离，边权也可以是随距离减小而增大的相似度，二者不能混淆。$k$近邻是为每点找最近的$k$个其他点；有向近邻关系不必对称，双方都选中才连边是交集规则，任一方选中就连边是并集规则。标准化例如把温度减去训练均值再除以训练标准差，避免温度的数值尺度压过振动；零标准差通道需要删除或另作约定。

\section{线性基线与测地距离}
\paragraph{PCA 与核 PCA。}
对中心化数据矩阵$X\in\R^{n\times d}$，PCA 求
\begin{equation}
 \min_{V^{\mathsf T}V=I}\|X-XVV^{\mathsf T}\|_F^2,
\end{equation}
其解由$X^{\mathsf T}X$的主特征向量组成。坐标为$Z=XV$，截断奇异值分解可避免显式构造大协方差矩阵。高方差方向不一定携带故障信息，低方差异常可能被截断。核 PCA 在核矩阵上中心化后求特征向量，允许非线性表示，但需要选择核与处理新样本映射\cite{scholkopf1998kpca}。

\paragraph{PCA 的重构推导与尺度选择。}
令$V\in\R^{d\times r}$列正交，$P=VV^{\mathsf T}$满足$P^2=P=P^{\mathsf T}$。因此
\begin{align}
 \|X-XP\|_F^2
 &=\operatorname{tr}[(I-P)X^{\mathsf T}X(I-P)]\\
 &=\operatorname{tr}(X^{\mathsf T}X)-\operatorname{tr}(V^{\mathsf T}X^{\mathsf T}XV).
\end{align}
第二行利用迹的循环性质$\operatorname{tr}(AB)=\operatorname{tr}(BA)$以及$(I-P)^2=I-P$，将两侧投影合并。最小化重构误差等价于最大化最后的迹。设$X^{\mathsf T}X=U\Lambda U^{\mathsf T}$，特征值降序排列，把$V$展开到$U$基底后，每个特征方向的被保留权重介于$0$与$1$，且总和为$r$；最大值于是为前$r$个特征值之和。若$X=U_X\Sigma V_X^{\mathsf T}$，坐标为$XV_{X,r}=U_{X,r}\Sigma_r$，最小误差为被舍弃奇异值平方之和。这也说明为何投影坐标有$r$列，而不是把样本特征向量直接当作载荷矩阵。

中心化均值和标准差仅由训练数据估计。若只中心化，新样本坐标为$z=V^{\mathsf T}(x-\bar x)$；若训练时还做了标准化，则先令$\widetilde x_j=(x_j-\bar x_j)/s_j$，再用在标准化训练数据上求得的载荷计算$z=V^{\mathsf T}\widetilde x$，不能遗漏除以训练标准差$s_j$的步骤。只中心化的 PCA 偏重数值方差大的传感器，标准化后的 PCA 则赋予低方差传感器更大相对权重；二者都不是天然正确。重特征值对应的基底可以旋转，不能把单条主轴的符号或方向变化解释成物理机制变化。

\paragraph{PCA 的每个矩阵各做什么。}
$V\in\R^{d\times r}$的列是保留方向，$Z=XV\in\R^{n\times r}$是每条记录的新坐标，$ZV^{\mathsf T}\in\R^{n\times d}$把它们重构回原特征数。$\|A\|_F^2$指矩阵所有元素平方之和，$\operatorname{tr}(A)$是方阵对角线之和；特征向量$v$满足$Av=\lambda v$，表示该方向只被矩阵缩放而不改变方向。取三条已中心化记录$(-1,0),(0,0),(1,0)$，则$X^{\mathsf T}X=\diag(2,0)$，选$r=1$时$V=(1,0)^{\mathsf T}$，得到坐标$(-1,0,1)^{\mathsf T}$，重构误差为零。若新点为$(0,1)$，投影却为$0$，第二个方向的异常完全丢失；训练重构好并不保证保留所有未来故障。

\paragraph{核 PCA 的特征空间与新样本坐标。}
设正半定核$k(x,x')=\langle\varphi(x),\varphi(x')\rangle$，$K_{ij}=k(x_i,x_j)$，$K\in\R^{n\times n}$。特征空间中心化对应$K_c=HKH$，其中$H=I-\mathbf1\mathbf1^{\mathsf T}/n$。若$K_cv_j=\lambda_jv_j$、$\|v_j\|=1$且$\lambda_j>0$，则单位主轴为
\begin{equation}
 w_j=\frac1{\sqrt{\lambda_j}}\sum_i(v_j)_i
 [\varphi(x_i)-\bar\varphi],\qquad
 z_{ij}=\sqrt{\lambda_j}(v_j)_i.
\end{equation}
归一化来自$\|w_j\|^2=v_j^{\mathsf T}K_cv_j/\lambda_j=1$。对新样本$x$，先形成交叉核向量
\begin{equation}
 (k_c(x))_i=k(x_i,x)-\frac1n\sum_a k(x_a,x)
 -\frac1n\sum_bK_{ib}+\frac1{n^2}\sum_{a,b}K_{ab},
\end{equation}
再计算$z_j(x)=v_j^{\mathsf T}k_c(x)/\sqrt{\lambda_j}$。新点必须使用训练均值，不能单独重新中心化。极小特征值会放大数值误差，核宽度会改变整个几何；核空间中的好重构也不等于可以唯一还原原始传感器读数。

\paragraph{Isomap。}
Isomap 用邻域图最短路径距离$D_{ij}$近似测地距离，再执行经典多维尺度分析\cite{tenenbaum2000isomap}：
\begin{equation}
 H=I-\frac1n\mathbf1\mathbf1^{\mathsf T},\qquad
 B=-\frac12H D^{\circ2}H,\qquad Z=U_r\Lambda_r^{1/2}.
\end{equation}
$D^{\circ2}$表示逐元素平方，$U_r,\Lambda_r$取正的主特征对。算法依次构图、计算最短路径、双中心化和谱分解。图断裂时路径距离无穷，不能直接进入 Multidimensional Scaling (MDS)；跨流形捷径则会系统性缩短距离。负特征值提示估计距离不能完全对应欧氏嵌入，不应直接对其开平方。

\paragraph{沿图走与直接穿过去的差别。}
设三点为$A=(0,0),B=(1,0),C=(1,1)$，图只连$A$--$B$和$B$--$C$，边长都为$1$。图上$D_{AC}=2$，而原空间直线距离是$\sqrt2$；Isomap 使用前者，因为假设允许的局部移动必须沿图边。如果这三点其实来自平面而非折线路径，稀疏图会高估距离，说明图距离只是估计，不自动等于真实测地距离。反过来，误加一条跨折叠短边会低估距离。把距离矩阵变成坐标时，$H\in\R^{n\times n}$执行去均值，$B\in\R^{n\times n}$存样本两两内积，$U_r\in\R^{n\times r}$是样本侧特征向量，与 PCA 的特征侧载荷$V$维度不同。

\paragraph{双中心化为什么把距离变成内积。}
暂设$D_{ij}$确实来自中心化坐标$Z\in\R^{n\times r}$，即$Z^{\mathsf T}\mathbf1=0$。令$B=ZZ^{\mathsf T}$、$b_i=B_{ii}$，平方距离展开为
\begin{equation}
 D_{ij}^2=B_{ii}+B_{jj}-2B_{ij},\qquad
 D^{\circ2}=b\mathbf1^{\mathsf T}+\mathbf1b^{\mathsf T}-2B.
\end{equation}
由于$H\mathbf1=0$且$HBH=B$，左右乘$H$便得到$B=-HD^{\circ2}H/2$。因此谱分解不是无来由的步骤，而是在从距离反求坐标内积。若$B$半正定，保留前$r$个正特征对得到$Z=U_r\Lambda_r^{1/2}$；坐标只确定到正交旋转，中心化则固定了平移自由度。

图最短路径得到的距离未必能精确嵌入欧氏空间。截断谱解最小化的是秩约束下的 Gram 矩阵误差$\|B-ZZ^{\mathsf T}\|_F^2$，不是一般意义下的$\sum_{ij}(D_{ij}-\|z_i-z_j\|)^2$。即便图构造正确，球面等具有内在曲率的流形也不能被无失真地摊平成同维欧氏平面。Isomap 尤其依赖采样足够密、邻域不跨越折叠以及测地结构近似可展平；“存在低维流形”本身还不足以保证成功。

经典实现需要存储$n\times n$距离矩阵。Landmark 变体仅用部分代表点近似全局几何，降低计算和存储，但代表点必须覆盖各工况；选择失衡会遗漏少数状态。

\section{局部重构与谱嵌入}
\paragraph{LLE。}
局部线性嵌入先固定每个样本的邻居，求局部重构权重\cite{roweis2000lle}：
\begin{equation}
 \min_W\sum_i\left\|x_i-\sum_{j\in\mathcal N_i}W_{ij}x_j\right\|_2^2,
 \qquad \sum_jW_{ij}=1.
\end{equation}
和为一的约束使权重对整体平移不变。邻域协方差奇异时需加小的对角正则。固定$W$后，再求
\begin{equation}
 \min_Z\|(I-W)Z\|_F^2,\qquad
 Z^{\mathsf T}\mathbf1=0,\quad Z^{\mathsf T}Z=I.
\end{equation}
取$(I-W)^{\mathsf T}(I-W)$最小非平凡特征值对应向量。第一步保留局部重构关系，第二步避免坐标全部坍缩。噪声、采样密度变化或过大的邻域会破坏局部线性假设。Hessian LLE 与局部切空间对齐等变体改变局部几何估计，并非仅更换一个可视化参数\cite{donoho2003hessian,zhang2004ltsa}。

\paragraph{LLE 的两个约束优化如何求解。}
对点$i$设邻居数为$m_i$，把差向量组成$A_i=[x_j-x_i]_{j\in\mathcal N_i}\in\R^{d\times m_i}$，局部权重$w_i\in\R^{m_i}$。利用$\mathbf1^{\mathsf T}w_i=1$，重构误差化为$\|A_iw_i\|^2=w_i^{\mathsf T}C_iw_i$，其中$C_i=A_i^{\mathsf T}A_i$。拉格朗日函数取
\begin{equation}
 \mathcal J(w_i,\nu)=w_i^{\mathsf T}C_iw_i
 -2\nu(\mathbf1^{\mathsf T}w_i-1).
\end{equation}
驻点满足$C_iw_i=\nu\mathbf1$，故当$C_i$可逆时
\begin{equation}
 w_i=\frac{C_i^{-1}\mathbf1}{\mathbf1^{\mathsf T}C_i^{-1}\mathbf1}.
\end{equation}
实现应求解线性系统而非显式求逆。若邻居近似落在低维切平面，$C_i$很容易奇异，可用$C_i+\gamma I$、$\gamma>0$替代。正则化提高稳定性，但同时改变重构关系。权重允许为负，LLE 的权重矩阵不是随机游走概率矩阵。

第二步设$M=(I-W)^{\mathsf T}(I-W)\in\R^{n\times n}$。因$W\mathbf1=\mathbf1$，有$M\mathbf1=0$，不加约束时$Z=0$就是平凡最优解。约束$Z^{\mathsf T}Z=I$固定尺度、$Z^{\mathsf T}\mathbf1=0$排除常数方向；Rayleigh 商最小化给出其余最小特征值对应向量。若零空间不止一维，还需诊断连通分量或重构退化，不能机械删除一个特征向量就结束。

对新点可以在训练点中求同样的和为一重构权重，并令$z(x)=\sum_jw_j(x)z_j$。这是一种已声明的邻域插值规则，不是重新求解训练谱问题。它在已覆盖的局部区域有意义，但无法可靠外推到新工况形成的流形分支。

\paragraph{局部重构不是按邻居投票。}
在一维上，点$1$由邻居$0,2$重构，权重$1/2,1/2$使$1=(1/2)0+(1/2)2$。若全体加$10$，仍有$11=(1/2)10+(1/2)12$，这就是权重和为一的平移不变性。点$3$若只能用邻居$0,2$重构，则权重为$-1/2,3/2$，和仍为一；负权重不是概率，也不能拿去做随机游走。嵌入的第二步寻找新的$z_i\in\R^r$，让同一组权重仍尽量重构它们，不是在新空间重新随意选权重。由于小邻域可能重构得过于精确，零误差本身不足以证明全局坐标可靠。

\paragraph{拉普拉斯特征映射。}
取对称非负邻接权$W_{ij}$，例如邻域内的高斯核权重；令$D_{ii}=\sum_jW_{ij}$、$L=D-W$。拉普拉斯特征映射最小化\cite{belkin2003laplacian}
\begin{equation}
 \frac12\sum_{i,j}W_{ij}\|z_i-z_j\|_2^2
 =\operatorname{tr}(Z^{\mathsf T}LZ),\qquad Z^{\mathsf T}DZ=I.
\end{equation}
排除常数方向后求$Lv=\lambda Dv$。图断裂会产生多个零特征值，因此非平凡方向的选择必须结合连通性。归一化拉普拉斯使用$D^{-1/2}LD^{-1/2}$，对孤立节点需单独处理，避免除零。

\paragraph{从边上的差异到广义特征值。}
本段$W\in\R^{n\times n}$是对称非负相似度，区别于上一方法允许负值的重构权重；本段的$D$是度矩阵，不是 Isomap 的距离矩阵。对单列坐标$v\in\R^n$逐项展开，有
\begin{align}
 \frac12\sum_{i,j}W_{ij}(v_i-v_j)^2
 &=\sum_iD_{ii}v_i^2-\sum_{i,j}W_{ij}v_iv_j\\
 &=v^{\mathsf T}(D-W)v\geq0.
\end{align}
在$v^{\mathsf T}Dv=1$下求驻点，$\nabla_v[v^{\mathsf T}Lv-\lambda(v^{\mathsf T}Dv-1)]=0$给出$Lv=\lambda Dv$。还需$v^{\mathsf T}D\mathbf1=0$排除常数方向。令$u=D^{1/2}v$就转成对称问题$D^{-1/2}LD^{-1/2}u=\lambda u$。低特征值表示沿强连接边缓慢变化，因而可能分开弱连接的设备群，但并不保证这些群等于故障类别。图密度、核带宽和正常工况的采样数量都可能决定这些方向。

\paragraph{图平滑可以在三个节点上手算。}
三节点链$1$--$2$--$3$，两条边权为$1$，则度为$(1,2,1)$，
\begin{equation}
 L=\begin{pmatrix}1&-1&0\\-1&2&-1\\0&-1&1\end{pmatrix},\qquad
 v^{\mathsf T}Lv=(v_1-v_2)^2+(v_2-v_3)^2.
\end{equation}
$v=(0,1,2)^{\mathsf T}$的能量为$2$，$v=(0,2,0)^{\mathsf T}$的能量为$8$，因此后者在边上变化更剧烈。所有节点取同一个值则能量为零，所以必须排除常数并固定尺度，否则求最平滑坐标只会得到重合的点。这里没有使用故障标签，故最平滑方向未必是分类最有用方向。

\paragraph{扩散坐标与扩散距离。}
Diffusion Maps 将邻接权归一化为随机游走转移矩阵，以特征值的时间次幂缩放坐标\cite{coifman2006diffusion}。设图连通、各度$d_i>0$，$P=D^{-1}W$，则$P_{ij}$表示一步从$i$到$j$的概率，平稳分布为$\pi_i=d_i/\sum_ad_a$。对称性给出细致平衡$\pi_iP_{ij}=\pi_jP_{ji}$，故可选右特征函数$\psi_\ell$，满足$\sum_i\pi_i\psi_\ell(i)\psi_m(i)=\mathbf1[\ell=m]$。第一项为$\lambda_0=1,\psi_0=1$。

走$t$步后的两行概率分布之差定义扩散距离：
\begin{align}
 D_t^2(i,j)&=\sum_a\frac{[P^t_{ia}-P^t_{ja}]^2}{\pi_a}\\
 &=\sum_{\ell\geq1}\lambda_\ell^{2t}
 [\psi_\ell(i)-\psi_\ell(j)]^2.
\end{align}
第二式由$P^t_{ia}=\pi_a\sum_\ell\lambda_\ell^t\psi_\ell(i)\psi_\ell(a)$代入，并用加权正交性消去交叉项得到。因此坐标$\Psi_t(i)=(\lambda_1^t\psi_1(i),\ldots,\lambda_r^t\psi_r(i))$的欧氏距离近似扩散距离，截断误差正是未保留项之和。它衡量的是多个路径共同形成的可达性，而非一条最短路径。

时间$t$越大，$|\lambda_\ell|<1$的方向衰减越强；短时间强调局部邻域，较长时间突出缓慢变化的结构，过长则使所有点近乎重合。二部图可能出现$-1$特征值和周期振荡，可采用懒惰随机游走$(I+P)/2$消除周期性。采样密度校正改变几何与密度的相对影响，应按任务解释，不能只选最美观的嵌入。

\paragraph{扩散为什么比较整行概率。}
在前面的三节点链中，从端点$1$出发一步必到$2$，从端点$3$出发一步也必到$2$，所以两行转移概率都为$(0,1,0)$，一步扩散距离为零，尽管两端的最短路径距离为$2$。扩散关心后续能到哪里，不保证每个原节点都可区分；这种对称结构会合并具有相同扩散行为的点。若使用懒惰随机游走，两端还有留在自身的概率，就能改变这一现象。选时间尺度和转移规则因此是在选择保留什么关系，而非无关紧要的绘图选项。

\section{概率邻域可视化}
\paragraph{t-SNE。}
t-SNE 在原空间构造对称邻域概率$p_{ij}$，在低维空间使用重尾分布\cite{vandermaaten2008tsne}：
\begin{equation}
 q_{ij}=\frac{(1+\|z_i-z_j\|^2)^{-1}}
 {\sum_{a\ne b}(1+\|z_a-z_b\|^2)^{-1}},\qquad
 \min_Z\sum_{i\ne j}p_{ij}\log\frac{p_{ij}}{q_{ij}}.
\end{equation}
原空间高斯带宽逐点调整，使条件分布熵对应指定 perplexity。重尾低维分布缓解拥挤问题，KL 方向更重视保留高概率邻居。优化非凸，不同初始化、perplexity 和学习率可能改变结果。图中簇间空隙、面积和距离不能直接解释为真实故障严重度。

\paragraph{t-SNE 的概率、熵与梯度。}
先在高维空间定义条件邻域概率：
\begin{equation}
 p_{j\mid i}=
 \frac{\exp\bigl(-\|x_i-x_j\|^2/(2\sigma_i^2)\bigr)}
 {\sum_{a\ne i}\exp\bigl(-\|x_i-x_a\|^2/(2\sigma_i^2)\bigr)},
 \qquad j\ne i,\quad p_{i\mid i}=0.
\end{equation}
带宽$\sigma_i$按$\exp[-\sum_jp_{j\mid i}\log p_{j\mid i}]$等于给定 perplexity 调整；它描述有效邻居数而不是类别数。再令$p_{ij}=(p_{j\mid i}+p_{i\mid j})/(2n)$，便有$\sum_{i\ne j}p_{ij}=1$。

令$a_{ij}=(1+\|z_i-z_j\|^2)^{-1}$、$A=\sum_{a\ne b}a_{ab}$。去掉不含$Z$的常数，目标为
\begin{equation}
 C=-\sum_{i\ne j}p_{ij}\log a_{ij}+\log A.\label{eq:manifold-tsne-expanded-loss}
\end{equation}
式\eqref{eq:manifold-tsne-expanded-loss}右端的第一项鼓励原空间高概率邻居靠近，第二项因归一化耦合所有点而产生排斥。利用$\nabla_{z_i}\log a_{ij}=-2a_{ij}(z_i-z_j)$，并计入有序点对的两个方向，可得
\begin{equation}
 \nabla_{z_i}C=4\sum_{j\ne i}(p_{ij}-q_{ij})
 \frac{z_i-z_j}{1+\|z_i-z_j\|^2}.
\end{equation}
当$p_{ij}>q_{ij}$时，该点对对下降方向的贡献是吸引，反之为排斥；总更新还叠加其他点的作用，不能据单个差值保证这对点本步实际靠近或远离。重尾使较远点仍能产生排斥，缓解把高维邻域挤到二维的拥挤问题，但不会恢复真实全局距离。平移、旋转不改变目标，初始化和优化路径却可能改变局部极小值；因此多次结果中稳定的邻域比某次图上的簇间距离更值得解释。

\paragraph{可视化概率不是分类概率。}
$p_{j\mid i}$表示给定样本$i$时把$j$视为邻居的相对权重，不是$j$属于某故障的概率。若一个点的两个邻居等概率，熵为$\log2$，perplexity 为$\exp(\log2)=2$；若有四个等概率邻居，则为$4$。这帮助解释有效邻居数，但不表示算法一定只连这么多条边。t-SNE 的$q_{ij}$在所有有序点对上归一化，UMAP 的$q_{ij}$则是单条边的连接强度，通常不要求全部边强度相加为一；两处使用相同字母不意味着概率空间相同。

\paragraph{UMAP。}
UMAP 用局部距离尺度构造模糊邻域图，再优化低维边连接强度\cite{mcinnes2018umap}。设原图权重为$w_{ij}$，低维相似度为$q_{ij}=(1+a\|z_i-z_j\|^{2b})^{-1}$，其图交叉熵包含
\begin{equation}
 -\sum_{i<j}\left[w_{ij}\log q_{ij}+(1-w_{ij})\log(1-q_{ij})\right],
\end{equation}
其中与$Z$无关的常数省略。实际实现用边采样与负采样近似优化，不能把采样目标不加区分地当成全对全精确计算。邻居数控制局部与较大尺度结构的折中，最小距离参数影响低维点的紧密程度，不是分类间隔的直接估计。

\paragraph{UMAP 的局部尺度与吸引排斥。}
一种典型有向邻域强度在邻居上定义为
\begin{equation}
 s_{j\mid i}=\exp\left[-\frac{\max\{0,d(x_i,x_j)-\rho_i\}}{\sigma_i}\right],
\end{equation}
非邻居取零。其中$\rho_i$给出局部连接距离，$\sigma_i>0$调节有效邻域规模。以模糊并集对称化得到
\begin{equation}
 w_{ij}=s_{j\mid i}+s_{i\mid j}-s_{j\mid i}s_{i\mid j},
\end{equation}
从而即便只有一侧视对方为近邻，也可能保留该边。这与对条件概率取平均不是同一构图过程。

令$r=\|z_i-z_j\|>0$、$q=(1+ar^{2b})^{-1}$，$a,b>0$。单对损失$\ell=-w\log q-(1-w)\log(1-q)$对距离的导数为
\begin{equation}
 \frac{dq}{dr}=-\frac{2b}{r}q(1-q),\qquad
 \frac{d\ell}{dr}=\frac{2b}{r}(w-q),\qquad
 \nabla_{z_i}\ell=\frac{2b(w-q)}{r^2}(z_i-z_j).\label{eq:manifold-umap-gradient}
\end{equation}
推导式\eqref{eq:manifold-umap-gradient}只需先对$q$求导，再乘$dq/dr$。当$w>q$时下降方向拉近点对，当$w<q$时推远点对；$r=0$附近需数值稳定处理。该式描述全对全理想目标，实际边采样和负采样的频率会改变有效权重，没有对应重要性校正时不能把随机更新称为该目标的严格无偏梯度。$a,b$控制低维相似度曲线，因而紧簇大小部分来自人为设定，不是原始数据中的物理尺寸。

参数化嵌入学习显式映射，便于处理新数据；普通谱嵌入与非参数可视化通常需要邻域插值、Nystr\"om 扩展或重新拟合。全数据联合嵌入属于传导式设定，不能据此声称对未见设备具备归纳泛化能力。

\section{流形正则与工业验证}
流形正则将未标注样本几何引入监督学习\cite{belkin2006regularization}：
\begin{equation}
 \mathcal L(f)=\frac1{n_l}\sum_{i=1}^{n_l}\ell(f(x_i),y_i)
 +\lambda_A\|f\|_{\mathcal H}^2+\lambda_I\bm f^{\mathsf T}L\bm f,
\end{equation}
其中$\bm f$为全部训练侧样本的预测向量。最后一项鼓励邻近样本预测平滑；若邻域跨越真实类别边界，它会损害少数故障。类别条件图可以利用已知标签约束邻接，但不能使用测试标签构图。

\paragraph{图上的平滑如何影响一个未知标签。}
沿三节点链，若两端预测固定为$0$和$1$，中间无标签，则图能量为$f_2^2+(f_2-1)^2$，导数$4f_2-2=0$给出$f_2=1/2$。这是沿邻域插值，不是发现中间点的真实类别；若第一条边跨越故障突变，这种平均会有害。一般式中的$\bm f\in\R^n$先按标量回归理解，多类预测需对各类分量求图惩罚或使用矩阵迹。$\|f\|_{\mathcal H}^2$约束核函数空间中的复杂度，$\bm f^{\mathsf T}L\bm f$约束训练点之间的变化，两种约束的对象不同，不能只保留后者就认为未标注孤立区域也能得到唯一答案。

\paragraph{平方损失下把函数问题化为线性系统。}
设训练侧共$n$个样本，其中$n_l$个有标签，取平方损失并假设$\lambda_A>0$。核展开$f(x)=\sum_{j=1}^na_jk(x_j,x)$足以表达最优解：任何与这些核截面正交的函数在训练点上取零，不改变经验损失和图项，却只会增加 Reproducing Kernel Hilbert Space (RKHS) 范数。因此令$a\in\R^n$、$\bm f=Ka$，$J\in\R^{n\times n}$为有标签位置为$1$的对角选择矩阵，$y$在无标签位置补零，目标化为
\begin{equation}
 \mathcal L(a)=\frac1{n_l}(Ka-y)^{\mathsf T}J(Ka-y)
 +\lambda_Aa^{\mathsf T}Ka+\lambda_Ia^{\mathsf T}KLKa.
\end{equation}
对$a$求导并除以$2$，得到
\begin{equation}
 \left(\frac1{n_l}KJK+\lambda_AK+\lambda_IKLK\right)a
 =\frac1{n_l}KJy.
\end{equation}
该系统的矩阵对称半正定。若$K$正定，$\lambda_A>0$保证唯一系数解；若$K$奇异，不能直接消去$K$，可在其非零特征子空间求解，系数虽可能不唯一，所表达的训练函数仍可相同。预测新点时使用核展开，不需要把测试标签或测试点加入训练图。

为理解图项的作用，也可暂时直接优化节点预测$f\in\R^n$。在$\frac12(f-y)^{\mathsf T}J(f-y)+\frac\gamma2f^{\mathsf T}Lf$中，驻点满足$(J+\gamma L)f=Jy$。对无标签节点$i$，有$d_if_i=\sum_jW_{ij}f_j$，即预测等于邻居的加权平均。若一个连通分量完全没有标签，其常数水平无法由此确定；若正常与故障之间有错误强边，平滑会把真实突变抹掉。几何提供的是可检验的先验，而不是无标签数据自动携带正确标签的保证。

对多工况振动数据，先统一转速、负载和通道量纲，再比较 PCA、Isomap 和图嵌入。随后改变邻居数、距离定义和随机种子，检查哪些邻近关系始终稳定。在归纳评估下，图构造和参数选择只能用训练侧数据；新设备上线时，要按事先写好的映射规则处理，不能把测试样本重新放回图中。若预先声明传导式评估，可联合使用无标签测试输入，但不得使用测试标签，也不能把结果称为逐个新点的归纳性能。除了邻域保持、重构误差或谱诊断，还要报告未见设备的故障召回、误报和计算成本。二维图看起来更整齐，却让异常召回下降时，它就不适合这个业务。

\paragraph{自测与答案。}
三节点链两端预测固定为$0,1$，若两条边权分别为$2,1$，中间点的最平滑预测是多少？答案：最小化$2f_2^2+(f_2-1)^2$，由$6f_2-2=0$得$f_2=1/3$，它更接近强边连接的左端。这个值是图假设下的插值，不是已验证的类别概率；若强边跨越真实故障边界，结果仍可能错误。

以下文档与讲解可用于比较降维目标，并检查低维图形的解释边界。
\section*{辅助学习材料}
\begingroup
\emergencystretch=3em
\begin{itemize}
\item \textbf{Manifold learning}；作者/平台：scikit-learn Developers；英文；官方文档。重点阅读 Isomap、LLE、Spectral Embedding、t-SNE 与 MDS 小节，帮助逐一对应本章的测地距离、局部重构、图平滑和概率邻域。\href{https://scikit-learn.org/stable/modules/manifold.html}{在线阅读}。
\item \textbf{Basic UMAP Parameters}；作者/平台：UMAP-learn；英文；官方教程。比较 n\_neighbors、min\_dist、n\_components 和 metric 的示例，理解局部与全局结构、嵌入紧密度和距离选择。\href{https://umap-learn.readthedocs.io/en/latest/parameters.html}{在线阅读}。
\item \textbf{【机器学习】数据降维（Dimensionality Reduction）}；知乎；中文解说。先建立PCA基线，再比较邻域和距离保持目标；标题与主题据必应索引，原页受限。\href{https://zhuanlan.zhihu.com/p/342129669}{在线阅读}。
\item \textbf{降维算法总结（超全！附代码）}；CSDN；中文博客。核对降维的目标与新样本映射；二维可视化不能直接证明类别结构或部署泛化；标题与主题据必应索引，原页受限。\href{https://blog.csdn.net/SeafyLiang/article/details/118701759}{在线阅读}。
\item \textbf{(选修)To Learn More - Unsupervised Learning - Neighbor Embedding}；李宏毅；bilibili；中文课程。选看第93集，对照本章邻域嵌入与可视化限制。2026年10月4日官方公开元数据显示整套课程累计播放约299.5万，并非本分集播放量；已核对目录与计数，未逐帧观看。\href{https://www.bilibili.com/video/BV1Wv411h7kN/?p=93}{观看分集}。
\end{itemize}
\endgroup
\paragraph{本章小结。}
本章从同一批高维观测得到了不同的数学对象：PCA 的正交投影、核 PCA 的特征空间内积、Isomap 的测地距离、LLE 的局部重构权重、谱方法的图平滑与扩散距离，以及 t-SNE、UMAP 的邻域强度。每一种谱分解或梯度优化都服务于一个明确目标；它们不能因最终都画出二维点图而互相替代。

流形假设真正有用的地方，是把局部连续变化与有限样本下的学习联系起来；真正危险的地方，是把错误距离、采样密度或跨类别边界的邻接当成真实结构。工业验证因此要同时检查图连通性、邻域稳定性、新点映射和故障指标。低维图形提供可探索的线索，而不是类别存在性、因果关系或部署泛化的证明。

本章利用样本之间的邻近关系，下一章则利用做过多个任务的经验。例如已经在多台设备上练习过少样本诊断，能否让下一台设备只需少量标签就完成适应？流形学习要检查样本邻近是否有意义，元学习则要检查以往任务的共同规律是否适用于新任务。换了故障定义或设备机理以后，两种方法的原有假设都可能需要重审。

\chapter{元学习}
\label{chap:meta}
\begin{quote}\small
\textit{本章摘要。}元学习在一组任务上优化学习过程，使新任务能够用少量数据完成适应。本章从支持集、查询集和任务分布出发，分别推导度量方法、梯度初始化与学习优化器，再比较一阶近似、隐式梯度和概率变体。少样本只是常见应用，任务级泛化才是评价核心。
\end{quote}

上一章利用样本几何学习表示，但即使表示合理，新设备仍可能只有极少标注数据。普通迁移学习先学一个模型再微调；元学习更进一步，把“看到少量样本之后如何学习”本身作为训练对象。这里被共享的不是所有任务的同一个最优参数，而是比较规则、初始化、更新器或先验，使不同任务能够以较低代价形成各自的预测器。

可以把每台新设备的适应过程看成一次练习：支持集是允许先看的少量已标注记录，查询集用来检查适应后的效果。训练时重复经历多个任务，才学出共享的元参数。下面先介绍不更新全部参数的度量方法，再用链式法则推导 Model-Agnostic Meta-Learning (MAML) 及其近似，最后讨论学习优化器和概率适应。作为全书最后一章，这里仍要回到最初的检查：测试任务是否真的没见过，查询标签是否参与了适应？这些条件不成立，少样本高分也不能说明模型会适应新任务；评价的对象必须是“用支持集适应后，在未见查询样本上的表现”。

\section{任务分布与双层目标}

% auxiliary-resource-guide
本章末的“辅助学习材料”列出与核心方法对应的讲解和实践入口，可在阅读相关推导后，按所列小节对照学习。

\begin{figure}[htbp]
\centering
\begin{tikzpicture}[node distance=0.75cm and 0.7cm]
 \node[idi input, minimum width=1.9cm] (task) {采样任务\\$\mathcal T_i$};
 \node[idi state, right=of task, minimum width=1.9cm] (support) {支持集\\$S_i$};
 \node[idi process, right=of support, minimum width=1.9cm] (adapt) {任务内适应\\$U_\phi(S_i)$};
 \node[idi output, right=of adapt, minimum width=1.9cm] (query) {查询集\\$Q_i$};
 \draw[idi arrow] (task)--(support); \draw[idi arrow] (support)--(adapt); \draw[idi arrow] (adapt)--(query);
 \draw[idi feedback] (query.south) .. controls +(0,-0.42) and +(0,-0.42) .. node[below,font=\scriptsize]{外循环更新 $\phi$} (task.south);
\end{tikzpicture}
\caption{元学习的双层结构：支持集用于适应，查询集评价适应后的模型并更新元参数。}
\label{fig:meta-support-query}
\end{figure}
图\ref{fig:meta-support-query}将支持集上的任务内适应与查询集上的外层更新分开；查询误差衡量的是适应后的泛化，而不是支持集拟合程度。
任务$\mathcal T\sim p(\mathcal T)$包含支持集$S$与查询集$Q$。支持集用于任务内适应，查询集用于评价适应结果。这里的“查询”不同于主动学习中向专家请求标签：$Q$是适应后要预测的样本集合，元测试时不能把其标签补入支持集。设适应算子为$U_\phi(S)$，元目标为
\begin{equation}
 \min_\phi\E_{\mathcal T}\left[\mathcal L_Q(U_\phi(S))\right].
\end{equation}
元参数可以是编码器、初始化、更新规则或概率先验\cite{hospedales2022meta}。元训练查询集提供外循环梯度，元测试查询标签则只能用于最终评价。把所有任务样本混合做监督训练是多任务或预训练基线，不等同于上述适应后目标。

$N$-way $K$-shot 表示每个任务有$N$个类别、每类$K$个支持样本；shot 只计支持集，不计查询集。闭集分类中的$1$-way只需输出唯一类别，不能作为有意义的分类难度；单类故障检测需要另设正常参考、未知类别和拒识阈值。工业任务应按设备、批次或工况定义，并明确新任务的类别是否与元训练重合。

\paragraph{样本、设备、任务与 episode 的具体对应。}
一段振动窗口及其故障标签是样本$(x,y)$，电机是物理设备；任务则是一个有明确类别和数据分布的预测问题，例如在设备甲上区分正常与磨损。客户端是联邦计算中的数据持有者，本章不要求客户端存在，也不要求设备之间交换模型。一次 episode 是从任务分布中抽出一个任务并组织支持、查询的一次练习，不是一条样本。若在设备甲上做$2$-way $2$-shot，支持集有正常两条、磨损两条，共四条；查询集可以另有每类三条，共六条，shot 不限制查询数。同一设备可形成多个 episode，但它们共享来源，不能被当作多个独立设备来夸大统计证据。

元训练先用支持标签适应，再用查询标签改进共享规则；元测试先固定已学到的共享规则，只用新任务支持标签适应，最后才揭示查询标签计分。度量方法的$U_\phi(S)$可以输出原型和编码器组成的预测器状态，不一定输出与$\phi$等长的网络权重。$\mathcal L_S$和$\mathcal L_Q$分别是在对应集合上平均的标量损失，内循环是任务内改变状态，外循环是跨任务改变共享参数，两个循环优化的对象不同。

\paragraph{为什么必须在任务内划分支持与查询。}
令元参数$\phi\in\R^p$，适应后参数$U_\phi(S)\in\R^d$；两者维度未必相同。训练批包含$B$个任务时，经验元目标为
\begin{equation}
 \widehat R(\phi)=\frac1B\sum_{i=1}^B
 \frac1{|Q_i|}\sum_{(x,y)\in Q_i}
 \ell(f_{U_\phi(S_i)}(x),y).
\end{equation}
先在每个查询集内平均、再在任务间平均，表达的是任务等权风险；把所有查询样本直接汇总则偏向查询集更大的任务。两者都可以定义，但必须与部署时抽到任务的分布对应。在任务独立同分布且支持查询按声明机制抽样时，该经验目标估计适应后的期望风险。高度相关的窗口不是独立任务，任务数量不能用窗口数量替代。

若把支持集也作为外循环评价集，目标会偏向快速记住已见样本，而不是快速泛化到任务内的新样本。相反，元训练查询标签用于更新$\phi$是算法本身所需，并非泄漏；真正禁止的是利用元测试查询标签选择初始化、步数、阈值或检查点。对于时间序列，支持与查询的时间关系也属于任务分布：训练时用随机窗口构造任务，部署时却要求从过去适应未来，会改变元目标所对应的条件分布。

\section{度量与关系方法}
\paragraph{Matching Networks 与原型网络。}
Matching Networks 根据查询与支持样本的相似度对支持标签加权\cite{vinyals2016matching}：
\begin{equation}
 p(y\mid x,S)=\sum_{i\in S}a(x,x_i)\mathbf1[y_i=y],
 \quad a(x,x_i)=\frac{\exp s(f(x),g(x_i))}{\sum_{j\in S}\exp s(f(x),g(x_j))}.
\end{equation}
适应主要表现为改变支持集，而非更新全部网络参数。支持样本增加会提高推理成本，噪声支持样本可能直接传播错误标签。

若查询编码$h=f(x)\in\R^r$、支持编码$e_i=g(x_i)\in\R^r$，可取$s(h,e_i)=-\|h-e_i\|^2/\tau$，$\tau>0$。对类别$k$，Matching Networks 等价于
\begin{equation}
 p(y=k\mid x,S)=
 \frac{\sum_{i:y_i=k}\exp[-\|h-e_i\|^2/\tau]}
 {\sum_j\exp[-\|h-e_j\|^2/\tau]}.
\end{equation}
它可以看作嵌入空间中的核加权投票：当$\tau\to0$且最近点唯一时接近最近邻，当$\tau\to\infty$时接近支持集类别频率。于是支持集不平衡会隐含改变类别先验，不能把“多放几个正常样本”视为不影响预测规则的改动。训练查询交叉熵时，梯度同时改变查询和支持编码器，使适用于任务的近邻关系被反复强化；若相似度本身受新工况干扰，少量错误近邻就可能主导结果。

原型网络以类均值压缩支持集\cite{snell2017proto}：
\begin{equation}
 c_k=\frac1{|S_k|}\sum_{(x_i,y_i)\in S_k}f_\phi(x_i),\qquad
 p(y=k\mid x,S)=\frac{\exp[-d(f_\phi(x),c_k)/\tau]}
 {\sum_j\exp[-d(f_\phi(x),c_j)/\tau]}.
\end{equation}
每个 episode 先编码支持集并计算原型，再计算查询交叉熵，最后反向更新共享编码器。类均值与平方欧氏距离配合有明确几何意义，换用任意距离不保证均值仍是最佳代表。温度$\tau$影响概率尖锐程度，应在元验证任务选择。

\paragraph{手算一次原型适应与预测。}
假设编码器已把支持样本映到一维，第一类为$0,2$，第二类为$4,6$。计算均值得到$c_1=1,c_2=5$；新查询编码$h=2$到两中心的平方距离为$1,9$。取$\tau=2$，第一类概率为$e^{-1/2}/(e^{-1/2}+e^{-9/2})=1/(1+e^{-4})\approx0.982$。看到支持集以后只重算原型就已经发生任务适应，无须先对编码器做梯度下降。若改把距离$|h-c_k|$而非平方距离代入，概率与相应概率假设都会改变；因此实现必须说明$d$是哪一种距离。支持集中每类至少有一个样本，否则类均值分母为零。

\paragraph{均值原型、概率模型与反向传播。}
设每个编码$h_i=f_\phi(x_i)\in\R^r$，类别$k$的原型解决$\min_c\sum_{i\in S_k}\|h_i-c\|^2$。展开并求导得到$2|S_k|c-2\sum_ih_i=0$，故最优解正是均值$c_k$。若再假设$h\mid y=k\sim\mathcal N(c_k,\sigma^2I_r)$且各类先验相同，贝叶斯公式给出平方欧氏原型的 softmax，温度$\tau=2\sigma^2$。先验不同则需在 logit 中加$\log\pi_k$。因此均值和距离的配合隐含了单峰、近似同尺度类别分布；同类多工况形成多个分支时，均值可能落在没有真实样本的位置。

对查询真标签$y$，令$\ell=-\log p_y$、$h=f_\phi(x)$。平方欧氏距离下
\begin{equation}
 \frac{\partial\ell}{\partial h}
 =\frac2\tau\left(\sum_kp_kc_k-c_y\right),\qquad
 \frac{\partial\ell}{\partial c_k}
 =\frac2\tau(p_k-\mathbf1[k=y])(h-c_k).
\end{equation}
这来自交叉熵对 logit 的导数$p_k-\mathbf1[k=y]$以及$\partial[-\|h-c_k\|^2/\tau]/\partial h=-2(h-c_k)/\tau$；对类别求和时，含$h$的项因$\sum_k(p_k-\mathbf1[k=y])=0$而抵消。固定原型、直接更新查询编码时，负梯度方向为$c_y-\sum_kp_kc_k$，即真类原型相对于概率加权原型的位置；它不保证每步都缩短到真类原型的欧氏距离。实际编码器更新还会同时改变支持与查询编码。又因$\partial c_k/\partial h_i=I_r/|S_k|$，查询误差会通过原型回传到每个支持编码，不能在元训练中默认将原型从计算图分离。推理时则只需计算支持原型和新查询距离。若支持噪声独立且协方差为$\Sigma_k$，均值方差为$\Sigma_k/|S_k|$；相关窗口会降低这种方差缩减收益，说明 shot 数必须按真正独立证据理解。

多原型变体用多个中心表示同类不同工况，但需要估计中心数和分配规则；类不平衡变体应避免支持样本数量被隐含当作类别先验。Relation Networks 用可学习关系模块替代固定距离\cite{sung2018relation}，表达能力增加，同时也更易记住训练任务的比较模式。传导式变体联合利用无标签查询集结构，须与逐查询归纳预测分开比较。

\paragraph{可学习关系与无标签查询的边界。}
关系模块可以取$r_\psi([h;c_k])\in\R$，其中拼接输入为$2r$维，令$p_k=\softmax_k(r_\psi([h;c_k]))$，再用查询交叉熵训练$\phi,\psi$。此时$\partial\ell/\partial r_k=p_k-\mathbf1[k=y]$，随后通过关系网络传回编码器。关系分数不一定对称，也不必满足三角不等式，所以它是任务比较函数而非严格距离；灵活性提高的同时失去了固定度量的一些约束。

若允许一次看到整个无标签查询批，可用软分配$r_{jk}\propto\exp[-\|h_j-c_k\|^2/\tau]$，再以$c_k^{new}=(\sum_{i\in S_k}h_i+\sum_{j\in Q}r_{jk}h_j)/(|S_k|+\sum_{j\in Q}r_{jk})$更新中心。它类似固定各向同性尺度混合模型的一次软聚类更新，不需要查询标签，但确实使用了查询分布。未知类、类别不平衡或初始原型错误可能形成自我强化，因此必须与逐查询预测分开报告，也不能在只到达一个样本的部署流程中假定整批查询可见。

\section{梯度初始化与一阶变体}
\paragraph{MAML 的内外循环。}
MAML 学习适于快速更新的初始化$\theta$\cite{finn2017maml}。单步内循环与外循环为
\begin{equation}
 \theta_i'=\theta-\alpha\nabla_\theta\mathcal L_{S_i}(\theta),
 \qquad\theta\leftarrow\theta-\beta\sum_i\nabla_\theta\mathcal L_{Q_i}(\theta_i').
\end{equation}
先采样多个任务，各自从同一初始化开始适应，再汇总查询梯度；不能让任务$i+1$直接继承任务$i$的适应结果。上式对任务求和，固定批量$B$时可把平均因子$1/B$并入外循环学习率；若要保持前述任务等权平均目标且批量大小变化，应显式除以$B$。链式法则给出
\begin{equation}
 \nabla_\theta\mathcal L_{Q_i}(\theta_i')
 =(I-\alpha H_{S_i})^{\mathsf T}\nabla_{\theta_i'}\mathcal L_{Q_i}(\theta_i').
\end{equation}
这里$i$索引本批任务，$\alpha,\beta>0$分别为内、外循环学习率，$H_{S_i}=\nabla_\theta^2\mathcal L_{S_i}(\theta)$是在当前初始化处求值的 Hessian，$I$与参数维数相同。
实现通常使用 Hessian--向量积，无需显式存储整个 Hessian，但多步内循环仍增加计算图与显存。适应步数、内循环学习率和是否更新归一化统计量都属于算法设定。

\paragraph{一维数值看清为什么要再乘一次导数。}
取初始化$\theta=0$、支持损失$(\theta-2)^2/2$、内步长$\alpha=1/2$。支持梯度为$-2$，适应后$\theta'=1$。若查询损失为$(\theta'-3)^2/2$，查询梯度对$\theta'$是$-2$，而$\partial\theta'/\partial\theta=1-\alpha=1/2$，故对初始化的元梯度为$-1$。外步长$\beta=0.1$把共享初始化更新至$0.1$，不是把它替换为本任务的$1$，也不是直接使用$-2$更新。查询标签参与元训练的这一步是合法的，元测试则不能执行这一步。Hessian 是损失的二阶偏导矩阵，用来描述梯度如何随参数改变；Jacobian 是向量映射的一阶偏导矩阵，用来把适应后参数的变化追溯到初始化。

\paragraph{多步链式法则与曲率直觉。}
设$\theta_{i,0}=\theta\in\R^d$，第$s$步$\theta_{i,s+1}=\theta_{i,s}-\alpha\nabla\mathcal L_{S_i}(\theta_{i,s})$。记$H_{i,s}\in\R^{d\times d}$为该步 Hessian、$A_{i,s}=I-\alpha H_{i,s}$，则
\begin{equation}
 \frac{\partial\theta_{i,m}}{\partial\theta}
 =A_{i,m-1}\cdots A_{i,0},\qquad
 \nabla_\theta\mathcal L_{Q_i}(\theta_{i,m})
 =A_{i,0}^{\mathsf T}\cdots A_{i,m-1}^{\mathsf T}g_{Q_i}.
\end{equation}
矩阵一般不可交换，顺序不能倒置。反向实现从查询梯度$v_m=g_{Q_i}$出发，逐步计算$v_s=v_{s+1}-\alpha H_{i,s}^{\mathsf T}v_{s+1}$。对固定向量$v$，$Hv=\nabla_\theta[(\nabla_\theta\mathcal L)^{\mathsf T}v]$，所以无须建立$d\times d$矩阵。但仍需保存或重算各步状态，而且内循环梯度必须保留可微路径；把更新后的参数无条件分离会改变算法。

用一维任务$\mathcal L_S(\theta)=a(\theta-b)^2/2$、$\mathcal L_Q(v)=c(v-b)^2/2$说明曲率作用，其中$a,c>0$。一步适应后$\theta'-b=(1-\alpha a)(\theta-b)$，外层损失及其梯度为
\begin{equation}
 \mathcal L_Q(\theta')=\frac c2(1-\alpha a)^2(\theta-b)^2,
 \qquad \frac{d\mathcal L_Q(\theta')}{d\theta}
 =c(1-\alpha a)^2(\theta-b).
\end{equation}
元目标评价的是适应后误差，而非原点误差。若支持曲率$a$允许一步恰好到达$b$，该任务对初始化的梯度可以为零；若步长使$|1-\alpha a|>1$，适应反而放大误差。这说明“容易微调”的初始化必须和步长、任务曲率共同讨论。实际支持与查询目标不同，有限支持样本产生的曲率还可能是噪声，不能把该示例的理想一步解推广到真实网络。

First-Order Model-Agnostic Meta-Learning (FOMAML) 忽略内循环 Jacobian 中的 Hessian 项，以降低开销；它是梯度近似，不是与 MAML 完全等价的目标。Reptile 在单任务上执行多步更新得到$\theta_i'$，再令\cite{nichol2018reptile}
\begin{equation}
 \theta\leftarrow\theta+\epsilon(\theta_i'-\theta).\label{eq:meta-reptile-update}
\end{equation}
式\eqref{eq:meta-reptile-update}不显式对查询损失求二阶梯度，多步任务内更新使算法偏向能在同一任务不同批次间泛化的初始化。一步更新时与多任务梯度训练的关系更接近，因此步数不是无关的实现细节。

\paragraph{一阶近似到底舍弃了什么。}
单步时 FOMAML 使用$g_Q$代替$(I-\alpha H_S)^{\mathsf T}g_Q$，故误差为$\alpha H_S^{\mathsf T}g_Q$，若$\|H_S\|\leq L$则范数不超过$\alpha L\|g_Q\|$。小步长且曲率受控时这种近似较合理；多步情况下误差来自一串 Jacobian 的省略，不能只凭单步界判断。前述一维例中，FOMAML 给出$c(1-\alpha a)(\theta-b)$，而真梯度多一个$1-\alpha a$因子，过大步长甚至会使两者方向相反。

对 Reptile，考虑同一任务上两个小批次损失，初始点处梯度为$g_1,g_2$、第二批 Hessian 为$H_2$。两步内更新用 Taylor 展开得到
\begin{align}
 \theta_1&=\theta-\alpha g_1,\\
 \theta_2&=\theta_1-\alpha\nabla\mathcal L_2(\theta_1)\notag\\
 &=\theta-\alpha(g_1+g_2)+\alpha^2H_2g_1+O(\alpha^3).
\end{align}
因此 Reptile 的位移除平均梯度项外，还含跨批次相互作用。由梯度内积的求导关系
\begin{equation*}
 \nabla(g_1^{\mathsf T}g_2)=H_1g_2+H_2g_1,
\end{equation*}
可知，若批次交换对称，二阶位移在期望上与提高同任务梯度一致性有关。这只是光滑、小步长下的局部解释，不是 Reptile 与 MAML 完全等价的证明。只有一步时没有该项，更新退化为任务采样下的普通梯度训练。

\paragraph{Meta-SGD、Almost No Inner Loop (ANIL) 与隐式梯度。}
Meta-SGD 同时学习初始化和逐参数更新尺度$\alpha$，内循环写为$\theta_i'=\theta-\alpha\odot\nabla\mathcal L_{S_i}$\cite{li2017metasgd}。更新尺度可改变方向和幅度，但也增加元参数与过拟合风险。ANIL 仅在内循环更新任务头，编码器由外循环训练\cite{raghu2020anil}；它检验快速适应是否主要依赖可复用表示，而非整个网络都必须变化。

对 Meta-SGD，$\alpha\in\R^d$也是元参数。令$g_S=\nabla\mathcal L_S(\theta)$、$g_Q=\nabla\mathcal L_Q(\theta')$，链式法则给出
\begin{equation}
 \nabla_\alpha\mathcal L_Q=-g_S\odot g_Q,\qquad
 \nabla_\theta\mathcal L_Q=
 [I-\diag(\alpha)H_S]^{\mathsf T}g_Q.
\end{equation}
当同一坐标的支持与查询梯度一致时，外循环倾向增大该坐标更新幅度；不一致时倾向减小。允许负尺度还可翻转方向，但这不再是传统的正学习率下降法，应监控任务外的稳定性。

对 ANIL，令参数为编码器$\omega$与任务头$w$，仅更新$w'=w-\alpha\nabla_w\mathcal L_S(\omega,w)$。编码器的外层梯度仍为
\begin{equation}
 \nabla_\omega\mathcal L_Q(\omega,w')
 =\partial_\omega\mathcal L_Q
 +\left(\frac{\partial w'}{\partial\omega}\right)^{\mathsf T}
 \nabla_{w'}\mathcal L_Q,\qquad
 \frac{\partial w'}{\partial\omega}
 =-\alpha\frac{\partial(\nabla_w\mathcal L_S)}{\partial\omega}.
\end{equation}
所以“内循环不更新编码器”不等于“外循环冻结编码器”，也不自动消除混合二阶导数；若额外舍弃该项，还应声明采用了一阶近似。

iMAML 在任务内目标加入到初始化的近端项\cite{rajeswaran2019imaml}：
\begin{equation}
 \theta_i^*=\arg\min_\vartheta\left[\mathcal L_{S_i}(\vartheta)
 +\frac\lambda2\|\vartheta-\theta\|^2\right].
\end{equation}
在局部可微、驻点充分求解且相关矩阵可逆时，隐式微分给出
\begin{equation}
 \frac{\partial\theta_i^*}{\partial\theta}
 =\lambda(H_{S_i}+\lambda I)^{-1}.\label{eq:meta-implicit-jacobian}
\end{equation}
此处$\lambda>0$，$H_{S_i}=\nabla_\vartheta^2\mathcal L_{S_i}(\vartheta)|_{\vartheta=\theta_i^*}$改在内层解处求值，不能沿用前面初始化处的数值。
通过迭代解线性系统获得向量积，可避免保存全部内循环轨迹；代价转为内层求解和线性系统精度。它不是任意少量梯度步下都成立的精确替代。

\paragraph{隐式微分的驻点方程与数值误差。}
令$G(v,\theta)=\nabla\mathcal L_S(v)+\lambda(v-\theta)$。内层解满足$G(\theta^*,\theta)=0$，对初始化求全导数得到
\begin{equation}
 (H_S+\lambda I)\frac{\partial\theta^*}{\partial\theta}
 -\lambda I=0.
\end{equation}
若$A=H_S+\lambda I$可逆，就得到式\eqref{eq:meta-implicit-jacobian}。对不直接依赖初始化的查询损失，先解$A^{\mathsf T}v=g_Q$，再返回$\lambda v$作为元梯度。若$A$对称正定可用共轭梯度法，并以 Hessian--向量积执行矩阵作用；单纯可逆但不定时，不能无条件使用正定系统的收敛保证。

若线性求解的残差为$r=g_Q-A^{\mathsf T}\widetilde v$，则$\|\widetilde v-v\|\leq\|A^{-\mathsf T}\|\|r\|$。小残差在病态系统中仍可能对应大误差；内层未达到驻点还会引入另一项偏差。对凸二次支持损失的特征方向$a\geq0$，初始化敏感度为$\lambda/(a+\lambda)$：证据曲率强的方向更由支持集决定，曲率弱的方向更依赖初始化。$\lambda$很大时稳定但适应受限，过小时可能病态，非凸情况下还可能切换局部解分支。隐式方法省去的是展开轨迹，而不是内层求解本身。

\paragraph{隐式解与一步适应不是同一个答案。}
取支持损失$(v-2)^2/2$、初始化$\theta=0$、$\lambda=1$。内层驻点由$(v-2)+(v-0)=0$得到$v^*=1$，一般初始化下$v^*=(2+\theta)/2$，所以对初始化的导数为$1/2$。若查询损失为$(v-3)^2/2$，在$v^*=1$时$g_Q=-2$，解$(1+1)w=-2$得$w=-1$，返回$\lambda w=-1$。如果只走一次步长$0.1$的内更新，得到$v=0.2$，它并非上述驻点，不能直接套用隐式导数当成该一步程序的精确梯度。

\section{学习更新规则与概率适应}
学习优化器将梯度与历史状态输入参数化更新器：
\begin{equation}
 h_{r+1}=g_\phi(h_r,\nabla\mathcal L_S),\qquad
 \theta_{r+1}=\theta_r+u_\phi(h_{r+1}).
\end{equation}
外循环根据优化后的损失训练$\phi$\cite{andrychowicz2016optimizer}。相比学习一个初始化，这一路线学习整个更新行为，但展开步数、参数尺度和网络结构变化会造成新的分布偏移。

\paragraph{学习优化器与直接学参数的区别。}
普通 SGD 规定更新器$u=-\eta g$，给定梯度$g=2$和步长$0.1$就移动$-0.2$。学习优化器把这个计算规则本身参数化；例如简单特例$u_\phi(g)=-\phi g$中，$\phi$是跨任务学出的步长，而$\theta$仍是本任务要优化的模型参数。历史状态$h_r$可以记住前几步梯度，类似动量，但不是输入样本的隐藏标签。训练它时必须检查经过多步后的查询损失，不能仅奖励第一步下降；过大更新可能第一步有效、随后发散。

\paragraph{学习优化器的状态展开。}
设$\theta_r\in\R^d$、隐藏状态$h_r\in\R^m$，合并状态$s_r=(\theta_r,h_r)\in\R^{d+m}$，则上述更新可以写为$s_{r+1}=T_\phi(s_r;S)$。对$\phi\in\R^p$，定义$J_r=\partial s_r/\partial\phi\in\R^{(d+m)\times p}$；状态转移的 Jacobian 为$(d+m)\times(d+m)$矩阵，故下面的乘积仍与$J_r$同形。若初始状态与$\phi$无关，有
\begin{equation}
 J_0=0,\qquad
 J_{r+1}=\frac{\partial T_\phi}{\partial s_r}J_r
 +\frac{\partial T_\phi}{\partial\phi},\qquad
 \nabla_\phi\mathcal L_Q(\theta_R)
 =J_R^{\mathsf T}\nabla_{s_R}\mathcal L_Q.
\end{equation}
若初始化也可学习，则$J_0$应包含相应导数。由于状态转移使用$\nabla\mathcal L_S(\theta_r)$，第一项通常含 Hessian；对输入梯度停止求导可降低开销，但会改变元梯度。反复相乘的状态 Jacobian 还会产生梯度消失或爆炸，截断反向传播因此是一种有偏近似，而非免费计算技巧。

逐坐标共享更新器可让同一个小网络应用于不同维度参数，然而梯度尺度、层结构、训练时长仍可能超出元训练范围。训练只展开十步的优化器不能凭短期损失下降就保证运行一千步稳定。验证时应报告完整轨迹、不同展开长度及普通 SGD、Adam 在相同计算预算下的结果。

概率元学习不只输出单个适应参数，而是在少量证据下保留多个合理任务解释。Probabilistic Model-Agnostic Meta-Learning (PLATIPUS) 等方法学习初始化的不确定性，并通过任务内更新得到预测分布\cite{finn2018platipus}。模型平均可表达任务歧义，但需要检查近似后验、采样成本与校准；支持样本少并不保证任何概率方法都能正确估计不确定性。

\paragraph{层次先验如何产生少样本适应。}
一个可明确计算的概率元学习模型是$\theta_i\sim p_\phi(\theta_i)$、$S_i\sim p(S_i\mid\theta_i)$。任务内适应对应
\begin{equation}
 p_\phi(\theta_i\mid S_i)
 =\frac{p(S_i\mid\theta_i)p_\phi(\theta_i)}
 {\int p(S_i\mid v)p_\phi(v)\,dv},\qquad
 p(y\mid x,S_i)=\int p(y\mid x,\theta_i)p_\phi(\theta_i\mid S_i)\,d\theta_i.
\end{equation}
元训练可最大化支持条件下查询数据的预测概率，从而学习使新任务容易适应的先验。这是层次贝叶斯的基本解释，不意味着所有概率元学习算法都使用相同后验或目标。

以线性高斯任务为例，$\theta_i\in\R^d$、$X_S\in\R^{m\times d}$、$y_S\in\R^m$，设$y_S\mid\theta_i\sim\mathcal N(X_S\theta_i,\sigma^2I)$，先验$\theta_i\sim\mathcal N(\mu,\Lambda^{-1})$、$\Lambda\succ0$。负对数后验中依赖$\theta_i$的部分为
\begin{equation}
 \frac1{2\sigma^2}\|y_S-X_S\theta_i\|^2
 +\frac12(\theta_i-\mu)^{\mathsf T}\Lambda(\theta_i-\mu).
\end{equation}
展开二次项并配方得到
\begin{equation}
 \Sigma_S=(\Lambda+\sigma^{-2}X_S^{\mathsf T}X_S)^{-1},\qquad
 m_S=\Sigma_S(\Lambda\mu+\sigma^{-2}X_S^{\mathsf T}y_S).
\end{equation}
因此新输入$x\in\R^d$的预测分布为$\mathcal N(x^{\mathsf T}m_S,\sigma^2+x^{\mathsf T}\Sigma_Sx)$。第一项方差来自观测噪声，第二项来自任务参数不确定性。样本未覆盖的方向主要保留先验，样本充分的方向更多服从似然，给出了“借助任务经验”的精确含义。先验均值及精度可由多个训练任务学习，但若新设备机制不在训练任务覆盖范围内，强先验也可能造成自信错误。

\paragraph{把概率适应缩成一次加权平均。}
令$d=1$，先验均值$\mu=0$、精度$\Lambda=1$，一条支持观测满足$x=1,y=2$且$\sigma^2=1$。后验方差为$\Sigma_S=(1+1)^{-1}=1/2$，后验均值为$m_S=(1/2)(0+2)=1$，介于先验的$0$与样本偏好的$2$之间。下一条$x=1$的预测方差为$1+1/2=1.5$，其中$1$是观测噪声，$1/2$是参数尚不确定。精度是方差的倒数，精度大表示先验更集中而不是更不确定。若相同噪声模型下增加独立支持观测，参数方差可下降，但观测噪声不会因此归零；重复复制同一记录不算独立证据。

非共轭模型常以$q_\psi(\theta\mid S)$近似后验，由 Jensen 不等式有
\begin{equation}
 \log p_\phi(S)\geq
 \E_{q_\psi}\log p(S\mid\theta)
 -\KL(q_\psi(\theta\mid S)\|p_\phi(\theta)).
\end{equation}
这个证据下界解释了数据拟合与接近共享先验的权衡；若元目标是条件查询预测，还需相应的外层训练，不能把优化支持集下界自动当作优化查询泛化。近似分布过于简单或采样不足会低估多种任务解释，因而概率输出仍须用覆盖率和校准检验。

记忆增强方法把任务内观测写入可寻址状态，以读取记忆完成预测。它与参数梯度适应的主要区别是任务状态的保存位置。元测试时应重置任务记忆，防止跨任务残留信息改变规则；持续学习的回放缓存则有意跨阶段保留，两者不能因都称为“记忆”而沿用同一重置规则。对于序列任务，还必须保证读取的记忆均早于当前预测时刻。

\paragraph{记忆读取也是一种适应算子。}
令任务记忆键$k_j\in\R^r$、值$v_j\in\R^q$，查询键$k(x)\in\R^r$。一个基本读取形式是
\begin{equation}
 a_j(x)=\frac{\exp(k(x)^{\mathsf T}k_j/\sqrt r)}
 {\sum_b\exp(k(x)^{\mathsf T}k_b/\sqrt r)},\qquad
 o(x)=\sum_ja_j(x)v_j.
\end{equation}
若值存标签编码，这与软近邻分类相联系；若值存任务历史状态，则输出还依赖写入与遗忘规则。分母使权重和为一，读取向量位于值向量的凸包内；要产生凸包外的复杂响应，还需后续预测网络。记忆容量有限、错误标签写入和相似键冲突都会限制适应。尤其在顺序标注规则中，当前样本标签必须在本次预测之后才能写入，否则模型学到的是读取答案而非预测。

\paragraph{一次读取的可核算结果。}
若两条记忆的键与查询相似度相同，softmax 权重各为$1/2$，存储的标签值为$v_1=(1,0)^{\mathsf T}$、$v_2=(0,1)^{\mathsf T}$，则读出$o=(1/2,1/2)^{\mathsf T}$，表示证据打平而非已确定类别。若先把当前查询真标签写入记忆再读取，模型可能轻易拿到答案，所以写入时序本身就是评价规则的一部分。切换到另一任务时既要重置任务记忆，也要重置任务内优化器的状态，除非明确研究允许跨任务保留状态的另一种设定。

\section{跨设备快速适应的完整设置}
在轴承诊断中，可以把“正常”和一种故障组成 $2$-way 任务；只有确实有五类可靠标签时，才定义 $5$-way 任务。应按设备划分元训练、元验证和元测试，同一条原始记录的相邻窗口不能一边进入支持集、一边进入查询集。元测试时，支持集只能用于预先规定的适应步骤，查询集再用来评价不同时间和工况下的表现。类别不够就修改任务设计，不能复制标签来凑 way 数。

\paragraph{任务内标签要与预测头同步。}
若某 episode 的两个类为正常与磨损，可以临时编码为$0,1$；下一 episode 换成裂纹与松动时也可以编码为$0,1$，但该任务的支持集、查询集和输出头必须使用同一映射。这些编号只表示本任务的类别位置，不能据其相同就认为两任务故障含义相同。基于固定头的梯度方法通常还要求任务输出维度相容，跨 way 数测试需另行规定头的构造，而不是直接沿用不匹配的参数。

每个 episode 的步骤是抽取任务与类别、划分支持查询、重置初始化或任务状态、执行内循环、在查询上计算损失；只有元训练阶段反向更新共享元参数。比较原型网络、MAML、一阶方法与普通预训练微调时，应匹配支持样本、预训练数据和适应计算预算。预训练基线很强时，元学习的额外收益需要单独验证\cite{chen2019closer}。

评价以任务而非高度相关的窗口为统计单元，报告平均性能、任务间区间、失败任务比例、适应前后增益、校准与延迟。改变支持集大小、标签噪声和目标设备类型，检查元训练任务分布之外的退化。不能用同一设备随机切片上的高准确率证明快速跨设备适应。

\paragraph{自测与答案。}
支持损失为$(\theta-2)^2/2$、初始化为$0$、内步长为$1/2$，查询损失为$(\theta'-3)^2/2$。一步后，MAML 与 FOMAML 对初始化给出的梯度各是多少？答案：$\theta'=1$，查询梯度为$-2$，内更新导数为$1/2$，所以分别为$-1$和$-2$。这些查询梯度可用于元训练外循环，不能用于元测试任务的适应或挑选步数。

以下课程、实现文档与入门讲解可用于复习任务划分、度量方法和可微适应。
\section*{辅助学习材料}
\begingroup
\emergencystretch=3em
\begin{itemize}
\item \textbf{CS330: Deep Multi-Task and Meta Learning}；作者/平台：Stanford University；英文；课程讲义、作业与项目。重点阅读 meta-learning、few-shot learning、MAML、metric learning 和 task distribution，帮助理解支持集/查询集与双层目标。\href{https://web.stanford.edu/class/cs330/}{课程主页}。
\item \textbf{higher documentation: Top-Level Functions}；作者/平台：Facebook Research / higher；英文；API 教程文档。阅读 innerloop\_ctx 的 copy\_initial\_weights 与 track\_higher\_grads 参数，核对可微内循环、快速权重及 MAML 的梯度路径。\href{https://higher.readthedocs.io/en/latest/toplevel.html}{在线阅读}。
\item \textbf{一文入门元学习（Meta-Learning）（附代码）}；知乎；中文解说。明确任务划分、支持集与查询集，比较度量学习和梯度初始化；标题与主题据必应索引，原页受限。\href{https://zhuanlan.zhihu.com/p/136975128}{在线阅读}。
\item \textbf{通俗讲解元学习（Meta-Learning）}；CSDN；中文博客。对照内外层更新，检查测试任务是否泄漏到元训练；标题与主题据必应索引，原页受限。\href{https://blog.csdn.net/bruce\_\_ray/article/details/131144371}{在线阅读}。
\item \textbf{第十五节 2021 - 元学习 Meta Learning (一) - 元学习和机器学习一样也是三個步骤}；李宏毅；bilibili；中文课程。选看第135集，对照本章跨任务训练与快速适应。2026年10月4日官方公开元数据显示整套课程累计播放约299.5万，并非本分集播放量；已核对目录与计数，未逐帧观看。\href{https://www.bilibili.com/video/BV1Wv411h7kN/?p=135}{观看分集}。
\end{itemize}
\endgroup
\paragraph{本章小结。}
元学习把一个预测器的训练问题改写为适应过程的训练问题。度量方法通过支持集重建比较规则，原型网络把平方距离下的类均值作为任务状态；MAML 通过内外循环链式法则学习初始化，一阶方法以梯度近似换取计算效率，隐式方法则以驻点方程替代轨迹展开。学习优化器和记忆方法把适应放在更新规则或显式状态中，概率方法进一步保留少量证据下的多种任务解释。这些路线的共同目标都是适应后的查询风险，而不是仅把支持样本拟合得更好。

本章最重要的限制同样位于任务层面：训练任务的共同性未必能延续到新设备，过强先验、错误近邻和不稳定更新都可能放大偏差。只有设备级隔离、合法的支持查询划分、匹配的适应预算以及强预训练基线，才能说明元学习带来了真实收益。对同一设备重复切窗，或者利用元测试查询表现调参，都会改变原本宣称的学习问题。

全书从怎样定义样本和损失开始，经过优化、概率推断和深度模型，最后讨论了不同的数据组织与学习方式。回到设备现场，仍应先问手上有什么记录、要提前作出什么判断，再选择模型和训练办法，最后检查误差、不确定性与运行成本。接下来的结语把这些检查整理为五个问题。开始新任务时，把记录来源、可用时间和评价方法写清楚，比先选一个热门模型更有帮助。


\backmatter
\chapter*{结语：如何开始一个工业数据智能问题}
面对新的工业任务，先写下五件事：现在能看到哪些变量，要提前多久给出什么判断，线上什么时候拿得到每条记录，设备和工况会怎样变化，以及漏报、误报各会带来什么代价。例如要预测未来二十四小时内的故障，就必须同时写清标签确认时间、预测窗口和允许的告警延迟。五件事没有写清之前，模型结构的比较缺少明确依据。

回看全书，第一部分把实际问题写成数据、假设空间、损失和风险之间的关系，并说明训练误差为什么不能代替部署表现；第二部分把未知量纳入概率推理，区分参数估计、预测分布与决策所需的不确定性；第三部分根据时间、空间、语言、跨模态关系和物理规律组织模型；第四部分进一步改变数据到达、知识迁移、标注获取和协作学习的方式。这条顺序反映了一个基本判断：先明确需要保留的结构与需要承担的错误，再选择模型，而不是从某个热门模型反推问题。

数学推导使这种选择具有可检查的依据。最小二乘的正规方程揭示了可辨识性与数值条件的重要性，概率模型的似然揭示了损失函数背后的观测假设，泛化分析揭示了样本数量与复杂度之间的约束，梯度与近似推断揭示了可计算性所付出的代价。深度网络并没有取消这些问题；它只是把表示和参数空间扩展得更大。迁移、联邦或元学习也没有自动消除分布变化，而是为利用已有知识增加了新的结构假设。公式的价值因此不在于显得复杂，而在于让假设、推理和限制能够被逐项追问。

真正开始一个项目时，可以先固定预测时刻和标签窗口，建立按设备与时间隔离的评估集，再实现可解释的简单基线。只有在误差分析指出明确瓶颈之后，才引入更丰富的表示、不确定性估计或新的学习方式。每次修改都应留下可复核的比较：使用了哪些新增信息，改变了什么假设，改善了哪类错误，增加了多少计算与维护成本。在本文的维护案例中，风险预测、缺陷定位和规程问答最终都必须回到同一件事：现有证据是否足以支持下一步维护动作；若不足，系统应请求复检或补充观测，而不是用更肯定的输出掩盖知识的缺口。

因此，学完本书并不意味着为每种工业任务记住一个标准答案，而是形成一套可以反复使用的工作方式：把问题写清楚，把推导做完整，把近似讲明白，把证据与结论对应起来，并在环境变化后重新检验。能够说明一个方法何时不应使用，与能够实现它同样重要。

\bibliographystyle{plain}
\bibliography{references}
\end{document}